%
% $Id: SANDtemplate.tex,v 1.3 2007-12-13 21:27:14 rolf Exp $
% A template to build SAND reports. See the examples for more details and
% formatting suggestions. A command reference is available at
% http://www.cs.sandia.gov/~rolf/SANDreport
%
\documentclass[pdf,12pt,report]{SANDreport}

% ---------------------------------------------------------------------------- %
% User-defined packages
%
\usepackage{tabularx} % for making acronym table \textwidth
\usepackage{mathptmx}    % Use the Postscript Times font
\usepackage{amsmath}
\usepackage[]{caption}
\usepackage[margin=0.25cm]{subcaption}
\usepackage{wrapfig}
\usepackage{pbox}
\usepackage{placeins}
\usepackage{multirow}
\usepackage{longtable}
\usepackage{pdflscape}
\usepackage{bm}
\usepackage{pgfplots}
\usepackage{array}
%\usepackage{chngcntr}
%\counterwithin{table}{section}
%\counterwithin{figure}{section}
\usepackage{alphalph}
\usepackage[final]{pdfpages}
\usepackage{cite}
\renewcommand\thesubfigure{\alphalph{\value{subfigure}}}
\usepackage{listings}
\usepackage{inconsolata}
\usepackage{acronym}
%\usepackage{notoccite}
\usepackage[notextcomp]{stix}
\usepackage{pgfgantt}
\usepackage{mathptmx}
\usepackage{setspace}
\usepackage{soul}
\usepackage{float}
\usepackage[hypertexnames=false]{hyperref}
\usepackage{xcolor}
\hypersetup{
    colorlinks=true,
    linkcolor=.,%black,%{red!50!black},
    citecolor={blue!50!black},
    urlcolor={blue!80!black}
}
\usepackage{color}
\usepackage{bigstrut}
\usepackage{attachfile}
\usepackage{enumitem}
\usepackage{bookmark}
\setlist[description]{%
    topsep=0pt,
    itemsep=-5pt,
    font={\ttfamily},
}
\pdfpageattr {/Group << /S /Transparency /I true /CS /DeviceRGB>>}
\definecolor{light-gray}{gray}{0.98}
\definecolor{number-gray}{gray}{0.5}
\definecolor{comment-green}{rgb}{0,0.6,0}
\definecolor{string-gray}{gray}{0.0}

\lstset{ %
    backgroundcolor=\color{light-gray},% choose the background color; you must add \usepackage{color} or \usepackage{xcolor}
    basicstyle=\ttfamily\scriptsize, % the size of the fonts that are used for the code
    breakatwhitespace=false,           % sets if automatic breaks should only happen at whitespace
    breaklines=true,                   % sets automatic line breaking
    captionpos=b,                      % sets the caption-position to bottom
    commentstyle=\color{comment-green},% comment style
    % deletekeywords={...},              % if you want to delete keywords from the given language
    % escapeinside={\%*}{*)},            % if you want to add LaTeX within your code
    extendedchars=true,                % lets you use non-ASCII characters; for 8-bits encodings only, does not work with UTF-8
    frame=single,                      % adds a frame around the code
    framesep=4pt,
    keepspaces=true,                   % keeps spaces in text, useful for keeping indentation of code (possibly needs columns=flexible)
    keywordstyle=\color{blue},         % keyword style
    % language=Python,                   % the language of the code
    morekeywords={*,...},              % if you want to add more keywords to the set
    numbers=left,                      % where to put the line-numbers; possible values are (none, left, right)
    numbersep=10pt,                     % how far the line-numbers are from the code
    numberstyle=\scriptsize\ttfamily,  % the style that is used for the line-numbers
    rulecolor=\color{black},           % if not set, the frame-color may be changed on line-breaks within not-black text (e.g. comments (green here))
    showspaces=false,                  % show spaces everywhere adding particular underscores; it overrides 'showstringspaces'
    showstringspaces=false,            % underline spaces within strings only
    showtabs=false,                    % show tabs within strings adding particular underscores
    stepnumber=1,                      % the step between two line-numbers. If it's 1, each line will be numbered
    stringstyle=\color{string-gray},   % string literal style
    tabsize=2                         % sets default tabsize to 2 spaces
}

\newcommand{\inlineCode}{\lstinline[basicstyle=\bfseries\ttfamily]}

\newcommand{\multilinecell}[2][l]{%
    \begin{tabular}[#1]{@{}c@{}}#2\end{tabular}}

\newcommand{\overbar}[1]{\mkern 1.5mu\overline{\mkern-1.5mu#1\mkern-1.5mu}\mkern 1.5mu}

\newcommand{\needsattention}[1]{\colorbox{red!50}{NEEDS ATTENTION: #1}}

\newcommand{\GitlabProject}{\href{https://github.com/sandialabs/rattlesnake-vibration-controller}{Github project}}
\newcommand{\GitlabReleasesPage}{\href{https://github.com/sandialabs/rattlesnake-vibration-controller/releases}{Releases}}
\newcommand{\GitlabMasterDownload}{\href{https://github.com/sandialabs/rattlesnake-vibration-controller/archive/refs/heads/main.zip}{Download}}
\newcommand{\GitlabIssueTracker}{\href{https://github.com/sandialabs/rattlesnake-vibration-controller/issues}{Github Issue Tracker}}

% ---------------------------------------------------------------------------- %
% Set the title, author, and date
%
    \title{Rattlesnake User's Manual\\Version 3}
    \author{Daniel P. Rohe}        % Use First, Middle, Initial
    \date{}        % Leave this here but empty


% ---------------------------------------------------------------------------- %
% These are mandatory
%
\SANDnum{SAND2024-14378}        % e.g. \SANDnum{SAND2006-0420}
\SANDprintDate{Oct, 2024}    % Month, year
\SANDauthor{Daniel P. Rohe, Ryan Schultz, and Norman Hunter}        % One line, separated by commas


% ---------------------------------------------------------------------------- %
% These are optional
%
%\SANDrePrintDate{}    % May be repeated for successive printings
%\SANDsupersed{}{}    % {Old SAND number}{Old date}


% ---------------------------------------------------------------------------- %
% Build your markings. See example files and SAND Report Guide
%
    %\SANDreleaseType{}
    %\SANDmarkTopBottomCoverBackTitle{}
    %\SANDmarkBottomCover{}
    %\SANDmarkTopBottomCoverTitle{}
    %\SANDmarkTop{}
    %\SANDmarkBottom{}
    %\SANDmarkTopBottom{}
    %\SANDmarkCover{}
    %\SANDmarkCoverTitle{}


% ---------------------------------------------------------------------------- %
% Start the document
%
\begin{document}
    \maketitle

    % ------------------------------------------------------------------------ %
    % An Abstract is required for SAND reports
    %
    \begin{abstract}
        Rattlesnake is a combined-environments, multiple input/multiple output control system for dynamic excitation of structures under test.  It provides capabilities to control multiple responses on the part using multiple exciters using various control strategies.  Rattlesnake is written in the Python programming language to facilitate multiple input/multiple output vibration research by allowing users to prescribe custom control laws to the controller.  Rattlesnake can target multiple hardware devices, or even perform synthetic control to simulate a test virtually.  Rattlesnake has been used to execute control problems with up to 200 response channels and 24 shaker drives.  This document describes the functionality, architecture, and usage of the Rattlesnake controller to perform combined environments testing.
    \end{abstract}


    % ------------------------------------------------------------------------ %
    % An Acknowledgment section is optional but important
    %
%    \clearpage
%    \chapter*{Acknowledgment}
%    

    % ------------------------------------------------------------------------ %
    % The table of contents and list of figures and tables
    %
    \cleardoublepage        % TOC needs to start on an odd page
    \tableofcontents
    \listoffigures
    \listoftables


    % ---------------------------------------------------------------------- %
    % An optional preface or Foreword
%    \clearpage
%    \chapter*{Preface}
%    \addcontentsline{toc}{chapter}{Preface}


    % ---------------------------------------------------------------------- %
    % An optional executive summary
%    \clearpage
%    \chapter*{Summary}
%    \addcontentsline{toc}{chapter}{Summary}


    % ---------------------------------------------------------------------- %
    % An optional glossary. We don't want it to be numbered or in the TOC
    \clearpage
    \chapter*{Nomenclature}
    \acresetall
    \begin{acronym}[xxxxxxxxxx]
        \acro{6DoF}{six degree-of-freedom}
        \acro{API}{application programming interface}
        \acro{APSD}{auto-power spectral density}
        \acro{CCLD}{constant current line drive}
        \acro{COLA}{constant overlap and add}
        \acro{CPSD}{cross-power spectral density}
        \acro{FEM}{finite element model}
        \acro{FFT}{fast Fourier transform}
        \acro{FRF}{frequency response function}
        \acro{GUI}{graphical user interface}
        \acro{ICP}{integrated circuit piezoelectric}
        \acro{IDE}{integrated development environment}
        \acro{IEPE}{integrated electronics piezoelectric}
        \acro{IFFT}{inverse fast Fourier transform}
        \acro{JSON}{javascript object notation}
        \acro{MIMO}{multiple input/multiple output}
        \acro{ReST}{representational state transfer}
        \acro{RMS}{root-mean-square}
        \acro{SVD}{singular value decomposition}
        \acro{TRAC}{time response assurance criterion}
        \acro{UI}{user interface}
    \end{acronym}

     \clearpage
    \chapter*{Notation}
    \subsection*{Accents}
    \begin{description}
        \item[$\mathbf{x}$ \normalfont (bold, non-italic typeface)] Matrix or vector
        \item[$x$ \normalfont (non-bold, italic typeface)] Scalar or array
    \end{description}
    \subsection*{Variables}
    \begin{description}
        \item[$n_f$] Number of frequency lines
        \item[$n_c$] Number of control channels
        \item[$n_o$] Number of output signals
        \item[$n_s$] Number of samples in a signal
        \item[$n_os$] Number of samples in an oversampled output signal
        \item[$\mathbf{G}_{vv}$] CPSD matrix for the output signals
        \item[$\mathbf{G}_{xx}$] CPSD matrix for the responses (i.e. the specification)
        \item[$\mathbf{H}_{xv}$] Transfer function matrix between the responses and the output signals
        \item[$\mathbf{X}$] Response spectra (FFT)
        \item[$\mathbf{V}$] Output spectra (FFT)
        \item[$\mathbf{T}$] Transformation matrix
        \item[$\mathbf{A}$] State or system matrix for state space equations of motion
        \item[$\mathbf{B}$] Input matrix for state space equations of motion
        \item[$\mathbf{C}$] Output matrix for state space equations of motion
        \item[$\mathbf{D}$] Feedthrough or feedforward matrix for state space equations of motion
        \item[$\mathbf{x}$] State vector for state space equations of motion
        \item[$\mathbf{y}$] Output vector for state space equations of motion
        \item[$\mathbf{u}$] Input vector for state space equations of motion
        \item[$\mathbf{M}$] Mass matrix for 2nd-order differential equations of motion
        \item[$\mathbf{C}$] Damping matrix for 2nd-order differential equations of motion
        \item[$\mathbf{K}$] Stiffness matrix for 2nd-order differential equations of motion
        \item[$\mathbf{I}$] Identity matrix
    \end{description}
    \subsection*{Operations}
    \begin{description}
        \item[$\dottedsquare^+$] Pseudoinverse of a matrix
        \item[$\dottedsquare^H$] Complex-conjugate transpose of a matrix (hermitian)
    \end{description}

    % ---------------------------------------------------------------------- %
    % This is where the body of the report begins; usually with an Introduction
    %
    \SANDmain        % Start the main part of the report

    \acresetall
    
    \chapter{Introduction}\label{sec:introduction}
    The field of \ac{MIMO} vibration testing has grown substantially in the past few years.  \Ac{MIMO} vibration testing provides the capability to match a field environment more accurately and at more locations on the test article than traditional single-axis vibration testing.  Unfortunately, many existing vibration control systems are proprietary, which makes it difficult to implement new \ac{MIMO} techniques.  Currently, \ac{MIMO} vibration practitioners must either develop a control system from scratch to implement their ideas, or alternatively convince an \ac{MIMO} vibration software vendor to implement their ideas into existing devices, and neither of these approaches are conducive to the rapid and iterative nature of research.  The Rattlesnake framework was developed to overcome these limitations and facilitate \ac{MIMO} vibration research.  Rattlesnake is a \ac{MIMO} control system that provides the user the ability to overcome testing challenges by providing a flexible framework that can be extended and modified to meet testing demands.  Rattlesnake can run multiple environments simultaneously, providing a combined-environments capability that does not yet exist in commercial software packages.  It can target multiple hardware devices, or even perform control virtually using a state space model, \ac{FEM} results, or a SDynPy System.

    \begin{figure}[!b]
        \centering
        \includegraphics[width=\linewidth]{figures/Rattlesnake_Logo_Banner}
        \caption{Rattlesnake Logo}
        \label{fig:rattlesnakelogobanner}
    \end{figure}
    
    \section{Document Overview}\label{sec:intro_document_overview}
    This document provides information about the functionality in the Rattlesnake software, as well as instructions for how to use that functionality.  The document is divided into Parts targeting different aspects of the software.
    \begin{itemize}
        \item Part \ref{sec:rattlesnake_overview} provides an overview of the software as well as instructions for how to acquire and run the software.  Chapter \ref{sec:intro_running_rattlesnake} includes instructions for setting up the Python ecosystem required to run the software, if necessary.  Chapter \ref{sec:using_rattlesnake} describes the Rattlesnake \ac{UI}.  Each main portion of the Rattlesnake interface is described along with the parameters that should be defined within that interface.
        \item Part \ref{sec:rattlesnake_hardware} describes the hardware devices available to the Rattlesnake software, as well as the hardware-specific considerations in the controller.  Each Chapter in this Part is dedicated to a specific hardware device.  Synthetic or virtual control is also discussed in this part, as well as instructions to extend Rattlesnake to additional hardware devices.
        \item Part \ref{sec:rattlesnake_environments} describes the various control environments contained within the Rattlesnake software.  The environments defined within Rattlesnake are where the next output data that will be sent to the exciters are computed based off the previously acquired data.  Each chapter in this Part is dedicated to a environment type within Rattlesnake.  This chapter also provides instructions to combine environments as well as extend Rattlesnake to additional environments.
        \item The Appendices contain several example problems.  Users with a reasonable understanding of the Rattlesnake workflow can use these chapters as a kind of ``Quick Start'' guide to the software.  Appendix \ref{sec:example_nidaqmx} demonstrates a series of tests on a simple beam using a NI cDAQ hardware device.  Appendices \ref{sec:example_sdynpy} and \ref{sec:example_state_space} demonstrate virtual control problems using SDynPy System and State Space models, respectively, which only require a computer to run.
    \end{itemize}
    
    \part{Rattlesnake Overview}\label{sec:rattlesnake_overview}
    
    This Part provides a general overview of the Rattlesnake software.  Chapter \ref{sec:intro_running_rattlesnake} provides an overview of how to acquire and run the software.  Chapter \ref{sec:using_rattlesnake} describes the general software workflow.
    
    This Part will only deal with ``global'' Rattlesnake parameters.  For hardware-specific parameters, see Part \ref{sec:rattlesnake_hardware}.  For environment-specific parameters, see Part \ref{sec:rattlesnake_environments}.
    
    \chapter{Acquiring and Running Rattlesnake}\label{sec:intro_running_rattlesnake}

    Two methods are provided to acquire the Rattlesnake software.  The software can be downloaded as an executable and run directly with no other dependencies.  Alternatively, the software can be downloaded in its Python script form and run using a Python interpreter.  The former approach is simpler, but results in a larger file size and longer software loading time due to the necessity to pack the Python ecosystem into the executable for distribution and unpack it prior to execution.  The latter approach is more suited to users who wish to utilize the full functionality of the Rattlesnake framework, which would include activities such as coding up custom control laws.  In this case, it will be advantageous to have a Python ecosystem installed on the user's computer, so simply downloading the source code and executing it similarly to other Python scripts will potentially be easier than the executable approach.
    
    \section{Running from an Executable}\label{sec:intro_running_executable}
    Python code can be compiled into a single executable file, which makes it easier to distribute Python code.  The user need not have a Python distribution installed on their computer to simply run the executable, as the executable will contain the required Python interpreter and libraries compiled within it.  The executable approach has a few disadvantages.  The file size of the executable will generally be larger than the raw source code.  Additionally, the executable will generally be slower to start due to the necessary unpacking of the Python ecosystem from it.  Still, if a user is not familiar with Python, the executable will be the easiest approach to run the software.
    
    \subsection{Downloading the Executable}
    Executables for Rattlesnake are generally stored in the \GitlabProject{} on the \GitlabReleasesPage{} page.  A user can simply download the executable corresponding to the user's operating system and save it to their computer.  No installation is necessary to run this executable.
    
    \subsection{Running the Executable}
    Running Rattlesnake from an executable is as simple as running any other program.  Simply double click on the executable (or otherwise execute it) and the program will run.  Note that there may be a significant delay between executing the executable and the program appearing, as the executable will unpack a Python distribution into the user's temporary space in order to run the included Python code.
    
    As with any other executable program, users may create links to the program to put on the desktop or in the start menu to make accessing the program easier.
    
    It may be beneficial to run the executable through the command terminal, as otherwise if an error occurs in the program, the program may simply disappear without the user being aware of any errors.  When running through the command terminal, the user will be able to view any error messages if the program unexpectedly exits, which will be useful in diagnosing the issue if submitted to the issue tracker in the \GitlabProject{} (See Section \ref{sec:intro_obtaining_support} for more information).
    
    \section{Running from the Python script}
    The alternative to running the Rattlesnake as an executable is to run it as a Python script, which users familiar with the Python programming language should be accustomed to.  This approach provides the user the ability to modify code directly without needing to recompile an executable.  Additionally, if the user plans on developing custom control laws for the Rattlesnake (see, e.g., Chapter \ref{sec:rattlesnake_environments_mimo_random} for more information), they will generally require a Python ecosystem installed on their computer, so the running of the code as a script is not a great burden.
    
    \subsection{Setting up a Python Ecosystem}
    The first step to running Rattlesnake from its Python script is to install Python.  This can be done in multiple ways.  Python can be downloaded and installed from the \href{https://www.python.org/}{Python Website} directly.  When installed this way, Python will not include any of the numeric or scientific libraries such as \inlineCode{numpy} or \inlineCode{scipy}.  For this reason, many users will prefer to download a scientific Python distribution which contains many numeric or scientific libraries.  \href{https://www.anaconda.com/}{Anaconda} or \href{https://winpython.github.io/}{WinPython} are popular distributions.
    
    Regardless of the distribution selected, users will need to install packages that Rattlesnake depends on.  The \GitlabProject{} repository contains a \inlineCode{requirements.txt} file that can be used by the Python package manager \inlineCode{pip} to install all required dependencies:
    \begin{center}
        \inlineCode{pip install -r requirements.txt}
    \end{center}
    
    Note that if running through a corporate or university firewall, the proxy may need to be specified in \inlineCode{pip}.  Additionally, on some networks the Python package repositories must be added as trusted hosts.  Such a command may look like
    \begin{center}
        \inlineCode{pip --proxy <proxy_address> install --trusted-host pypi.org}\\ \inlineCode{ --trusted-host files.pythonhosted.org -r requirements.txt}
    \end{center}
    where \inlineCode{<proxy_address>} is the address of the proxy.
    
    \subsection{Downloading the Python Code}
    
    With Python installed the Rattlesnake code can be downloaded from the \GitlabProject{} repository.  For users who aim to develop the Rattlesnake software, the preferred approach to acquire the software is to use Git to clone the repository.  For users who only wish to use the software, the code can be \GitlabMasterDownload ed in a \inlineCode{zip} file or other archive format and extracted to the user's computer.  
    
    \subsection{Running Rattlesnake}
    If Python is on the user's path, the user can simply call
    \begin{center}
        \inlineCode{python rattlesnake.py}
    \end{center}
    from a command line in the directory in which Rattlesnake was downloaded to execute the software.  If Python is not on the user's path, it will be necessary to provide the command line the full path to the Python executable.
    
    Many users will find it more comfortable to forgo the command line and launch Rattlesnake directly from their favorite \ac{IDE}.  Note that the Rattlesnake uses the \inlineCode{multiprocessing} Python package to spawn several processes, and these processes sometimes do not play nice with \ac{IDE} consoles.  Therefore if running from an \ac{IDE}, the \ac{IDE} should be instructed to run the \inlineCode{rattlesnake.py} script from an external system terminal rather than a terminal inside the \ac{IDE}.  When executing in an external terminal, it is again useful to keep the terminal active after execution in case of errors occurring.  The error message and traceback displayed in the terminal will be instrumental in debugging the source of the error.  In Spyder, this can be configured per file from the \inlineCode{Run} menu, so the run settings for \inlineCode{rattlesnake.py} should be set as shown in Figure \ref{fig:spyderrunconfiguration}.
    
    \begin{figure}
        \centering
        \includegraphics[width=\linewidth]{figures/spyder_run_configuration}
        \caption{Spyder run configuration showing execution in an external system terminal as well as allowing interaction with the Python console after execution.}
        \label{fig:spyderrunconfiguration}
    \end{figure}
    
    \section{Computational Requirements}
    Rattlesnake is a process-heavy software; it spawns processes for each environment, as well as for various portions of the controller that should operate in parallel.  In general, the core controller utilizes 2-3 full processes.  Various environments will also utilize multiple processes; for example, the \ac{MIMO} Random Vibration uses 3 main processes to compute spectral quantities (FRFs, CPSDs, etc.), perform control calculations, and generate time histories simultaneously.  If running virtual control, keep in mind that acquisition processes will be more fully subscribed due to the need to integrate the equations of motion rather than just read data off the data acquisition hardware.
    
    While the exact computational requirements will depend on the channel count of the test and size of the control computations, the authors have had success using a 6-core processor with 32 GB RAM for multi-environment control approximately 20 control channels and 4 outputs.  For a 200 acquisition channel test with 50 control channels and 8 outputs, the authors needed to upgrade to a 16 core, 32 GB RAM computer.
    
    \section{Obtaining Support}\label{sec:intro_obtaining_support}
    Rattlesnake was developed by a small team as a research tool, and as such will not be as polished as commercial vibration software.  Therefore, it should be expected that bugs and errors will occur from time to time.  If a bug occurs, please report it by creating a new issue in the \GitlabIssueTracker.  This will notify the development team of the issue and they can work to solve it.  If users have a feature request, these can also be submitted to the \GitlabIssueTracker.
    
    The issue tracker provides templates for Bugs, Feature Requests, and Questions.  The more thoroughly that these can be filled out, the greater the chance that the development team can solve the issue.  For bug reports, files can also be attached to the issue, so the users can include the \inlineCode{Rattlesnake.log} file that is generated for each run of the Rattlesnake, as well as any screenshot of error messages that get shown in dialog boxes or the command window.  The question template can be used for basic support questions, which will be answered by the developers as their time permits.  Users should certainly consult this User's Manual and the Rattlesnake Source Code first to ensure their question is not covered by it.
    
    \chapter{Using Rattlesnake}\label{sec:using_rattlesnake}

    This chapter will describe how to use Rattlesnake through its \ac{GUI}.  Rattlesnake is capable of running several different types of control, therefore the \ac{GUI} may look different for different tests.  In general, the \ac{GUI} consists of a tabbed interface across the top of the main window, and users must complete each tab before proceeding to the next.  The tabs that exist in a given test will depend on which control type is being run.  For example, in a combined environments test (see Section \ref{sec:rattlesnake_environments_combining_environments}) such as the one shown in Figure \ref{fig:rattlesnakemainguicombinedenvironments}, there is a \inlineCode{Test Profile} tab that allows the user to define a testing timeline.  Additionally, environments such as the \Ac{MIMO} Random Vibration environment (see Section \ref{sec:rattlesnake_environments_mimo_random}) require a system identification phase where the controller identifies relationships between the output signals and the control degrees of freedom.  Therefore, tests using the \Ac{MIMO} Random Vibration environment will also have a \inlineCode{System Identification} and \inlineCode{Test Predictions} tab.  Figure \ref{fig:rattlesnakemainguitimegeneration}, on the other hand, shows the \ac{GUI} for a test that only utilizes the Time History environment (see Section \ref{sec:rattlesnake_environments_time_generator}) so these optional tabs are not displayed.
    
    Users of Rattlesnake must be aware that depending on their test configuration, their \ac{GUI} may not appear identical to images shown in this User's Manual.  Additionally, users should be aware that the \ac{GUI} library used by this software will inherit stylistic features from the operating system.  There may therefore be cosmetic differences between the images of the \ac{GUI} shown in this document and the \ac{GUI} seen by the user.  All images in this document were created using Microsoft Windows 10 or Windows 11 operating systems, so users with Mac or Linux operating systems will note a difference in \ac{GUI} appearance.
    
    Note that the Rattlesnake enforces an order to operations when defining a particular test by enabling and disabling tabs in the \ac{GUI}.  Initially, only the first tab will be enabled.  As the users complete each tab, the next tab will become available.  In Figures \ref{fig:rattlesnakemainguicombinedenvironments} and \ref{fig:rattlesnakemainguitimegeneration}, it can be seen that only the initial tab is enabled, and subsequent tabs are disabled.
    
    \begin{figure}
        \centering
        \includegraphics[width=\linewidth]{figures/Rattlesnake_Main_GUI_Combined_Environments}
        \caption{Rattlesnake \ac{GUI} tabs when running a combined environments test with an environment that requires a system identification.}
        \label{fig:rattlesnakemainguicombinedenvironments}
    \end{figure}
    
    \begin{figure}
        \centering
        \includegraphics[width=\linewidth]{figures/Rattlesnake_Main_GUI_Time_Generation}
        \caption[Rattlesnake \ac{GUI} tabs when running a single environment with no system identification phase.]{Rattlesnake \ac{GUI} when running a single environment with no system identification phase.}
        \label{fig:rattlesnakemainguitimegeneration}
    \end{figure}
    
    \section{Environment Selection}
    
    When Rattlesnake is opened, the first \ac{GUI} window that the user will see allows the user to select the environment that they will run (Figure \ref{fig:environmentselection}).  Users can select a single environment, or alternatively select a combined environments test (see Section \ref{sec:rattlesnake_environments_combining_environments}).  The selection made in this dialog box will determine which tabs are set up in the main \ac{GUI}.
    
    \begin{figure}
        \centering
        \includegraphics[width=0.4\linewidth]{figures/environment_selection}
        \caption{Initial Rattlesnake dialog to select the type of control that will be run.}
        \label{fig:environmentselection}
    \end{figure}
    
    
    \section{Global Data Acquisition Settings}\label{sec:using_rattlesnake_data_acq_settings}
    The \inlineCode{Data Acquisition Setup} tab of the Rattlesnake \ac{GUI} specifies the global test parameters that the controller will use.  Parameters are determined to be global when they affect all environments or the controller itself.  The three main sections of this portion of the interface are the Channel Table, Environment Table, and Global Data Acquisition Parameters.  Figure \ref{fig:dataacquisitionsetup} shows this 

    \begin{figure}
        \centering
        \includegraphics[width=\linewidth]{figures/data_acquisition_setup}
        \caption{Data Acqisition Setup tab in the Rattlesnake Controller where the Channel Table, Environment Table, and Data Acquisition Parameters are specified.}
        \label{fig:dataacquisitionsetup}
    \end{figure}
    
    \subsection{Channel Table}\label{sec:channel_table}
    The channel table specifies how the instrument channels in a given test are connected to the data acquisition hardware, as well as how the data read from those channels are used by the software.
    
    In general, for a given test there will be a set of excitation devices that use the output signals from Rattlesnake as well as instrumentation to record the test article's responses to those exciters.  Rattlesnake requires each instrument (or each channel on each instrument for multi-axial instruments) as well as each excitation device to have a row in the channel table.  This is perhaps contrary to other control software where only the response channels need to be set up in the channel table.  However, to maintain the flexibility to run multiple types of hardware devices, some of which having limitations to their triggering capabilities, Rattlesnake must read in the signal from its output directly in order to be able to synchronize its outputs and the responses to those outputs.  Therefore, for all Rattlesnake test setups, the output signal should be split using a tee to the exciter and the corresponding input channel.  Because of this requirement, one should keep in mind that the number of acquisition channels required on the hardware device for a given test is actually the number of responses plus the number of outputs.  Figure \ref{fig:lanxisourceteelabelled} shows a schematic of a four acquisition channel, two output channel LAN-XI module set up for use with Rattlesnake.
    
    \begin{figure}
    \centering
    \includegraphics[width=0.7\linewidth]{figures/lanxi_source_tee_labelled}
    \caption{Output channels teed to acquisition channels so they can be read by the controller.}
    \label{fig:lanxisourceteelabelled}
    \end{figure}
   
    The required data input into the channel table varies with the physical or virtual hardware used for the test.  For device-specific channel table requirements, see the appropriate section of Part \ref{sec:rattlesnake_hardware}.  In general, the entries to the channel table are as follows:
    \begin{description}
        \item[Node Number] Determines the instrumentation position on the test article.  While not used directly by the controller except to label plots, it is important for book-keeping and test documentation.  The node number will generally correspond to a node in a test geometry or \ac{FEM}.
        \item[Node Direction] Determines the instrumentation direction on the test article at the position specified by the Node Number.  Again, this is not used directly by the controller except to label plots, but it is important for book-keeping.  The Node Direction will generally correspond to the node's local coordinate system if one exists in the test geometry.
        \item[Comment] Provides space for additional information about a channel that may not be captured by the Node Number and Node Direction.
        \item[Serial Number] The serial number of the instrument used for the given channel.  This field is not used by the controller but will be stored with the test data and is important for data traceability to know which instruments were used to measure which channels.
        \item[Triax DoF] The degree of freedom on a given instrument corresponding to the given channel.  This is primarily used to distinguish between the three axes of a triaxial accelerometer, but has the potential to be used for other multi-axis instrumentation types such as strain gauge rosettes.
        \item[Sensitivity] The sensitivity of the instrument in millivolts per Engineering Unit.  This is used to transform the acquired data from a raw voltage to a engineering quantity such as acceleration or force.
        \item[Engineering Unit] The unit in which the measured signal for the given instrument will be reported.  Certain hardware will limit the units that can be specified: see Part \ref{sec:rattlesnake_hardware} for more information.
        \item[Make] The name of the instrument's manufacturer, used for data traceability.
        \item[Model] The product name or model number of the instrument, used for data traceability.
        \item[Expiration] The expiration date of the instrument's calibration certificate.  Note that this is only for data traceability; no checking of this date with the current data to ensure a valid calibration is performed by the software.
        \item[Physical Device] The reference to a physical device attached to the computer.  The entries in this field will be specific to the acquisition hardware being used for a given test.  For virtual control, this column must be filled to specify that a given channel is active.  See Part \ref{sec:rattlesnake_hardware} for more information.
        \item[Physical Channel] The reference to a channel on a physical device attached to the computer.  The entries in this column will be specific to the acquisition hardware being used for a given test.  See Part \ref{sec:rattlesnake_hardware} for more information.
        \item[Channel Type] The type of the channel being used for a given test, such as Acceleration, Force, or Voltage.  The allowable entries in this column will be specific to the acquisition hardware being used for a given test.  See Part \ref{sec:rattlesnake_hardware} for more information.
        \item[Minimum Value (V)] The minimum voltage that the data acquisition system can handle during a test.  This is used to set the range on the data acquisition system.  For hardware devices with symmetric ranges (e.g. $\pm$10V), this column can be left blank.
        \item[Maximum Value (V)]  The maximum voltage that the data acquisition system can handle during a test.  This is used to set the range on the data acquisition system.  For hardware devices with symmetric ranges (e.g. $\pm$10V), this column is used to set the maximum and minimum voltage values.
        \item[Coupling] The coupling used by the data acquisition system.  This may include filtering in addition to AC/DC coupling, which is dependent on the hardware being used for a given test.  See Part \ref{sec:rattlesnake_hardware} for more information.
        \item[Excitation Source] Used to specify the signal conditioning that is required by the instrument.  This column is generally where the \ac{CCLD}/\ac{IEPE}/\ac{ICP} is specified for a given hardware device.  See Part \ref{sec:rattlesnake_hardware} for more information.
        \item[Current Excitation (A)] Used to specify the excitation current sent to the device for signal conditioning.  Depending on whether the device has a fixed or variable excitation current, this field may be left empty.  This can also be left empty if no signal conditioning is provided by the data acquisition system.  See Part \ref{sec:rattlesnake_hardware} for more information.
        \item[Feedback Device] For output channels, this is the reference to the output or excitation device that is being fed back into the current channel's Physical device.  If the current channel is not an output channel, it should be left empty.  A populated Feedback Device column tells the controller that the given channel is an output channel.
        \item[Feedback Channel] For output channels, this is the reference to the output channel on the output or excitation device that is being fed back into the current channel's Physical Device.  As an example using generic device and channel names, if ``Channel 2'' on ``Generator 1'' is teed off to ``Channel 3'' on ``Acquisition Card 2'', the corresponding row in the channel table would have ``Acquisition Card 2'' specified as the Physical Device, ``Channel 3'' specified as the Physical Channel, ``Generator 1'' specified as the Feedback Device'' and ``Channel 2'' specified as the feedback channel.
        \item[Warning Level] A warning level can be implemented for each channel.  The warning level is specified in the same units as the Engineering Unit column.  When a channel hits the warning limit, it will be flagged as Yellow in the Channel Monitor (see Section \ref{sec:channel_monitor}).  The warning level can be left blank if no warning is desired.
        \item[Abort Level] An abort level can be implemented for each channel.  The abort level is specified in the same units as the Engineering Unit column.  When a channel hits the abort limit, it will be flagged as Red in the Channel Monitor (see Section \ref{sec:channel_monitor}).  The controller will also shut down if an abort level is reached.  The abort level can be left blank if no abort is desired.
    \end{description}

    To limit the tediousness of inputting channel table information into the \ac{GUI} by hand, the channel table can be loaded from an Excel spreadsheet or Comma-separated-value file.  A channel table can be loaded by clicking the \inlineCode{Load Channel Table} button under the channel table, which will bring up a file selection dialog, enabling the user to select a file to load.  For convenience, a template Excel spreadsheet is attached to this PDF: \attachfile{attachments/channel_table_template.xlsx}.  A template Excel file can also be generated by creating a test in Rattlesnake and saving the empty channel table by clicking the \inlineCode{Save Channel Table} button under the channel table.  If a channel table is filled out in Rattlesnake's \ac{GUI}, its contents will be saved to the file as well.
    
    A complete test can be loaded by clicking the \inlineCode{Load Test from File} button.  See Section \ref{sec:load_rattlesnake_test} for more details.

    \subsection{Environment Table}
    
    For combined environments tests, an environment table is also provided to the right of the standard Channel Table.  This table specifies which channels are used by which environments.  A channel can be used for multiple environments, a single environment, or no environments.  Channels used by no environments will still be measured and streamed to disk, but will not be sent to any environment for use in the respective control approaches.  The environment table is also used to specify which excitation devices are used by which environment.
    
    For single environment tests, the environment table is not visible, and the software assumes that all channels in the channel table are used by the single environment.
    
    If importing a channel table from an Excel spreadsheet for a combined environments test, the Environment Table can be specified as the columns after the main Channel Table information (starting in Column X with one column for each environment) with the environment name specified in row 2 and an entry (e.g. an \inlineCode{X} or some other mark) in the row corresponding to a given channel if that channel is used for the given environment.
    
    \subsection{Data Acquisition Parameters}
    The final portion of the \inlineCode{Data Acquisition Setup} tab specifies data acquisition parameters.  These parameters may change depending on the hardware selected.
    
    \begin{description}
        \item[Hardware Selector] The physical or virtual hardware used for the test.  See Part \ref{sec:rattlesnake_hardware} for hardware specific details of the controller.  For some devices, a file selector window will appear will appear when the device is selected, as that device needs more information to operate.  This is primarily the case for virtual hardware where some model of the test article must be loaded.  This is also used when a specific hardware device needs to access external functionality in a library such as a \inlineCode{dll} file.
        \item[Sample Rate] The sample rate of all hardware devices used for the test.  Some devices will have arbitrary sample rates, and some devices have fixed sample rates, so the options available will depend on the acquisition hardware being used.
        \item[Time per Read] The amount of data that the acquisition system will acquire with each read from the hardware.  By reading data in chunks, hardware input/output operations with relatively large overhead can be limited, and the buffer gives the controller time to catch up if e.g. the operating system decides to start a computationally intensive task in the background of the computer.  Note that specifying large numbers for this quantity (e.g. 10) will reduce the responsiveness of the controller, because the controller will potentially not receive the acquired data until 10 seconds after it was acquired.  Note also that this does not need to correspond to the Samples per Analysis Frame or any other signal processing parameter used by an environment.  Each environment should be buffered such that it creates appropriately sized analysis windows from the differently sized acquisition chunks.
        \item[Time per Write] The amount of data that the output system will write with each write to the hardware.  By writing data in chunks, hardware input/output operations with relatively large overhead can be limited, and the buffer gives the control time to catch up if e.g. the operating system decides to start a computationally intensive task in the background of the computer.   Note that this output is also buffered, so it does not need to be equal to the size of the data that will be created during each control loop of the controller.
        \item[Maximum Acquisition Processes] For specific hardware devices with large channel count tests, it can be difficult to pull down large quantities of data fast enough for the controller to keep up using a single process.  This option allows the user to specify how many processes can be given to the acquisition system to be used to stream data off the hardware.  Note that too many processes will bog down the computer, and too few will result in the controller falling behind.  Generally about 20-40 channels per processor is sufficient, but this will depend on the sample rate.  For higher sample rates, more processors may be needed.
        \item[Integration Oversample] For synthetic hardware devices that integrate equations of motion, an integration oversample factor can be specified.  This factor will be applied to the sample rate to determine the time step for the integration.  Generally a factor of 10 is sufficient for reasonably accurate data without significant computational expense.
    \end{description}
    
    \subsection{Initialize Data Acquisition}
    With the Data Acquisition Settings specified in the \ac{GUI}, the Data Acquisition can be initialized by pressing the \inlineCode{Initialize Data Acquisition} button.  At this point, the controller will go through and create the programming interfaces to the hardware device, specify the sampling parameters, and create the channels on the devices.  The software will then proceed to the next tab.
    
    Figure \ref{fig:dataacquisitionsetuptab} shows a completed \inlineCode{Data Acquisition Setup} tab with three response channels and one output channel for a test with two environments ``A'' and ``B''.  The first response and output channels are used by both environments, while the second response channel is used only by environment ``A'' and the third response channel is used only by environment ``B''.
    
    \begin{figure}
        \centering
        \includegraphics[width=\linewidth]{figures/data_acquisition_setup_tab}
        \caption{Example of a completed \inlineCode{Data Acquisition Setup} tab with three response channels and one output channel.}
        \label{fig:dataacquisitionsetuptab}
    \end{figure}
    
    \section{Environment Definition}\label{sec:using_rattlesnake_environment_definition}
    The \inlineCode{Environment Definition} tab is the second tab in the Rattlesnake software.  It is in this tab that the various environments are defined.  The main tab will have one sub-tab for each environment, as shown in Figure \ref{fig:environmentdefinitionsubtabs}.
    
    \begin{figure}
        \centering
        \includegraphics[width=0.4\linewidth]{figures/environment_definition_subtabs}
        \caption{Sub-tabs for environments ``A'' and ``B'' in the \inlineCode{Environment Definition} tab.}
        \label{fig:environmentdefinitionsubtabs}
    \end{figure}
    
    Different environment types will have different parameters that can be set.  See Part \ref{sec:rattlesnake_environments} for a description of each environment type in Rattlesnake and the parameters that define it.
    
    When all environments are defined, the \inlineCode{Initialize Environments} button can be pressed to proceed to the next portion of the controller.
    
    \section{System Identification}\label{sec:using_rattlesnake_system_identification}
    With the environments defined, the controller proceeds to the \inlineCode{System Identification} tab if required by any environment, shown in Figure \ref{fig:systemidentification}.  During this phase of the controller, the controller will develop relationships between the excitation signals and the responses of the test article to those excitation signals.  It will also make a measurement of the noise floor of the test.
    
    
    \begin{figure}
        \centering
        \includegraphics[width=\linewidth]{figures/system_identification}
        \caption{System identification tab showing various signals and spectral quantities that can be used to evaluate the test.}
        \label{fig:systemidentification}
    \end{figure}
    
    Not all environment types will require a system identification.  For environments that simply stream output data, a system identification will generally not be required.  However for any environment that aims to produce an output that creates some response on the test article, a system identification will be required to understand the relationships between the excitation signals and the response signals.
    
    There will be one sub-tab for each environment that requires a System Identification.  System identification must be run for each sub-tab before the test can be run.  When system identification is performed, the software will first perform a noise floor measurement, where all channels are recorded, but no excitation signal is provided.  After the noise floor calculation completes, the system identification will begin.
    
    The \inlineCode{System Identification} tab has been significantly overhauled since the previous version of controller.  The system identification now has a number of dedicated parameters on its tab that the user can select.  These are:
    
    \begin{description}
        \item[Samples per Frame] The number of samples used in each measurement frame
        \item[Averaging Type] The type of averaging used to compute the spectral quantities.  Linear averaging weights each measurement frame equally.  Exponential averaging weights more recent frames more heavily.
        \item[Noise Averages] The number of averages used in the noise characterization.
        \item[System ID Averages] The number of averages used in the System Identification characterization.
        \item[Averaging Coefficient] If Exponential Averaging is used, this is the weighting of the most recent frame compared to the weighting of the previous frames.  If the averaging coefficient is $\alpha$, then the most recent frame will be weighted $\alpha$, the frame before that will be weighted $\alpha(1-\alpha)$, the frame before that will be $\alpha(1-\alpha)^2$, etc.
        \item[Estimator] The estimator used to compute transfer functions between voltage signals and responses.
        \item[Level] The RMS voltage level used for the system identification
        \item[Level Ramp Time] The startup and shutdown time of the system identification.
    \end{description}

    The new system identification tab also gives the option to select the signal to use for system identification.
    
    \begin{description}
        \item[Signal Type] The type of signal that will be used for System Identification.  This can be Random, Burst Random, Chirp, or Pseudorandom.  Random is the most flexible, but requires a Hann window which can distort data.  Burst Random does not require a window, but the response signal must decay within the measurement frame.  Chirp and Pseudorandom do not require windows, and do not need to decay, but they are only useful for environments with a single excitation device.
        \item[Window] The window function used for the system identification signal
        \item[Overlap] The overlap percentage between measurement frames used in System Identification
        \item[On Fraction] The fraction of the frame that the Burst Random signal is active for
        \item[Pretrigger] The fraction of the frame before the Burst Random signal starts
        \item[Ramp Fraction] The fraction of the Burst Random On Fraction that is used to ramp up and ramp down.
    \end{description}

    The system identification phase can stream time data to disk by selecting a streaming file and clicking the \inlineCode{Stream Time Data} checkbox.  If streaming time data, the noise measurement will be saved to the variable name \inlineCode{time_data} and the system identification measurement will be saved to the variable name \inlineCode{time_data_1} (see Section \ref{sec:using_rattlesnake_output_files} for more information on the structure of this file).  The spectral data from the system identification can be saved to disk by clicking the \inlineCode{Save System Identification Data} button and selecting the file.
    
    To run the system identification, there are buttons to Preview the Noise or System ID characterizations.  When ready, the \inlineCode{Start} button can be clicked.  It will run a Noise Characterization for the specified number of \inlineCode{Noise Averages}, and then subsequently run the System Identification characterization for the specified number of \inlineCode{System ID Averages}.  If the user wishes to run either the noise or system identification phases continuously, they can click the \inlineCode{Preview Noise} or \inlineCode{Preview System ID} buttons.
    
    Data will be plotted as the system identification proceeds.  The signals to visualize can be selected by clicking one or more of the \inlineCode{References} or \inlineCode{Responses} channels on the right side of the screen.  In the bottom right corner, there are options to show or hide various quantities of interest.  The \inlineCode{System Identification} tab can show the following:
    
    \begin{description}
        \item[Time Data] Raw Time Data as it is streamed from the data acquisition system.  Only data used in spectral computations is shown, so the users shouldn't see any data that is ramping up or down as if the controller is starting or stopping.
        \item[Transfer Function] These are the transfer functions between the References (e.g. voltage signals) and the Responses (e.g. Accelerations or Forces).  The controller will use these transfer functions to develop excitation signals that will be played to the shaker to achieve a desired response.
        \item[Impulse Response]  The impulse response of the system can be visualized.  This can be helpful to debug issues with transient control.
        \item[Coherence and Conditioning] Coherence will be displayed so the user can judge how satisfactory the input/output relationships that are developed are.  If coherence is poor, it could suggest that the controller won't be able to control the structure properly.
        The condition number of the Transfer Function Matrix is also displayed.  This can be useful to determine what level of regularization a control law might need to implement.
        \item[Levels] The autopower spectral density of each signal will be displayed both for the system identification as well as for the noise characterization.  This can help identify if the system identification is high enough out of the noise floor.
    \end{description}

    \section{Test Predictions}\label{sec:using_rattlesnake_test_predictions}
    Once the system identification phase completes, the controller will compute a test prediction for each environment where system identification was completed.  This prediction will be based on the measured transfer functions between output signals and measured responses, as well as the environment parameters specified on the \inlineCode{Environment Definition} tab.  Predictions will be made both for outputs required as well as response accuracy.  These predictions will be displayed on the \inlineCode{Test Predictions} tab.
    
    \section{Test Profiles}\label{sec:using_rattlesnake_test_profiles}
    The \inlineCode{Test Profile} tab gives the user the ability to set up a test timeline for complex combined environments tests.  The user can add a list of events that will be executed at certain times during the test.  The tab will also display a graphical representation of the test timeline.
    
    Events can be added or removed from the test timeline by clicking the \inlineCode{Add Event} or \inlineCode{Remove Event} buttons.  Users can also load a series of events from an Excel spreadsheet.
    
    For each event, the following parameters are defined:
    \begin{description}
        \item[Timestamp (s)] The time in seconds after the timeline has started that the event will be executed.
        \item[Environment] The environment in which the event will occur.
        \item[Operation] The operation that will occur to the event.  Each environment defines its own set of operations that can be executed through the test profile interface.
        \item[Data] Any additional data that the operation requires.  For example, if a ``Set Test Level'' event is chosen, the Data field should specify the value that the test level is set to.
    \end{description}

    Figure \ref{fig:testprofile} shows an example of a test profile that ramps up the test level of environment ``A'' from -6 to 0 dB, and then subsequently starts environment ``B''.

    \begin{figure}
        \centering
        \includegraphics[width=\linewidth]{figures/test_profile}
        \caption{Example test profile showing a ramp up of test level for environment ``A'' and subsequently starting environment ``B''}
        \label{fig:testprofile}
    \end{figure}
    
    \section{Run Test}\label{sec:using_rattlesnake_run_test}
    The \inlineCode{Run Test} tab is where Rattlesnake finally runs the test.  This tab again has sub-tabs for the different environments in the test, however these sub-tabs will not be enabled until the data acquisition system is armed.
    
    Rattlesnake gives the user many options to save data to the disk through a set of Radio buttons at the top of this tab.  These options are:
    
    \begin{description}
        \item[No Streaming] Do not save data to disk, just run the test.
        \item[Start Streaming from Test Profile Instruction] Selecting this option allows data to be saved to disk after a ``Global Start Streaming'' event from the test profile is executed.  This allows the user to fine tune at which point in the test data is acquired.
        \item[Start Streaming at <environment> Target Test Level] Selecting this option starts streaming data when the selected environment hits its target test level.  This can be useful if, for example in a random environment, the user wishes to start at a low level and slowly creep up to the target test level.  If all data is saved, it might require a large amount of file space, so instead only the data at the test level of interest can be saved.
        \item[Start Streaming Immediately] Saves all data from the time the first environment starts until the data acquisition system is disarmed.
        \item[Manually Start/Stop Streaming] Allows the user to start and stop the measurement periodically throughout the test.  A \inlineCode{Start Streaming} button will appear when this option is selected.  When clicked, the button will change to \inlineCode{Stop Streaming}.  Multiple data streams can be saved in a given test.  These will be stored to separate variables in the output NetCDF4 file (see Section \ref{sec:using_rattlesnake_output_files}).
    \end{description}

    When streaming data, it is important to note that the software does not stop streaming until the data acquisition system is disarmed by pressing the \inlineCode{Disarm Data Acquisition} button.  This is because for a combined-environments test, the environments may have down-time between them where no environment is running, and that data should still be saved.
    
    The \inlineCode{Run Test} tab contains global \inlineCode{Arm Data Acquisition} and \inlineCode{Disarm Data Acquisition} buttons that start and stop the data acquisition system.  When the data acquisition system is armed, the user can no longer change streaming options, and the sub-tabs for each environment are enabled.  The user can then start or stop each environment manually using the \inlineCode{Start Environment} or \inlineCode{Stop Environment} buttons on each environment's sub-tab.  The sub-tab for each environment is described more thoroughly in Part \ref{sec:rattlesnake_environments}.
    
    Alternatively, the user can start or stop the test profile by clicking on the \inlineCode{Start Profile} or \inlineCode{Stop Profile} buttons respectively.  The profile capability also includes the option to switch the active environment sub-tab when an event is executed so the user can see the results.  Note that the profile options only appear when a profile has been defined on the \inlineCode{Test Profile} tab.
    
    Figure \ref{fig:runtesttab} shows an example \inlineCode{Run Test} tab with test profile events.
    
    \begin{figure}
        \centering
        \includegraphics[width=\linewidth]{figures/run_test_tab}
        \caption{Run Test Tab}
        \label{fig:runtesttab}
    \end{figure}

    \section{Rattlesnake Output Files}\label{sec:using_rattlesnake_output_files}
    After data is acquired, the user may wish to analyze or plot the data acquired for a given test report.  Rattlesnake stores data in a self-documenting netCDF file \cite{unidata2019_netcdf}, which can be read by multiple platforms.  The output file is described as self-documenting because it contains all parameters necessary to reconstruct a given test using the Rattlesnake controller.  Any parameter that is set by the user in the \ac{GUI} is stored to the netCDF file.
    
    A full description of the netCDF file format is out of this document's scope, but the important points are briefly described here.  NetCDF files have a number of data structures.  Variables are multi-dimensional arrays of data.  Dimensions describe the axes of the variable arrays.  Attributes are used to store small data such as scalars or 1D arrays.  NetCDF files can be separated into different groups, and each group can have its own variables, dimensions, and attributes.
    
    The Rattlesnake output files contain the following data members
    \subsection{NetCDF Dimensions}
    \begin{description}
        \item[{response\_channels}] The number of response channels in a given test
        \item[{output\_channels}] The number of output channels in a given test
        \item[{time\_samples}] The number of time samples measured in the file, this dimension can expand as more data is acquired.
        \item[{time\_samples\_$X$}] If manual streaming is used and streaming is started multiple times, each subsequent stream will have the \inlineCode{time_samples} name with an underscore and appended number (e.g. \inlineCode{time_samples_1}, \inlineCode{time_samples_2})
        \item[{num\_environments}] The total number of environments in the test
    \end{description}
    \subsection{NetCDF Attributes}
    \begin{description}
        \item[sample\_rate] The global sample rate of the data acquisition system
        \item[time\_per\_write] The amount of data put to the output hardware per write operation, in seconds
        \item[time\_to\_read] The amount of data read from the acquisition hardware per read operation, in seconds
        \item[hardware] The hardware index used for the test.  0 -- National Instruments NI-DAQmx, 1 -- HBK LAN-XI Open API, 2 -- Data Physics Quattro, 3 -- Data Physics 900 Series, 4 -- Virtual Control defined by Exodus Modal Solution, 5 -- Virtual Control defined by State Space Matrices, 6 -- Virtual Control defined with a SDynPy System
        \item[hardware\_file] The path to the file used to define the Virtual test article, or the path to the external code library used by the data acquisition hardware.  Otherwise, it will be \inlineCode{None}
        \item[maximum\_acquisition\_processes] The maximum number of processes that the LAN-XI hardware can use for acquisition
        \item[output\_oversample] The oversample used either due to sample rate restrictions on the data acquisition system, or due to oversampling the integration
    \end{description}
    \subsection{NetCDF Variables}
    \begin{description}
        \item[time\_data] The measured data from the test. Type: 64-bit float; Dimensions: \inlineCode{response_channels} $\times$ \inlineCode{time_samples}
        \item[time\_data\_$X$] If manual streaming is used and streaming is started multiple times, each subsequent stream will have the \inlineCode{time_data} name with an underscore and appended number (e.g. \inlineCode{time_data_1}, \inlineCode{time_data_2})
        \item[environment\_names] The name of each environment. Type: string; Dimensions: \inlineCode{num_environments}
        \item[environment\_active\_channels] The channels active in each environment.  1 if active, 0 if not. Type: 8-bit int; Dimensions: \inlineCode{response_channels} $\times$ \inlineCode{num_environments}
    \end{description}
    \subsection{Channels Group}
    The netCDF files from Rattlesnake store all channel information into a separate group called \inlineCode{channels}.  Inside the \inlineCode{channels} group, there is a variable for each column of the channel table.  See Section \ref{sec:channel_table} for more complete descriptions of each channel variable.
    \begin{description}
        \item[/channels/node\_number] The node number of each channel.  Type: str; Dimensions: \inlineCode{response_channels}
        \item[/channels/node\_direction] The instrument direction of each channel.  Type: str; Dimensions: \inlineCode{response_channels}
        \item[/channels/comment] The commend for each channel.  Type: str; Dimensions: \inlineCode{response_channels}
        \item[/channels/serial\_number] The serial number of the instrument for each channel.  Type: str; Dimensions: \inlineCode{response_channels}
        \item[/channels/triax\_dof] The sensor degree of freedom for each channel.  Type: str; Dimensions: \inlineCode{response_channels}
        \item[/channels/sensitivity] The sensitivity of the instrument for each channel.  Type: str; Dimensions: \inlineCode{response_channels}
        \item[/channels/unit] The engineering unit of the instrument for each channel. Type: str; Dimensions: \inlineCode{response_channels}
        \item[/channels/make] The manufacturer of the instrument for each channel.  Type: str; Dimensions: \inlineCode{response_channels}
        \item[/channels/model] The model number or product name of the instrument for each channel.  Type: str; Dimensions: \inlineCode{response_channels}
        \item[/channels/expiration] The expiration date of the instrument's calibration for each channel.  Type: str; Dimensions: \inlineCode{response_channels}
        \item[/channels/physical\_device] The physical device that the instrument is connected to for each channel.  Type: str; Dimensions: \inlineCode{response_channels}
        \item[/channels/physical\_channel] The channel in the physical device that the instrument is attached to for each channel.  Type: str; Dimensions: \inlineCode{response_channels}
        \item[/channels/channel\_type] The type of quantity that is measured by the channel.  Type: str; Dimensions: \inlineCode{response_channels}
        \item[/channels/minimum\_value] The minimum voltage that the channel can accept.  Type: str; Dimensions: \inlineCode{response_channels}
        \item[/channels/maximum\_value] The maximum voltage that the channel can accept.  Type: str; Dimensions: \inlineCode{response_channels}
        \item[/channels/coupling] The coupling type used by each channel (AC/DC/filter/etc.).  Type: str; Dimensions: \inlineCode{response_channels}
        \item[/channels/excitation\_source] The excitation source for each channel, used to specify \ac{CCLD}/\ac{ICP}/\ac{IEPE}.  Type: str; Dimensions: \inlineCode{response_channels}
        \item[/channels/excitation] The excitation current value used in the signal conditioning for each channel.  Type: str; Dimensions: \inlineCode{response_channels}
        \item[/channels/feedback\_device] The device that the channel's generator originates from if the channel is an output channel.  Type: str; Dimensions: \inlineCode{response_channels}
        \item[/channels/feedback\_channel] The channel that the channel's generator originates from if the channel is an output channel.  Type: str; Dimensions: \inlineCode{response_channels}
        \item[/channels/warning\_level] The warning level of each channel.  Type: str; Dimensions: \inlineCode{response_channels}
        \item[/channels/abort\_level] The abort level of each channel.  Type: str; Dimensions: \inlineCode{response_channels}
    \end{description}

    \subsection{Environment Groups}
    Environment-specific attributes, dimensions, and variables are also stored within a group corresponding to each environment.  For example, in the case where there were two environments ``A'' and ``B'', parameters specific to environment ``A'' would be stored within the group ``A'' in the netCDF file, and similarly for ``B''.  See Part \ref{sec:rattlesnake_environments} for more information on environment-specific parameters.
    
    \subsection{Reading Rattlesnake Output Files using Python}
    To read data from a netCDF using Python, it is recommended to use the \inlineCode{netCDF4} Python package.  This library is a dependency of Rattlesnake, so if the user is not running Rattlesnake via an executable, this package should already be installed in the user's Python ecosystem.
    
    \inlineCode{netCDF4} provides a sleek Python interface into the data of a netCDF4 file.  This section will assume the command \inlineCode{import netCDF4 as nc4} was used to import the package, so \inlineCode{nc4} is used as a shorter alias.
    
    A netCDF4 dataset can be opened using the following command:
    \begin{lstlisting}[language=Python,numbers=none]
dataset = nc4.Dataset('path/to/netcdf4/file.nc4')
\end{lstlisting}
    after which all data can be accessed through the \inlineCode{dataset} object.
    
    Attribute names can be queried using the \inlineCode{dataset.ncattrs()} function and the attribute values can be accessed directly from the \inlineCode{dataset} object using that name.
    \begin{lstlisting}[language=Python,numbers=none]
>>> dataset.ncattrs()
['sample_rate',
'samples_per_write',
'samples_per_read',
'hardware',
'hardware_file']

>>> dataset.sample_rate
2048
\end{lstlisting}

    Dimensions can be accessed using the \inlineCode{dataset.dimensions} property, which gives a Python \inlineCode{dict} where the keys are the dimension names and the values are references to the dimension.  The size of the dimension can be accessed using the \inlineCode{size} parameter in each dimension object.
    
\begin{lstlisting}[language=Python,numbers=none]
>>> dataset.dimensions
{'response_channels': <class 'netCDF4._netCDF4.Dimension'>: name = 'response_channels', size = 30,
'output_channels': <class 'netCDF4._netCDF4.Dimension'>: name = 'output_channels', size = 3,
'time_samples': <class 'netCDF4._netCDF4.Dimension'> (unlimited): name = 'time_samples', size = 31745,
'num_environments': <class 'netCDF4._netCDF4.Dimension'>: name = 'num_environments', size = 2}

>>> dataset.dimensions['response_channels'].size
30
\end{lstlisting}

    Variables can be accessed similarly to dimensions using the \inlineCode{dataset.variables} property.  Variables have many properties that may be interesting to the users, including the netCDF dimensions that were used to size the variable (accessible with the \inlineCode{dimensions} parameter) or the actual shape of the array (accessible with the \inlineCode{shape} parameter).  The data inside the dimension can be accessed by slicing or indexing the array, or simply passing it to a \inlineCode{numpy} array.  Note that slicing or indexing the variable returns the data in a \inlineCode{numpy} masked array which allows data to potentially to be missing from the array.  Rattlesnake does not use the missing data capabilities of the netCDF file, so data can safely be transformed directly to a regular \inlineCode{numpy} array.
    
\begin{lstlisting}[language=Python,numbers=none]
>>> dataset.variables
{'time_data': <class 'netCDF4._netCDF4.Variable'>
float64 time_data(response_channels, time_samples)
unlimited dimensions: time_samples
current shape = (30, 31745)
filling on, default _FillValue of 9.969209968386869e+36 used,
'environment_names': <class 'netCDF4._netCDF4.Variable'>
vlen environment_names(num_environments)
vlen data type: <class 'str'>
unlimited dimensions: 
current shape = (2,),
'environment_active_channels': <class 'netCDF4._netCDF4.Variable'>
int8 environment_active_channels(response_channels, num_environments)
unlimited dimensions: 
current shape = (30, 2)
filling on, default _FillValue of -127 ignored}

# Get the dimensions used by the variable
>>> dataset.variables['time_data'].dimensions
('response_channels', 'time_samples')

# Get the shape of the variable
>>> dataset.variables['time_data'].shape
(30, 31745)

# Access via slice returns a masked array
>>> dataset.variables['time_data'][0,0]
masked_array(data=-0.00312098,
mask=False,
fill_value=1e+20)

# Can pass directly to a numpy array to get the full variable data
>>> np.array(dataset.variables['time_data'])
array([[-3.12098493e-03,  4.26820006e-03,  3.77395182e-03, ...,
         2.00690958e-01,  3.38505511e-01,  0.00000000e+00],
        [-6.10438702e-03,  1.50628999e-02, -1.50619535e-02, ...,
         2.67639515e-01,  5.50047023e-01,  0.00000000e+00],
        [-3.42732089e-03,  7.76593927e-03, -2.66239267e-03, ...,
         2.05434816e-01,  3.21815820e-01,  0.00000000e+00],
         ...,
        [ 3.71743658e-06, -8.77497995e-08, -8.80558595e-06, ...,
         -5.26559214e-06,  0.00000000e+00,  0.00000000e+00],
        [-1.32020650e-05,  2.74453772e-05, -2.01551409e-05, ...,
         -1.02501347e-05,  0.00000000e+00,  0.00000000e+00],
        [-5.96816619e-07, -1.47868461e-05, -5.24157875e-05, ...,
         -1.61722666e-05,  0.00000000e+00,  0.00000000e+00]])
\end{lstlisting}
    
    Group names in the netCDF dataset can be queried using \inlineCode{dataset.groups}, which returns a dictionary similar to the dimensions and variables.  Groups can also be accessed by indexing the dataset directly with the group name.  A group object can be treated exactly the same as the root-level dataset, and will have its own set of attributes, dimensions, and variables.
    
\begin{lstlisting}[language=Python,numbers=none]
>>> dataset['channels'].variables['node_number']
<class 'netCDF4._netCDF4.Variable'>
vlen control(response_channels)
vlen data type: <class 'str'>
path = /channels
unlimited dimensions: 
current shape = (30,)
\end{lstlisting}
    
    \subsection{Reading Rattlesnake Output Files using Matlab}
    Matlab can also be used to read netCDF files from Rattlesnake.  The Matlab \inlineCode{ncdisp} function can be used to quickly determine which parameters are in a file.
    
\begin{lstlisting}[language=Matlab,numbers=none]
>>> ncdisp('path/to/netcdf/file.nc4')
Source:
           path/to/netcdf/file.nc4
Format:
           netcdf4
Global Attributes:
           sample_rate       = 2048
           samples_per_write = 512
           samples_per_read  = 512
           hardware          = 2
           hardware_file     = 'path/to/hardware/file.exo'
Dimensions:
           response_channels = 30
           output_channels   = 3
           time_samples      = 31745 (UNLIMITED)
           num_environments  = 2
Variables:
    time_data                  
           Size:       31745x30
           Dimensions: time_samples,response_channels
           Datatype:   double
    environment_names          
           Size:       2x1
           Dimensions: num_environments
           Datatype:   UNSUPPORTED DATATYPE
    environment_active_channels
           Size:       2x30
           Dimensions: num_environments,response_channels
           Datatype:   int8
Groups:
    /channels/
        Variables:
            node_number      
                   Size:       30x1
                   Dimensions: /response_channels
                   Datatype:   UNSUPPORTED DATATYPE
       .            
       .
       .
\end{lstlisting}

    Attributes, dimensions, and other metadata can be read into Matlab using the \inlineCode{ncinfo} function.  Variables information must be read using the \inlineCode{ncread} function.
\begin{lstlisting}[language=matlab,numbers=none]
>>> finfo = ncinfo('path/to/netcdf/file.nc4')

finfo = 

  struct with fields:

      Filename: 'C:\Users\dprohe\Documents\Local_Respositories\Combined_Environments_Controller\test_data\BARC_Exodus_Test\barc_combined.nc4'
          Name: '/'
    Dimensions: [1x4 struct]
     Variables: [1x3 struct]
    Attributes: [1x5 struct]
        Groups: [1x3 struct]
        Format: 'netcdf4'
        
>>> finfo.Dimensions(1)

ans = 

  struct with fields:

         Name: 'response_channels'
       Length: 30
    Unlimited: 0
    
>>> time_data = ncread('path/to/netcdf/file.nc4','time_data')
\end{lstlisting}

    One issue with the Matlab interface is that string variables are unsupported.  This means that the majority of the channel information cannot be read through the Matlab netCDF interface.  However, they can be read using the lower level \inlineCode{h5read} function.

\begin{lstlisting}[language=matlab,numbers=none]
>>> ncread(file,'channels/node_number')
Error using netcdf.getVar (line 137)
12 is not a recognized netCDF datatype.

Error in internal.matlab.imagesci.nc/read (line 605)
                data = netcdf.getVar(gid, varid);

Error in ncread (line 66)
vardata = ncObj.read(varName, varargin{:});

>>> h5read(file,'/channels/node_number')

ans =

  30x1 cell array
\end{lstlisting}

    \section{Loading Rattlesnake Tests}\label{sec:load_rattlesnake_test}
    
    It can be tedious to set up a test from scratch each time a test is to be run, so Rattlesnake offers two ways to load test settings from files.
    
    On the \inlineCode{Data Acquisition Setup} page, selecting the \inlineCode{Load Test From File} button allows the user to load in a netCDF data file that was output from Rattlesnake.  As all the test metadata is stored to this file, Rattlesnake can read the file and set itself up accordingly to reproduce a given test.  Note that difficulties may arise using this approach if parameters specified by file paths are no longer valid.  For example, if the control law is read from a given file on one computer, but the file is in a different place on a separate computer, Rattlesnake will not be able to find the file.
    
    The second way to load in an entire test is by using the Test Profile functionality in the Combined Environments mode.  While this capability was designed to make it easier to load in complex multi-environment test setups, it can be used just as effectively for single environment tests.  See Chapter \ref{sec:rattlesnake_environments_combining_environments} for more information.

    \section{Channel Monitor}\label{sec:channel_monitor}
    
    To aid with understanding the test levels and headroom available for the sensors in the test, a Channel Monitor is available where the levels are shown for each channel.  The channel monitor is displayed by clicking on the \inlineCode{Channel Monitor} button on the lower left side of the GUI.  The display shows both an instantaneous level (green) as well as a running historical maximum (blue).  If a channel reaches the Warning or Abort level, it will be flagged with a yellow or red tint, respectively.  These warnings ``latch''; once the level is reached, it will stay highlighted in the channel monitor until the \inlineCode{Clear Alerts} button is clicked.  Figure \ref{fig:channelmonitor} shows an example channel monitor.
    
    \begin{figure}
        \centering
        \includegraphics[width=0.6\linewidth]{figures/channel_monitor}
        \caption{View of the Channel Monitor dialog box showing several channels that have reached the ``warning'' level (highlighted yellow) and one channel that has reached the ``abort'' level (highlighted red).}
        \label{fig:channelmonitor}
    \end{figure}
    
    The aspect ratio of the Channel Monitor can be customized to different sizes modifying the \inlineCode{Channels per Row}.

    \part{Rattlesnake Hardware Devices}\label{sec:rattlesnake_hardware}

%    \chapter{Rattlesnake Hardware}\label{sec:rattlesnake_hardware}
    
    Designed for flexibility, Rattlesnake can be used with multiple hardware devices and even perform virtual control using a synthetic data acquisition system.  This Part will cover the hardware-specific implementation details that must be considered when running Rattlesnake with each hardware device.
    
    Rattlesnake is designed so there is minimal differences in software workflow when using different hardware devices.  Nonetheless, there are some slight differences in how channels and devices must be specified, and these differences are primarily found on the \inlineCode{Data Acquisition Setup} tab in the Rattlesnake software.
    
    The Chapters in this Part document the hardware devices that are able to be used by Rattlesnake, as well as the virtual devices that can be used to simulate control.  Chapter \ref{sec:rattlesnake_hardware_nidaqmx} describes the implementation and utilization of NI-DAQmx devices.  Chapter \ref{sec:rattlesnake_hardware_lanxi} describes the HBK LAN-XI hardware.  Chapters \ref{sec:rattlesnake_hardware_dp_quattro} and \ref{sec:rattlesnake_hardware_dp900} describe the Data Physics Quattro and 900-series hardware devices, respectively.  Chapters \ref{sec:rattlesnake_hardware_state_space}, \ref{sec:rattlesnake_hardware_virtual_control}, and \ref{sec:rattlesnake_hardware_sdynpy_system} describe virtual or synthetic hardware devices that are defined using state space matrices, eigensolution results stored in an exodus file, or a SDynPy System object, respectively.
    
    If a user is interested in implementing a new hardware device, Chapter \ref{sec:rattlesnake_hardware_implementation} describes some of the things to be aware of.  Implementation of new hardware devices will require a good amount of knowledge of the Rattlesnake architecture.
    
    \chapter{NI-DAQmx Devices}\label{sec:rattlesnake_hardware_nidaqmx}
    Rattlesnake is able to run National Instruments devices that utilize the NI-DAQmx programming interface.  See the \href{https://www.ni.com/pdf/manuals/374768ae.html#MXSupport}{NI-DAQmx Documentation} for a list of supported devices for this programming interface.  Note that users must have the proper drivers installed\footnote{Drivers can be found at \url{https://www.ni.com/en-us/support/downloads/drivers.html}} in order to communicate with their devices, though users need not have LabView or other NI software installed.  See the instruction manual or online documentation for the specific device in use.
    
    Appendix \ref{sec:example_nidaqmx} shows a complete example problem using a NI data acquisition system.
    
    \section{Setting up the Channel Table for NI-DAQmx Devices}
    This section lists the channel table requirements specific to NI-DAQmx.
    
     NI-DAQmx channels are identified by a device name and a channel name.  Device names vary depending on the type of device used.  For example, USB devices may simply be called \inlineCode{Dev#}, where \inlineCode{#} is a device number.  cDAQ devices might be called \inlineCode{cDAQ#Mod#} where the first \inlineCode{#} denotes the data acquisition system number and the second represents the module within the data acquisition system.  PXI/PXIe chassis devices will be similarly named \inlineCode{PXI#Slot#} where the first \inlineCode{#} is the chassis number and the second is the card number on the chassis.  In general, the names of NI-DAQmx devices that are attached to a given computer can be found using the National Instruments Measurement Automation Explorer (NI MAX).
        
     Channel names for NI-DAQmx devices generally are called \inlineCode{ai} for analog input and \inlineCode{ao} for analog output.  A four acquisition channel, one output device might have acquisition channels \inlineCode{ai0}, \inlineCode{ai1}, \inlineCode{ai2}, and \inlineCode{ai3} and output channel \inlineCode{ao0}.  Again, see the Measurement Automation Explorer to determine the channels that exist on each device.
    
    The channel type of a given channel determines what parameters are used and required for that channel.  Valid channel types are \inlineCode{Acceleration}, \inlineCode{Force}, or \inlineCode{Voltage}.
    
    \inlineCode{Acceleration} channel types must have a sensitivity specified in mV/G in the channel table, as the only valid Engineering Unit for an acceleration channel is \inlineCode{G}.
    
    \inlineCode{Force} channel types must also have a sensitivity specified.  Valid units for a force channel are pounds (specified by \inlineCode{lb}, \inlineCode{pound}, \inlineCode{pounds}, \inlineCode{lbf}, \inlineCode{lbs}, or \inlineCode{lbfs}, case insensitive) or newtons (specified by \inlineCode{n}, \inlineCode{newton}, \inlineCode{newtons}, or \inlineCode{ns}, case insensitive).
    
    \inlineCode{Voltage} channel types need not have a sensitivity or unit specified, as they will always be reported in volts.  A best practice is to fill out these columns anyways with the correct values \inlineCode{V} for engineering unit and sensitivity of 1000 mV/V.  \textbf{Note that specifying a different sensitivity for a voltage channel WILL NOT result in the voltage being scaled.  The NI-DAQmx implementation does not have the ability to scale voltage channels.  Users should specify 1000 mV/V as the sensitivity}.  Note that if a sensitivity is not specified, the Warning and Abort limits may not be correctly computed, as they rely on sensitivity information to convert between a measured raw voltage and the correct sensitivity unit.
    
    The \inlineCode{Physical Device} and \inlineCode{Physical Channel} (as well as \inlineCode{Feedback Device} and \inlineCode{Feedback Channel} for output channels) should correspond to the device and channel names in the Measurement Automation Explorer.  For example, channel \inlineCode{ai3} on device \inlineCode{PXI1Slot4} would have \inlineCode{Physical Device} set to \inlineCode{PXI1Slot4} and \inlineCode{Physical Channel} set to \inlineCode{ai3}.
    
    Maximum and Minimum Values should be specified based on the levels expected for the given test, taking into account that they are not outside the range acceptable for the device.
    
    The Coupling column in the channel table is not currently used by the NI-DAQmx system, rather the coupling is specified automatically by the channel type (Acceleration and Force are AC coupled, Voltage is DC coupled).
    
    Excitation Source should be set to either \inlineCode{Internal} or \inlineCode{None}.  If set to \inlineCode{Internal}, the device will generate the \ac{ICP}/\ac{CCLD}/\ac{IEPE} signal conditioning required by the sensor.  If set to \inlineCode{None}, no signal conditioning will be provided by the hardware device.
    
    Current Excitation should typically be set to 0.004 A (4 mA) unless the sensor requires a different current to be provided.  The Current Excitation value is only used if the Excitation Source is set to \inlineCode{Internal}.  If set to \inlineCode{None}, no current is generated.
    
    \section{Hardware Parameters}
    Besides the sample rate, no additional hardware-specific parameters must be specified for NI hardware.  Figure \ref{fig:nidaqmxdataacquisitionparameters} shows the parameters for NI-DAQmx hardware devices.
    
    \textbf{WARNING:} Some NI-DAQmx devices have discrete allowable sample rates, while others can have any sample rate that is desired up to some maximum value.  Please refer to the documentation of your device to determine which sample rates are allowable for your device.  The NI-DAQmx drivers, when provided an incompatible sample rate, often simply select the closest available rate or the next highest rate, which can result in data being acquired at a rate that is not equivalent to the rate selected in the software.  Additionally, further issues can occur when the sample rate of an output device is not compatible with the sample rate of an acquisition device, meaning the output is being delivered at a different rate than it is being measured, resulting in inconsistent data.  \textbf{Currently Rattlesnake does not do very rigorous checks to determine if the specified sample rate is allowable, so it falls on the user to ensure that it is!}
    
    \begin{figure}
    \centering
    \includegraphics[width=0.7\linewidth]{figures/nidaqmx_data_acquisition_parameters}
    \caption{NI-DAQmx Data Acquisition Parameters}
    \label{fig:nidaqmxdataacquisitionparameters}
    \end{figure}
    
    \section{Implementation Details}
    This section contains details on the NI-DAQmx implementation in Rattlesnake, which may be helpful for users when diagnosing issues that arise in the controller.
    
    \subsection{NI-DAQmx Tasks}
    NI-DAQmx hardware interfaces are defined within Tasks.  The acquisition exists within one task.  The output exists within one or more more separate tasks with one task per type of output card.  This multi-task output enables, for example, two different types of output cards to be used in the same chassis.  
    
    \subsection{Sampling Parameters}
    Both acquisition and output operate in continuous mode.  This means that if the controller cannot output signals fast enough and the hardware runs out of samples to generate, the controller will throw an error and stop abruptly.  A buffer size of three times the number of samples per write is specified in the source code which gives a cushion against the controller falling behind.  The output will accept a set of samples into the output buffer whenever there is less than two writes worth of samples remaining in the buffer.  The buffer size is computed by subtracting the total samples generated by the output task from the current write position in the output tasks buffer.
    
    \subsection{Starting the Hardware}
    When a measurement is started, the output is set up and started first.  The output task uses a start trigger from an analog input task (\inlineCode{/<physical_device_name>/ai/StartTrigger}) as its trigger, so once it starts, it effectively waits for the acquisition to be set up and started.  To determine which start trigger is used, it checks if the device is a cDAQ device.  If the device is a cDAQ device, it uses the start trigger from the cDAQ chassis.  If it is not a cDAQ device, it utilizes the start trigger from the first analog input channel.  This configuration has been tested with both cDAQ and PXI devices.  When the acquisition is started, its start trigger will trigger the output to start outputting signals simultaneously.  Because the system tries to use the analog input start trigger for all outputs, users may have trouble daisy-chaining multiple chassis together if the trigger signal is not able to be passed between all of the chassis.
    
    \chapter{LAN-XI Devices}\label{sec:rattlesnake_hardware_lanxi}
    Rattlesnake is able to run HBK LAN-XI devices using the hardware's \href{https://www.bksv.com/en/instruments/daq-data-acquisition/lan-xi-daq-system/open-api}{OpenAPI}.  Rattlesnake communicates with the LAN-XI hardware via Ethernet, so for best results, the computer's Ethernet card should be connected with the LAN-XI hardware rather than using a USB Ethernet dongle, which in the author's experience can lead to lower network speeds and more time required to pull data off the hardware.
    
    \section{Setting up the Channel Table for LAN-XI Devices}
    This section describes the process to set a channel table in Rattlesnake for the LAN-XI hardware.
    
    LAN-XI devices are defined by an IP address, which is displayed on each module when it is plugged into a computer or a LAN-XI frame.  This IP address should be specified as the \inlineCode{Physical Device} or \inlineCode{Feedback Device} in the channel table for a acquisition or output channel, respectively.  The \inlineCode{Physical Channel} or \inlineCode{Feedback Channel} range from 1 to the number of channels on the module.  Figure \ref{fig:lanxiphysicaldevice} shows an example case for setting up a LAN-XI test.
    
    \begin{figure}
    \centering
    \includegraphics[width=0.7\linewidth]{figures/lanxi_physical_device}
    \caption[Physical Channel and Feedback Channels for LAN-XI modules.]{Physical Channel and Feedback Channels for LAN-XI modules.  The left device would have the Physical Device 169.254.113.201 and Physical Channels 1, 2, 3, 4, 5, and 6.  The right device would have Physical Device 169.254.211.152 and Physical Channels 1, 2, 3, and 4.  The right device would also have Feedback Device 169.254.211.152 and Feedback Channels 1 and 2}
    \label{fig:lanxiphysicaldevice}
    \end{figure}
    
    The \inlineCode{Maximum Value} column in the channel table is used to set the range on the LAN-XI hardware.  The only two valid options for LAN-XI hardware ranges are \inlineCode{10} and \inlineCode{31.6}.  No other range is allowable.  The Minimum Value column is not used and can be left blank.
    
    The Coupling column in the channel table is used to specify the filter used in the LAN-XI hardware.  Valid values for coupling are \inlineCode{DC}, \inlineCode{0.7 Hz}, \inlineCode{7.0 Hz}, \inlineCode{22.4 Hz}, or \inlineCode{Intensity}.
    
    Current Excitation Source is used to specify if a channel uses \ac{CCLD} or not.  If \ac{CCLD} is to be used on a given channel the Current Excitation Source column should contain \inlineCode{CCLD}.  If \ac{CCLD} is not to be used for that channel, the column can be left blank.
    
    \section{Hardware Parameters}
    LAN-XI hardware devices have discrete sample rates, which are powers of 2 staring with 4096 samples per second.  The minimum sample rate of the generator is 16,384 Hz, so the output must be over-sampled if the acquisition sample rate is less than 16,384 samples per second.  Environments defined in Rattlesnake must be able to handle output oversampling when required by the hardware device.
    
    For large channel count tests, Rattlesnake struggles to pull data off the acquisition device fast enough using just one process.  Therefore, a maximum number of acquisition processes can be specified.  If too few processes are used, it will take longer to read data off the hardware than it took to acquire the data.  This will result in the controller falling behind and the data buffer on the hardware slowly filling.  Alternatively if too many processes are used, the computer running Rattlesnake will get bogged down swapping between processes, and other parts of the controller (particularly the GUI) may become slow.  Generally, 20-40 channels per process is a reasonable rule of thumb, though this will depend on the Sample Rate of the hardware.
    
    \begin{figure}
    \centering
    \includegraphics[width=\linewidth]{figures/lanxi_data_acquisition_parameters}
    \caption{LAN-XI data acquisition parameters.}
    \label{fig:lanxidataacquisitionparameters}
    \end{figure}
    
    \section{Implementation Details}
    This section contains details on the LAN-XI implementation in Rattlesnake, which may be helpful for users when diagnosing issues that arise in the controller.
    
    \subsection{ReST API}
    Rattlesnake communicates with the LAN-XI using a \ac{ReST} interface.  Rattlesnake sends and receives \ac{JSON} packages that define state transitions in the hardware using HTTP commands.
    
    \subsection{Setting up the LAN-XI}
    The LAN-XI hardware is set up primarily by the Output process of the controller.  The LAN-XI hardware can be used in either single or multi-module mode.  It will generally take a while to set up the LAN-XI configuration.  Lights will flash on the LAN-XI hardware while the setup is occurring, and information will be printed to the command terminal that appears when Rattlesnake is run.  When the LAN-XI setup is completed, the statement \inlineCode{Data Acquisition Ready for Acquire} will be printed to the command terminal.  Do not start a test prior to seeing this message.
    
    Various network issues and firewall settings can block or slow down data transfer between the LAN-XI hardware and the computer running Rattlesnake.  These issues are out of the scope of this document to diagnose and correct; users are encouraged to contact their network administrator if such issues arise.
    
    It has been found that for larger channel count tests (typically three or more 11-card frames used simultaneously) that individual cards can hang during this setup process.  The authors currently do not know why this occurs, as the exact same setup configuration causes no issues with lower channel counts.  Cycling power on the hardware devices has been found to resolve this issue.
    
    \subsection{Acquisition Processes}
    The LAN-XI interface used by Rattlesnake uses multiple processes to acquire data from the hardware.  When starting to acquire data, Rattlesnake starts a number of processes less than or equal to the maximum number of acquisition processes allowed on the \inlineCode{Data Acquisition Setup} tab.  Each process will generally handle multiple hardware modules, with each module having one socket over which the data communication occurs.  After all sockets are connected, the measurement is started.
    
    Raw data is read from each module by the various acquisition processes and put into queues that are read by the main acquisition process.  The main acquisition process takes these data and assembles them into the data array that is required by the controller.
    
    When an acquisition activity ends, the Rattlesnake process will attempt to recover the LAN-XI subprocesses.  When all processes have been recovered, the LAN-XI interface will print \inlineCode{All processes recovered, ready for next acquire.} to the command terminal.  Periodically, one of these processes may not be recovered successfully, which will cause the acquisition process to hang.  It is currently unknown why this occurs.
    
    \subsection{Decoding Acquisition Data}
    Raw acquisition data is provided to the controller from the hardware as bytes that must be interpreted correctly to be meaningful.  A data stream from the hardware consists of messages transmitted sequentially.  The message header is always 28 bytes long and specifies what type of message is being transmitted as well as the total length of the message.  This then allows the socket reader to receive the rest of the message.  Message types read by the Rattlesnake can either be ``interpretations'' or ``signals''. Interpretation messages specify how the subsequent signal data is to be interpreted and are sent whenever a signal changes.  Interpretations provide an offset and a scale factor to apply to the raw acquisition data to create the physical measurements.  Signal messages carry the raw acquisition data.  Rattlesnake uses Kaitai Struct to decode the binary data stream into Python objects which are accessed by the controller.
    
    \subsection{Encoding Output Data}
    Data is provided to the hardware for output in bytes as well.  LAN-XI hardware accepts output data as the 32-bit integers with only the upper 24 bits used.  Floating point signals that are desired to be output are divided by 10 and multiplied by 8,372,224.  These values are then truncated to integers and converted to bytes.  These bytes are then sent via socket communications to the generator on the LAN-XI module.
    
    \subsection{Output Oversampling}
    The generator on the LAN-XI module always runs at its full rate of 131,072 samples per second.  If the acquisition sample rate is equal to or greater than 16,384 samples per second, the LAN-XI hardware performs over-sampling automatically.  However, if the sample rate is less than 16,384, Rattlesnake must over-sample its output.  It is left to each environment within Rattlesnake to determine how to handle the oversampling of its output data.  For \ac{FFT}-based environments, this can be as simple as padding the FFT with zeros prior to creating the signal.  For time-based environments, more thought must be given to ensure signals are not made to be discontinuous by the up-sampling procedure.

    \subsection{Starting Up and Running the LAN-XI Acquisition}
    When an acquisition is started, the LAN-XI will immediately start acquiring data.  Rattlesnake may take some time to set up all of the required network connections to the device, so there may be a delay between starting the acquisition and receiving data in the Rattlesnake software.  During this time, the buffer on the hardware will start to fill up.  Once Rattlesnake starts pulling data off the LAN-XI device, the buffer will generally empty, as Rattlesnake ideally can pull data off the device faster than it is acquired.  For each read, Rattlesnake will print the amount of time it took to read to the command terminal.  Initially, the value reported will be smaller than the \inlineCode{Time per Read} value that is specified on the \inlineCode{Data Acquisition Setup} tab, because the software will be pulling data off of the hardware buffer that has accumulated during the delay in starting up the measurement.  After the buffer is emptied, users should see that the value reported will hover around \inlineCode{Time per Read} value that is specified on the \inlineCode{Data Acquisition Setup} tab, because that is the amount of time it takes to fill the buffer enough for the next read.  The value may be lower or higher for any given read, but on average, it should be approximately equal to the \inlineCode{Time per Read} value.  For maximal controller responsiveness, we want the hardware buffer to be empty so we can react to new data as quickly as possible, so it is therefore a good practice to not start any environment until the time that it is taking to read data becomes approximately equal to the specified \inlineCode{Time per Read} value.
    
    If the user notices that it is consistently taking more time to perform a single read than the \inlineCode{Time per Read} value that was specified, it typically means that the software cannot pull data off the hardware quickly enough.  This could be due to poor network speeds or a firewall interfering with the transfer.  It could also mean that the user should increase the \inlineCode{Maximum Acquisition Processes} value specified on the \inlineCode{Data Acquisition Setup} page.  If it consistently takes more time per read than the value specified in the \inlineCode{Time per Read} field, this means that the hardware buffer is filling up.  The controller will be responding to increasingly older data, and will not be responsive.  Finally, if the hardware buffer becomes full, the acquisition may simply stop entirely as there is no room to put more data.

    \chapter{Data Physics Quattro Devices}\label{sec:rattlesnake_hardware_dp_quattro}
    
    Rattlesnake is able to run Data Physics Quattro devices through the programming interface exposed by \inlineCode{DPQuattro.dll}, which is provided with installation materials for the Quattro device.  Users must have the proper drivers installed in order to communicate with the Quattro device.  Issues with device installation and licensing are outside the scope of the document; users are encouraged to contact Data Physics support for issues communicating with the hardware.  If users can run the examples provided by Data Physics for the Quattro device, it should also work with Rattlesnake, as the same library calls are used.
    
    \section{Setting up the Channel Table for Quattro Devices}
    
    This section lists the channel table requirements specific to the Data Physics Quattro device.
    
    The Quattro is a single standalone device, and therefore needs not be referenced by a specific name.  However, Rattlesnake requires an entry in the \inlineCode{Physical Device} column for each channel to show that the channel is active.  It is recommended to put \inlineCode{Quattro} in that space to make it clear that the connected device is a Quattro device.  The \inlineCode{Physical Channel} column should contain the channel number from the device that the row in the channel table represents.  For a Quattro device, this should range from 1--4.  Channels in Rattlesnake do not need to be in the same order as channels on the data acquisition system for Quattro devices.
    
    The \inlineCode{Minimum Value (V)} column is not used.  The \inlineCode{Maximum Value (V)} column specifies both the minimum and maximum frequency.  Valid Ranges for the Quattro device are 0.1, 1.0, 10.0, and 20.0 V for acquisition channels and 2.0 and 10.0 V for output channels.
    
    The \inlineCode{Coupling} column is where the channel coupling is specified.  Valid coupling values are
    \begin{description}
        \item[AC Differential] specified by the entry \inlineCode{AC}, \inlineCode{AC Differential}, or \inlineCode{AC Diff} (case insensitive)
        \item[DC Differential] specified by the entry \inlineCode{DC}, \inlineCode{DC Differential}, or \inlineCode{DC Diff} (case insensitive)
        \item[AC Single-Ended] specified by the entry \inlineCode{AC Single}, \inlineCode{AC Single Ended}, or \inlineCode{AC Single-Ended} (case insensitive)
        \item[DC Single-Ended] specified by the entry \inlineCode{DC Single}, \inlineCode{DC Single Ended}, or \inlineCode{DC Single-Ended} (case insensitive)
        \item[AC IEPE] specified by the entry \inlineCode{IEPE}, \inlineCode{ICP}, \inlineCode{AC ICP}, or \inlineCode{CCLD} (case insensitive)
    \end{description}

    The Quattro device handles all signal conditioning so there doesn't need to be anything entered in the \inlineCode{Current Excitation Source} or \inlineCode{Current Excitation Value} columns.  

    Recall for output channels, users of Rattlesnake must tee the output signal back into an acquisition channel.  This means using Quattro's two outputs with Rattlesnake will leave only two acquisition channels remaining.  The \inlineCode{Feedback Device} column must be populated in the row corresponding to the acquisition channel that the output is teed into, again using \inlineCode{Quattro} for convention, but otherwise any entry will suffice.  The \inlineCode{Feedback Channel} will be the output channel on the Quattro that is used.  This should range from 1--2.  Outputs in Rattlesnake do not need to be in the same order as outputs on the data acquisition system for Quattro Devices.
    
    \section{Hardware Parameters}
    When the Quattro device is selected from the \inlineCode{Hardware Selector}, Rattlesnake will bring up a file dialog asking for a Data Physics \ac{API} file.  Users should navigate to the file \inlineCode{DpQuattro.dll}, which should have been delivered with the Quattro \ac{API} media.  This \inlineCode{dll} file contains the Quattro programming interface that Rattlesnake uses to communicate with the Quattro device.  If users cannot find this file, please contact Data Physics support, as Rattlesnake maintainers cannot provide this file.  Each file is licensed separately by Data Physics and cannot be shared.  The \inlineCode{dll} file should be alongside a \inlineCode{signal.001.qrt} file, which is the license file specific to each Quattro device.  If this license file is missing or is not the one paired to the Quattro hardware, the interface will not run.
    
    The only other parameter specific to the Quattro device is the \inlineCode{Sample Rate}.  Quattro devices used with Rattlesnake have discrete sampling rates available.  These are 16, 20, 25, 32, 40, 50, 64, 80, 100, 128, 160, 200, 256, 320, 400, 512, 640, 800, 1024, 1280, 1600, 2048, 2560, 3200, 4096, 5120, 6400, 8192, 10240, 12800, 20480, 25600, 40960, 51200, 102400, or 204800 samples per second.  This is a subset of the full set of Quattro sample rates, because Rattlesnake only deals with integer numbers of samples per second.
    
    \section{Implementation Details}
    This section contains details on the Quattro implementation in Rattlesnake, which may be helpful for users when diagnosing issues that arise in the controller.
    
    \subsection{Drivers}
    The Data Physics Quattro device connects to the computer via USB, and must have drivers installed to correctly connect.  Drivers are found in the \inlineCode{\Driver} folder of the DpQuattro \ac{API} package.  If you have trouble finding or installing the drivers, please contact Data Physics support, as the Rattlesnake authors cannot provide these files.
    
    \subsection{Programming API}
    The Data Physics DpQuattro \ac{API} is a 64-bit programming interface written in C/C++.  It is compiled to a Windows \inlineCode{dll}, so it can be used by any other environment that can use a Windows-based C library.  Therefore the Quattro device as implemented in Rattlesnake is only usable on 64-bit Windows operating systems, not Mac or Linux.  Additionally, if you are not using a 64-bit Python executable, you will not be able to connect to the 64-bit \inlineCode{dll}, and will not be able to use the Quattro device.
    
    The \ac{API} is accessed using the Python \inlineCode{ctypes} module, by calling its \inlineCode{WinDLL} method on the path to the \inlineCode{dll}.  The file \inlineCode{components/data_physics_interface.py} in the main Rattlesnake directory constructs an object-oriented Python interface from the functions in this \inlineCode{dll}.  The file \inlineCode{components/data_physics_hardware.py} then defines the acquisition and output hardware objects with the correct methods that Rattlesnake will call.
    
    \subsection{Debugging}
    If issues with the Quattro devices are encountered, users can look at the Rattlesnake command prompt that appears when the software is loaded, or the command prompt from which Rattlesnake is run.  The DpQuattro \ac{API} will print messages to this prompt periodically that can be useful in debugging issues.  These will generally be prefixed by \inlineCode{CAceSys} and provide some information on the issue encountered.  Figure \ref{fig:quattroissuecommandline} shows an example that occurred if the Quattro devices is not connected to the computer.  Other diagnostic messages are also printed, so a message appearing does not necessarily mean an error has occurred.
    
    \begin{figure}
        \centering
        \includegraphics[width=0.5\linewidth]{figures/quattro_issue_command_line}
        \caption{An example issue where the Quattro fails to connect.}
        \label{fig:quattroissuecommandline}
    \end{figure}
    
    For more debugging information, there is a \inlineCode{DEBUG} flag in the file \inlineCode{components/data_physics_interface.py} that is by default set to \inlineCode{False}.  If users change that to \inlineCode{DEBUG = True}, then the Quattro interface will output a log file with more information about its activities.
    
    \chapter{Data Physics 900-Series Devices}\label{sec:rattlesnake_hardware_dp900}
    
    The Data Physics DP900 series device is not yet implemented with Rattlesnake.  This will be implemented pending the release of the DP900 \ac{API} by Data Physics.  This chapter serves as a placeholder to describe that implementation.  The implementation is expected to be very similar to the Quattro device.

    \chapter{Virtual Control using State Space Matrices}\label{sec:rattlesnake_hardware_state_space}
    
    If no data acquisition hardware is available, it can still be advantageous to run Rattlesnake in a ``virtual'' mode by simulating responses to virtual forces using some kind of model.  This would allow users of the software to develop new control laws and environments without risking damage to potentially expensive test hardware if the control law or environment is not implemented correctly.  Additionally, synthetic control can be used to determine proper test parameters prior to the actual test occurring, which can help determine if a given test is feasible or not.
    
    Rattlesnake is able to integrate linear time-invariant equations of motion using general state space matrices.  The general form for state space equations of motion is 
    
    \begin{equation}\label{eq:state_eoms}
        \dot{\mathbf{x}} = \mathbf{A}\mathbf{x} + \mathbf{B}\mathbf{u}
    \end{equation}
    \begin{equation}\label{eq:output_eoms}
        \mathbf{y} = \mathbf{C}\mathbf{x} + \mathbf{D}\mathbf{u}
    \end{equation}
    
    where $\mathbf{A}$ is the state matrix, $\mathbf{B}$ is the input matrix, $\mathbf{C}$ is the output matrix, $\mathbf{D}$ is the feedthrough matrix, $\mathbf{x}$ is the state vector, $\mathbf{u}$ is the input vector, and $\mathbf{y}$ is the output vector.  Equation \eqref{eq:state_eoms} defines how the state changes with time $\dot{\mathbf{x}}$ given the current state $\mathbf{x}$ and inputs $\mathbf{u}$ to the system.  Equation \eqref{eq:output_eoms} describes how the output degrees of freedom $\mathbf{y}$ are constructed from the current state and inputs to the system.  When used within Rattlesnake, the output signals from Rattlesnake will be supplied to the system as $\mathbf{u}$, and the output degrees of freedom $\mathbf{y}$ will be acquired by Rattlesnake.  A complete example for running this type of test can be found in Appendix \ref{sec:example_state_space}.
    
    For many structural dynamic analyses, the user will have mass, stiffness, and damping matrices $\mathbf{M}$, $\mathbf{K}$, and $\mathbf{C}$, respectively, in the typical 2nd-order differential equations of motion
    
    \begin{equation}\label{eq:second_order_eoms}
    \mathbf{M}\ddot{\mathbf{z}} + \mathbf{C}\dot{\mathbf{z}} + \mathbf{K}\mathbf{z} = \mathbf{u}
    \end{equation}
    
    where $\mathbf{z}$ is treated as the displacement degrees of freedom of the system and $\mathbf{u}$ are the input forces.  The state of the system is represented by 
    
    \begin{equation}
        \mathbf{x} = \begin{bmatrix}
        \mathbf{z} \\
        \dot{\mathbf{z}}
        \end{bmatrix}
    \end{equation}
    
    and therefore the derivative of the state is 
    
    \begin{equation}
        \dot{\mathbf{x}} = \begin{bmatrix}
        \dot{\mathbf{z}} \\
        \ddot{\mathbf{z}}
        \end{bmatrix}
    \end{equation}
    
    Given the state $\mathbf{x}$ and the inputs $\mathbf{u}$, equation \eqref{eq:second_order_eoms} can be transformed into equation \eqref{eq:state_eoms} with
    
    \begin{equation}\label{eq:state_space_a}
        \mathbf{A} = \begin{bmatrix}
        \mathbf{0} & \mathbf{I} \\
        -\mathbf{M}^{-1}\mathbf{K} & -\mathbf{M}^{-1}\mathbf{C}
        \end{bmatrix}
    \end{equation}
    \begin{equation}\label{eq:state_space_b}
        \mathbf{B} = \begin{bmatrix}
        \mathbf{0} \\ \mathbf{M}^{-1}
        \end{bmatrix}
    \end{equation}
    
    where the first row partition is the trivial equation $\dot{\mathbf{z}} = \dot{\mathbf{z}}$.
    
    The output degrees of freedom will then depend on which types of quantities are desired from the analysis.  Displacements, velocities, accelerations, and forces can be recovered in the output degrees of freedom $\mathbf{y}$ using
    
    \begin{equation}\label{eq:state_space_c}
        \mathbf{C} = \begin{bmatrix}
        \mathbf{I} & \mathbf{0} \\
        \mathbf{0} & \mathbf{I} \\
        -\mathbf{M}^{-1}\mathbf{K} & -\mathbf{M}^{-1}\mathbf{C} \\
        \mathbf{0} & \mathbf{0}
        \end{bmatrix}
    \end{equation}
    \begin{equation}\label{eq:state_space_d}
        \mathbf{D} = \begin{bmatrix}
        \mathbf{0} \\\mathbf{0} \\ \mathbf{M}^{-1} \\ \mathbf{I}
        \end{bmatrix}
    \end{equation}
    
    where the first row partition recovers displacements $\mathbf{z}$, the second row partition recovers velocities $\dot{\mathbf{z}}$, the third row partition recovers accelerations $\ddot{\mathbf{z}}$, and the fourth row partition recovers forces $\mathbf{u}$.  Note that Rattlesnake requires its output signals (corresponding to the input signals to the system $\mathbf{u}$) to be measured in order to function (see Section \ref{sec:channel_table}) so the input signals to the system $\mathbf{u}$ must be passed through directly to output degrees of freedom $\mathbf{y}$ which are the acquired data in Rattlesnake.  Practically, this means that there should always be at least one row partition of $\mathbf{C}$ containing zeros, and the same row partition of $\mathbf{D}$ should contain the identity matrix.
    
    This chapter describes the setup and implementation details for running a synthetic control problem using state space matrices.  For a complete example of using this type of control, see Appendix \ref{sec:example_state_space}.
    
    \section{Setting up the Channel Table for Virtual Control using State Space Matrices}
    Required channel table inputs for virtual control using state space matrices are minimal, as the $\mathbf{C}$ and $\mathbf{D}$ matrices effectively define the measured degrees of freedom.  Still users are encouraged to fill out the channel table as much as possible for test documentation's sake.  The only required parameters for virtual control are have an entry in the Physical Device column of the channel table.  The authors recommend simply using the word ``Virtual'' to reinforce the fact that these are not real hardware channels, but instead virtual channels that exist only in software.  The excitation channels must also have an entry in the Feedback Device column.  The authors recommend simply using the word ``Input'' to make it clear that these channels are input channels to the system.  One important consideration is that the number of channels in the channel table must be the same as the number of output degrees of freedom in $\mathbf{y}$ (and therefore the number of rows of matrices $\mathbf{C}$ and $\mathbf{D}$), and they must also be ordered identically.  Similarly, the number of input channels to the system specified in the channel table must be the same as the number of input degrees of freedom in $\mathbf{u}$ (and therefore the number of columns of matrices $\mathbf{B}$ and $\mathbf{D}$), also ordered identically.
    
    \section{Hardware Parameters}
    Hardware parameters are similar to those found in Chapter \ref{sec:rattlesnake_hardware_virtual_control}.  When the \inlineCode{State Space Integration...} hardware device is selected, Rattlesnake will bring up a file dialog box that will allow the user to specify a Numpy or Matlab file from which the state space matrices will be loaded.  This file should contain fields \inlineCode{A}, \inlineCode{B}, \inlineCode{C}, and \inlineCode{D} containing the appropriately sized matrices.  An additional parameter will appear in the \inlineCode{Data Acquisition Parameters} portion of the \inlineCode{Data Acquisition Setup} tab that allows the user to specify an \inlineCode{Integration Oversample} amount, as the integration time steps should generally be finer than the sample rate of the controller.  A value of 10 is typically sufficient for integration of a linear system using \inlineCode{lsim} from \inlineCode{scipy.signal}.
    
    \section{Implementation Details}
    \label{sec:state_space_integration}
    A \inlineCode{scipy.signal.StateSpace} model is created from the loaded matrices which are used for integration using the \inlineCode{scipy.signal.lsim} function.
    
    \subsection{Output Signal Force Buffering}
    Excitation signals coming into the state space virtual hardware are buffered to enable different read and write sizes.  The excitation signals will be delivered to the hardware based on the specified \inlineCode{Time per Write} on the \inlineCode{Data Acquisition Setup} tab of the controller.  The virtual hardware will then remove an amount of excitation signal from the buffer based on the specified \inlineCode{Time per Read}, as the acquisition task will perform the integration over time blocks equivalent to \inlineCode{Time per Read}.
        
    \subsection{Integration using LSIM}
    The \inlineCode{scipy.signal.lsim} function is used to integrate the state space model.  When the acquisition is started, the state of the integration is initialized to zero.  For subsequent integrations, the initial state is set to the final state of the previous integration.  The excitation signals are taken from the force buffer discussed previously.  Note that integration is performed at the rate specified by the \inlineCode{Sample Rate} of the data acquisition system multiplied by the \inlineCode{Integration Oversample} amount.  After the integration is performed, the output is downsampled by the \inlineCode{Integration Oversample} amount to return to the sample rate of the controller.  For subsequent integration blocks, the initial conditions of the integration are equivalent to the final value of the last integration block.
    
    \chapter{Virtual Control using Finite Element Results in Exodus Files}\label{sec:rattlesnake_hardware_virtual_control}
    
    In addition to virtual control using state space matrices, Rattlesnake can also perform virtual control using results from \ac{FEM} analyses.  This is perhaps less complex to set up than the virtual control using state space matrices described in \ref{sec:rattlesnake_hardware_state_space}, but it offers less flexibility to the user and relies on the user having results in the Exodus file format.
    
    Modern \ac{FEM}s can contain thousands or even millions of degrees of freedom, so it is not realistic to be able to integrate the full finite element model's equations of motion in real time.  Some model reduction is necessary to reduce the number of degrees of freedom that must be integrated.  Rattlesnake instead integrates modal equations of motion rather than the full set of equations of motion.  The modal transformation has the advantage of generally only requiring a small number of modal degrees of freedom to characterize the part over some frequency bandwidth.  Additionally, modal degrees of freedom are uncoupled from one another, which makes for simpler integration strategies.
    
    This chapter describes the setup and implementation details for running a synthetic control problem using Exodus \ac{FEM} results.
    
    \section{Setting up the Channel Table for Virtual Control using a \ac{FEM}}
    
    Setting up the channel table is very straightforward for a virtual control problem.  The node number that is specified in the channel table corresponds to the node in the finite element model.  Note that this corresponds to the node \textit{number} not the node \textit{index}, so users should be aware when picking out nodes for virtual control using visualization software whether or not the software is reporting node number or node index.  The node direction can be specified as either a principal direction (X+, Y+, Z+, X-, Y-, or Z-) or a comma separated 3-vector with unit magnitude that specifies the measurement direction.
    
    The modal damping value that will be applied to all modes in the model should be placed in the Comment column of the first channel.  This should be specified as a fraction of critical damping (e.g. 0.01) rather than a percentage (e.g. 1\%).  For more complex damping capabilities, a State Space formulation should be used instead, see Chapter \ref{sec:rattlesnake_hardware_state_space}.
    
    All active measurement degrees of freedom and output degrees of freedom must have an entry in the \inlineCode{Physical Device} column of the channel table.  The authors recommend simply using the word \inlineCode{Virtual} to reinforce the fact that these are virtual channels.  Any excitation channel must have an entry in the \inlineCode{Feedback Device} column.  The authors here recommend simply using the word \inlineCode{Input} to make it clear that these channels are excitation channels.
    
    All other columns of the channel table can be left empty, though it may be useful to put entries into them to document the virtual test better.
    
    \section{Hardware Parameters}
    When the \inlineCode{Exodus Modal Solution...} hardware device is selected, Rattlesnake will bring up a file dialog box that will allow the user to specify the Exodus file from which the modal solution will be loaded.  An additional parameter will appear in the \inlineCode{Data Acquisition Parameters} portion of the \inlineCode{Data Acquisition Setup} tab that allows the user to specify an \inlineCode{Integration Oversample} amount, as the integration time steps should generally be finer than the sample rate of the controller.  A value of 10 is typically sufficient for integration of a linear system using \inlineCode{lsim} from \inlineCode{scipy.signal}.
    
    \section{Implementation Details}
    This section contains implementation details of the virtual control using modal analysis solutions found in exodus files. 
    
    \subsection{Creating State Space Matrices from Modal EOMs}
    To assemble the state space equations of motion, the mode shape matrix is used to transform between modal and physical degrees of freedom.  A 3D displacement array is constructed from the \inlineCode{DispX}, \inlineCode{DispY}, and \inlineCode{DispZ} node variables in the Exodus file.  A mode shape matrix is computed by taking the dot product of the direction vector specified by the direction in the channel table with the displacement vector for the node specified by the channel table for each mode shape.
    
    \subsubsection{State Degrees of Freedom}
    States for the virtual control using an Exodus modal solution are the modal displacements and velocities for the modes in the exodus file below 1.5$\times$ the Nyquist frequency of the controller.  Mass normalized mode shapes are assumed, so the modal mass matrix is assumed to be the identify matrix, modal damping is assumed to be a diagonal matrix of damping coefficients $2\zeta_n\omega_n$, and modal stiffness is assumed to be a diagonal matrix of stiffness coefficients ${\omega_n}^2$.  The state equation matrices $\mathbf{A}$ and $\mathbf{B}$ are then defined by equations \eqref{eq:state_space_a} and \eqref{eq:state_space_b}.
    
    \subsubsection{Response Degrees of Freedom}
    Response degrees of freedom are defined using the $\mathbf{C}$ and $\mathbf{D}$ matrices defined by equations \eqref{eq:state_space_c} and \eqref{eq:state_space_d}.  This hardware device assumes responses are accelerations, so the third row partition of equations \eqref{eq:state_space_c} and \eqref{eq:state_space_d} are used.  These are modal accelerations, so they are multiplied by the mode shape matrix to produce physical accelerations.
    
    \subsubsection{Excitation Degrees of Freedom}
    Excitation to the system is specified as modal forces, so the physical excitation forces provided by Rattlesnake are multiplied by the mode shape matrix transposed to create modal forces.  To accommodate potentially different read and write sizes in the controller, physical forces to be output are appended to a force buffer and one write worth of forces are pulled from the buffer to be transformed to modal forces and applied to the system.
    
    Rattlesnake's measurement strategy requires that excitation signals also be measured by the data acquisition, and this is also true for virtual control.  Therefore, excitation signals must be returned as a measured response.
    
    \subsection{Integration of Equations of Motion}
    Once state space matrices are formed, the implementation of this virtual hardware is largely the same as the virtual control described in Section \ref{sec:rattlesnake_hardware_state_space}.  See Section \ref{sec:state_space_integration} for details on the integration scheme.
    
    \chapter{Virtual Control using SDynPy System objects}\label{sec:rattlesnake_hardware_sdynpy_system}
    
    The final option for synthetic control in Rattlesnake is to load in a SDynPy \href{https://sandialabs.github.io/sdynpy/_autosummary/sdynpy.core.sdynpy_system.System.html}{\inlineCode{System}} object, which gets stored from SDynPy using the \href{https://sandialabs.github.io/sdynpy/_autosummary/sdynpy.core.sdynpy_system.System.html#sdynpy.core.sdynpy_system.System.save}{\inlineCode{sdynpy.System.save}} method.  SDynPy \inlineCode{System} objects store mass, stiffness, and damping matrices.  They also store transformation matrices which transform its internal system space into the physical space.  This allows SDynPy \inlineCode{System} objects to represent both ``full'' physical systems as well as ``reduced'' systems (e.g., Craig-Bampton or Modal systems).  The final part of the SDynPy \inlineCode{System} is the degree of freedom information that is stored with the matrices.  This maps rows of the transformation matrix (or rows of the mass, stiffness, and damping matrices if the transformation is the identity matrix) to physical degrees of freedom.
    
    A complete example problem using a SDynPy \inlineCode{System} object can be found in Appendix \ref{sec:example_sdynpy}.
    
    \section{Setting Up the Channel Table}
    
    Setting up the channel table is very straightforward for a SDynPy System.  Degrees of freedom in the SDynPy system are selected in the Rattlesnake test by specifying the node and direction of the degree of freedom in the \inlineCode{Node} and \inlineCode{Direction} of the \inlineCode{coordinate} field of the SDynPy system.  The node is specified as an integer, and the direction is specified as a string (X+, Y+, Z+, X-, Y-, or Z-).
    
    All active measurement degrees of freedom and output degrees of freedom must have an entry in the \inlineCode{Physical Device} column of the channel table.  The authors recommend simply using the word \inlineCode{Virtual} to reinforce the fact that these are virtual channels.  Any excitation channel must have an entry in the \inlineCode{Feedback Device} column.  The authors here recommend simply using the word \inlineCode{Input} to make it clear that these channels are excitation channels.
    
    The \inlineCode{Channel Type} must also be specified.  It is used to specify the derivative of the degree of freedom that will be acquired in the response channels.  The \inlineCode{Channel Type} can be specified as \inlineCode{Acceleration}, \inlineCode{Velocity}, or \inlineCode{Displacement}, and the output matrices of the state space system that will be integrated will be assembled such that it will return the specified channel type.  If an unknown channel type is specified or if none are specified, then displacement will be assumed.  Note excitation degrees of freedom are not modified by a \inlineCode{Channel Type} modifier, they are applied directly as-is to the SDynPy system.
    
    All other columns of the channel table can be left empty, though it may be useful to put entries into them to document the virtual test better.
    
    \section{Hardware Parameters}
    
    When the \inlineCode{SDynPy System Integration...} hardware device is selected, Rattlesnake will bring up a file dialog box that will allow the user to specify the file from which the SDynPy system will be loaded.  An additional parameter will appear in the \inlineCode{Data Acquisition Parameters} portion of the \inlineCode{Data Acquisition Setup} tab that allows the user to specify an \inlineCode{Integration Oversample} amount, as the integration time steps should generally be finer than the sample rate of the controller.  A value of 10 is typically sufficient for integration of a linear system using \inlineCode{lsim} from \inlineCode{scipy.signal}.
    
    \section{Implementation Details}
    
    This section describes the implementation details for the SDynPy System integration.
    
    \subsection{Setting Up the State Space System}
    
    Rattlesnake is able to integrate linear time-invariant equations of motion using general state space matrices.  The general form for state space equations of motion is given in Equations \eqref{eq:state_eoms} and \eqref{eq:output_eoms}.
    
    Section \ref{sec:rattlesnake_hardware_state_space} describes how the state space equations can be constructed from mass, stiffness, and damping matrices.  The state space formulation can be extended to allow for a transformation matrix between the mass, stiffness, and damping matrices, and the physical degrees of freedom.
    
    \begin{equation}\label{eq:state_space_c_w_transform}
        \mathbf{C} = \begin{bmatrix}
            \mathbf{\Phi}_r & \mathbf{0} \\
            \mathbf{0} & \mathbf{\Phi}_r \\
            -\mathbf{\Phi}_r\mathbf{M}^{-1}\mathbf{K} & -\mathbf{\Phi}_r\mathbf{M}^{-1}\mathbf{C} \\
            \mathbf{0} & \mathbf{0}
        \end{bmatrix}
    \end{equation}
    \begin{equation}\label{eq:state_space_d_w_transform}
        \mathbf{D} = \begin{bmatrix}
            \mathbf{0} \\\mathbf{0} \\ \mathbf{\Phi}_r\mathbf{M}^{-1}{\mathbf{\Phi}_e}^T \\ \mathbf{I}
        \end{bmatrix}
    \end{equation}

    Here the system transformation $\mathbf{\Phi}$ is partitioned into the response degrees of freedom $\mathbf{\Phi}_r$ and excitation degrees of freedom $\mathbf{\Phi}_e$.
    
    Matrices $\mathbf{A}$ and $\mathbf{B}$ are unchanged and defined in equations \eqref{eq:state_space_a} and \eqref{eq:state_space_b}.
    
    The state space matrices are then integrated using \inlineCode{scipy.signal.lsim}, similar to the State Space formulation described in Section \ref{sec:rattlesnake_hardware_state_space}.
    
    \chapter{Implementing New Hardware Devices with Rattlesnake}\label{sec:rattlesnake_hardware_implementation}
    \textbf{WARNING:} This chapter contains advanced Rattlesnake operations that require a reasonably good understanding of the Rattlesnake software architecture.  Understanding this chapter is not required to run the software successfully, and therefore this chapter can be skipped if the user is not interested in implementing a new hardware device.
    
    Rattlesnake allows users to implement new hardware devices with minimal modifications to the controller.  Within the \GitlabProject{} repository, there is a source file \inlineCode{components/abstract_hardware.py} that defines abstract base classes \inlineCode{HardwareAcquisition} and \inlineCode{HardwareOutput}.  All hardware devices used by Rattlesnake must implement acquisition and output classes that inherit from these abstract base classes and define functions that overwrite the abstract methods in these base classes.  The acquisition and output will generally be run on separate processes in Rattlesnake to ensure that output can be streamed to the device as acquired data is simultaneously read from the device.  A flowchart of the operations performed by each hardware process is shown in Figure \ref{fig:hardwareflowchart}.
    
    \begin{figure}
        \centering
        \includegraphics[height=0.9\textheight]{figures/hardware_flowchart}
        \caption[Flowchart of Hardware Operations.]{Flowchart of Hardware Operations. Note that each process will proceed at its own pace, so it is generally not possible to ensure that an Acquisition function is called before or after an Output function. The only place where order of operations is enforced is at startup, where the Output \inlineCode{start} function finishes before the Acquisition \inlineCode{start} function is called.}
        \label{fig:hardwareflowchart}
    \end{figure}
    
    The hardware setup in Rattlesnake generally assumes that separate processes are required for acquisition and output.  However, this may not be the case for all hardware devices.  If a device can only run on a single process, users can utilize a \inlineCode{multiprocessing} queue to pass the output data from the output process to the acquisition process and perform all hardware processes there.  See, for example, the hardware implementation in \inlineCode{components/state_space_virtual_hardware.py} for more details on this approach; the excitation signal is obtained by the \inlineCode{write} method and immediately passed to a queue to deliver to the acquisition process.
    
    Users implementing new hardware devices are encouraged to look at existing hardware implementations in the \GitlabProject{} repository (\inlineCode{components/nidaqmx_hardware_multitask.py}, \inlineCode{components/lanxi_hardware.py}, \inlineCode{components/data_physics_hardware.py},
    \inlineCode{components/state_space_virtual_hardware.py},
    \inlineCode{components/sdynpy_system_virtual_hardware.py}, and \inlineCode{components/exodus_modal_solution_hardware.py}) to use as examples.
    
    \section{Defining a \inlineCode{HardwareAcquisition} Class}
    The acquisition portion of the data acquisition hardware is defined using a class that inherits the abstract base class \inlineCode{HardwareAcquisition} in \inlineCode{components/abstract_hardware.py}.  A class inheriting from \inlineCode{HardwareAcquisition} must implement the following functions.
    \begin{description}
        \item[\_\_init\_\_] The class constructor that initializes the hardware device.  Here default parameters can be set up.
        \item[set\_up\_data\_acquisition\_parameters\_and\_channels] This function is called when the \inlineCode{Data Acquisition Setup} tab is initialized.  This function will nominally set up the channels on the hardware as well as any other acquisition parameters such as sample rate or triggering that are required by the hardware.  As arguments, this function receives an object defining the data acquisition parameters as well as a list of channels.  Note that some hardware devices require separate setup for output and acquisition, while others may only need to be set up once for both hardware and acquisition.  For the latter case, users should evaluate whether or not it makes sense to do the setup in the \inlineCode{HardwareAcquisition} or \inlineCode{HardwareOutput} class.  Note also that there is no assurance made by the controller that operations in the \inlineCode{set_up_data_acquisition_parameters_and_channels} function will be called before or after operations in the the \inlineCode{set_up_data_output_parameters_and_channels} function from the \inlineCode{HardwareOutput} class, as they will both be called simultaneously by separate processes.  If a specific order of operations for hardware setup must be enforced, those operations should be contained entirely in either the \inlineCode{set_up_data_acquisition_parameters_and_channels} function or the \inlineCode{set_up_data_output_parameters_and_channels} function from the \inlineCode{HardwareOutput} class.
        \item[start] This function is called when the acquisition system starts acquiring and should start the device recording.  This function will always be called after the \inlineCode{start} function from the \inlineCode{HardwareOutput} class, so users can rely on that ordering for order of operations.
        \item[read] This function reads a set of data from the device.  The returned data should be a Numpy array with number of rows equal to the number of measurement channels and number of columns equal to the number of samples read.  The number of samples read is determined by the \inlineCode{Sample Rate} and \inlineCode{Time per Read} parameters that get specified on the \inlineCode{Data Acquisition Setup} tab.
        \item[read\_remaining] This function reads the remaining data on the device, as it is called before the device is stopped.  It returns a Numpy array with number of rows equal to the number of measurement channels and a number of columns equal to the number of samples remaining to be read on the device.
        \item[get\_acquisition\_delay]  This function returns an integer number of samples that defines how much longer the acquisition should acquire after output shuts down.  This is designed to accommodate any output buffering present in the hardware device.
        \item[stop] This function is called when the acquisition system should stop recording, so it should stop any data acquisition on the device.  The controller makes no guarantee that the \inlineCode{HardwareAcquisition} \inlineCode{stop} function will be called before or after the \inlineCode{HardwareOutput} \inlineCode{stop} function, so the user should not rely on one stopping prior to the other.
        \item[close] This function is called when the controller is shutting down the hardware for good.  It should gracefully shut down the hardware and release any memory or references related to it.  The controller makes no guarantee that the \inlineCode{HardwareAcquisition} \inlineCode{close} function will be called before or after the \inlineCode{HardwareOutput} \inlineCode{close} function, so the user should not rely on one stopping prior to the other.
    \end{description}
    
    \section{Defining a \inlineCode{HardwareOutput} Class}
    The acquisition portion of the data acquisition hardware is defined using a class that inherits the abstract base class \inlineCode{HardwareOutput} in \inlineCode{components/abstract_hardware.py}.  A class inheriting from \inlineCode{HardwareOutput} must implement the following functions.
    \begin{description}
        \item[\_\_init\_\_] The class constructor that initializes the hardware device.  Here default parameters can be set up.
        \item[set\_up\_data\_output\_parameters\_and\_channels] This function is called when the \inlineCode{Data Acquisition Setup} tab is initialized.  This function will nominally set up the channels on the hardware as well as any other output parameters such as sample rate or triggering that are required by the hardware.  As arguments, this function receives an object defining the data acquisition parameters as well as a list of channels.  Note that some hardware devices require separate setup for output and acquisition, while others may only need to be set up once for both hardware and acquisition.  For the latter case, users should evaluate whether or not it makes sense to do the setup in the \inlineCode{HardwareAcquisition} or \inlineCode{HardwareOutput} class.  Note also that there is no assurance made by the controller that operations in the \inlineCode{HardwareAcquisition} \inlineCode{set_up_data_acquisition_parameters_and_channels} function will be called before or after operations in the the \inlineCode{set_up_data_output_parameters_and_channels} function from the \inlineCode{HardwareOutput} class, as they will both be called simultaneously by separate processes.  If a specific order of operations for hardware setup must be enforced, those operations should be contained entirely in either the \inlineCode{set_up_data_acquisition_parameters_and_channels} function from the \inlineCode{HardwareAcquisition} class or the \inlineCode{set_up_data_output_parameters_and_channels} function from the \inlineCode{HardwareOutput} class.
        \item[start] This function is called when the acquisition system starts acquiring and should start the device outputting signals.  This function will always be called before the \inlineCode{start} function from the \inlineCode{HardwareAcquisition} class, so users can rely on that ordering for order of operations.
        \item[ready\_for\_next\_output] This function is called prior to writing any data to the hardware to check whether or not the device is ready for new outputs.  This function will return \inlineCode{True} if the device should accept new output signals, and \inlineCode{False} if it should not.  If the controller was able to write data to the hardware as fast as possible, it might fill up the hardware buffer.  This would make the controller very unresponsive as the buffer would need to be exhausted by outputting all stored samples prior to any new data being output.  In general, the function should be set up such that the minimum amount of data is stored on the device (i.e. written to the device but not yet output) at any given time, while also ensuring that variations in the controller's performance due to other tasks being performed simultaneously do not result in the buffer running out of samples to output before the next set of samples is written to the device.  Generally a buffer of two or three writes worth of samples is sufficient to ensure that the buffer does not exhaust itself.
        \item[write] This function accepts data as a Numpy array and writes it to the hardware.  The array will be sized such that the number of rows is equal to the number of output signals and the number of columns is equal to the number of samples that are written per channel, which is determined by the \inlineCode{Sample Rate} and \inlineCode{Time per Write} parameters that get specified on the \inlineCode{Data Acquisition Setup} tab.
        \item[stop] This function is called when the acquisition system should stop outputting data, so it should stop any data output on the device.  The controller makes no guarantee that the \inlineCode{HardwareAcquisition} \inlineCode{stop} function will be called before or after the \inlineCode{HardwareOutput} \inlineCode{stop} function, so the user should not rely on one stopping prior to the other.
        \item[close] This function is called when the controller is shutting down the hardware for good.  It should gracefully shut down the hardware and release any memory or references related to it.  The controller makes no guarantee that the \inlineCode{HardwareAcquisition} \inlineCode{close} function will be called before or after the \inlineCode{HardwareOutput} \inlineCode{close} function, so the user should not rely on one shutting down prior to the other.
    \end{description}
    
    \section{Controller Modifications to Recognize New Hardware Devices}
    With the new hardware device implemented, the controller should be modified to find the new hardware device and allow the user to run it.
    
    \subsection{\Ac{GUI} Modifications}
    In order for the user to select the newly-implemented hardware device, the hardware must be added to the \inlineCode{Hardware Selector} drop-down menu on the \inlineCode{Data Acquisition Setup} tab of the controller.  This requires editing the \inlineCode{components/combined_environments_controller.ui} file.  Editing of \inlineCode{*.ui} files is most easily done through the Qt Designer software, but can also be done by editing the \inlineCode{*.ui} file directly, as it contains XML.  The widget that needs to be modified is a \inlineCode{QComboBox} widget with the name \inlineCode{hardware_selector}.  An additional item needs to be added to the widget corresponding to the hardware device.
    
    If the hardware requires additional user parameters to be specified, then further modifications to the \ac{GUI} might be required.  Such modifications will require a thorough understanding of the Rattlesnake source code.  These modifications will be highly specific to the hardware being implemented, and are therefore out of scope for this User's Manual.  Such modifications might include adding a dropdown to select a trigger channel to synchronize acquisition and output or adding integration parameters such as minimum or maximum timestep for a virtual hardware device that utilizes a nonlinear integrator.
    
    \subsection{Modifications to \Ac{UI} Callback Functions}
    With the \ac{GUI} modified to enable selecting the hardware interface, the user might need to modify the \inlineCode{Hardware Selector} callback function in order to ensure the proper options are visible to the user.  These callbacks are defined in the \inlineCode{components/user_interface.py} file in the \inlineCode{Ui} class.  The \inlineCode{hardware_update} function may need to be modified to handle hiding or showing specific widgets when the new hardware device is selected.
    
    If the hardware interface requires additional user parameters to be specified, then further modifications to the code will be required.  These modifications will be highly specific to the hardware being implemented, so this User's Manual will only discuss where the changes will need to be made, but not necessarily what those changes need to be.  If additional parameters are required, users may need to modify the \inlineCode{DataAcquisitionParameters} class found in \inlineCode{components/utilities.py} to add these additional parameters to the global parameter set.  The user will also need to modify the code where the class is instantiated to include these new parameters, particularly in the \inlineCode{initialize_data_acquisition} function of the \inlineCode{Ui} class in \inlineCode{components/user_interface.py}.
    
    \subsection{Modifications to the Acquisition and Output Processes}
    Finally, the \inlineCode{initialize_data_acquisition} functions in both the \inlineCode{AcquisitionProcess} class in \inlineCode{components/acquisition.py} and the \inlineCode{OutputProcess} class in \inlineCode{components/output.py} will need to be modified to correctly initialize the hardware.  A new \inlineCode{elif} statement should be added to the main \inlineCode{if} statement inside that function to accommodate the new hardware device index that was added to the \inlineCode{Hardware Selector} dropdown menu.  The \inlineCode{*.py} file that contains the hardware implementation should be imported and the respective Acquisition or Output class should be instantiated.  It is at this point that any additional arguments to the classes' \inlineCode{__init__} functions are passed.  The object instantiated from the class should be stored to the \inlineCode{self.hardware} property of the class.

    \part{Rattlesnake Environments}\label{sec:rattlesnake_environments}
    
    The key to Rattlesnake's flexibility is its implementation of various environment types.  These environments can be run stand-alone or in combination with other environments.  A general environment will receive response data from the controller and then perform some analysis on that data to create the next set of output signals that will be required by the controller.  This Part will describe the existing environments in the controller and provide information required to implement additional environments.
    
    Chapter \ref{sec:rattlesnake_environments_mimo_random} describes the \Ac{MIMO} Random Vibration environment where users will control to vibration spectra in the form of \Ac{CPSD} matrices.  Chapter \ref{sec:rattlesnake_environments_mimo_transient} describes the \Ac{MIMO} Transient Vibration environment, where users will control directly to time responses.  Chapter \ref{sec:rattlesnake_environments_time_generator} describes the Time History Generator environment, which allows users to specify the signal played to the excitation devices.  Chapter \ref{sec:rattlesnake_environments_modal} describes the Modal environment, where users can perform dynamic characterization tests using shaker or impact hammer excitation.  Chapter \ref{sec:rattlesnake_environments_combining_environments} describes the Combined Environments functionality in Rattlesnake which allows users to run multiple environments simultaneously.
    
    Chapter \ref{sec:rattlesnake_environments_implementation} describes the general process to implement new environments.  Implementing new environments in Rattlesnake is a very involved process that requires a good amount of knowledge of the Rattlesnake software architecture, as well as advanced programming concepts including object-oriented programming and inheritance, multiprocessing, and \Ac{GUI} design and implementation.
    
    \chapter{Multiple Input/Multiple Output Random Vibration}\label{sec:rattlesnake_environments_mimo_random}
    
    The first environment implemented in the Rattlesnake controller was the \Ac{MIMO} Random Vibration environment.  This environment aims to control the vibration response of a component to specified levels by creating output signals with the correct levels, coherence, and phase at each frequency line.  The governing equation for \Ac{MIMO} Random Vibration is 
    
    \begin{equation}\label{eq:forward_mimo_vibration}
    \mathbf{G}_{xx} = {\mathbf{H}_{xv}}\mathbf{G}_{vv}{{\mathbf{H}_{xv}}}^H
    \end{equation}
    
    where the \ac{CPSD} matrix of the responses $\mathbf{G}_{xx}$ result from some signals $\mathbf{G}_{vv}$ exciting the structure represented by transfer function matrices $\mathbf{H}_{xv}$.  In a typical vibration control problem, the control system tries to compute the signal matrix $\mathbf{G}_{vv}$ that best reproduces the desired response $\mathbf{G}_{xx}$.
    
    \section{Specification Definition}
    The first step in defining a Random Vibration control problem is the definition of the vibration response that is desired.  This vibration specification can be derived using various approaches, perhaps from test data from some environment test, predictions from a model, or derivations from a standard.  Regardless of its source, the specification defines the response levels, coherence, and phase of each control channel at each frequency line in the test.
    
    Rattlesnake accepts the specification in the form of a 3D array consisting of a complex \ac{CPSD} matrix defined at each frequency line.  Specification \ac{CPSD} matrices can be loaded from Numpy \inlineCode{*.npz} files or Matlab \inlineCode{*.mat} files.  For each of these files, Rattlesnake respects the natural dimension ordering of a dataset consisting of ``stacks'' of matrices that the specification can be visualized to be.  For Matlab, which customarily uses the third dimension as the ``stacking'' dimension for 3D datasets, the specification dimensions should be $n_c \times n_c \times n_f$ where $n_c$ is the number of control channels and $n_f$ is the number of frequency lines.  For Numpy/Python the more natural ordering is $n_f \times n_c \times n_c$, essentially taking the last dimension of the Matlab array and moving it to the first dimension in the Numpy array.  Both Matlab \inlineCode{*.mat} and Numpy \inlineCode{*.npz} files should contain the following data fields:
    
    \begin{description}
        \item[cpsd] A $n_f \times n_c \times n_c$ (for \inlineCode{*.npz} files) or $n_c \times n_c \times n_f$ (for \inlineCode{*.mat} files) complex array containing the \ac{CPSD} matrix at each frequency defined in \inlineCode{f}.
        \item[f] A $n_f$ array of frequencies corresponding to the frequency lines in the \inlineCode{CPSD} matrix.
    \end{description}

    For example, for a test consisting of three control channels has a given specification is defined from 10 Hz to 100 Hz with 2 Hz spacing, the variable \inlineCode{f} in the specification file would be length 46 and have values \inlineCode{[10, 12, 14, ... 98, 100]} and the variable \inlineCode{cpsd} would be size $46\times3\times3$ in a \inlineCode{*.npz} file or $3\times3\times46$ in a \inlineCode{*.mat} file.
    
    The ordering of the rows and columns of the \ac{CPSD} matrices defining the specification are the same order as the control channels in the Channel Table on the \inlineCode{Data Acquisition Setup} tab.  This means that the first row and column of the \ac{CPSD} matrix will correspond to the first channel that is selected as a control channel in the \inlineCode{Control Channels} list on the \inlineCode{Environment Definitions} tab.  The second row and column to the second channel selected as a control channel, and so on.  Note that if Control transformations are specified, then the first row and column of the specification will correspond to the first virtual control channel, which is the first row of the control transformation matrix.  The second row and column will correspond to the second virtual control channel.
    
    The specification is defined in units of $EU^2/Hz$ where $EU$ is the engineering unit specified by the \inlineCode{Engineering Unit} column of the channel table for the control channels.
    
    \textbf{Please note that Rattlesnake will not interpolate a specification for you!}  Any frequency line that is not defined in the specification will be set to zero.  This allows a user to specify the specification only over certain bandwidths of interest.  This also means that if a user provides a specification with 2 Hz frequency spacing but runs a test with parameters that result in 1 Hz frequency spacing, every second frequency line will end up being set to zero, which will generally result in very poor control.
    
    Rattlesnake \Ac{MIMO} Random Vibration specification files can also contain optional warning and abort limits.  Note that these limits only operate on the \Ac{APSD} portion (i.e. the diagonal) of the \Ac{CPSD} matrices.  It is not currently possible to set a limit based on, for example, the coherence between two channels in Rattlesnake.  These are defined in the specification files in fields:
    
    \begin{description}
        \item[warning\_upper] A $n_f \times n_c$ (for \inlineCode{*.npz} files) or $n_c \times n_f$ (for \inlineCode{*.mat} files) array containing an upper warning level at each frequency defined in \inlineCode{f} for each control channel.
        \item[warning\_lower] A $n_f \times n_c$ (for \inlineCode{*.npz} files) or $n_c \times n_f$ (for \inlineCode{*.mat} files) array containing a lower warning level at each frequency defined in \inlineCode{f} for each control channel.
        \item[abort\_upper] A $n_f \times n_c$ (for \inlineCode{*.npz} files) or $n_c \times n_f$ (for \inlineCode{*.mat} files) array containing an upper abort level at each frequency defined in \inlineCode{f} for each control channel.
        \item[abort\_lower] A $n_f \times n_c$ (for \inlineCode{*.npz} files) or $n_c \times n_f$ (for \inlineCode{*.mat} files) array containing a lower abort level at each frequency defined in \inlineCode{f} for each control channel.
    \end{description}
    
    Any combination of the above fields can be specified.  For example, a lower limit can be defined without an equivalent upper limit.  An abort limit can be defined without a warning limit.  However, if the field is defined in the specification file, it must have the correct shape, which means that the limit must be defined for all frequency lines and for all control channels.  If a user does not want to limit on specific frequency ranges or specific channels, the limit can be set to a value of \inlineCode{NaN}.  Rattlesnake will ignore portions of the limit specifications that contain \inlineCode{NaN} values.
    
    Throughout the \Ac{MIMO} Random Vibration environment, channels will be flagged as yellow if they cross a warning limit, and flagged as red if they cross an abort limit.  Additionally, if the \inlineCode{Allow Automatic Aborts?} checkbox is checked on the \inlineCode{Environment Definition} tab, the environment will automatically stop if the abort limit is crossed.
    
    \section{Defining the \Ac{MIMO} Random Vibration Environment in Rattlesnake}
    
    In addition to the specification, there are a number of signal processing parameters that are used by the \Ac{MIMO} Random Vibration environment.  These, along with the specification, are defined on the \inlineCode{Environment Definition} tab in the Rattlesnake controller on a sub-tab corresponding to a \Ac{MIMO} Random Vibration environment.  Figure \ref{fig:randomvibrationenvironmentdefinition} shows a \Ac{MIMO} Random Vibration sub-tab.  The following subsections describe the parameters that can be specified, as well as their effects on the analysis.
    
    \begin{figure}
        \centering
        \includegraphics[width=\linewidth]{figures/random_vibration_environment_definition}
        \caption{\Ac{GUI} used to define a \Ac{MIMO} Random Vibration environment.}
        \label{fig:randomvibrationenvironmentdefinition}
    \end{figure}

    \subsection{Sampling Parameters}
    The \inlineCode{Sampling Parameters} section of the \Ac{MIMO} Random Vibration definition sub-tab consists of the following parameters:
    \begin{description}
        \item[Sample Rate] The global sample rate of the data acquisition system.  This is set on the \inlineCode{Data Acquisition Setup} tab, and displayed here for convenience as a read-only value.
        \item[Samples Per Frame] The number of samples used when computing the \ac{FFT}.  Making this number larger will increase the frequency resolution of the measurement while also making each measurement frame longer.  Note that this value should never be set so that the frequency resolution of the \ac{FFT} is finer than the frequency resolution of the specification, otherwise the frequency lines without a corresponding value in the specification will be controlled to zero.  Because adjusting this parameter will adjust the frequency resolution of the measurement, any loaded specification will need to be re-loaded with the new frequency resolution.
        \item[Samples Per Aquire] The number of additional samples required per measurement frame when taking into account overlapping of the signals that will occur in the system identification and \ac{CPSD} calculations.  This is a computed quantity presented for convenience, so the user cannot modify it directly.
        \item[Frame Time] The amount of time it takes to measure a frame of data, in seconds.  This is a computed quantity presented for convenience, so the user cannot modify it directly.
        \item[Nyquist Frequency] The maximum bandwidth of the measurement given the sampling rate.  This is the largest frequency value in the \ac{FFT}.  This is a computed quantity presented for convenience, so the user cannot modify it directly.
        \item[FFT Lines] The number of frequency lines in the \ac{FFT}.  This is a computed quantity presented for convenience, so the user cannot modify it directly.
        \item[Frequency Spacing] The frequency resolution of the \ac{FFT}.  This is a computed quantity presented for convenience, so the user cannot modify it directly.
        \item[Test Level Ramp Time] The time it takes the environment to ramp between test levels.  This is done to prevent sharp adjustments to the controller which may damage test hardware.  This is also the time it takes for the controller to ramp up to the starting level from zero when the test is started, as well as to ramp back to zero from the final test level when testing is completed.
    \end{description}

    \subsection{Signal Generation Parameters}
    The \inlineCode{Signal Generation Parameters} section of the \Ac{MIMO} Random Vibration definition sub-tab consists of the following parameters:
    \begin{description}
        \item[COLA Window] The window function used to smooth between individual time signal realizations during the \ac{COLA} process.
        \item[COLA Overlap] The amount of overlap between signals when smoothing individual time signal realizations during the \ac{COLA} process.
        \item[Window Exponent] The exponent used on the \ac{COLA} window function.  This should generally be left at 0.5, which will result in a signal with constant variance.
        \item[Samples per Output] The number of additional samples output per measurement frame when taking into account overlapping of the signals that will occur \ac{COLA} calculations.  This is a computed quantity presented for convenience, so the user cannot modify it directly.
    \end{description}
    
    \subsection{Control Channels}
    The \inlineCode{Control Channels} list allows users to select the channels in the test that will be used by the environment to perform control calculations.  These are the channels that will match the rows and columns of the specification file.
    
    \subsection{Control Parameters}
    The \inlineCode{Control Parameters} section of the \Ac{MIMO} Random Vibration definition sub-tab consists of the following parameters:
    \begin{description}
        \item[Input Channels] The total number of channels being measured by the Rattlesnake, including response channels and output channels.  This is a computed quantity presented for convenience, so the user cannot modify it directly.
        \item[Output Channels] The number of excitation signals being used to control the vibration response of this environment.  This is a computed quantity presented for convenience, so the user cannot modify it directly.
        \item[Control Channels] The number of response channels being used in the control.  This is a computed quantity presented for convenience, so the user cannot modify it directly.
        \item[Update System ID During Control?]  Selecting this option allows Rattlesnake to continue to update \Ac{FRF} during control.  This can be useful for nonlinear systems as the structure may behave differently at higher levels of excitation.  One should be aware that if signals become too correlated during a test, the \ac{FRF} calculations may not be able to separate responses by exciter and the \ac{FRF}s may be incorrect leading to poor control.  Additionally, if something in the test setup breaks or comes loose, the controller may see that the responses are decreasing for a given excitation level and continually try to increase the excitation level, which could result damage to test equipment or the test article.
        \item[Frames in CPSD] The number of measurement frames used when computing the \ac{CPSD} matrices of the responses and output signals.  Note that if \inlineCode{Allow Automatic Aborts?} is active, the environment will not shut down the number of frames specified by this setting has been acquired.  If fewer frames have been acquired than are specified by this setting, the \ac{CPSD} matrix might be noisy due to only having acquired a small number of measurement frames.
        \item[CPSD Window]  The window function applied during the \ac{CPSD} matrix calculations.
        \item[CPSD Overlap] The percentage overlap between measurements frames used to compute the CPSD matrix.
        \item[Allow Automatic Aborts?] Selecting this option will allow the environment to shut itself down if an abort limit is reached as defined in the specification file.  If not selected, the channel will be flagged red, but the environment will still continue to run.
    \end{description}
    The \inlineCode{Control Parameters} section of the \Ac{MIMO} Random Vibration definition sub-tab also includes functionality for loading in custom control laws.  See Section \ref{sec:rattlesnake_environments_custom_control_law} for information on defining a custom control law.
    \begin{description}
        \item[Control Python Script] The Python script containing a custom control law that is currently loaded into the Rattlesnake controller.
        \item[Load] Pressing this button will bring up a file selection dialog to load in a new Python script containing a custom control law.
        \item[Control Python Function] This selector presents the functions, generators, and classes within the loaded Python script that can be used as custom control laws.
        \item[Control Parameters] This text box allows arbitrary input to be passed to custom control laws as a string.  It is up to the control law to specify what this extra input must be and parse whatever input the user gives it.
    \end{description}

    \subsection{Control and Drive Transforms}
    The \inlineCode{Control and Drive Transforms} section of the \Ac{MIMO} Random Vibration definition sub-tab consists of the following parameters:
    \begin{description}
        \item[Transformation Matrices...] Selecting this button will bring up the transformation matrices dialog box, which allows the user to specify linear transformations between the physical responses and excitation signals and virtual responses and excitation signals.  See Section \ref{sec:rattlesnake_environments_transformation_matrices} for more information.
        \item[Transform Channels] The number of virtual control degrees of freedom after applying transformation matrices.  This is a computed quantity presented for convenience, so the user cannot modify it directly.
        \item[Transform Outputs] The number of virtual excitation devices after applying transformation matrices.  This is a computed quantity presented for convenience, so the user cannot modify it directly.
    \end{description}
    Note that if Transformation matrices are defined, the number of control channels ends up being the number of rows of the \inlineCode{Response Transformation Matrix}, rather than the number of physical control channels.  The number of physical control channels will be equal to the number of columns of the transformation matrix.  The number of rows and columns of the specification loaded should be equal to the number of rows in the transformation.
    
    \subsection{Test Specification}
    The test specification is loaded into the environment in the \inlineCode{Test Specification} section of the \Ac{MIMO} Random Vibration definition:
    
    \begin{description}
        \item[Row] A drop-down menu consisting of all response channels used to visualize a certain row and column of the specification.
        \item[Column] A drop-down menu consisting of all the response channels used to visualize a certain row and column of the specification.
        \item[Load Spec] A button that when pressed will bring up a file selection dialog to read in a specification from a Numpy \inlineCode{*.npz} or Matlab \inlineCode{*.mat} file.
        \item[Specification: Single Entry] A visualization of all frequency lines of a single row and column of the \ac{CPSD} matrix in the specification.
        \item[Specification: Sum of ASDs] The sum of the diagonal entries of the \ac{CPSD} matrix at each frequency line, which is commonly used as a test metric.
    \end{description}
    
    \section{System Identification for the \Ac{MIMO} Random Vibration Environment}
    
    When all environments are defined and the \inlineCode{Initialize Environments} button is pressed, Rattlesnake will proceed to the next phase of the test, which is defined on the \inlineCode{System Identification} tab.
    
    \Ac{MIMO} Random Vibration requires a system identification phase to compute the matrices $\mathbf{H}_{xv}$ used in the control calculations of equation \eqref{eq:forward_mimo_vibration}.  Figure \ref{fig:randomvibrationenvironmentsystemidentification} shows the \ac{GUI} used to perform this phase of the test.
    
    \begin{figure}
        \centering
        \includegraphics[width=\linewidth]{figures/random_vibration_environment_system_identification}
        \caption{System identification \ac{GUI} used by the \Ac{MIMO} Random Vibration environment.}
        \label{fig:randomvibrationenvironmentsystemidentification}
    \end{figure}
    
    Rattlesnake's system identification phase will start with a noise floor check, where the data acquisition records data on all the channels without specifying an output signal.  After the noise floor is computed, the system identification phase will play out the specified signals to the excitation devices, and transfer functions will be computed using the responses of the control channels to those excitation signals.  Section \ref{sec:using_rattlesnake_system_identification} describes the System Identification tab and its various parameters and capabilities.
    
    \section{Test Predictions for the \Ac{MIMO} Random Vibration Environment}
    
    Once the system identification is performed, a test prediction will be performed and results displayed on the \inlineCode{Test Predictions} tab, shown in Figure \ref{fig:randomvibrationenvironmenttestprediction}.  This is meant to give the user an idea of the test feasibility.  The left side of the window displays excitation information, including \ac{RMS} signal levels required as well as the excitation spectra expected.  The right side of the window displays the predicted responses compared to the specification as well as the predicted \ac{RMS} dB error.  This figure will also show any abort or warning limits imposed.  Channels will be highlighted in yellow if they cross a warning level and will be highlighted in red if they cross an abort level.  For example in the test in Figure \ref{fig:randomvibrationenvironmenttestprediction}, all channels are predicted to cross the warning threshold, and a handful are predicted to cross the abort threshold.
    
    \begin{figure}
        \centering
        \includegraphics[width=\linewidth]{figures/random_vibration_environment_test_prediction}
        \caption{Test prediction \ac{GUI} which gives the user some idea of the test feasibility.}
        \label{fig:randomvibrationenvironmenttestprediction}
    \end{figure}
    
    \section{Running the \Ac{MIMO} Random Vibration Environment}
    The \Ac{MIMO} Random Vibration environment is then run on the \inlineCode{Run Test} tab of the controller.
    
    With the data acquisition system armed, the environment can be started manually with the \inlineCode{Start Environment} button.  Once running, it can be stopped manually with the \inlineCode{Stop Environment} button.  With the data acquisition system armed and the environment running, the \ac{GUI} looks like Figure \ref{fig:randomvibrationenvironmentruntest}.  
    
    \begin{figure}
        \centering
        \includegraphics[width=\linewidth]{figures/random_vibration_environment_run_test}
        \caption{\Ac{GUI} for running the \Ac{MIMO} Random Vibration environment.}
        \label{fig:randomvibrationenvironmentruntest}
    \end{figure}
    
    There are various operations that can be performed when setting up and running the \Ac{MIMO} Random Vibration environment, and many visualization operations as well.
    
    \subsection{Test Level}
    Two test levels exist in the \Ac{MIMO} Random Vibration Environment.  The \inlineCode{Current Test Level} specifies the current level of the control in decibels relative to the specification level, which is 0 dB.  Note that all data and visualizations on the \inlineCode{Run Test} window are scaled back to full level, so users should not be surprised if for example the values reported in the \inlineCode{Output Voltages (RMS)} table do not change significantly with test level.  See Chapter \ref{sec:rattlesnake_environments_generating_time_histories} for more information on this implementation detail.
    
    The second test level is the \inlineCode{Target Test Level}.  This option can be used to specify a level at which data starts streaming to the disk if the user does not wish to save low level data.  Additionally, the controller can be made to stop controlling automatically after a certain time at a target test level.  This is done to ensure that the controller does not spend too much time at a level that could eventually damage a part.
    
    \subsection{Test Timing}
    The \Ac{MIMO} Random Vibration environment has multiple options for test timing.  If \inlineCode{Continuous Run} is selected, the environment will continue until it is manually stopped.  A specific run time can be specified using the \inlineCode{Run for h:mm:ss} option and specifying a time in the \inlineCode{h:mm:ss} selector.  The \inlineCode{at Target Test Level} checkbox specifies whether or not to activate the timer at any test level or only when the test is at the target test level.
    
    The \inlineCode{MIMO} Random Vibration environment will constantly update the \inlineCode{Total Test Time} and \inlineCode{Time At Level} time displays when the environment is active.  A progress bar will be displayed when the controller is set to only run for a specified time.  When the progress bar reaches 100\%, the environment will shut down automatically.
    
    \subsection{Test Metrics and Visualizations}
    The \Ac{MIMO} Random Vibration environment displays a number of global metrics to help evaluate the success of a test.  \Ac{RMS} signal voltage values are displayed in the \inlineCode{Output Voltages (RMS)} table.  \Ac{RMS} dB errors for each control channel are displayed in the \inlineCode{Response Error (dB)} table.  These errors will also be colored yellow or red if the given channel is crossing a warning or abort level.  If an abort level is reached and the \inlineCode{Allow Automatic Aborts?} option is selected on the \inlineCode{Environment Definition} page, then the environment will shut down automatically.
    
    The \inlineCode{Run Test} tab for the \Ac{MIMO} Random Vibration environment displays the sum of \ac{APSD} functions of the response \ac{CPSD} matrix compared to the sum of \ac{APSD} functions of the specification in a large plot in the middle of the main window, which can be seen in Figure \ref{fig:randomvibrationenvironmentruntest}.  This can be considered an ``Average'' response level for the test compared to the ``Average'' specification level.
    
    To interrogate specific channels, the \inlineCode{Data Display} section of the \inlineCode{Run Test} window offers several options.  The row and column of the \ac{CPSD} matrix can be selected using \inlineCode{Control Channel 1} and \inlineCode{Control Channel 2} selectors.  The \inlineCode{Data Type} of the plot can be specified as \inlineCode{Magnitude}, \inlineCode{Phase}, \inlineCode{Coherence}, \inlineCode{Real}, or \inlineCode{Imaginary}.  Pressing the \inlineCode{Create Window} button then creates the specified plot.
    
    Some convenience operations are also included to visualize all channels.  In the \inlineCode{Show all:} section, pressing the \inlineCode{Autospectral Densities} button will bring up one window per control channel and display the \ac{APSD} function for each.  Pressing the \inlineCode{Spectral Densities (phase/coh)} or \inlineCode{Spectral Densities (real/imag)} buttons will attempt to display the entire \ac{CPSD} matrix, displaying either the phase and coherence or real and imaginary parts in the upper and lower triangular portions of the matrix.  Be aware that these convenience functions with a large number of control channels can easily overwhelm the computer with plotting operations causing the \ac{GUI} to become unresponsive, so use these operations with caution.  Figure \ref{fig:randomvibrationenvironmentchannelvisualizations} shows an example displaying the full \ac{CPSD} matrix with coherence and phase for a test with six control degrees of freedom. 
    
    \begin{figure}
        \centering
        \includegraphics[width=\linewidth]{figures/random_vibration_environment_channel_visualizations}
        \caption{Visualizing individual channels (magnitude, coherence, and phase).}
        \label{fig:randomvibrationenvironmentchannelvisualizations}
    \end{figure}
    
    If a specification has warning and abort limits defined, these will also be plotted, as shown in Figure \ref{fig:randomvibrationenvironmentchannelvisualizationswlimits}.
    
    
    \begin{figure}
        \centering
        \includegraphics[width=\linewidth]{figures/random_vibration_environment_channel_visualizations_w_limits}
        \caption{Figure showing the \ac{APSD} data for each channel, as well as the warning and abort limits.}
        \label{fig:randomvibrationenvironmentchannelvisualizationswlimits}
    \end{figure}
    
    Further convenience operations are available in the \inlineCode{Window Operations:} section.  Pressing \inlineCode{Tile All Windows} will rearrange all channel windows neatly across the screen.  Pressing \inlineCode{Close All Windows} will close all open channel windows.
    
    \subsection{Saving Data from the \Ac{MIMO} Random Environment}
    
    Time data can be saved from the \Ac{MIMO} random vibration environment through Rattlesnake's streaming functionality, described in Section \ref{sec:using_rattlesnake_run_test}.
    
    Users can also directly write the spectral data from the environment to a file.  This will also result in a netCDF file, however the fields will be slightly different.  This is described more fully in Section \ref{sec:rattlesnake_environments_store_to_netcdf}.
    
    \section{Output NetCDF File Structure}\label{sec:rattlesnake_environments_store_to_netcdf}
    When Rattlesnake streams time data to a netCDF file, environment-specific parameters are stored in a netCDF group with the same name as the environment name.  Similar to the root netCDF structure described in Section \ref{sec:using_rattlesnake_output_files}, this group will have its own attributes, dimensions, and variables, which are described here.
    
    \subsection{NetCDF Dimensions}
    \begin{description}
        \item[fft\_lines] The number of frequency lines in the \ac{FFT}.
        \item[two] A dimension of size 2, which is required for the warning and abort variables (there are two limits, upper and lower).
        \item[specification\_channels] The number of channels defined in the specification provided to the \Ac{MIMO} Random Vibration environment.
        \item[response\_transformation\_rows] The number of rows in the response channel transformation (See Section \ref{sec:rattlesnake_environments_transformation_matrices}).  This is not defined if no response transformation is used.
        \item[response\_transformation\_cols] The number of columns in the response channel transformation (See Section \ref{sec:rattlesnake_environments_transformation_matrices}).  This is not defined if no response transformation is used.
        \item[reference\_transformation\_rows] The number of rows in the output transformation (See Section \ref{sec:rattlesnake_environments_transformation_matrices}).  This is not defined if no output transformation is used.
        \item[reference\_transformation\_cols] The number of columns in the output transformation (See Section \ref{sec:rattlesnake_environments_transformation_matrices}).  This is not defined if no output transformation is used.
        \item[control\_channels] The number of physical channels used for control.  Note that this may be different from the \inlineCode{specification_channels} due to the presence of a transformation matrix.
    \end{description}
    
    \subsection{NetCDF Attributes}
    \begin{description}
        \item[sysid\_frame\_size] The number of samples per measurement frame in the system identification
        \item[sysid\_averaging\_type] The type of averaging used in the system identification, linear or exponential
        \item[sysid\_noise\_averages] The number of measurement frames acquired for the noise floor calculation
        \item[sysid\_averages] The number of measurement frames acquired for the system identification calculation
        \item[sysid\_exponential\_averaging\_coefficient] The weighting coefficient used for new frames in the exponential averaging scheme
        \item[sysid\_estimator] The \Ac{FRF} estimator used to compute the transfer functions during the system identification
        \item[sysid\_level] The level used by the system identification in volts RMS.
        \item[sysid\_level\_ramp\_time] The time to ramp up to the test level when starting and ramp back to zero when stopping the system identification
        \item[sysid\_signal\_type] The signal type used by the system identification
        \item[sysid\_window] The window function applied to the time data during the system identification
        \item[sysid\_overlap] The overlap fraction between measurement frames used for system identification
        \item[sysid\_burst\_on] The fraction of a measurement frame that a burst is active for burst random excitation during system identification
        \item[sysid\_pretrigger] The fraction of a measurement used as a pre-trigger for burst random excitation during system identification
        \item[sysid\_burst\_ramp\_fraction] The fraction of a measurement frame used to ramp the burst up to full level and back to zero
        \item[samples\_per\_frame] The number of samples per measurement frame used in the \ac{FFT}
        \item[test\_level\_ramp\_time] The time to ramp between test levels
        \item[cpsd\_overlap] The percentage overlap used when computing \ac{FRF} and \ac{CPSD} matrices
        \item[update\_tf\_during\_control] 1 if transfer functions were updated during control, 0 otherwise
        \item[cola\_window] The window function used by the \ac{COLA} process
        \item[cola\_overlap] The overlap between realizations of excitation signals used during the \ac{COLA} process
        \item[cola\_window\_exponent] The exponent on the \ac{COLA} window function
        \item[frames\_in\_cpsd] The number of frames used to compute \ac{CPSD} matrices
        \item[cpsd\_window] The window function used to compute \ac{CPSD} matrices
        \item[control\_python\_script] The path to the Python script used to control the \Ac{MIMO} Random Vibration environment
        \item[control\_python\_function] The function (or class or generator function) in the Python script used to control the \Ac{MIMO} Random Vibration environment
        \item[control\_python\_function\_type] The type of the object used for the control law (function, generator, or class)
        \item[control\_python\_function\_parameters] The extra parameters passed to the control law.
    \end{description}
    
    \subsection{NetCDF Variables}
    \begin{description}
        \item[specification\_frequency\_lines] The frequency values in the specification associated with each frequency line.  Type: 64-bit float; Dimensions: \inlineCode{fft_lines}
        
        \item[specification\_cpsd\_matrix\_real] The real part of the \Ac{MIMO} Random Vibration specification.  NetCDF files cannot handle complex data types so real and imaginary parts are split into two variables.  Type: 64-bit float; Dimensions: \inlineCode{fft_lines} $\times$ \inlineCode{specification_channels} $\times$ \inlineCode{specification_channels}
        
        \item[specification\_cpsd\_matrix\_imag] The imaginary part of the \Ac{MIMO} Random Vibration specification.  NetCDF files cannot handle complex data types so real and imaginary parts are split into two variables.  Type: 64-bit float; Dimensions: \inlineCode{fft_lines} $\times$ \inlineCode{specification_channels} $\times$ \inlineCode{specification_channels}
        
        \item[specification\_warning\_matrix] The data used to define the warning limits in the specification.  The first index in the first dimension defines the lower limit, and the second index in the first dimension defines the upper limit. Type: 64-bit float; Dimensions: \inlineCode{two} $\times$ \inlineCode{specification_channels}
        
        \item[specification\_abort\_matrix] The data used to define the abort limits in the specification.  The first index in the first dimension defines the lower limit, and the second index in the first dimension defines the upper limit. Type: 64-bit float; Dimensions: \inlineCode{two} $\times$ \inlineCode{specification_channels}
        
        \item[response\_transformation\_matrix] The response transformation matrix (See Section \ref{sec:rattlesnake_environments_transformation_matrices}).  This is not defined if no response transformation is used.  Type: 64-bit float; Dimensions: \inlineCode{response_transformation_rows} $\times$ \inlineCode{response_transformation_cols}
        
        \item[output\_transformation\_matrix] The output transformation matrix (See Section \ref{sec:rattlesnake_environments_transformation_matrices}).  This is not defined if no output transformation is used.  Type: 64-bit float; Dimensions: \inlineCode{output_transformation_rows} $\times$ \inlineCode{output_transformation_cols}
        
        \item[control\_channel\_indices] The indices of the active control channels in the environment.  Type: 32-bit int; Dimensions: \inlineCode{control_channels}
        
    \end{description}
    
    \subsection{Saving Spectral Data}
    
    In addition to time streaming, Rattlesnake's \Ac{MIMO} Random Vibration environment can also save the current realization of spectral data directly to the disk by clicking the \inlineCode{Save Current Spectral Data} button.  The spectral data is stored in a NetCDF file similar to the time streaming data; however, it has additional dimensions and variables to store the spectral data.
    
    The single additional dimension is:
    
    \begin{description}
        \item[drive\_channels] The number of drive channels active in the environment.
    \end{description}
    
    There are also several additional variables to store the spectral data:
    
    \begin{description}
        \item[frf\_data\_real] The real part of the most recently computed value for the transfer functions between the excitation signals and the control response signals.  NetCDF files cannot handle complex data types so real and imaginary parts are split into two variables.  Type: 64-bit float; Dimensions: \inlineCode{fft_lines} $\times$ \inlineCode{specification_channels} $\times$ \inlineCode{drive_channels}
        
        \item[frf\_data\_imag] The imaginary part of the most recently computed value for the transfer functions between the excitation signals and the control response signals.  NetCDF files cannot handle complex data types so real and imaginary parts are split into two variables.  Type: 64-bit float; Dimensions: \inlineCode{fft_lines} $\times$ \inlineCode{specification_channels} $\times$ \inlineCode{drive_channels}
        
        \item[frf\_coherence] The multiple coherence of the control channels computed during the test.  Type: 64-bit float; Dimensions: \inlineCode{fft_lines} $\times$ \inlineCode{specification_channels}
        
        \item[response\_cpsd\_real] The real part of the most recently computed value for the \ac{CPSD} matrix at the control channels.  NetCDF files cannot handle complex data types so real and imaginary parts are split into two variables.  Type: 64-bit float; Dimensions: \inlineCode{fft_lines} $\times$ \inlineCode{specification_channels} $\times$ \inlineCode{specification_channels}
        
        \item[response\_cpsd\_imag] The imaginary part of the most recently computed value for the \ac{CPSD} matrix at the control channels.  NetCDF files cannot handle complex data types so real and imaginary parts are split into two variables.  Type: 64-bit float; Dimensions: \inlineCode{fft_lines} $\times$ \inlineCode{specification_channels} $\times$ \inlineCode{specification_channels}
        
        \item[drive\_cpsd\_real] The real part of the most recently computed value for the \ac{CPSD} matrix at the excitation channels.  NetCDF files cannot handle complex data types so real and imaginary parts are split into two variables.  Type: 64-bit float; Dimensions: \inlineCode{fft_lines} $\times$ \inlineCode{drive_channels} $\times$ \inlineCode{drive_channels}
        
        \item[drive\_cpsd\_imag] The imaginary part of the most recently computed value for the \ac{CPSD} matrix at the excitation channels.  NetCDF files cannot handle complex data types so real and imaginary parts are split into two variables.  Type: 64-bit float; Dimensions: \inlineCode{fft_lines} $\times$ \inlineCode{drive_channels} $\times$ \inlineCode{drive_channels}
        
        \item[response\_noise\_cpsd\_real] The real part of the \ac{CPSD} matrix at the control channels during the noise floor measurement that occurred during system identification.  NetCDF files cannot handle complex data types so real and imaginary parts are split into two variables.  Type: 64-bit float; Dimensions: \inlineCode{fft_lines} $\times$ \inlineCode{specification_channels} $\times$ \inlineCode{specification_channels}
        
        \item[response\_noise\_cpsd\_imag] The imaginary part of the \ac{CPSD} matrix at the control channels during the noise floor measurement that occurred during system identification.  NetCDF files cannot handle complex data types so real and imaginary parts are split into two variables.  Type: 64-bit float; Dimensions: \inlineCode{fft_lines} $\times$ \inlineCode{specification_channels} $\times$ \inlineCode{specification_channels}
        
        \item[drive\_noise\_cpsd\_real] The real part of the \ac{CPSD} matrix at the excitation channels during the noise floor measurement that occurred during system identification.  NetCDF files cannot handle complex data types so real and imaginary parts are split into two variables.  Type: 64-bit float; Dimensions: \inlineCode{fft_lines} $\times$ \inlineCode{drive_channels} $\times$ \inlineCode{drive_channels}
        
        \item[drive\_noise\_cpsd\_imag] The imaginary part of the \ac{CPSD} matrix at the excitation channels during the noise floor measurement that occurred during system identification.  NetCDF files cannot handle complex data types so real and imaginary parts are split into two variables.  Type: 64-bit float; Dimensions: \inlineCode{fft_lines} $\times$ \inlineCode{drive_channels} $\times$ \inlineCode{drive_channels}
    \end{description}
    
    \section{Writing a Custom Control Law}\label{sec:rattlesnake_environments_custom_control_law}
    The flexibility of the Rattlesnake framework is highlighted by the ease in which users can implement and iterate on their own ideas.  For the \Ac{MIMO} Random Vibration control type, users can implement custom control laws using a custom Python function, or alternatively a generator function or class which allow state to be maintained between function calls.  This section will provide instructions and examples for implementing a custom control law.
    
    The controller will provide various data types to the control law functions which are:
    \begin{itemize}
        \item \inlineCode{specification} -- The target \Ac{CPSD} matrix for the control channels; complex 3D array ($n_f \times n_c \times n_c$)
        
        \item \inlineCode{warning_levels} -- The warning levels provided with the specification; complex 2D array ($2 \times n_f \times n_c$)
        
        \item \inlineCode{abort_levels} -- The abort levels provided with the specification; complex 2D array ($2 \times n_f \times n_c$)
        
        \item \inlineCode{transfer_function} -- The current estimate of the transfer function between the control responses and the excitation voltages; complex 3D array ($n_f \times n_c \times n_o$)
        
        \item \inlineCode{noise_response_cpsd} -- The levels and correlation of the noise floor measurement on the control channels obtained during system identification; complex 3D array ($n_{f} \times n_{c}\times n_{c}$)
        
        \item \inlineCode{noise_reference_cpsd} -- The levels and correlation of the noise floor measurement on the excitation channels obtained during system identification; complex 3D array ($n_{f} \times n_{o}\times n_{o}$)
        
        \item \inlineCode{sysid_response_cpsd} -- The levels and correlation of the control channels obtained during system identification; complex 3D array ($n_{f} \times n_{c}\times n_{c}$)

        \item \inlineCode{sysid_reference_cpsd} -- The levels and correlation of the noise floor measurement on the excitation channels obtained during system identification; complex 3D array ($n_{f} \times n_{o}\times n_{o}$)
        
        \item \inlineCode{multiple_coherence} -- The multiple coherence for each control channel; real 2D array ($n_f \times n_c$)
        
        \item \inlineCode{frames} -- The number of measurement frames acquired so far, used to compute various parameters in the control law.  This can be compared to \inlineCode{total_frames} to determine if a full set of measurement frames has been acquired, or if the estimation of the various parameters could improve with continued averaging; scalar integer
        
        \item \inlineCode{total_frames} -- The total number of frames used to compute the \Ac{CPSD} and \Ac{FRF} matrices; scalar integer
 
        \item \inlineCode{extra_parameters} -- Extra parameters provided to the controller.  The control law can parse this value to allow extra arguments to be passed to the control law; string
        
        \item \inlineCode{last_response_cpsd} -- The most recent control \Ac{CPSD}, which can be used for error-based control; complex 3D array ($n_{f} \times n_{c}\times n_{c}$)
        
        \item \inlineCode{last_output_cpsd} -- The most recent excitation \Ac{CPSD}, which can be used for drive-based control; complex 3D array ($n_{f} \times n_{o}\times n_{o}$)
    \end{itemize}
    
    where size $n_f$ is the number of frequency lines, $n_c$ is the number of control channels, and $n_o$ is the number of output signals.  Note that the values passed into the function may be defined using arbitrary variable names (e.g. \inlineCode{transfer_function} may be instead called \inlineCode{H}, or \inlineCode{specification} may be instead called \inlineCode{spec} or \inlineCode{Syy}); however, the order of the variables passed into each function will always be consistent.

    \subsection{Defining a control law using a Python function}
    Python functions are the simplest approach to define a custom control law that can be used with the Rattlesnake software; however, they are limited in that a function's state is completely lost when a function returns.  Still, they can be used to implement relatively complex control laws as long as no state persistence is required.
    
    A Python function used to define a \Ac{MIMO} Random Vibration control law in Rattlesnake would have the following general structure within a Python script.
    
    \begin{lstlisting}[language=Python,caption={General Python function structure for defining a custom Random Vibration control law called \inlineCode{control_law} in Rattlesnake},label=lst:control_function_structure]
# Any module imports, initialization code, or helper functions would go here

# Now we define the control law.  It always receives the same arguments from the controller.
def control_law(specification, # Specifications
                warning_levels, # Warning levels
                abort_levels, # Abort Levels
                transfer_function,  # Transfer Functions
                noise_response_cpsd,  # Noise levels and correlation 
                noise_reference_cpsd, # from the system identification
                sysid_response_cpsd,  # Response levels and correlation
                sysid_reference_cpsd, # from the system identification
                multiple_coherence, # Coherence from the system identification
                frames, # Number of frames in the CPSD and FRF matrices
                total_frames, # Total frames that could be in the CPSD and FRF matrices
                extra_parameters = '', # Extra parameters for the control law
                last_response_cpsd = None, # Last Control Response for Error Correction
                last_output_cpsd = None, # Last Control Excitation for Drive-based control
                ):

    # Code to perform the control would go here
    output_cpsd = ...

    # Finally, we need to return an output CPSD matrix
    return output_cpsd\end{lstlisting}
        
    The function must return an \inlineCode{output_cpsd}, which is a complex 3D array with size ($n_{f} \times n_{o}\times n_{o}$).
    
    Three examples are presented to illustrate how a control function may be created.
    
    \subsubsection{Pseudoinverse Control}
    Perhaps the simplest strategy to perform \ac{MIMO} control is to simply invert the transfer function matrix to recover the least-squares solution of the optimal output signal from the desired responses.  This first example will demonstrate that approach.
    
    The mathematics for this control strategy are relatively simple; pre- and post-multiply the specification $\mathbf{G}_{xx}$ by the pseudoinverse ($^+$) of the transfer function matrix $\mathbf{H}_{xv}$, noting that the post-multiplicand is complex-conjugate transposed ($^H$).  This calculation is performed for each frequency line.
    
    \begin{equation}
        \mathbf{G}_{vv} = {\mathbf{H}_{xv}}^+\mathbf{G}_{xx}{{\mathbf{H}_{xv}}^+}^H
    \end{equation}
    
    In Python code, the above mathematics would look like
    \begin{lstlisting}[language=Python,caption={Computing the pseudoinverse calculation to solve for a least-squares output \ac{CPSD} matrix},label=lst:pseudoinverse_computation]
import numpy as np # Import numpy to get access to the pseudoinverse (pinv) function
H_pinv = np.linalg.pinv(H_xv) # Invert the transfer function and assign to a variable so we don't have to invert twice
G_vv = H_pinv@G_xx@H_pinv.conjugate().transpose(0,2,1) # Perform the mathematics described above.\end{lstlisting}
    
    For users not familiar with Python and its numeric library \inlineCode{numpy}, the following points are clarified
    \begin{itemize}
        \item \inlineCode{numpy} is imported and assigned to the alias \inlineCode{np}, which lets us just type in \inlineCode{np} rather than the longer name \inlineCode{numpy} when we want to access \inlineCode{numpy} functions.
        \item The \inlineCode{numpy} pseudoinverse function \inlineCode{pinv} is stored in the linear algebra package \inlineCode{linalg} within \inlineCode{numpy}, therefore to access \inlineCode{pinv}, we need to call \inlineCode{np.linalg.pinv}
        \item The \inlineCode{pinv} can perform a pseudoinverse on "stacks" of matrices, so even though we are only calling the \inlineCode{pinv} function once, it is actually performing the pseudoinverse over all frequency lines
        \item The \inlineCode{@} symbol in Python is the matrix multiplication operation.  Unlike Matlab, Python doesn't support the syntax \inlineCode{.*} to differentiate elementwise and matrix multiplcation.  In Python, \inlineCode{*} is elementwise and \inlineCode{@} is matrix multiplication.  This operation also works over stacks of matrices, so \inlineCode{G_vv} is computed over all frequency lines.
        \item The \inlineCode{transpose} function of a \inlineCode{numpy} array accepts as its arguments the new ordering of the indices.  Recalling that in Python, the first index is index 0, the second is index 1, etc., essentially what this command is doing is taking the existing indices \inlineCode{(0,1,2)} and re-ordering them as \inlineCode{(0,2,1)}, or said another way \inlineCode{(frequency_line,row,column)} re-ordered as \inlineCode{(frequency_line,column,row)}, effectively transposing each matrix in the stack without modifying the 0-index corresponding to frequency line.
    \end{itemize}
    
    Wrapping the above mathematics into the function definition from Listing \ref{lst:control_function_structure}, the control law can be defined as
    
    \begin{lstlisting}[language=Python,caption={A pseudoinverse control law that can be loaded into Rattlesnake},label=lst:pseudoinverse_control]
import numpy as np

def pseudoinverse_control(
        specification, # Specifications
        warning_levels, # Warning levels
        abort_levels, # Abort Levels
        transfer_function,  # Transfer Functions
        noise_response_cpsd,  # Noise levels and correlation 
        noise_reference_cpsd, # from the system identification
        sysid_response_cpsd,  # Response levels and correlation
        sysid_reference_cpsd, # from the system identification
        multiple_coherence, # Coherence from the system identification
        frames, # Number of frames in the CPSD and FRF matrices
        total_frames, # Total frames that could be in the CPSD and FRF matrices
        extra_parameters = '', # Extra parameters for the control law
        last_response_cpsd = None, # Last Control Response for Error Correction
        last_output_cpsd = None, # Last Control Excitation for Drive-based control
        ):
    # Invert the transfer function using the pseudoinverse
    tf_pinv = np.linalg.pinv(transfer_function)
    # Return the least squares solution for the new output CPSD
    return tf_pinv@specification@tf_pinv.conjugate().transpose(0,2,1)\end{lstlisting}
    
    where the variables have been renamed from single letters (\inlineCode{G}, \inlineCode{H}) to something more meaningful (\inlineCode{specification}, \inlineCode{transfer_function}).
    
    This example shows that a control law can be implemented in only two lines of code in addition to the function boilerplate code.  Therefore even users not familiar with the Python programming language should not be intimidated by the coding required to implement custom control laws.
    
    \subsubsection{Adding Optional Arguments}
    
    When performing the pseudoinverse, users may be wary of having an ill-conditioned transfer function matrix.  This can be due to the fact that at a resonance of the structure, all responses tend to look like the mode shape at that resonance.  Therefore the condition number of the \Ac{FRF} matrix can be quite high.  \inlineCode{numpy}'s \inlineCode{pinv} function can accept an optional argument \inlineCode{rcond} which performs singular value truncation for very small singular values.  We can allow users to enter an \inlineCode{rcond} value through the \inlineCode{extra_parameters} argument.
    
    In this implementation, we try to convert the data passed as an extra parameter to a floating point number.  If we can, we use that as the \inlineCode{rcond} value.  If we can't we just use a default value of \inlineCode{rcond}.
    
    \begin{lstlisting}[language=Python,caption={A pseudoinverse control law that can be loaded into Rattlesnake that utilizes extra parameters},label=lst:pseudoinverse_control_w_rcond]
import numpy as np

def pseudoinverse_control(
        specification, # Specifications
        warning_levels, # Warning levels
        abort_levels, # Abort Levels
        transfer_function,  # Transfer Functions
        noise_response_cpsd,  # Noise levels and correlation 
        noise_reference_cpsd, # from the system identification
        sysid_response_cpsd,  # Response levels and correlation
        sysid_reference_cpsd, # from the system identification
        multiple_coherence, # Coherence from the system identification
        frames, # Number of frames in the CPSD and FRF matrices
        total_frames, # Total frames that could be in the CPSD and FRF matrices
        extra_parameters = '', # Extra parameters for the control law
        last_response_cpsd = None, # Last Control Response for Error Correction
        last_output_cpsd = None, # Last Control Excitation for Drive-based control
        ):
    try:
        rcond = float(extra_parameters)
    except ValueError:
        rcond = 1e-15
    # Invert the transfer function using the pseudoinverse
    tf_pinv = np.linalg.pinv(transfer_function,rcond)
    # Return the least squares solution for the new output CPSD
    output = tf_pinv@specification@tf_pinv.conjugate().transpose(0,2,1)
    return output\end{lstlisting}
    
    \subsubsection{Trace-matching Pseudoinverse Control}
    
    While the previous example showed that a simple control law could be implemented in a few lines of code, users may argue that this simple control scheme is not representative of a control law that one might use in practice.  Therefore, the next example will illustrate the transformation of the first example into a closed-loop control law that corrects for error at each frequency line.  This is the essence of a closed-loop controller: the controller is able to respond to errors in the response and modify the output to accommodate.
    
    This control strategy is implemented by computing the trace (the sum of the diagonal of the matrix) of the specification and the trace of the last response, and then multiplying the last output \ac{CPSD} by the ratio of the two at each frequency line.  The trace can be computed efficiently using the \inlineCode{numpy} \inlineCode{einsum} function.
    
    \begin{lstlisting}[language=Python,caption={A short function to compute the trace of a \ac{CPSD} matrix in Python}]
def trace(cpsd):
    return np.einsum('ijj->i',cpsd)\end{lstlisting}
    
    The first time through the control law, when there is no previous data to use for error correction, the control strategy will perform a simple pseudoinverse control scheme.
    
    \begin{lstlisting}[language=Python]
tf_pinv = np.linalg.pinv(transfer_function)
output = tf_pinv@specification@tf_pinv.conjugate().transpose(0,2,1)\end{lstlisting}
    
    Subsequent times through the control law, the trace ratio is computed from the previous responses, and the ratio is multiplied by the previous output.  The trace ratio is also checked for \inlineCode{nan} quantities to ensure that there are no divide-by-zero errors.
    
    \begin{lstlisting}[language=Python]
trace_ratio = trace(specification)/trace(last_response_cpsd)
trace_ratio[np.isnan(trace_ratio)] = 0
output = last_output_cpsd*trace_ratio[:,np.newaxis,np.newaxis]\end{lstlisting}
    
    The final code for the closed-loop control law shown in Listing \ref{lst:match_trace_pseudoinverse_control}.  On the first run-through, the \inlineCode{last_response_cpsd} and \inlineCode{last_output_cpsd} are set to \inlineCode{None} which is how the function knows whether or not to compute the output using the pseudoinverse control or by updating the trace.
    
    \begin{lstlisting}[language=Python,caption={A closed-loop control law to match the trace of the \ac{CPSD} matrix at each frequency line.},label=lst:match_trace_pseudoinverse_control]
def match_trace_pseudoinverse(
        specification, # Specifications
        warning_levels, # Warning levels
        abort_levels, # Abort Levels
        transfer_function,  # Transfer Functions
        noise_response_cpsd,  # Noise levels and correlation 
        noise_reference_cpsd, # from the system identification
        sysid_response_cpsd,  # Response levels and correlation
        sysid_reference_cpsd, # from the system identification
        multiple_coherence, # Coherence from the system identification
        frames, # Number of frames in the CPSD and FRF matrices
        total_frames, # Total frames that could be in the CPSD and FRF matrices
        extra_parameters = '', # Extra parameters for the control law
        last_response_cpsd = None, # Last Control Response for Error Correction
        last_output_cpsd = None, # Last Control Excitation for Drive-based control
        ):
    try:
        rcond = float(extra_parameters)
    except ValueError:
        rcond = 1e-15
    # If it's the first time through, do the actual control
    if last_output_cpsd is None:
        # Invert the transfer function using the pseudoinverse
        tf_pinv = np.linalg.pinv(transfer_function,rcond)
        # Return the least squares solution for the new output CPSD
        output = tf_pinv@specification@tf_pinv.conjugate().transpose(0,2,1)
    else:
        # Scale the last output cpsd by the trace ratio between spec and last response
        trace_ratio = trace(specification)/trace(last_response_cpsd)
        trace_ratio[np.isnan(trace_ratio)] = 0
        output =  last_output_cpsd*trace_ratio[:,np.newaxis,np.newaxis]
    return output\end{lstlisting}
    
    As can be seen in the previous Listing, the simple pseudoinverse control law can be extended to a closed-loop, error-correcting control law simply by the addition of perhaps 10 more lines of code.  This again shows that even relatively complex control strategies can be implemented easily within the Rattlesnake framework.
    
    \subsubsection{Shape-Constrained Control}
    
    The final example control law that will be shown in this section is a more complex control law that constrains the exciters to work together to reduce the force required in a given test \cite{schultz2020_shape_constrained_input_estimation_efficient_multishaker_vibration_testing}.  This shape-constrained approach utilizes a singular value decomposition of the transfer function matrix to determine the constraints to apply to the shakers as well as how many shapes to keep at each frequency line.
    
    A set of shapes used as a constraint can be defined by a matrix $\mathbf{C}$ to form a constrained transfer function matrix $\mathbf{H}_{c}$
    
    \begin{equation}
    \mathbf{H}_c = \mathbf{H}_{xv}\mathbf{C}
    \end{equation}
    
    where $\mathbf{C}$ will generally have fewer columns than rows.  This matrix effectively reduces the number of control degrees of freedom at a frequency line.  The control equation then looks like 
    
    \begin{equation}
    \hat{\mathbf{G}}_{vv} = [\mathbf{H}_{xv}\mathbf{C}]^+\mathbf{G}_{xx}{[\mathbf{H}_{xv}\mathbf{C}]^+}^H
    \end{equation}
    
    The \ac{CPSD} matrix $\hat{\mathbf{G}}_{vv}$ is defined using the constrained control degrees of freedom.  The true physical degrees of freedom can be computed from the constrained set by 
    
    \begin{equation}
    \mathbf{G}_{vv} = \mathbf{C}\hat{\mathbf{G}}_{vv}\mathbf{C}^H
    \end{equation}
    
    To select the constraint shapes $\mathbf{C}$, the right singular vectors $\mathbf{V}$ of the singular value decomposition of the transfer function matrix are used.  A singular value threshold is used to only keep the right singular vectors $\mathbf{V}_1$ corresponding to large singular values, and discarding the right singular vectors $\mathbf{V}_2$ corresponding to the small singular values.
    
    \begin{equation}
    \mathbf{H}_{xv} = \mathbf{U}\mathbf{\Sigma}\mathbf{V}^H
    \end{equation}
    \begin{equation}
    \mathbf{V} = [\mathbf{V}_1 \; \mathbf{V}_2]
    \end{equation}
    \begin{equation}
    \mathbf{H}_{xv}\mathbf{C} = \mathbf{U}\mathbf{\Sigma}\mathbf{V}^H\mathbf{V}_1
    \end{equation}
    
    Converting this control strategy into a Python control law is reasonably straightforward.  The first approach is to perform the \ac{SVD} on the transfer function matrix.
    
    \begin{lstlisting}[language=Python]
[U,S,Vh] = np.linalg.svd(H,full_matrices=False)
V = Vh.conjugate().transpose(0,2,1)\end{lstlisting}
    
    Here, we use the \inlineCode{numpy} singular value decomposition function \inlineCode{svd}.  Again, like many \inlineCode{numpy} functions, this function behaves correctly on stacks of matrices, so the \inlineCode{svd} function need be called only once to perform the operation over all frequency lines.  The \inlineCode{full_matrices} argument essentially asks whether or not the null space of the larger singular vector matrix is computed (returning an $m \times m$ \inlineCode{U} matrix rather than an $m \times k$ matrix where $k$ is the number of singular values).  That isn't required for this operation, so it is set to \inlineCode{False}.  The output from the \inlineCode{svd} function returns $\mathbf{V}^H$, so it is complex-conjugate transposed to get $\mathbf{V}$.
    
    The next step is to compute the singular values to keep based off the singular value ratios.  Singular values are kept if they are above a certain ratio to the primary singular value.
    
    \begin{lstlisting}[language=Python]
singular_value_ratios = S/S[:,0,np.newaxis]
num_shape_vectors = np.sum(singular_value_ratios >= shape_constraint_threshold,axis=1)\end{lstlisting}
    
    At this point, we perform the shape constrained control.  A \inlineCode{for} loop is required to iterate through the frequency lines because a different number of vectors is used for each frequency line.  The constraint matrix is computed using the right singular vectors corresponding to the singular values that are above the threshold.  The transfer function matrix is then constrained and the control problem is solved using the constrained transfer function matrix.  The constrained output response is then transformed back to the physical space using the constraint matrix.
    
    \begin{lstlisting}[language=Python]
output = np.empty((transfer_function.shape[0],transfer_function.shape[2],transfer_function.shape[2]),dtype=complex)
for i_f,(V_f,spec_f,H_f,num_shape_vectors_f) in enumerate(zip(V,specification,transfer_function,num_shape_vectors)):
    # Form constraint matrix
    constraint_matrix = V_f[:,:num_shape_vectors_f]
    # Constraint FRF matrix
    HC = H_f@constraint_matrix
    HC_pinv = np.linalg.pinv(HC)
    # Estimate inputs (constrained)
    SxxC = HC_pinv@spec_f@HC_pinv.conjugate().T
    # Convert to full inputs
    output[i_f] = constraint_matrix@SxxC@constraint_matrix.conjugate().T\end{lstlisting}
    
    The entire script is then shown below.  Note that the singular value threshold is passed to the function in the \inlineCode{extra_parameters} string, which is converted from a string to a floating point number.
    
    \begin{lstlisting}[language=Python,caption={A shape-constrained control law that can be used in the Rattlesnake controller}]
import numpy as np
    
def shape_constrained_pseudoinverse(
        specification, # Specifications
        warning_levels, # Warning levels
        abort_levels, # Abort Levels
        transfer_function,  # Transfer Functions
        noise_response_cpsd,  # Noise levels and correlation 
        noise_reference_cpsd, # from the system identification
        sysid_response_cpsd,  # Response levels and correlation
        sysid_reference_cpsd, # from the system identification
        multiple_coherence, # Coherence from the system identification
        frames, # Number of frames in the CPSD and FRF matrices
        total_frames, # Total frames that could be in the CPSD and FRF matrices
        extra_parameters = '', # Extra parameters for the control law
        last_response_cpsd = None, # Last Control Response for Error Correction
        last_output_cpsd = None, # Last Control Excitation for Drive-based control
        ):
    shape_constraint_threshold = float(extra_parameters)
    # Perform SVD on transfer function
    [U,S,Vh] = np.linalg.svd(transfer_function,full_matrices=False)
    V = Vh.conjugate().transpose(0,2,1)
    singular_value_ratios = S/S[:,0,np.newaxis]
    # Determine number of constraint vectors to use
    num_shape_vectors = np.sum(singular_value_ratios >= shape_constraint_threshold,axis=1)
    # We have to go into a For Loop here because V changes size on each iteration
    output = np.empty((transfer_function.shape[0],transfer_function.shape[2],transfer_function.shape[2]),dtype=complex)
    for i_f,(V_f,spec_f,H_f,num_shape_vectors_f) in enumerate(zip(V,specification,transfer_function,num_shape_vectors)):
        # Form constraint matrix
        constraint_matrix = V_f[:,:num_shape_vectors_f]
        # Constraint FRF matrix
        HC = H_f@constraint_matrix
        HC_pinv = np.linalg.pinv(HC)
        # Estimate inputs (constrained)
        SxxC = HC_pinv@spec_f@HC_pinv.conjugate().T
        # Convert to full inputs
        output[i_f] = constraint_matrix@SxxC@constraint_matrix.conjugate().T
    return output\end{lstlisting}
    
    This example shows that even complex control laws can be written in less than 100 lines of code.
    
    \subsection{Defining a control law using state-persistent approaches}
    
    While the function approach is useful in its simplicity there are certain applications where it is not sufficient.  These primarily revolve around cases where there is a significant amount of setup computations or response history that must be tracked.  To demonstrate, the Buzz Test approach by Daborn is used for illustration \cite{daborn2014_smarter_dynamic_testing_critical_structures}.
    
    The Buzz Test control strategy uses a flat random ``buzz'' test of the part to determine preferred phasing and coherence between the control degrees of freedom.  This ``buzz'' comes from the system identification phase of the controller, which is one of the inputs to the control law.  The specification is then modified so the coherence and phase of the buzz test are matched by the specification.  The same pseudoinverse control is then performed as described above, except now with the modified specification.
    
    For this case, it is helpful to define several smaller functions.  The first function will compute the coherence of each entry in a \ac{CPSD} matrix
    
    \begin{equation}
    {\gamma^2}_{ij}=\frac{{\|G_{ij}\|}^2}{G_{ii}G_{jj}}
    \end{equation}
    
    A vectorized Python implementation of this function is 
    
    \begin{lstlisting}[language=Python]
def cpsd_coherence(cpsd):
    num = np.abs(cpsd)**2
    den = (cpsd[:,np.newaxis,np.arange(cpsd.shape[1]),np.arange(cpsd.shape[2])]*
           cpsd[:,np.arange(cpsd.shape[1]),np.arange(cpsd.shape[2]),np.newaxis])
    den[den==0.0] = 1 # This prevents divide-by-zero errors from ruining the matrix for frequency lines where the specification is zero
    return np.real(num/den)\end{lstlisting}
    
    Similarly, a second function is defined that computes the phase of each entry in a \ac{CPSD} matrix.
    
    \begin{equation}
    {\phi_{ij}} = \angle{G_{ij}}
    \end{equation}
    
    and the vectorized Python function is 
    
    \begin{lstlisting}[language=Python]
def cpsd_phase(cpsd):
    return np.angle(cpsd)
    \end{lstlisting}
    
    We also need a function that can get the \ac{APSD} functions (diagonal terms) from a \ac{CPSD} matrix.  This can be done very efficiently with the \inlineCode{numpy} Einstein Summation function \inlineCode{einsum}.
    
    \begin{lstlisting}[language=Python]
def cpsd_autospectra(cpsd):
    return np.einsum('ijj->ij',cpsd)
    \end{lstlisting}
    
    Now that functions are defined to extract the various parts of a \ac{CPSD} matrix, a function is defined that assembles a \ac{CPSD} matrix from those parts.  This will look like
    \begin{equation}
    G_{jk} = e^{i\phi_{jk}}\sqrt{{\gamma^2}_{jk}G_{jj}G_{kk}}
    \end{equation}
    
    which in vectorized \inlineCode{numpy} Python looks like
    \begin{lstlisting}[language=Python]
def cpsd_from_coh_phs(asd,coh,phs):
    return np.exp(phs*1j)*np.sqrt(coh*asd[:,:,np.newaxis]*asd[:,np.newaxis,:])\end{lstlisting}
    
    Then finally, function is defined that will extract the autospectra from one \ac{CPSD} matrix and assemble a new \ac{CPSD} matrix using the coherence and phase from a second \ac{CPSD} matrix.
    
    \begin{lstlisting}[language=Python]
def match_coherence_phase(cpsd_to_modify,cpsd_to_match):
    coh = cpsd_coherence(cpsd_to_match)
    phs = cpsd_phase(cpsd_to_match)
    asd = cpsd_autospectra(cpsd_to_modify)
    return cpsd_from_coh_phs(asd,coh,phs)\end{lstlisting}
    
    In a Buzz Test control law defined using a Python function, the specification is updated using the phase and coherence of the CPSD from the system identification phase, and control is performed using pseudoinverse control to the updated specification.
    
    \begin{lstlisting}[language=Python]
def buzz_control(
        specification, # Specifications
        warning_levels, # Warning levels
        abort_levels, # Abort Levels
        transfer_function,  # Transfer Functions
        noise_response_cpsd,  # Noise levels and correlation 
        noise_reference_cpsd, # from the system identification
        sysid_response_cpsd,  # Response levels and correlation
        sysid_reference_cpsd, # from the system identification
        multiple_coherence, # Coherence from the system identification
        frames, # Number of frames in the CPSD and FRF matrices
        total_frames, # Total frames that could be in the CPSD and FRF matrices
        extra_parameters = '', # Extra parameters for the control law
        last_response_cpsd = None, # Last Control Response for Error Correction
        last_output_cpsd = None, # Last Control Excitation for Drive-based control
        ):
    # Create a new specification using the autospectra from the original and
    # phase and coherence of the buzz_cpsd
    spec = match_coherence_phase(specification,sysid_response_cpsd)
    # Invert the transfer function using the pseudoinverse
    tf_pinv = np.linalg.pinv(transfer_function)
    # Return the least squares solution for the new output CPSD
    return tf_pinv@spec@tf_pinv.conjugate().transpose(0,2,1)\end{lstlisting}
    
    The entire control script that is loaded into the Rattlesnake controller is then here:
    
    \begin{lstlisting}[language=Python,caption={A Buzz Test control law defined using a Python function that can be used with the Rattlesnake software.},label = lst:buzz_test_function]
import numpy as np

# Helper functions
def cpsd_coherence(cpsd):
    num = np.abs(cpsd)**2
    den = (cpsd[:,np.newaxis,np.arange(cpsd.shape[1]),np.arange(cpsd.shape[2])]*
           cpsd[:,np.arange(cpsd.shape[1]),np.arange(cpsd.shape[2]),np.newaxis])
    den[den==0.0] = 1 # This prevents divide-by-zero errors from ruining the matrix for frequency lines where the specification is zero
    return np.real(num/
    den)
    
def cpsd_phase(cpsd):
    return np.angle(cpsd)
    
def cpsd_autospectra(cpsd):
    return np.einsum('ijj->ij',cpsd)
    
def cpsd_from_coh_phs(asd,coh,phs):
    return np.exp(phs*1j)*np.sqrt(coh*asd[:,:,np.newaxis]*asd[:,np.newaxis,:])
    
def match_coherence_phase(cpsd_to_modify,cpsd_to_match):
    coh = cpsd_coherence(cpsd_to_match)
    phs = cpsd_phase(cpsd_to_match)
    asd = cpsd_autospectra(cpsd_to_modify)
    return cpsd_from_coh_phs(asd,coh,phs)
    
def buzz_control(
        specification, # Specifications
        warning_levels, # Warning levels
        abort_levels, # Abort Levels
        transfer_function,  # Transfer Functions
        noise_response_cpsd,  # Noise levels and correlation 
        noise_reference_cpsd, # from the system identification
        sysid_response_cpsd,  # Response levels and correlation
        sysid_reference_cpsd, # from the system identification
        multiple_coherence, # Coherence from the system identification
        frames, # Number of frames in the CPSD and FRF matrices
        total_frames, # Total frames that could be in the CPSD and FRF matrices
        extra_parameters = '', # Extra parameters for the control law
        last_response_cpsd = None, # Last Control Response for Error Correction
        last_output_cpsd = None, # Last Control Excitation for Drive-based control
        ):
    # Create a new specification using the autospectra from the original and
    # phase and coherence of the buzz_cpsd
    spec = match_coherence_phase(specification,sysid_response_cpsd)
    # Invert the transfer function using the pseudoinverse
    tf_pinv = np.linalg.pinv(transfer_function)
    # Return the least squares solution for the new output CPSD
    return tf_pinv@spec@tf_pinv.conjugate().transpose(0,2,1)\end{lstlisting}
    
    \subsubsection{Defining Control Laws with Generator Functions}
    
    Note that while the above is a usable control law, it is not optimal.  Note that every time the function is called, the modified specification is recomputed, which can result in a significant amount of computation for large control problems.  Even if we could tell the control law to only compute it the first time (for example, by checking if \inlineCode{last_response_cpsd} or \inlineCode{last_output_cpsd} is \inlineCode{None} as was done previously) there is no way to store the modified specification as all variables local to the function are lost when the function returns.
    
    For this reason, Rattlesnake has alternative approaches to defining control laws that allow state persistence.  The second strategy to define a control law in Rattlesnake is to use a Generator Function.  A generator function is simply a function that maintains its internal state between function calls.
    
    Listing \ref{lst:buzz_test_generator} shows the Buzz test approach described above in a generator format.
    
    \begin{lstlisting}[language=Python, caption = {A Buzz Test control law defined using a Python generator function to allow for state persistence},label=lst:buzz_test_generator]
import numpy as np

def cpsd_coherence(cpsd):
    num = np.abs(cpsd)**2
    den = (cpsd[:,np.newaxis,np.arange(cpsd.shape[1]),np.arange(cpsd.shape[2])]*
           cpsd[:,np.arange(cpsd.shape[1]),np.arange(cpsd.shape[2]),np.newaxis])
    den[den==0.0] = 1 # Set to 1
    return np.real(num/
                   den)

def cpsd_phase(cpsd):
    return np.angle(cpsd)

def cpsd_from_coh_phs(asd,coh,phs):
    return np.exp(phs*1j)*np.sqrt(coh*asd[:,:,np.newaxis]*asd[:,np.newaxis,:])

def cpsd_autospectra(cpsd):
    return np.einsum('ijj->ij',cpsd)

def match_coherence_phase(cpsd_original,cpsd_to_match):
    coh = cpsd_coherence(cpsd_to_match)
    phs = cpsd_phase(cpsd_to_match)
    asd = cpsd_autospectra(cpsd_original)
    return cpsd_from_coh_phs(asd,coh,phs)

def buzz_control_generator():
    output_cpsd = None
    modified_spec = None
    while True:
        (specification, # Specifications
         warning_levels, # Warning levels
         abort_levels, # Abort Levels
         transfer_function,  # Transfer Functions
         noise_response_cpsd,  # Noise levels and correlation 
         noise_reference_cpsd, # from the system identification
         sysid_response_cpsd,  # Response levels and correlation
         sysid_reference_cpsd, # from the system identification
         multiple_coherence, # Coherence from the system identification
         frames, # Number of frames in the CPSD and FRF matrices
         total_frames, # Total frames that could be in the CPSD and FRF matrices
         extra_parameters, # Extra parameters for the control law
         last_response_cpsd, # Last Control Response for Error Correction
         last_output_cpsd, # Last Control Excitation for Drive-based control
            ) = yield output_cpsd
        # Only comput the modified spec if it hasn't been yet.
        if modified_spec is None:
            modified_spec = match_coherence_phase(specification,sysid_response_cpsd)
         # Invert the transfer function using the pseudoinverse
        tf_pinv = np.linalg.pinv(transfer_function)
        # Assign the output_cpsd so it is yielded next time through the loop
        output_cpsd = tf_pinv@modified_spec@tf_pinv.conjugate().transpose(0,2,1)\end{lstlisting}
        
    Note that the generator function itself is not called with any arguments, as the initial function call simply starts up the generator.  Also note that there is no \inlineCode{return} statement in a generator function, only a \inlineCode{yield} statement.   When a program requests the next value from a generator, the generator code proceeds until it hits a \inlineCode{yield} statement, at which time it pauses and waits for the next value to be requested.  During this pause, all internal data is maintained inside the generator function.  The \inlineCode{yield} statement also accepts new data into the generator function, so this is where the same arguments used to define a control law using a Python function are passed in to the generator control law.  Therefore, by creating a while loop inside a generator function, the generator can be called infinitely many times to deliver data to the controller.
    
    To implement the Buzz test as shown in Listing \ref{lst:buzz_test_generator} the modified specification is initialized as \inlineCode{None}, which enables the generator to check whether or not it has been computed yet.  If it has, then there is no need to compute it again.
    
    \subsubsection{Defining Control Laws using Classes}
    
    A final way to implement more complex control laws is using a Python class.  This approach allows for the near infinite flexibility of Python's object-oriented programming at the expense of more complex syntax.  Users not familiar with Python's object-oriented programming paradigms are encouraged to \href{https://docs.python.org/3/tutorial/classes.html}{learn more} about the topic prior to reading this section.
    
    A class in Python is effectively a container that can have functions and properties stored inside of it, so it provides a good way to encapsulate all the parameters and helper functions associated with a given control law into one place.  A class allows for arbitrary properties to be stored within it, so arbitrary data can be made persistent between control function calls.
    
    For the Rattlesnake implementation of a control law, the class must have at a minimum of three functions defined.  These are the class constructor \inlineCode{__init__} that is called when the class is instantiated, a \inlineCode{system_id_update} function that is called upon completion of the System Identification portion of the controller, and a \inlineCode{control} function that actually computes the output \ac{CPSD} matrix.  A general class structure is shown below
    
    \begin{lstlisting}[language=Python,caption={Structure for a class defining a control law in Rattlesnake}]
# Any module imports or constants would go here    

class ControlLawClass:
    def __init__(
            self,
            specification : np.ndarray, # Specifications
            warning_levels  : np.ndarray, # Warning levels
            abort_levels  : np.ndarray, # Abort Levels
            extra_parameters : str, # Extra parameters for the control law
            transfer_function : np.ndarray = None,  # Transfer Functions
            noise_response_cpsd : np.ndarray = None,  # Noise levels and correlation 
            noise_reference_cpsd : np.ndarray = None, # from the system identification
            sysid_response_cpsd : np.ndarray = None,  # Response levels and correlation
            sysid_reference_cpsd : np.ndarray = None, # from the system identification
            multiple_coherence : np.ndarray = None, # Coherence from the system identification
            frames = None, # Number of frames in the CPSD and FRF matrices
            total_frames = None, # Total frames that could be in the CPSD and FRF matrices
            last_response_cpsd : np.ndarray = None, # Last Control Response for Error Correction
            last_output_cpsd : np.ndarray = None, # Last Control Excitation for Drive-based control
            ):
        # Code to initialize the control law would go here
    
    def system_id_update(
            self,
            transfer_function : np.ndarray = None,  # Transfer Functions
            noise_response_cpsd : np.ndarray = None,  # Noise levels and correlation 
            noise_reference_cpsd : np.ndarray = None, # from the system identification
            sysid_response_cpsd : np.ndarray = None,  # Response levels and correlation
            sysid_reference_cpsd : np.ndarray = None, # from the system identification
            multiple_coherence : np.ndarray = None, # Coherence from the system identification
            frames = None, # Number of frames in the CPSD and FRF matrices
            total_frames = None, # Total frames that could be in the CPSD and FRF matrices
            ):
        # Code to update the control law with system identification information would go here
    
    def control(
            self,
            transfer_function : np.ndarray = None,  # Transfer Functions
            multiple_coherence : np.ndarray = None, # Coherence from the system identification
            frames = None, # Number of frames in the CPSD and FRF matrices
            total_frames = None, # Total frames that could be in the CPSD and FRF matrices
            last_response_cpsd : np.ndarray = None, # Last Control Response for Error Correction
            last_output_cpsd : np.ndarray = None) -> np.ndarray:
        # Code to perform the actual control operations would go here
        
    # Any helper functions or properties that belong with the class could go here
    \end{lstlisting}
    
    The class's \inlineCode{__init__} constructor function is called whenever a class is instantiated (e.g. \inlineCode{control_law_object = ControlLawClass(specification,warning_levels,...)}).  Note that the the constructor function accepts as arguments not only the data available at the time (e.g. the specification and any extra control parameters) but also any parameters that will eventually exist.  This is because it needs to be able to seamlessly transition in case the control law is changed during control when there is already a transfer function, buzz \ac{CPSD}, etc.
    
    The \inlineCode{system_id_function} is called after the system identification is complete, so inside this function is where all setup calculations that require a transfer function or buzz \ac{CPSD} would go.  In the case of the buzz test approach currently under consideration, this function is where the modified specification would be computed.
    
    The \inlineCode{control} function is then the function that actually performs the control operations to compute the output \ac{CPSD} matrix using the updated transfer functions or last response or output \ac{CPSD} matrices in addition to any data that had been stored inside the class.
    
    The class implementation of the buzz test control is shown in Listing \ref{lst:buzz_test_class}.  Note how all the helper functions can be stored directly within the class.
    
    \begin{lstlisting}[language=Python,caption = {Class implementation of the buzz test approach},label=lst:buzz_test_class]
import numpy as np

class buzz_control_class:
    def __init__(
            self,
            specification : np.ndarray, # Specifications
            warning_levels  : np.ndarray, # Warning levels
            abort_levels  : np.ndarray, # Abort Levels
            extra_parameters : str, # Extra parameters for the control law
            transfer_function : np.ndarray = None,  # Transfer Functions
            noise_response_cpsd : np.ndarray = None,  # Noise levels and correlation 
            noise_reference_cpsd : np.ndarray = None, # from the system identification
            sysid_response_cpsd : np.ndarray = None,  # Response levels and correlation
            sysid_reference_cpsd : np.ndarray = None, # from the system identification
            multiple_coherence : np.ndarray = None, # Coherence from the system identification
            frames = None, # Number of frames in the CPSD and FRF matrices
            total_frames = None, # Total frames that could be in the CPSD and FRF matrices
            last_response_cpsd : np.ndarray = None, # Last Control Response for Error Correction
            last_output_cpsd : np.ndarray = None, # Last Control Excitation for Drive-based control
            ):
        # Store the specification to the class
        if sysid_response_cpsd is None: # If it's the first time through we won't have a buzz test yet
            self.specification = specification
        else: # Otherwise we can compute the modified spec right away
            self.specification = self.match_coherence_phase(specification, sysid_response_cpsd)
            
    def system_id_update(
            self,
            transfer_function : np.ndarray = None,  # Transfer Functions
            noise_response_cpsd : np.ndarray = None,  # Noise levels and correlation 
            noise_reference_cpsd : np.ndarray = None, # from the system identification
            sysid_response_cpsd : np.ndarray = None,  # Response levels and correlation
            sysid_reference_cpsd : np.ndarray = None, # from the system identification
            multiple_coherence : np.ndarray = None, # Coherence from the system identification
            frames = None, # Number of frames in the CPSD and FRF matrices
            total_frames = None, # Total frames that could be in the CPSD and FRF matrices
            ):
        # Update the specification with the buzz_cpsd
        self.specification = self.match_coherence_phase(self.specification,sysid_response_cpsd)

    def control(
            self,
            transfer_function : np.ndarray = None,  # Transfer Functions
            multiple_coherence : np.ndarray = None, # Coherence from the system identification
            frames = None, # Number of frames in the CPSD and FRF matrices
            total_frames = None, # Total frames that could be in the CPSD and FRF matrices
            last_response_cpsd : np.ndarray = None, # Last Control Response for Error Correction
            last_output_cpsd : np.ndarray = None) -> np.ndarray:
        # Perform the control
        tf_pinv = np.linalg.pinv(transfer_function)
        return tf_pinv @ self.specification @ tf_pinv.conjugate().transpose(0,2,1)
        
    def cpsd_coherence(self,cpsd):
        num = np.abs(cpsd)**2
        den = (cpsd[:,np.newaxis,np.arange(cpsd.shape[1]),np.arange(cpsd.shape[2])]*
               cpsd[:,np.arange(cpsd.shape[1]),np.arange(cpsd.shape[2]),np.newaxis])
        den[den==0.0] = 1 # Set to 1
        return np.real(num/
                       den)
    
    def cpsd_phase(self,cpsd):
        return np.angle(cpsd)
    
    def cpsd_from_coh_phs(self,asd,coh,phs):
        return np.exp(phs*1j)*np.sqrt(coh*asd[:,:,np.newaxis]*asd[:,np.newaxis,:])
    
    def cpsd_autospectra(self,cpsd):
        return np.einsum('ijj->ij',cpsd)
    
    def match_coherence_phase(self,cpsd_original,cpsd_to_match):
        coh = self.cpsd_coherence(cpsd_to_match)
        phs = self.cpsd_phase(cpsd_to_match)
        asd = self.cpsd_autospectra(cpsd_original)
        return self.cpsd_from_coh_phs(asd,coh,phs)\end{lstlisting}
        
    Note that the number of lines of code for a class implementation of the buzz test approach is not significantly more than the simpler function implementation, therefore users should not immediately discard the class-based approach as too difficult to implement in favor of the simpler function implementation.  Each approach has its merits and limitations, so it is up to the user to decide the best approach for their control law.
    
    \section{Using Transformation Matrices}\label{sec:rattlesnake_environments_transformation_matrices}
    
    The \Ac{MIMO} Random Vibration and Transient (see Section \ref{sec:rattlesnake_environments_mimo_transient}) environments allow the use of transformation matrices to constrain or transform response measurements or output signals into more favorable degrees of freedom.  The basic usage of the transformation matrix is
    \begin{equation}
        \hat{\mathbf{x}} = \mathbf{T}\mathbf{x}
    \end{equation}
    where $\mathbf{x}$ is the measurement or signal in the physical degrees of freedom, $\mathbf{T}$ is the transformation matrix, and $\hat{\mathbf{x}}$ is the transformed quantity.
    
    By this definition, a response transformation must have the same number of columns as there are control channels in the environment.  Similarly, an output transformation must have the same number of columns as excitation signals in the environment.  The number of rows in these matrices will then be the number of virtual control degrees of freedom or virtual excitation signals that will be used by the environment.
    
    A common application of transformation matrices is for so-called \ac{6DoF} testing, where shakers are constrained to excite the six rigid body motions of a rigid table.  A representative \ac{6DoF} configuration is shown in Figure \ref{fig:6dof} with 12 shakers exciting a rigid table on which a test article is mounted.  To measure the table response, four triaxial gauges are positioned symmetrically across the table.

    \begin{figure}
        \centering
        \includegraphics[width=\linewidth]{figures/6dof}
        \caption[Representative \ac{6DoF} setup]{Representative \ac{6DoF} setup showing 12 shakers attached to a rigid table with four triaxial accelerometers measuring the table's response. Note that shaker 11 is occluded by the table in this figure.}
        \label{fig:6dof}
    \end{figure}

    To set up the output transformation, one should examine the geometry of the test setup to identify which shakers excite which motions.  For example, to move the table vertically in the Y+ direction, the four bottom shakers (9, 10, 11, and 12) should excite with a positive signal (here, a positive signal is assumed to be compressive, pushing the table away from the shaker; this may vary depending on shaker wiring).  For this case, the row of the output transformation matrix should contain a 1 for shaker signals 9, 10, 11, and 12, and a zero for other shakers.  Similar reasoning can be used for rotations: to rotate about the positive Z+ direction, shakers 11 and 12 should push with a positive signal and shakers 9 and 10 should pull with a negative signal.  Therefore, the row corresponding to this degree of freedom would have -1 for shakers 9 and 10 and +1 for 11 and 12.  The full output transformation for this case can be seen in Table \ref{tab:6DoF_Output_Transform}.  This matrix transforms the 12 physical shaker signals into 6 rigid body motions of the table.

    \begin{table}
        \centering
        \caption{Output transformation for the \ac{6DoF} test shown in Figure \ref{fig:6dof}}
        \label{tab:6DoF_Output_Transform}
        \begin{tabular}{c|cccccccccccc}
            Shaker & 1 & 2 & 3 & 4 & 5 & 6 & 7 & 8 & 9 & 10 & 11 & 12 \\
            \hline
            DX+ & 1 & 1 & 0 & 0 & -1 & -1 & 0 & 0 & 0 & 0 & 0 & 0 \\
            DY+ & 0 & 0 & 0 & 0 & 0 & 0 & 0 & 0 & 1 & 1 & 1 & 1 \\
            DZ+ & 0 & 0 & 1 & 1 & 0 & 0 & -1 & -1 & 0 & 0 & 0 & 0 \\
            RX+ & 0 & 0 & 0 & 0 & 0 & 0 & 0 & 0 & -1 & 1 & 1 & -1 \\
            RY+ & 1 & -1 & 1 & -1 & 1 & -1 & 1 & -1 & 0 & 0 & 0 & 0 \\
            RZ+ & 0 & 0 & 0 & 0 & 0 & 0 & 0 & 0 & -1 & -1 & 1 & 1 \\
        \end{tabular}
    \end{table}

    To generate the response transformation matrix, it is often helpful to construct the inverse transformation and then invert it to recover the response transformation matrix.  As an example, we will investigate the measured signals by the accelerometers if the table is translated one unit in the X+ direction.  In this case, all accelerometers pointing in the X+ direction would see a one unit signal.  Therefore the column of the inverse transformation corresponding to a X+ direction motion would have 1 for all channels pointing in that direction.  If the table is rotated about the Y+ direction, the 1X+, 1Z+, 2X+, and 4Z+ channels will see positive motion, and the 2Z+, 3X+, 3Z+ and 4X+ channels will see negative motions, so the rows corresponding to those channels should be populated with +1 or -1, respectively.  Note that there technically should be a moment arm computed by the distance from the accelerometer to the center-line of the table in the real rotation calculation; however, if all accelerometers are equidistant, this term can be dropped, and the system identification will compensate accordingly.  Table \ref{tab:6DoF_Response_Transform_inverse} shows the constructed inverse transformation, and Table \ref{tab:6DoF_Response_Transform} shows the response transform matrix that should be provided to Rattlesnake.
    
    
    \begin{table}
        \centering
        \caption{Inverse response transformation for the \ac{6DoF} test shown in Figure \ref{fig:6dof}}
        \label{tab:6DoF_Response_Transform_inverse}
        \begin{tabular}{c|cccccc}
            Channel & DX+ & DY+ & DZ+ & RX+ & RY+ & RZ+ \\
            \hline
            1X+ & 1 & 0 & 0 & 0 & 1 & 0 \\
            1Y+ & 0 & 1 & 0 & -1 & 0 & -1 \\
            1Z+ & 0 & 0 & 1 & 0 & 1 & 0 \\
            2X+ & 1 & 0 & 0 & 0 & 1 & 0 \\
            2Y+ & 0 & 1 & 0 & -1 & 0 & 1 \\
            2Z+ & 0 & 0 & 1 & 0 & -1 & 0 \\
            3X+ & 1 & 0 & 0 & 0 & -1 & 0 \\
            3Y+ & 0 & 1 & 0 & 1 & 0 & 1 \\
            3Z+ & 0 & 0 & 1 & 0 & -1 & 0 \\
            4X+ & 1 & 0 & 0 & 0 & -1 & 0 \\
            4Y+ & 0 & 1 & 0 & 1 & 0 & -1 \\
            4Z+ & 0 & 0 & 1 & 0 & 1 & 0 \\
            \hline
        \end{tabular}
    \end{table}    
    
    \begin{table}
        \centering
        \caption{Response transformation for the \ac{6DoF} test shown in Figure \ref{fig:6dof}}
        \label{tab:6DoF_Response_Transform}
        \resizebox{\linewidth}{!}{
        \begin{tabular}{c|cccccccccccc}
            Channel &1X+ & 1Y+ & 1Z+ & 2X+ & 2Y+ & 2Z+ & 3X+ & 3Y+ & 3Z+ & 4X+ & 4Y+ & 4Z+ \\
            \hline
            DX+ & 0.250 & 0.000 & 0.000 & 0.250 & 0.000 & 0.000 & 0.250 & 0.000 & 0.000 & 0.250 & 0.000 & 0.000 \\
            DY+ & 0.000 & 0.250 & 0.000 & 0.000 & 0.250 & 0.000 & 0.000 & 0.250 & 0.000 & 0.000 & 0.250 & 0.000 \\
            DZ+ & 0.000 & 0.000 & 0.250 & 0.000 & 0.000 & 0.250 & 0.000 & 0.000 & 0.250 & 0.000 & 0.000 & 0.250 \\
            RX+ & 0.000 & -0.250 & 0.000 & 0.000 & -0.250 & 0.000 & 0.000 & 0.250 & 0.000 & 0.000 & 0.250 & 0.000 \\
            RY+ & 0.125 & -0.000 & 0.125 & 0.125 & 0.000 & -0.125 & -0.125 & -0.000 & -0.125 & -0.125 & -0.000 & 0.125 \\
            RZ+ & -0.000 & -0.250 & -0.000 & 0.000 & 0.250 & -0.000 & 0.000 & 0.250 & 0.000 & 0.000 & -0.250 & 0.000 \\
        \end{tabular}
    }
    \end{table}

    Note that when a response transformation matrix is used, the vibration specification must be delivered to Rattlesnake in terms of the transformed coordinates.  Note that the specification in terms of physical measurement can be transformed into the virtual quantities in the same way that the physical quantities themselves are transformed.
    \begin{equation}
        \hat{\mathbf{G}}_{xx} = \mathbf{T}\mathbf{G}_{xx}\mathbf{T}^H
    \end{equation}

    Transformation matrices can be used for more general coordinate system transformations.  Perhaps the vibration specification has been derived from a finite element model, but the test article has instrumentation mounted obliquely to the global finite element coordinate system.  A transformation matrix consisting a rotation matrices can be used to transform between coordinate systems in the model and coordinate systems in the test.  Additionally, transformation matrices can be used to perform spatial filtering for example to control to specific modes of the structure $\mathbf{q}$.  Given the modal transformation
    \begin{equation}
        \mathbf{x} = \mathbf{\Phi}\mathbf{q}
    \end{equation}
    the modal quantities $\mathbf{q}$ can be computed from physical quantities $\mathbf{x}$ as 
    \begin{equation}
        \mathbf{q} = \mathbf{\Phi}^+\mathbf{x}
    \end{equation}
    In this case, the transformation matrix is simply the inverse mode shape matrix $\mathbf{\Phi}$, which acts as a modal filter.

    \section{Generation of Time Histories}\label{sec:rattlesnake_environments_generating_time_histories}
    Rattlesnake's \Ac{MIMO} Random Vibration environment allows users to run custom control laws that produce \ac{CPSD} matrices.  Rattlesnake must then create time histories from those \ac{CPSD} matrices to send to the shakers.  This section will discuss some of the implementation details of this process that may be useful for users to understand.
    
    \subsection{Creating Time Signal Realizations from \ac{CPSD} Matrices}
    To start the signal generation process, the signal generation routine receives an output \ac{CPSD} matrix from the custom control law provided by the user.  The controller uses the process described in \cite{schultz2020_input_signal_synthesis_openloop_multipleinput_multipleoutput_testing} to create a time history realization from the \ac{CPSD} matrix.  Note that for large control problems, it might take more than one measurement frame to perform the control problem, so there might not be an output \ac{CPSD} matrix ready when the controller requires it.  In that case, the previous output \ac{CPSD} matrix is used.
    
    To generate a time history realization, the square root of the \ac{CPSD} matrix is taken using a \ac{SVD} approach.  This creates linear spectra $\mathbf{L}$ from the power spectra.
    \begin{equation}
        \mathbf{G}_{vv} = \mathbf{U}\mathbf{S}\mathbf{V}^H
    \end{equation}
    \begin{equation}
        \mathbf{L} = \mathbf{U}\sqrt{\mathbf{S}}\mathbf{V}^H
    \end{equation}

    A random process vector is computed using Gaussian real and imaginary components $\mathbf{A}_k$ and $\mathbf{B}_k$.
    \begin{equation}
        \mathbf{\Psi} = \frac{1}{{df}\sqrt{2}}(\mathbf{A}_k+i\mathbf{B}_k)
    \end{equation}
    where $df$ is the frequency spacing and $i$ is the complex variable.
    
    A realization of the linear spectra is computed by multiplying the random process vector by the linear spectra
    \begin{equation}
        \mathbf{X}_v = \mathbf{L}\mathbf{\Psi}
    \end{equation}

    Taking the inverse \ac{FFT} then transforms the linear spectra into time signals that can be sent to the shakers.
    \begin{equation}
        \mathbf{x}_v = \mathcal{F}^{-1}{\mathbf{X}_v}
    \end{equation}

    Note that this signal has the length of one \ac{FFT} frame, which is generally not long enough to be immediately useful for vibration testing.  Multiple signals can be generated by taking additional realizations of the random process vector; however, simply concatenating multiple realizations will result in discontinuous signals, as there is no guarantee that these signals will start and end at the same levels.
    
    \subsection{Synthesizing Continuous Time Histories using \ac{COLA}}
    To handle the synthesis of longer signals, a \ac{COLA} process is used.  This process is also described in detail in \cite{schultz2020_input_signal_synthesis_openloop_multipleinput_multipleoutput_testing}.  A \ac{COLA} window function is specified that allows each time history realization to start and end at zero.  To ensure the proper variance is maintained, the square root of the window function is applied to the signal (specified using the \inlineCode{Window Exponent} parameter in the \inlineCode{Environment Definition} tab).  The overlap is specified such that the original (un-square-rooted) window function sums to one over the signal.  The windowed signals are then added together to form a continuous signal.
    
    Rattlesnake cannot realistically compute a large number of realizations all at once to create one long signal, as this would impede its ability to respond to changing test conditions, effectively making it an open-loop control system.  Instead, Rattlesnake must compute this \ac{COLA} operation on the fly.  It does this by keeping two output signals in memory at once.  It overlaps the last part of the first signal with the first part of the next signal and sends that overlapped portion of the signal to the output of the controller.  Then, the previous next signal becomes the next first signal and a new realization replaces the previous next signal and the process continues.
    
    \subsection{Setting the Test Level}
    The final operation performed on the output signal is to adjust for test level.  Initial versions of the Rattlesnake simply scaled the vibration specification so the controller would naturally scale its output to match the scaled specification.  This was found to be lacking especially for large control problems where it could take a second or more to perform a control calculation.  This combined with the \ac{COLA} approach meant that it could take several seconds for a test level change to be realized, which was more than enough time to damage test hardware.  Another issue with this approach was for error-based control approaches.  The controller would immediately recognize that it was in error due to the change in test level and produce a modified output.  However, there could be 20 or more frames of data being used in the computation of \ac{CPSD} matrices that would need to be overwritten prior to the controller recovering.  In practice, this meant that the controller would severely overshoot the desired level as it further and further modified the output with very little change in the response \ac{CPSD}.  One final issue is that it relied on the \ac{COLA} process to smooth jumps between test levels.  As the \ac{COLA} overlap regions could be quite short, this lead to fairly severe level changes that could damage test hardware.
    
    Recent versions of the Rattlesnake have moved away from scaling the vibration specification and started scaling the output signal directly.  The acquired data is then scaled back to full level prior to control calculations taking place.  For this reason, all information written to the \inlineCode{Run Test} tab in the \ac{GUI} is presented as if it were at full level.  Because the time histories are modified directly, it allows ramping between test levels to be implemented directly and at whatever speed is required for hardware safety.  Data acquired during the ramp between test levels is generally ignored; only data acquired while the test level is constant is used for spectral computations.
    
    \chapter{Multiple Input/Multiple Output Transient Control}\label{sec:rattlesnake_environments_mimo_transient}
    
    The \Ac{MIMO} Transient environment aims to control the vibration response of a component to a specific time history by creating output signals with the correct levels and phasing.  The governing equation for \Ac{MIMO} Transient control is 
    
    \begin{equation}\label{eq:forward_mimo_transient}
    \mathbf{X} = {\mathbf{H}_{xv}}\mathbf{V}
    \end{equation}
    
    where the frequency spectrum matrix of the responses $\mathbf{X}$ result from some signals $\mathbf{V}$ exciting the structure represented by transfer function matrices $\mathbf{H}_{xv}$.  In a typical transient control problem, the control system tries to compute the a signal with spectrum matrix $\mathbf{V}$ that best reproduces the desired response with frequency spectrum $\mathbf{X}$.
    
    \section{Signal Definition}
    The first step in defining a transient control problem is the definition of the response signal that is desired.  Rattlesnake accepts the specification in the form of a 2D array consisting of a time response of one or more channels in the test.  Signals can be loaded from Numpy \inlineCode{*.npy} or \inlineCode{*.npz} files or Matlab \inlineCode{*.mat} files.  Both Matlab \inlineCode{*.mat} and Numpy \inlineCode{*.npz} files should contain the following data members:
    
    \begin{description}
           \item[signal] A $n_c \times n_s$ array containing the time signal the controller will try to reproduce on the test article.  For \inlineCode{*.npy} files which do not have fields, the signal array is stored directly to the file.
           \item[t] A $n_s$ array of times corresponding to columns of the \inlineCode{signal} field.  If not specified (either by not including a \inlineCode{t} field in a \inlineCode{*.mat} or \inlineCode{*.npz} file or by using an \inlineCode{*.npy} file), the \inlineCode{signal} will be assumed to be at the sample rate defined in the controller.  If \inlineCode{t} is specified but is not at the sample rate of the controller, the signal will be linearly interpolated to be at the sample rate of the controller.
    \end{description}
      
    The ordering of the rows of the array defining the signal is the same order as the control channels in the Channel Table on the \inlineCode{Data Acquisition Setup} tab that are selected as control channels on the \inlineCode{Environment Definition} tab.  The specification is defined in the engineering units specified by the \inlineCode{Engineering Unit} column of the channel table for the control channels.
    
    \section{Defining the \Ac{MIMO} Transient Environment in Rattlesnake}
    
    In addition to the specification, there are a number of signal processing parameters that are used by the \Ac{MIMO} Transient environment.  These, along with the specification, are defined on the \inlineCode{Environment Definition} tab in the Rattlesnake controller on a sub-tab corresponding to a \Ac{MIMO} Transient environment.  Figure \ref{fig:transientenvironmentdefinition} shows a \Ac{MIMO} Transient sub-tab.  The following subsections describe the parameters that can be specified, as well as their effects on the analysis.
    
    \begin{figure}
           \centering
           \includegraphics[width=\linewidth]{figures/transient_environment_definition}
           \caption{\Ac{GUI} used to define a \Ac{MIMO} Transient environment.}
           \label{fig:transientenvironmentdefinition}
    \end{figure}

    \subsection{Signal Parameters}
    The \inlineCode{Signal Parameters Parameters} section of the \Ac{MIMO} Transient definition sub-tab consists of the following parameters:
    \begin{description}
           \item[Sample Rate] The global sample rate of the data acquisition system.  This is set on the \inlineCode{Data Acquisition Setup} tab, and displayed here for convenience as a read-only value.
           \item[Signal Time] The total time it will take to play the signal that the controller will attempt to reproduce on the test article.
           \item[Signal Samples] The number of time steps in the signal that the controller will attempt to reproduce on the test article.
           \item[Ramp Time] The time taken to ramp the signal to zero when the signal is stopped, preventing ``Hard Stops'' that could damage equipment if the test is aborted midway through.
    \end{description}
    
    \subsection{Control Channels}
    The \inlineCode{Control Channels} list allows users to select the channels in the test that will be used by the environment to perform control calculations.  These are the channels that will match the rows and columns of the specification file.
    
    \subsection{Control Parameters}
    The \inlineCode{Control Parameters} section of the \Ac{MIMO} Transient definition sub-tab consists of the following parameters:
    \begin{description}
           \item[Input Channels] The total number of channels being measured by the Rattlesnake, including response channels and output channels.  This is a computed quantity presented for convenience, so the user cannot modify it directly.
           
           \item[Output Channels] The number of excitation signals being used to control the vibration response of this environment.  This is a computed quantity presented for convenience, so the user cannot modify it directly.
           
           \item[Control Channels] The number of response channels being used in the control.  This is a computed quantity presented for convenience, so the user cannot modify it directly.

    \end{description}
       The \inlineCode{Control Parameters} section of the \Ac{MIMO} Transient definition sub-tab also includes functionality for loading in custom control laws.  See Section \ref{sec:rattlesnake_environments_custom_control_law_transient} for information on defining a custom transient control law.
       \begin{description}
           \item[Control Python Script] The Python script containing a custom control law that is currently loaded into the Rattlesnake controller.
           \item[Load] Pressing this button will bring up a file selection dialog to load in a new Python script containing a custom control law.
           \item[Control Python Function] This selector presents the functions, generators, and classes within the loaded Python script that can be used as custom control laws.
           \item[Control Parameters] This text box allows arbitrary input to be passed to custom control laws as a string.  It is up to the control law to specify what this extra input must be and parse whatever input the user gives it.
       \end{description}

    \subsection{Control and Drive Transforms}
    The \inlineCode{Control and Drive Transforms} section of the \Ac{MIMO} Random Vibration definition sub-tab consists of the following parameters:
    \begin{description}
        \item[Transformation Matrices...] Selecting this button will bring up the transformation matrices dialog box, which allows the user to specify linear transformations between the physical responses and excitation signals and virtual responses and excitation signals.  See Section \ref{sec:rattlesnake_environments_transformation_matrices} for more information.
        \item[Transform Channels] The number of virtual control degrees of freedom after applying transformation matrices.  This is a computed quantity presented for convenience, so the user cannot modify it directly.
        \item[Transform Outputs] The number of virtual excitation devices after applying transformation matrices.  This is a computed quantity presented for convenience, so the user cannot modify it directly.
    \end{description}
    Note that if Transformation matrices are defined, the number of control channels ends up being the number of rows of the \inlineCode{Response Transformation Matrix}, rather than the number of physical control channels.  The number of physical control channels will be equal to the number of columns of the transformation matrix.  The number of rows of the specification loaded should be equal to the number of rows in the transformation.

   \subsection{Signal Specification}
   The control signal is loaded and displayed on the right side of the \Ac{MIMO} Transient definition tab:
   
   \begin{description}
       \item[Load Signal] A button that when pressed will bring up a file selection dialog to read in a specification from a Numpy \inlineCode{*.npz} or \inlineCode{*.npy} or Matlab \inlineCode{*.mat} file.
       \item[Signal Plot] A visualization of the signals loaded into the controller.
       \item[Signal Table] A table showing statistics of each of the signals loaded into the controller.  The user can also select which signals to plot using the checkboxes in the table.
   \end{description}
    
    \section{System Identification for the \Ac{MIMO} Transient Environment}
    
    When all environments are defined and the \inlineCode{Initialize Environments} button is pressed, Rattlesnake will proceed to the next phase of the test, which is defined on the \inlineCode{System Identification} tab.
    
    \Ac{MIMO} Transient requires a system identification phase to compute the matrix $\mathbf{H}_{xv}$ used in the control calculations of equation \eqref{eq:forward_mimo_vibration}.  Unlike the \Ac{MIMO} Random Vibration environment, the \Ac{MIMO} Transient environment's system identification phase can have a number of samples per measurement frame that is not equal to the length of the signal.  Control laws written for this environment must be able to handle this, either by doing a \ac{COLA} scheme or interpolating the \Ac{FRF}.  Figure \ref{fig:transientenvironmentsystemidentification} shows the \ac{GUI} used to perform this phase of the test.  Given the nature of the \Ac{MIMO} Transient environment, it may be useful to visualize impulse response functions as the system identification is proceeding.
    
    \begin{figure}
           \centering
           \includegraphics[width=\linewidth]{figures/transient_environment_system_identification}
           \caption{System identification \ac{GUI} used by the \Ac{MIMO} Transient environment.}
           \label{fig:transientenvironmentsystemidentification}
    \end{figure}
    
    Rattlesnake's system identification phase will start with a noise floor check, where the data acquisition records data on all the channels without specifying an output signal.  After the noise floor is computed, the system identification phase will play out the specified signals to the excitation devices, and transfer functions will be computed using the responses of the control channels to those excitation signals.  Section \ref{sec:using_rattlesnake_system_identification} describes the System Identification tab and its various parameters and capabilities.
    
    \section{Test Predictions for the \Ac{MIMO} Transient Environment}
    
    Once the system identification is performed, a test prediction will be performed and results displayed on the \inlineCode{Test Predictions} tab, shown in Figure \ref{fig:transientenvironmenttestprediction}.  This is meant to give the user an idea of the test feasibility.  The top portion of the window displays excitation information, including peak signal levels required as well as the excitation time history that will be output.  The bottom portion of the window displays the predicted responses compared to the specification as well as the \ac{TRAC} between the signals.  The \ac{TRAC} is a metric that compares two time signals, and has a value of 1 if the signals are identical or 0 if the signals are not related.
    
    \begin{figure}
        \centering
        \includegraphics[width=\linewidth]{figures/transient_environment_test_prediction}
        \caption{Test prediction \ac{GUI} which gives the user some idea of the test feasibility.}
        \label{fig:transientenvironmenttestprediction}
    \end{figure}
    
    \section{Running the \Ac{MIMO} Transient Environment}
    The \Ac{MIMO} Transient environment is then run on the \inlineCode{Run Test} tab of the controller.
    
    With the data acquisition system armed, the environment can be started manually with the \inlineCode{Start Environment} button.  Once running, it can be stopped manually with the \inlineCode{Stop Environment} button.  With the data acquisition system armed and the environment run, the \ac{GUI} looks like Figure \ref{fig:transientenvironmentruntest}.
    
    \begin{figure}
        \centering
        \includegraphics[width=\linewidth]{figures/transient_environment_run_test}
        \caption{\Ac{GUI} for running the \Ac{MIMO} Transient environment.}
        \label{fig:transientenvironmentruntest}
    \end{figure}
    
    There are operations that can be performed when setting up and running the \Ac{MIMO} Transient environment, and many visualization operations as well.
    
    \subsection{Test Level}
    The \Ac{MIMO} Transient Environment allows scaling of the test by modifying the \inlineCode{Signal Level} selector on the \inlineCode{Run Test} tab.  The \inlineCode{Signal} specifies the scaling in decibels relative to the specification level, which is 0 dB.  Note that all data and visualizations on the \inlineCode{Run Test} window are scaled back to full level, so users should not be surprised if for example the signals shown in the \inlineCode{Outputs} or \inlineCode{Responses} plots do not change significantly with test level.  Unlike the \ac{MIMO} Random Vibration environment, the \Ac{MIMO} Transient environment does not allow changing test level while the environment is active.
    
    \subsection{Repeating a Signal}
    The \Ac{MIMO} Transient environment also has the capability to repeat a signal continuously if the \inlineCode{Repeat Signal} checkbox is checked.  If this is the case, users will be required to stop the environment by clicking the \inlineCode{Stop Environment} button.  If the controller is not set to repeat, the environment will stop automatically after the signal has been output one time.
    
    \subsection{Test Metrics and Visualizations}
    The \Ac{MIMO} Transient environment displays the output signals in the \inlineCode{Outputs} plot and the responses to those signals at the control channels in the \inlineCode{Responses} plot.  The controller will identify the position of the control signal in the response signal by using a correlation, and will draw a black box around that portion of the response signal, as shown in Figure \ref{fig:transientenvironmentruntest}.  The errors in the response channels are reported as a \Ac{TRAC} value between the signal measured and the specification desired.

    To interrogate specific channels, the \inlineCode{Data Display} section of the \inlineCode{Run Test} can be used.  The specific control channel to visualize can be selected using \inlineCode{Control Channel} selector.  Pressing the \inlineCode{Create Window} button then creates the specified plot.
    
    Some convenience operations are also included to visualize all channels.  The \inlineCode{Show All Channels} button will bring up one window per control channel.  Be aware that showing a large number of channels can easily overwhelm the computer with plotting operations causing the \ac{GUI} to become unresponsive, so use this operations with caution.  Figure \ref{fig:transientenvironmentchannelvisualizations} shows an example displaying all channels for a test with six control degrees of freedom. 
    
    \begin{figure}
        \centering
        \includegraphics[width=\linewidth]{figures/transient_environment_channel_visualizations}
        \caption{Visualizing individual channels' time response.}
        \label{fig:transientenvironmentchannelvisualizations}
    \end{figure}
    
    Further convenience operations are available in the \inlineCode{Window Operations:} section.  Pressing \inlineCode{Tile All Windows} will rearrange all channel windows neatly across the screen.  Pressing \inlineCode{Close All Windows} will close all open channel windows.

    \section{Output NetCDF File Structure}
    When Rattlesnake saves data to a netCDF file, environment-specific parameters are stored in a netCDF group with the same name as the environment name.  Similar to the root netCDF structure described in Section \ref{sec:using_rattlesnake_output_files}, this group will have its own attributes, dimensions, and variables, which are described here.
    
    \subsection{NetCDF Dimensions}
    \begin{description}
        \item[signal\_samples] The number of time samples in the specification signal provided to the \Ac{MIMO} Transient environment.
       \item[specification\_channels] The number of channels in the specification signal provided to the \Ac{MIMO} Transient environment.
       \item[control\_channels] The number of physical channels used for control.  Note that this may be different from the \inlineCode{specification_channels} due to the presence of a transformation matrix.
       \item[response\_transformation\_rows] The number of rows in the response channel transformation.  This is not defined if no response transformation is used.
       \item[response\_transformation\_cols] The number of columns in the response channel transformation.  This is not defined if no response transformation is used.
       \item[output\_transformation\_rows] The number of rows in the output transformation.  This is not defined if no output transformation is used.
       \item[output\_transformation\_cols] The number of columns in the output transformation.  This is not defined if no output transformation is used.
    \end{description}
    
    \subsection{NetCDF Attributes}
    \begin{description}
        \item[sysid\_frame\_size] The number of samples per measurement frame in the system identification
        \item[sysid\_averaging\_type] The type of averaging used in the system identification, linear or exponential
        \item[sysid\_noise\_averages] The number of measurement frames acquired for the noise floor calculation
        \item[sysid\_averages] The number of measurement frames acquired for the system identification calculation
        \item[sysid\_exponential\_averaging\_coefficient] The weighting coefficient used for new frames in the exponential averaging scheme
        \item[sysid\_estimator] The \Ac{FRF} estimator used to compute the transfer functions during the system identification
        \item[sysid\_level] The level used by the system identification in volts RMS.
        \item[sysid\_level\_ramp\_time] The time to ramp up to the test level when starting and ramp back to zero when stopping the system identification
        \item[sysid\_signal\_type] The signal type used by the system identification
        \item[sysid\_window] The window function applied to the time data during the system identification
        \item[sysid\_overlap] The overlap fraction between measurement frames used for system identification
        \item[sysid\_burst\_on] The fraction of a measurement frame that a burst is active for burst random excitation during system identification
        \item[sysid\_pretrigger] The fraction of a measurement used as a pre-trigger for burst random excitation during system identification
        \item[sysid\_burst\_ramp\_fraction] The fraction of a measurement frame used to ramp the burst up to full level and back to zero
        \item[test\_level\_ramp\_time] The time to ramp to the test level and back to zero
        \item[control\_python\_script] The path to the Python script used to control the \Ac{MIMO} Random Vibration environment
        \item[control\_python\_function] The function (or class or generator function) in the Python script used to control the \Ac{MIMO} Transient environment
        \item[control\_python\_function\_type] The type of the object used for the control law (function, generator, or class)
        \item[control\_python\_function\_parameters] The extra parameters passed to the control law.
       \end{description}
    
    \subsection{NetCDF Variables}
    \begin{description}
        \item[control\_signal] The control signal used by the \Ac{MIMO} Transient environment  Type: 64-bit float; Dimensions: \inlineCode{specification_channels} $\times$ \inlineCode{signal_samples}
        \item[response\_transformation\_matrix] The response transformation matrix.  This is not defined if no response transformation is used.  Type: 64-bit float; Dimensions: \inlineCode{response_transformation_rows} $\times$ \inlineCode{response_transformation_cols}
        \item[output\_transformation\_matrix] The output transformation matrix.  This is not defined if no output transformation is used.  Type: 64-bit float; Dimensions: \inlineCode{output_transformation_rows} $\times$ \inlineCode{output_transformation_cols}
        \item[control\_channel\_indices] The indices of the active control channels in the environment.  Type: 32-bit int; Dimensions: \inlineCode{control_channels}
       \end{description}
    
    \subsection{Saving Control Data}
    
    In addition to time streaming, Rattlesnake's \Ac{MIMO} Transient environment can also save the previous control data directly to the disk by clicking the \inlineCode{Save Current Control Data}.  The control data is stored in a NetCDF file similar to the time streaming data.  The primary difference is that the control data is time-aligned to the specification, so users don't need to worry about extracting the signal from an a time stream and trying to align it with specification.  The control data also saves out the information from the system identification, including \Ac{FRF} and \Ac{CPSD} matrices.
    
    The two additional dimensions are:
    
    \begin{description}
        \item[drive\_channels] The number of drive channels active in the environment.
        \item[fft\_lines] The number of frequency lines in the spectral quantities.
    \end{description}
    
    There are also several additional variables to store the spectral data:
    
    \begin{description}
        \item[frf\_data\_real] The real part of the most recently computed value for the transfer functions between the excitation signals and the control response signals.  NetCDF files cannot handle complex data types so real and imaginary parts are split into two variables.  Type: 64-bit float; Dimensions: \inlineCode{fft_lines} $\times$ \inlineCode{specification_channels} $\times$ \inlineCode{drive_channels}
        
        \item[frf\_data\_imag] The imaginary part of the most recently computed value for the transfer functions between the excitation signals and the control response signals.  NetCDF files cannot handle complex data types so real and imaginary parts are split into two variables.  Type: 64-bit float; Dimensions: \inlineCode{fft_lines} $\times$ \inlineCode{specification_channels} $\times$ \inlineCode{drive_channels}
        
        \item[frf\_coherence] The multiple coherence of the control channels computed during the test.  Type: 64-bit float; Dimensions: \inlineCode{fft_lines} $\times$ \inlineCode{specification_channels}
        
        \item[response\_cpsd\_real] The real part of the \ac{CPSD} matrix at the control channels from the system identification.  NetCDF files cannot handle complex data types so real and imaginary parts are split into two variables.  Type: 64-bit float; Dimensions: \inlineCode{fft_lines} $\times$ \inlineCode{specification_channels} $\times$ \inlineCode{specification_channels}
        
        \item[response\_cpsd\_imag] The imaginary part of the \ac{CPSD} matrix at the control channels from the system identification.  NetCDF files cannot handle complex data types so real and imaginary parts are split into two variables.  Type: 64-bit float; Dimensions: \inlineCode{fft_lines} $\times$ \inlineCode{specification_channels} $\times$ \inlineCode{specification_channels}
        
        \item[drive\_cpsd\_real] The real part of the \ac{CPSD} matrix at the excitation channels from the system identification.  NetCDF files cannot handle complex data types so real and imaginary parts are split into two variables.  Type: 64-bit float; Dimensions: \inlineCode{fft_lines} $\times$ \inlineCode{drive_channels} $\times$ \inlineCode{drive_channels}
        
        \item[drive\_cpsd\_imag] The imaginary part of the \ac{CPSD} matrix at the excitation channels from the system identifications.  NetCDF files cannot handle complex data types so real and imaginary parts are split into two variables.  Type: 64-bit float; Dimensions: \inlineCode{fft_lines} $\times$ \inlineCode{drive_channels} $\times$ \inlineCode{drive_channels}
        
        \item[response\_noise\_cpsd\_real] The real part of the \ac{CPSD} matrix at the control channels during the noise floor measurement that occurred during system identification.  NetCDF files cannot handle complex data types so real and imaginary parts are split into two variables.  Type: 64-bit float; Dimensions: \inlineCode{fft_lines} $\times$ \inlineCode{specification_channels} $\times$ \inlineCode{specification_channels}
        
        \item[response\_noise\_cpsd\_imag] The imaginary part of the \ac{CPSD} matrix at the control channels during the noise floor measurement that occurred during system identification.  NetCDF files cannot handle complex data types so real and imaginary parts are split into two variables.  Type: 64-bit float; Dimensions: \inlineCode{fft_lines} $\times$ \inlineCode{specification_channels} $\times$ \inlineCode{specification_channels}
        
        \item[drive\_noise\_cpsd\_real] The real part of the \ac{CPSD} matrix at the excitation channels during the noise floor measurement that occurred during system identification.  NetCDF files cannot handle complex data types so real and imaginary parts are split into two variables.  Type: 64-bit float; Dimensions: \inlineCode{fft_lines} $\times$ \inlineCode{drive_channels} $\times$ \inlineCode{drive_channels}
        
        \item[drive\_noise\_cpsd\_imag] The imaginary part of the \ac{CPSD} matrix at the excitation channels during the noise floor measurement that occurred during system identification.  NetCDF files cannot handle complex data types so real and imaginary parts are split into two variables.  Type: 64-bit float; Dimensions: \inlineCode{fft_lines} $\times$ \inlineCode{drive_channels} $\times$ \inlineCode{drive_channels}
        
        \item[control\_response] The time response of the system aligned in time to the specification. Type: 64-bit float; Dimensions: \inlineCode{specification_channels} $\times$ \inlineCode{signal_samples}
        
        \item[control\_drives] The drive signals aligned in time to the specification.  Type: 64-bit float; Dimensions: \inlineCode{drive_channels} $\times$ \inlineCode{signal_samples}
    \end{description}
    
    \section{Writing a Custom Control Law}\label{sec:rattlesnake_environments_custom_control_law_transient}
    The flexibility of the Rattlesnake framework is highlighted by the ease in which users can implement and iterate on their own ideas.  For the \Ac{MIMO} Transient control type, users can implement custom control laws using a custom Python function, or alternatively a generator function or class which allow state to be maintained between function calls.  This section will provide instructions and examples for implementing a custom control law.
    
    The controller will provide various data types to the control law functions which are:
    \begin{itemize}
        \item \inlineCode{sample_rate} -- The sample rate of the acquisition portion of the controller; real scalar integer
        
        \item \inlineCode{specification_signals} -- The target signal for the control channels; real 2D array ($n_c \times n_s$)
        
        \item \inlineCode{frequency_spacing} -- The frequency spacing of the transfer function array; real scalar float
        
        \item \inlineCode{transfer_function} -- The current estimate of the transfer function between the control responses and the output voltages; complex 3D array ($n_f \times n_c \times n_o$)
        
        \item \inlineCode{noise_response_cpsd} -- The levels and correlation of the noise floor measurement on the control channels obtained during system identification; complex 3D array ($n_{f} \times n_{c}\times n_{c}$)
        
        \item \inlineCode{noise_reference_cpsd} -- The levels and correlation of the noise floor measurement on the excitation channels obtained during system identification; complex 3D array ($n_{f} \times n_{o}\times n_{o}$)
        
        \item \inlineCode{sysid_response_cpsd} -- The levels and correlation of the control channels obtained during system identification; complex 3D array ($n_{f} \times n_{c}\times n_{c}$)
        
        \item \inlineCode{sysid_reference_cpsd} -- The levels and correlation of the noise floor measurement on the excitation channels obtained during system identification; complex 3D array ($n_{f} \times n_{o}\times n_{o}$)
        
        \item \inlineCode{multiple_coherence} -- The multiple coherence for each control channel; real 2D array ($n_f \times n_c$)
        
        \item \inlineCode{frames} -- The number of measurement frames acquired so far, used to compute various parameters in the control law.  This can be compared to \inlineCode{total_frames} to determine if a full set of measurement frames has been acquired, or if the estimation of the various parameters could improve with continued averaging; real scalar integer
        
        \item \inlineCode{total_frames} -- The total number of frames used to compute the \Ac{CPSD} and \Ac{FRF} matrices; real scalar integer
        
        \item \inlineCode{output_oversample_factor} -- Some hardware devices and settings within Rattlesnake require output signals from Rattlesnake to be oversampled compared to the input samples.  For example, on the LAN-XI hardware, the slowest output rate is four times larger than the slowest acquisition rate.  This argument tells the control law the factor by which the output returned from the controller is oversampled compared to the acquisition; real scalar integer
        
        \item \inlineCode{extra_parameters} -- Extra parameters provided to the controller; string
        
        \item \inlineCode{last_excitation_signals} -- The most recent output signal, which can be used for error-based control; real 2D array ($n_{o} \times n_{s}$)
        
        \item \inlineCode{last_response_signals} -- The most recent responses to the last output signals, which can be used for error-based control; real 2D array ($n_{c} \times n_{s}$)
    \end{itemize}
    
    where size $n_s$ is the number of time samples in the control signal, $n_f$ is the number of frequency lines in the transfer function matrix, $n_c$ is the number of control channels, and $n_o$ is the number of output signals.  Note that the values passed into the function may be defined using arbitrary variable names (e.g. \inlineCode{transfer_function} may be instead called \inlineCode{H}, or any other valid variable name); however, the order of the variables passed into each function will always be constant.

    \subsection{Defining a control law using a Python function}
    Python functions are the simplest approach to define a custom control law that can be used with the Rattlesnake software; however, they are limited in that a function's state is completely lost when a function returns.  Still, they can be used to implement relatively complex control laws as long as no state persistence is required.
    
    A Python function used to define a \Ac{MIMO} Transient control law in Rattlesnake would have the following general structure within a Python script.
    
    \begin{lstlisting}[language=Python,caption={General Python function structure for defining a custom Transient control law called \inlineCode{control_law} in Rattlesnake},label=lst:control_function_structure_transient]
# Any module imports, initialization code, or helper functions would go here

# Now we define the control law.  It always receives the same arguments from the controller.
def control_law(
    sample_rate,
    specification_signals, # Signal to try to reproduce
    frequency_spacing,
    transfer_function, # Transfer Functions
    noise_response_cpsd,  # Noise levels and correlation 
    noise_reference_cpsd, # from the system identification
    sysid_response_cpsd,  # Response levels and correlation
    sysid_reference_cpsd, # from the system identification
    multiple_coherence, # Coherence from the system identification
    frames, # Number of frames in the CPSD and FRF matrices
    total_frames, # Total frames that could be in the CPSD and FRF matrices
    output_oversample_factor, # Oversample factor to output
    extra_parameters = '', # Extra parameters for the control law
    last_excitation_signals = None, # Last excitation signal for drive-based control
    last_response_signals = None, # Last response signal for error correction
    ):

    # Code to perform the control would go here
    # output_signal = ...

    # Finally, we need to return an output signal matrix
    return output_signal\end{lstlisting}
        
    The function must return an \inlineCode{output_signal}, which is a 2D array with size ($n_{o} \times n_{os}$) where $n_{os}$ is the number of samples in the specification signal times the \inlineCode{output_oversample_factor}.
    
    One example is shown to demonstrate how a Transient control law may be written.
    
    \subsubsection{Pseudoinverse Control}
    Perhaps the simplest strategy to perform \ac{MIMO} Transient control is to simply invert the transfer function matrix to recover the least-squares solution of the optimal output signal from the desired responses.  This example will demonstrate that approach with some additional arguments that may improve the control.
    
    The mathematics for this control strategy are relatively simple; pre-multiply the spectra of the desired responses $\mathbf{X}$ by the pseudoinverse ($^+$) of the transfer function matrix $\mathbf{H}_{xv}$.  This calculation is performed for each frequency line.
    
    \begin{equation}
        \mathbf{V} = {\mathbf{H}_{xv}}^+\mathbf{X}
    \end{equation}
    
    The implementation details are a bit more complex than this formula would reveal, as the controller gives and wants returned the time signals and not the spectra.  This means the \ac{FFT} must be computed to transform to the frequency domain for the calculations to be performed, and the \ac{IFFT} must be computed to return to the time domain.  Additionally, to handle the potential for oversampling the output, the drive signal is zero-padded to up-sample the signal appropriately.  We must also handle the case where the number of samples in the system identification is not the same as the number of samples in the control signal.  We will finally add optional arguments that may be specified using key-value pairs separated by a colon (e.g. the \inlineCode{rcond} parameter could be specified as \inlineCode{rcond:1e-5} to specify a singular value threshold of 10,000 in the inverse).  In Python code, the above mathematics and implementation details would look like
    \begin{lstlisting}[language=Python,caption={Computing the pseudoinverse calculation to solve for a least-squares output spectrum},label=lst:pseudoinverse_computation_transient]
import numpy as np # Import numpy to get access to its function
# Parse the input arguments in extra_parameters
rcond = 1e-15
zero_impulse_after = None
# Split it up into lines
for entry in extra_parameters.split('\n'):
    try:
        # For each entry, split the key from the value using the colon
        field,value = entry.split(':')
        # Strip any whitespace
        field = field.strip()
        # Check the field to figure out which value to assign
        if field == 'rcond':
            rcond = float(value)
        elif field == 'zero_impulse_after':
            zero_impulse_after = float(value)
        else:
            # Report if we cannot understand the parameter
            print('Unrecognized Parameter: {:}, skipping...'.format(field))
    except ValueError:
        # Report if we cannot parse the line
        print('Unable to Parse Line {:}, skipping...'.format(entry))

# Compute impulse responses using the IFFT of the transfer function
# We will zero pad the IFFT to do interpolation in the frequency domain
# to match the length of the required signal
impulse_response = np.fft.irfft(transfer_function,axis=0)

# The impulse response should be going to zero at the end of the frame,
# but practically there may be some gibbs phenomenon effects that make the
# impulse response noncausal.  If we zero pad, this might be wrong.  We
# therefore give the use the ability to zero out this non-causal poriton of
# the impulse response.
if zero_impulse_after is not None:
    # Remove noncausal portion
    impulse_response_abscissa = np.arange(impulse_response.shape[0])/sample_rate
    zero_indices = impulse_response_abscissa > zero_impulse_after
    impulse_response[zero_indices] = 0
    
# Zero pad the impulse response to create a signal that is long enough for
# the specification signal
added_zeros = np.zeros((specification_signals.shape[-1]-impulse_response.shape[0],) 
                       + impulse_response.shape[1:])
full_impulse_response = np.concatenate((impulse_response,added_zeros),axis=0)

# Compute FRFs using the FFT from the impulse response.  This is now
# interpolated such that it matches the frequency spacing of the specification
# signal
interpolated_transfer_function = np.fft.rfft(full_impulse_response,axis=0)

# Perform convolution by frequency domain multiplication
signal_fft = np.fft.rfft(specification_signals,axis=-1)
# Invert the FRF matrix using the specified rcond parameter
inverted_frf = np.linalg.pinv(interpolated_transfer_function,rcond=rcond)
# Multiply the inverted FRFs by the response spectra to get the drive spectra
drive_signals_fft = np.einsum('ijk,ki->ij',inverted_frf,signal_fft)

# Zero pad the drive FFT to oversample to the output_oversample_factor
drive_signals_fft_zero_padded = np.concatenate((drive_signals_fft[:-1],
    np.zeros((drive_signals_fft[:-1].shape[0]*(output_oversample_factor-1)+1,)
             +drive_signals_fft.shape[1:])),axis=0)

# Finally, take the IFFT to get the time domain signal.  We need to scale
# by the output_oversample_factor due to how the IFFT is normalized.
drive_signals_oversampled = np.fft.irfft(
    drive_signals_fft_zero_padded.T,axis=-1)*output_oversample_factor
return drive_signals_oversampled\end{lstlisting}
    
    For users not familiar with Python and its numeric library \inlineCode{numpy}, the following points are clarified
    \begin{itemize}
        \item \inlineCode{numpy} is imported and assigned to the alias \inlineCode{np}, which lets us just type in \inlineCode{np} rather than the longer name \inlineCode{numpy} when we want to access \inlineCode{numpy} functions.
        \item To interpolate a function in the frequency domain, we can zero-pad its time response, and to interpolate a function in the time domain, we can zero-pad its frequency response.
        \item The \inlineCode{numpy} module has a function \inlineCode{zeros} to create an array of zeros and a function \inlineCode{concatenate} to combine two arrays together.  These two functions allow us to zero-pad an array.
        \item \Ac{FFT} calculations can be performed using the Fourier Transform package \inlineCode{fft} within \inlineCode{numpy}.  We will be using real signals, so we use \inlineCode{rfft} and \inlineCode{irfft} to do these transforms.
        \item The \inlineCode{numpy} pseudoinverse function \inlineCode{pinv} is stored in the linear algebra package \inlineCode{linalg} within \inlineCode{numpy}, therefore to access \inlineCode{pinv}, we need to call \inlineCode{np.linalg.pinv}
        \item The \inlineCode{pinv} can perform a pseudoinverse on "stacks" of matrices, so even though we are only calling the \inlineCode{pinv} function once, it is actually performing the pseudoinverse over all frequency lines
        \item The \inlineCode{numpy} \inlineCode{einsum} function utilizes a syntax similar to Einstein Summation Notation to perform multiplication and summation over certain dimensions of the arrays.
    \end{itemize}
    
    Wrapping the above code into the function definition from Listing \ref{lst:control_function_structure_transient}, the control law can be defined as
    
    \begin{lstlisting}[language=Python,caption={A pseudoinverse transient control law that can be loaded into Rattlesnake},label=lst:pseudoinverse_transient_control]
import numpy as np

def pseudoinverse_control(
        sample_rate,
        specification_signals, # Signal to try to reproduce
        frequency_spacing,
        transfer_function, # Transfer Functions
        noise_response_cpsd,  # Noise levels and correlation 
        noise_reference_cpsd, # from the system identification
        sysid_response_cpsd,  # Response levels and correlation
        sysid_reference_cpsd, # from the system identification
        multiple_coherence, # Coherence from the system identification
        frames, # Number of frames in the CPSD and FRF matrices
        total_frames, # Total frames that could be in the CPSD and FRF matrices
        output_oversample_factor, # Oversample factor to output
        extra_parameters = '', # Extra parameters for the control law
        last_excitation_signals = None, # Last excitation signal for drive-based control
        last_response_signals = None, # Last response signal for error correction
        ):
    # Parse the input arguments in extra_parameters
    rcond = 1e-15
    zero_impulse_after = None
    # Split it up into lines
    for entry in extra_parameters.split('\n'):
        try:
            # For each entry, split the key from the value using the colon
            field,value = entry.split(':')
            # Strip any whitespace
            field = field.strip()
            # Check the field to figure out which value to assign
            if field == 'rcond':
                rcond = float(value)
            elif field == 'zero_impulse_after':
                zero_impulse_after = float(value)
            else:
                # Report if we cannot understand the parameter
                print('Unrecognized Parameter: {:}, skipping...'.format(field))
        except ValueError:
            # Report if we cannot parse the line
            print('Unable to Parse Line {:}, skipping...'.format(entry))

    # Compute impulse responses using the IFFT of the transfer function
    # We will zero pad the IFFT to do interpolation in the frequency domain
    # to match the length of the required signal
    impulse_response = np.fft.irfft(transfer_function,axis=0)

    # The impulse response should be going to zero at the end of the frame,
    # but practically there may be some gibbs phenomenon effects that make the
    # impulse response noncausal.  If we zero pad, this might be wrong.  We
    # therefore give the use the ability to zero out this non-causal poriton of
    # the impulse response.
    if zero_impulse_after is not None:
        # Remove noncausal portion
        impulse_response_abscissa = np.arange(impulse_response.shape[0])/sample_rate
        zero_indices = impulse_response_abscissa > zero_impulse_after
        impulse_response[zero_indices] = 0
        
    # Zero pad the impulse response to create a signal that is long enough for
    # the specification signal
    added_zeros = np.zeros((specification_signals.shape[-1]-impulse_response.shape[0],) 
                           + impulse_response.shape[1:])
    full_impulse_response = np.concatenate((impulse_response,added_zeros),axis=0)

    # Compute FRFs using the FFT from the impulse response.  This is now
    # interpolated such that it matches the frequency spacing of the specification
    # signal
    interpolated_transfer_function = np.fft.rfft(full_impulse_response,axis=0)

    # Perform convolution by frequency domain multiplication
    signal_fft = np.fft.rfft(specification_signals,axis=-1)
    # Invert the FRF matrix using the specified rcond parameter
    inverted_frf = np.linalg.pinv(interpolated_transfer_function,rcond=rcond)
    # Multiply the inverted FRFs by the response spectra to get the drive spectra
    drive_signals_fft = np.einsum('ijk,ki->ij',inverted_frf,signal_fft)

    # Zero pad the drive FFT to oversample to the output_oversample_factor
    drive_signals_fft_zero_padded = np.concatenate((drive_signals_fft[:-1],
        np.zeros((drive_signals_fft[:-1].shape[0]*(output_oversample_factor-1)+1,)
                 +drive_signals_fft.shape[1:])),axis=0)

    # Finally, take the IFFT to get the time domain signal.  We need to scale
    # by the output_oversample_factor due to how the IFFT is normalized.
    drive_signals_oversampled = np.fft.irfft(
        drive_signals_fft_zero_padded.T,axis=-1)*output_oversample_factor
    return drive_signals_oversampled\end{lstlisting}
    
    The requirement that transient control laws are required to be able to oversample their output and may require switching between frequency and time domains mean that the Transient control laws will generally be more complex than the Random Vibration control laws.  Note that these computations to oversample the output are still performed by the controller for the Random Vibration environment described in Chapter \ref{sec:rattlesnake_environments_mimo_random}; however, the user does not need to handle them explicitly.  In the Transient environment, the user-defined control laws are actually creating the time signals that will be sent to the exciters, so user-defined control laws must handle these computations.
    
    \subsection{Defining a control law using state-persistent approaches}
    
    Similarly to the \Ac{MIMO} Random Vibration environment, the Transient environment may also use state-persistent approaches such as generator functions or classes.  This can be handy, for example, to compute the \ac{FFT} of the desired response signals one time, as opposed to computing them every single control iteration, because the desired response signals generally do not change throughout a test.
    
    \subsubsection{Defining Control Laws with Generator Functions}
    
    The second strategy to define a control law in Rattlesnake is to use a Generator Function.  A generator function is simply a function that maintains its internal state between function calls.  Listing \ref{lst:pseudoinverse_transient_generator} shows the Pseudoinverse Transient control law implemented as a generator function.
    
    \begin{lstlisting}[language=Python, caption = {A Pseudoinverse Transient control law defined using a Python generator function to allow for state persistence},label=lst:pseudoinverse_transient_generator]
import numpy as np

def pseudoinverse_control_generator():
    signal_fft = None
    inverted_frf = None
    drive_signals_oversampled = None
    while True:
        (sample_rate,
         specification_signals, # Signal to try to reproduce
         frequency_spacing,
         transfer_function, # Transfer Functions
         noise_response_cpsd,  # Noise levels and correlation 
         noise_reference_cpsd, # from the system identification
         sysid_response_cpsd,  # Response levels and correlation
         sysid_reference_cpsd, # from the system identification
         multiple_coherence, # Coherence from the system identification
         frames, # Number of frames in the CPSD and FRF matrices
         total_frames, # Total frames that could be in the CPSD and FRF matrices
         output_oversample_factor, # Oversample factor to output
         extra_parameters, # Extra parameters for the control law
         last_excitation_signals, # Last excitation signal for drive-based control
         last_response_signals, # Last response signal for error correction
         ) = yield drive_signals_oversampled
        if signal_fft is None:
            # Compute the FFT of the spec if it hasn't been done yet
            signal_fft = np.fft.rfft(specification_signals).T
        # Get a tolerance if specified
        rcond = 1e-15
        zero_impulse_after = None
        for entry in extra_parameters.split('\n'):
            field,value = entry.split(':')
            field = field.strip()
            if field == 'rcond':
                rcond = float(value)
            elif field == 'zero_impulse_after':
                zero_impulse_after = float(value)
            else:
                print('Unrecognized Parameter: {:}'.format(field))
        if inverted_frf is None:
            # Compute impulse responses
            impulse_response = np.fft.irfft(transfer_function,axis=0)

            if zero_impulse_after is not None:
                # Remove noncausal portion
                impulse_response_abscissa = np.arange(impulse_response.shape[0])/sample_rate
                zero_indices = impulse_response_abscissa > zero_impulse_after
                impulse_response[zero_indices] = 0
                
            # Zero pad the impulse response to create a signal that is long enough
            added_zeros = np.zeros((specification_signals.shape[-1]-impulse_response.shape[0],) + impulse_response.shape[1:])
            full_impulse_response = np.concatenate((impulse_response,added_zeros),axis=0)

            # Compute FRFs
            interpolated_transfer_function = np.fft.rfft(full_impulse_response,axis=0)

            # Perform convolution in frequency domain
            inverted_frf = np.linalg.pinv(interpolated_transfer_function,rcond=rcond)
            
        drive_signals_fft = np.einsum('ijk,ki->ij',inverted_frf,signal_fft)

        # Zero pad the FFT to oversample
        drive_signals_fft_zero_padded = np.concatenate((drive_signals_fft[:-1],
            np.zeros((drive_signals_fft[:-1].shape[0]*(output_oversample_factor-1)+1,)+drive_signals_fft.shape[1:])),axis=0)

        drive_signals_oversampled = np.fft.irfft(drive_signals_fft_zero_padded.T,axis=-1)*output_oversample_factor\end{lstlisting}
    
    Note that the generator function itself is not called with any arguments, as the initial function call simply starts up the generator.  Also note that there is no \inlineCode{return} statement in a generator function, only a \inlineCode{yield} statement.   When a program requests the next value from a generator, the generator code proceeds until it hits a \inlineCode{yield} statement, at which time it pauses and waits for the next value to be requested.  During this pause, all internal data is maintained inside the generator function.  The \inlineCode{yield} statement also accepts new data into the generator function, so this is where the same arguments used to define a control law using a Python function are passed in to the generator control law.  Therefore, by creating a while loop inside a generator function, the generator can be called infinitely many times to deliver data to the controller.
        
    To implement the Pseudoinverse control as shown in Listing \ref{lst:pseudoinverse_transient_generator} several parameters are initialized as \inlineCode{None}, which enables the generator to check whether or not they have been computed yet.  If they have, then there is no need to compute them again.
    
    \subsubsection{Defining Control Laws using Classes}
    
    A final way to implement more complex control laws is using a Python class.  This approach allows for the near infinite flexibility of Python's object-oriented programming at the expense of more complex syntax.  Users not familiar with Python's object-oriented programming paradigms are encouraged to \href{https://docs.python.org/3/tutorial/classes.html}{learn more} about the topic prior to reading this section.
    
    A class in Python is effectively a container that can have functions and properties stored inside of it, so it provides a good way to encapsulate all the parameters and helper functions associated with a given control law into one place.  A class allows for arbitrary properties to be stored within it, so arbitrary data can be made persistent between function calls.
    
    For the Rattlesnake implementation of a control law, the class must have at a minimum three functions defined.  These are the class constructor \inlineCode{__init__} that is called when the class is instantiated, a \inlineCode{system_id_update} function that is called upon completion of the System Identification portion of the controller, and a \inlineCode{control} function that actually computes the output \ac{CPSD} matrix.  A general class structure is shown below
    
    \begin{lstlisting}[language=Python,caption={Structure for a class defining a transient control law in Rattlesnake}]
# Any module imports or constants would go here    

class ControlLawClass:
    def __init__(self,
                 sample_rate,
                 specification_signals, # Signal to try to reproduce
                 output_oversample_factor, # Oversample factor to output
                 extra_parameters, # Extra parameters for the control law
                 frequency_spacing,
                 transfer_function, # Transfer Functions
                 noise_response_cpsd,  # Noise levels and correlation 
                 noise_reference_cpsd, # from the system identification
                 sysid_response_cpsd,  # Response levels and correlation
                 sysid_reference_cpsd, # from the system identification
                 multiple_coherence, # Coherence from the system identification
                 frames, # Number of frames in the CPSD and FRF matrices
                 total_frames, # Total frames that could be in the CPSD and FRF matrices
                 last_excitation_signals = None, # Last excitation signal for drive-based control
                 last_response_signals = None, # Last response signal for error correction
                 ):
        # Code to initialize the control law would go here
    
    def system_id_update(self,
                         frequency_spacing,
                         transfer_function, # Transfer Functions
                         noise_response_cpsd,  # Noise levels and correlation 
                         noise_reference_cpsd, # from the system identification
                         sysid_response_cpsd,  # Response levels and correlation
                         sysid_reference_cpsd, # from the system identification
                         multiple_coherence, # Coherence from the system identification
                         frames, # Number of frames in the CPSD and FRF matrices
                         total_frames, # Total frames that could be in the CPSD and FRF matrices
                         ):
        # Code to update the control law with system identification information would go here
    
    def control(self,
                last_excitation_signals = None, # Last excitation signal for drive-based control
                last_response_signals = None, # Last response signal for error correction
                ) -> np.ndarray:
        # Code to perform the actual control operations would go here
        
    # Any helper functions or properties that belong with the class would go here
    \end{lstlisting}
    
    The class's \inlineCode{__init__} constructor function is called whenever a class is instantiated (e.g. \inlineCode{control_law_object = ControlLawClass(sample_rate,specification_signals,output_oversample_factor,...)}). Note that the the constructor function accepts as arguments not only the data available at the time (e.g. the specification and any extra control parameters) but also any parameters that will eventually exist.  This is because it needs to be able to seamlessly transition in case the control law is changed during control when there is already a transfer function, last signal, etc.
    
    The \inlineCode{system_id_function} is called after the system identification is complete, so inside this function is where all setup calculations that require a transfer function would go.  In the case of a typical pseudoinverse control law, the system identification will not change as the test proceeds, so the inversion can take place in this function.
    
    The \inlineCode{control} function is then the function that actually performs the control operations to compute the output time history using the last response or output time histories in addition to any data that had been stored inside the class.
    
    The class implementation of the pseudoinverse control is shown in Listing \ref{lst:transient_pseudoinverse_class}.
    
    \begin{lstlisting}[language=Python,caption = {Class implementation of the buzz test approach},label=lst:transient_pseudoinverse_class]
import numpy as np

class pseudoinverse_control_class:
    def __init__(self,
                 sample_rate,
                 specification_signals, # Signal to try to reproduce
                 output_oversample_factor, # Oversample factor to output
                 extra_parameters, # Extra parameters for the control law
                 frequency_spacing,
                 transfer_function, # Transfer Functions
                 noise_response_cpsd,  # Noise levels and correlation 
                 noise_reference_cpsd, # from the system identification
                 sysid_response_cpsd,  # Response levels and correlation
                 sysid_reference_cpsd, # from the system identification
                 multiple_coherence, # Coherence from the system identification
                 frames, # Number of frames in the CPSD and FRF matrices
                 total_frames, # Total frames that could be in the CPSD and FRF matrices
                 last_excitation_signals = None, # Last excitation signal for drive-based control
                 last_response_signals = None, # Last response signal for error correction
                 ):
        self.rcond = 1e-15
        self.zero_impulse_after = None
        for entry in extra_parameters.split('\n'):
            field,value = entry.split(':')
            field = field.strip()
            if field == 'rcond':
                self.rcond = float(value)
            elif field == 'zero_impulse_after':
                self.zero_impulse_after = float(value)
            else:
                print('Unrecognized Parameter: {:}'.format(field))
        self.sample_rate = sample_rate
        self.specification_signals = specification_signals
        self.signal_fft = np.fft.rfft(specification_signals).T
        if self.transfer_function is not None:
            self.system_id_update(
                frequency_spacing,
                transfer_function, # Transfer Functions
                noise_response_cpsd,  # Noise levels and correlation 
                noise_reference_cpsd, # from the system identification
                sysid_response_cpsd,  # Response levels and correlation
                sysid_reference_cpsd, # from the system identification
                multiple_coherence, # Coherence from the system identification
                frames, # Number of frames in the CPSD and FRF matrices
                total_frames, # Total frames that could be in the CPSD and FRF matrices
                )
        
    def system_id_update(self,
                         frequency_spacing,
                         transfer_function, # Transfer Functions
                         noise_response_cpsd,  # Noise levels and correlation 
                         noise_reference_cpsd, # from the system identification
                         sysid_response_cpsd,  # Response levels and correlation
                         sysid_reference_cpsd, # from the system identification
                         multiple_coherence, # Coherence from the system identification
                         frames, # Number of frames in the CPSD and FRF matrices
                         total_frames, # Total frames that could be in the CPSD and FRF matrices
                         ):
        # Compute impulse responses
        impulse_response = np.fft.irfft(transfer_function,axis=0)

        if self.zero_impulse_after is not None:
            # Remove noncausal portion
            impulse_response_abscissa = np.arange(impulse_response.shape[0])/self.sample_rate
            zero_indices = impulse_response_abscissa > self.zero_impulse_after
            impulse_response[zero_indices] = 0
            
        # Zero pad the impulse response to create a signal that is long enough
        added_zeros = np.zeros((self.specification_signals.shape[-1]-impulse_response.shape[0],) 
                               + impulse_response.shape[1:])
        full_impulse_response = np.concatenate((impulse_response,added_zeros),axis=0)

        # Compute FRFs
        interpolated_transfer_function = np.fft.rfft(full_impulse_response,axis=0)

        # Perform convolution in frequency domain
        self.inverted_frf = np.linalg.pinv(interpolated_transfer_function,rcond=self.rcond)
        
        drive_signals_fft = np.einsum('ijk,ki->ij',self.inverted_frf,self.signal_fft)

        # Zero pad the FFT to oversample
        drive_signals_fft_zero_padded = np.concatenate((drive_signals_fft[:-1],
            np.zeros((drive_signals_fft[:-1].shape[0]*(self.output_oversample_factor-1)+1,)
                     +drive_signals_fft.shape[1:])),axis=0)

        self.drive_signals_oversampled = np.fft.irfft(
            drive_signals_fft_zero_padded.T,axis=-1)*self.output_oversample_factor

    def control(self,
                last_excitation_signals = None, # Last excitation signal for drive-based control
                last_response_signals = None, # Last response signal for error correction
                ) -> np.ndarray:       
        # We could modify the output signal based on new data that we obtained
        # Otherwise just output the same
        
        return self.drive_signals_oversampled\end{lstlisting}
        
    Note that the number of lines of code for a class implementation of the pseudoinverse control approach is not significantly more than the simpler function implementation, therefore users should not immediately discard the class-based approach as too difficult to implement in favor of the simpler function implementation.  Each approach has its merits and limitations, so it is up to the user to decide the best approach for their control law.
    
    \section{Using Transformation Matrices}
    Transformation matrices in the Transient environment behave identically to the the Random Vibration environment.  See Section \ref{sec:rattlesnake_environments_transformation_matrices} for more information.
    
    \section{Identifying the Signal in the Acquired Data}
    As the environment may be started or stopped arbitrarily during a given test, the Transient environment will need to identify where the signal that was desired actually occurs in the acquired signals.
    
    The Transient environment keeps a buffer twice the length of the desired control signal, so at some point during the acquisition, the entire signal should be within the buffer.
    
    To identify the position of the signal in the buffer, a correlation is performed between the signal that was output to the exciters and the acquired output data.  Correlation is performed on the output signal as it is generally read directly back into the controller and does not depend on an accurate prediction of the part's response.  The index of the maximum value of the correlation determines the sample closest to the starting point of the signal.
    
    A second, sub-sample alignment is then performed using phases from the \ac{FFT} of the signal truncated to the starting point and signal length compared to the phases of the specified output.  The slope of the phase change vs frequency line is proportional to the sub-sample signal shift.
    
    With the sample and subsample shift of the signal computed for the output signal, the same sample and subsample shift can be can be applied to the signal from the control gauges.  This portion of the control channel signals can then be compared directly to the desired signal to judge how well the environment is controlling.

    \chapter{Time History Generator}\label{sec:rattlesnake_environments_time_generator}
    
    Rattlesnake's Time History Generator environment provides the ability to simply stream a user-created signal to an output device.  While this is a relatively simple environment, it can be used to create simple shocks or run open-loop excitation devices such as a centrifuge with voltage input being proportional to speed.  There is no system identification, and therefore no test predictions that are made for this environment.
    
    \section{Signal Definition}
    
    The first step to defining a Time History Generator environment is to create the signal that will be output.  Rattlesnake accepts the signal in the form of a 2D array consisting of an output sample for each excitation signal for each time step.  Signals can be loaded from Numpy \inlineCode{*.npy} or \inlineCode{*.npz} files or Matlab \inlineCode{*.mat} files.  For \inlineCode{*.npy} files, the stored array defines the signal directly, and it is assumed that the signal uses the sample rate specified in Rattlesnake.  Note that if a hardware device over-samples the output (e.g. LAN-XI, see Section \ref{sec:rattlesnake_hardware_lanxi}), it is the over-sampled output sample rate that is used rather than the acquisition sample rate.  Matlab \inlineCode{*.mat} and Numpy \inlineCode{*.npz} files allow the users to specify a time vector as well as a signal, and should contain the following data members:
    \begin{description}
        \item[signal] A $n_o \times n_s$ array containing the signal for each exciter for each time step in \inlineCode{t}.
        \item[t] A $n_s$ array of times corresponding to the samples in the \inlineCode{signal} matrix.
    \end{description}
    where $n_o$ is the number of exciters and $n_s$ is the number of samples in the signal.  If the time vector specified by \inlineCode{t} does not match the sample rate specified in Rattlesnake, the \inlineCode{signal} data will be linearly interpolated to provide the correct sample rate.  If \inlineCode{t} is not provided in the \inlineCode{*.mat} or \inlineCode{*.npz} file, Rattlesnake will treat \inlineCode{signal} as if it were defined at the output sample rate of the controller.
    
    The ordering of the signals in the signal file is the same as the ordering of the excitation devices in the channel table that are active in the current environment.  The first signal will be played to the first excitation device, and so on.
    
    Note that the environment will play the signal as-is, so it is up to the user to implement graceful startup and shutdown at the start and end of the signal if the test configuration requires it.
    
    \section{Defining the Time History Generator Environment in Rattlesnake}
    In addition to the signal that will be played to the excitation devices, there is only one parameter that needs to be defined in the Time History Generator.  Figure \ref{fig:timehistorygeneratorenvironmentdefinition} shows a Time History Generator sub-tab in the \inlineCode{Environment Definition} tab of Rattlesnake.
    
    \begin{figure}
        \centering
        \includegraphics[width=\linewidth]{figures/time_history_generator_environment_definition}
        \caption{\Ac{GUI} defining the Time History Generation environment}
        \label{fig:timehistorygeneratorenvironmentdefinition}
    \end{figure}
    
    Pressing the \inlineCode{Load Signal} button brings up a file dialog from which the signal file can be loaded.  Once the file is loaded, it is displayed in the main plot window.  Signal statistics are also displayed in the adjacent table.  The checkbox in the \inlineCode{Show?} column of the table can be used to show or hide individual signals.  The signal name in the \inlineCode{Signal} column is constructed from the node number and direction in the channel table.  The \inlineCode{Max} and \inlineCode{RMS} value of the signal is also displayed.
    
    At the bottom of the window, there are various computed parameters, and one user defined parameter.
    \begin{description}
        \item[Sample Rate]  The global sample rate of the data acquisition system.  This is set on the \inlineCode{Data Acquisition Setup} tab, and displayed here for convenience as a read-only value.
        \item[Output Sample Rate]  This is the output sample rate, which might be different than \inlineCode{Sample Rate} if the hardware over-samples the output device.  This is set on the \inlineCode{Data Acquisition Setup} tab, and displayed here for convenience as a read-only value.
        \item[Output Channels]  The number of excitation signals being used by this environment.  This is a computed quantity presented for convenience, so the user cannot modify it directly.
        \item[Signal Samples] The number of samples in the loaded signal.  This is a computed quantity presented for convenience, so the use cannot modify it directly.
        \item[Signal Time] The amount of time it will take to play the signal.  This is a computed quantity presented for convenience, so the user cannot modify it directly.
        \item[Cancel Rampdown Time] The amount of time the environment will take to ramp to zero if the environment is stopped prior to the signal being completely played.
    \end{description}
    
    \section{Running the Time History Generator Environment}
    The Time History Generator environment is then run on the \inlineCode{Run Test} tab of the controller.  With the data acquisition system armed, the \ac{GUI} looks like Figure \ref{fig:timehistorygeneratorruntest}.
    
    \begin{figure}
        \centering
        \includegraphics[width=\linewidth]{figures/time_history_generator_run_test}
        \caption{\Ac{GUI} for running the Time History Generator Environment}
        \label{fig:timehistorygeneratorruntest}
    \end{figure}
    
    Two parameters can be defined prior to starting the environment.
    \begin{description}
        \item[Signal Level] Allows the user to scale the signal prior to output by a number of dB.
        \item[Repeat Signal] Checking this checkbox makes the signal repeat once started until the environment is stopped manually.  Alternatively, the environment will stop automatically once the signal has been played.
    \end{description}

    Similar to other environments, there are \inlineCode{Start Environment} and \inlineCode{Stop Environment} buttons to control when the environment occurs.  If the \inlineCode{Stop Environment} button is clicked, the signal will continue to play for the specified \inlineCode{Cancel Rampdown Time} while the environment ramps the signal level to zero.
    
    \section{Output NetCDF File Structure}
    When Rattlesnake saves data to a netCDF file, environment-specific parameters are stored in a netCDF group with the same name as the environment name.  Similar to the root netCDF structure described in Section \ref{sec:using_rattlesnake_output_files}, this group will have its own attributes, dimensions, and variables, which are described here.
    
    \subsection{NetCDF Dimensions}
    \begin{description}
        \item[output\_channels] The number of output signals used by the environment
        \item[signal\_samples] The number of samples in the output signals
    \end{description}
    
    \subsection{NetCDF Attributes}
    \begin{description}
        \item[cancel\_rampdown\_time] The time to ramp to zero if the environment is stopped.
    \end{description}
    
    \subsection{NetCDF Variables}
    \begin{description}
        \item[output\_signal] The signals that are played to the excitation devices  Type: 64-bit float; Dimensions: \inlineCode{output_channels} $\times$ \inlineCode{signal_samples}
    \end{description}
    
    \chapter{Modal Testing}\label{sec:rattlesnake_environments_modal}
    Rattlesnake's most recent environment is the Modal Testing environment, which is designed for dynamic characterization of a test article.  The environment computes transfer functions between a subset of reference channels and the remaining response channels.  Reference channels are traditionally measurements of the forces applied to the structure using either a modal impact hammer or modal shaker.  The modal testing environment can generate many standard modal signals such as chirp or burst random.  However, shaker signals can also be generated from other environments using Rattlesnake's combined environments capabilities.
    
    \section{Defining the Modal Testing Environment in Rattlesnake}
    
    The Rattlesnake Modal Testing environment has many signal processing parameters to specify when setting up the modal test.  These are defined on the \inlineCode{Environment Definition} tab in the Rattlesnake controller on a sub-tab corresponding to the name of the modal testing environment.  Figure \ref{fig:modalenvironmentdefinition} shows a Modal Testing sub-tab.  The following subsections describe the parameters that can be specified, as well as their effects on the analysis.
    
    \begin{figure}
        \centering
        \includegraphics[width=\linewidth]{figures/modal_environment_definition}
        \caption{\Ac{GUI} used to define a Modal Testing Environment}
        \label{fig:modalenvironmentdefinition}
    \end{figure}
    
    \subsection{Sampling Parameters}
    The \inlineCode{Sampling Parameters} section contains information and settings that pertain to the samping in the modal test.  It consists of the following parameters:
    
    \begin{description}
        \item[Sample Rate] The global sample rate of the data acquisition system.  This is set on the \inlineCode{Data Acquisition Setup} tab, and displayed here for convenience as a read-only value.
        \item[Samples Per Frame] The number of time samples used in each measurement frame.  Modifying this value will change the length of the measurement, as well as the frequency spacing of the spectral data.
        \item[Samples Per Acquire] The number of additional samples required per measurement frame when taking into account overlapping of the measurement frames.  This is a computed quantity presented for convenience, so the user cannot modify it directly.
        \item[Frame Time] The amount of time it takes to measure a frame of data, in seconds.  This is a computed quantity presented for convenience, so the user cannot modify it directly.
        \item[Nyquist Frequency] The maximum bandwidth of the measurement given the sample rate.  This is the largest frequency value in the \Ac{FFT}.  This is a computed quantity presented for convenience, so the user cannot modify it directly.
        \item[FFT Lines] The number of frequency lines in the \Ac{FFT}.  This is a computed quantity presented for convenience, so the user cannot modify it directly.
        \item[Frequency Spacing] The frequency resolution of the \Ac{FFT}.  This is a computed quantity presented for convenience, so the user cannot modify it directly.
    \end{description}
    
    \subsection{System ID Parameters}
    The \inlineCode{System ID Parameters} section contains information and settings that pertain to the computation of spectral quantities in the Modal Testing Environment.  This section contains the following parameters:
    
    \begin{description}
        \item[Averaging Type] The type of averaging used to compute the spectral quantities.  Linear averaging weights each measurement frame equally.  Exponential averaging weights more recent frames more heavily.
        \item[Number of Averages] The number of averaged measurement frames used to compute the spectral quantities.
        \item[Averaging Coefficient] If Exponential Averaging is used, this is the weighting of the most recent frame compared to the weighting of the previous frames.  If the averaging coefficient is $\alpha$, then the most recent frame will be weighted $\alpha$, the frame before that will be weighted $\alpha(1-\alpha)$, the frame before that will be $\alpha(1-\alpha)^2$, etc.
        \item[FRF Technique] The estimator used to compute transfer functions between voltage signals and responses.
        \item[FRF Window] The window function used to reduce leakage in the \Ac{FFT} calculation.
        \item[Overlap Percentage] The overlap percentage between measurement frames when computing spectral quantities.
        \item[Window value at frame end] For an exponential window, the value of the window function at the end of the measurement frame.  This value will determine how quickly the exponential window decays.  This parameter is only visible when an exponential window is selected.
    \end{description}
    
    \subsection{Triggering}
    The \inlineCode{Triggering} section contains information and settings that pertain to triggering each measurement frame using a measured signal.  The parameters are:
    
    \begin{description}
        \item[Triggering Type] The type of triggering used for a test.  If \inlineCode{Free Run} is selected, all data will be acquired immediately.  This is useful for Random, Pseudorandom, or Chirp signals.  If \inlineCode{First Frame} is selected, the software will wait for a trigger to occur before the first frame is measured, but subsequent frames will be measured immediately.  This setting is useful for Burst Random excitation, where the first frame must be synchronized with the excitation signal, but because the excitation is the same length as the frame, subsequent frames will automatically be aligned with the excitation.  If \inlineCode{Every Frame} is selected, a trigger will be required every measurement frame.  This is useful for Impact Hammer testing, where there may be an unknown amount of time between hammer impacts.
        
        \item[Acceptance] If \inlineCode{Accept All} is selected, all measurement frames will be automatically accepted.  If \inlineCode{Manual} is selected, the user will have the option to accept or reject measurements.  If \inlineCode{Autoreject...} is selected, the user will be prompted to load in a Python function that will accept the time history measurement for the given measurement frame and return a True or False depending on whether or not the frame should be accepted or rejected.  This could be used to automatically reject a signal if, for example, there were a double hit or if the peak level were not correct.
        
        \item[Wait for Steady State] This parameter specifies a time to wait before acquiring measurements to allow the system to come to steady state.  This is useful for signals that assume steady state response, such as Sinusoid, Pseudorandom, or Chirp excitation.
        
        \item[Trigger Channel] This parameter selects the channel used to trigger the measurement.
        
        \item[Pretrigger] This parameter specifies how much of the measurement frame occurs prior to the trigger.  This is useful to allow the entire signal to be measured, as a portion of the trigger signal often occurs before the trigger actually occurs.
        
        \item[Trigger Slope] This parameter specifies the slope of the signal used to find the trigger.  A value of \inlineCode{Positive} means that the trigger will occur when the signal is below the trigger level and then goes above the trigger level.  A value of \inlineCode{Negative} means that the trigger will occur when the signal is above the trigger level and then goes below the trigger level.
        
        \item[Trigger Level] The level of the trigger as a percentage of the full range of the channel.
        
        \item[Trigger Level (V)] The level of the trigger in volts.
        
        \item[Trigger Level (EU)] The level of the trigger in Engineering Units.  These are the units specified in the channel table based on the sensitivity of the channel specified in the channel table.
        
        \item[Hysteresis Level] The Hysteresis level is the level below which the trigger must return prior to another trigger being able to be accepted.  This can be used to ensure that the trigger signal returns to some nominal state (e.g. zero force) before another trigger can be acquired.
        
        \item[Hysteresis Level (V)] The Hysteresis Level of the trigger signal in volts.
        
        \item[Hysteresis Level (EU)] The Hysteresis level of the trigger signal in Engineering Units.  These are the units specified in the channel table based on the sensitivity of the channel specified in the channel table.
        
        \item[Hysteresis Frame \%] The percentage of the measurement frame that must be below the Hysteresis level before another trigger can be accepted.  This is useful for Burst Random or Hammer excitation where the signal should return to zero for some portion of the frame prior to a trigger being acquired.  Without this parameter, a trigger could occur in the middle of a Burst, depending on when the acquisition starts up relative to the output.
        
        \item[Hysteresis Samples] The number of samples that the trigger signal must be below the Hysteresis Level.
        
        \item[Hysteresis Time] The amount of time the trigger signal must be below the Hysteresis Level.
    \end{description}
    
    \subsection{Channel Information}
    The \inlineCode{Channel Information} section contains information regarding how many channels are active for the given test.  These quantities are read-only to inform the user about the number of channels the current setup will contain.
    
    \begin{description}
        \item[Response Channels] The number of response channels in the measurement.  These will be compared against the reference channels to compute frequency response functions.
        \item[Reference Channels] The number of reference channels in the measurement.
        \item[Output Channels] The number of output channels that will be present in the test.  If used as shaker drives, this value should generally match the number of \inlineCode{Reference Channels}.  For a hammer test, this value may be zero.
    \end{description}
        
    \subsection{Channel Selections}
    The \inlineCode{Channel Selections} section contains a list of all of the channels in the test and the ability to select which channels are references or responses, as well as to enable or disable channels.
    
    To disable a channel, simply uncheck the checkbox in the \inlineCode{Enabled} column associated with a given channel.  Disabling a voltage channel associated with a \inlineCode{Feedback Device} and \inlineCode{Feedback Channel} on the channel table will also disable that output.
    
    To make a channel a reference channel, simply check the checkbox in the \inlineCode{Reference} column associated with a given channel.  This will turn the channel from a Response channel to a Reference channel.
    
    Multiple rows of the \inlineCode{Channel Selctions} table can be selected at once.  The \inlineCode{Enable Selected} and \inlineCode{Disable Selected} buttons will enable or disable all selected channels, respectively.  Similarly, the \inlineCode{Check Selected References} and \inlineCode{Uncheck Selected References} buttons will turn the selected channels into Reference Channels or Response Channels, respectively.
    
    \subsection{Signal Generator}
    The \inlineCode{Signal Generator} section contains the parameters to determine what signal will be generated for the modal test.
    The tabs at the top of the \inlineCode{Signal Generator} section determine the type of signal that will be applied.  Each signal type may have different parameters to define it.
    
    \begin{description}
        \item[No Signal Generation] No signal will be generated if this is selected.
        \item[Random] A True Random signal will be generated.  This signal is continually generated, and being random, will not repeat for each measurement frame.  Generally a window function such as a Hann window should be used with a Random signal.
        \begin{description}
            \item[RMS] Sets the RMS voltage level for the output signal
            \item[Frequency Range] Sets the minimum and maximum content of the signal
        \end{description}
        \item[Burst Random] A Burst Random signal will be generated.  A Burst Random signal is essentially a Random signal that shuts off part of the way through a measurement frame.  This signal is continually generated, and being random, will not repeat for each measurement frame.  When using a Burst Random excitation, users will generally want to set up a \inlineCode{First Frame} trigger on a voltage channel with a \inlineCode{Hysteresis Frame %} of approximately half the time the burst will be ``off''.
        \begin{description}
            \item[RMS] Sets the RMS voltage level for the output signal.  The RMS level here describes the level when the burst is active.
            \item[Frequency Range] Sets the minimum and maximum content of the signal
            \item[On Percentage] Sets the fraction of the measurement frame that the burst will be active for.  The rest of the measurement frame, the signal will be zero.
        \end{description}
        \item[Pseudorandom] A Pseudorandom signal will be generated.  This is a signal that appears random, but is actually deterministically constructed from the frequency lines of the FFT with randomized phases.  This signal will repeat for each measurement frame.  When using Pseudorandom excitation, users should set the \inlineCode{Wait For Steady State} parameter to ensure the system has reached steady state prior to acquiring data.  Pseudorandom excitation will generally not work well with multiple shaker excitation.
        \begin{description}
            \item[RMS] Sets the RMS voltage level for the output signal
            \item[Frequency Range] Sets the minimum and maximum content of the signal
        \end{description}
        \item[Chirp] A Chirp signal will be generated.  A Chirp is a fast sine sweep between two frequency values.  This signal will repeat for each measurement frame.  When using Chirp excitation, users should set the \inlineCode{Wait For Steady State} parameter to ensure the system has reached steady state prior to acquiring data.  Chirp excitation will generally not work well with multiple shaker excitation.
        \begin{description}
            \item[Peak Level] Sets the peak voltage level for the output signal
            \item[Frequency Range] Sets the start and end frequencies of the Chirp
        \end{description}
        \item[Square Pulse] A square pulse wave will be generated.  This wave is useful for triggering other test hardware, such as an automatic hammer.
        \begin{description}
            \item[Peak Level] Sets the peak voltage level for the output signal
            \item[Frequency] Sets the frequency of the square pulse.  The period of the square pulse can be longer than one measurement frame.
            \item[Percent On] Sets the percentage of the period of the square wave that the signal is ``up''.
        \end{description}
        \item[Sine] A sine pulse wave will be generated.  This is useful for characterizing a single frequency line or for performing rigid body checkouts.  Be careful when selecting a frequency to ensure there is no leakage at that frequency line.
        \begin{description}
            \item[Peak Level] Sets the peak voltage level for the output signal
            \item[Frequency] Sets the frequency of the sine wave
        \end{description}
    \end{description}
    
    The selected signal will be plotted in the \inlineCode{Output Signal} plot, which will give the user an idea of what the signal looks like.  If the \inlineCode{Automatically Regenerate Signal on Parameter Change} checkbox is checked, the signal should update automatically when different parameters are selected.  If not, the user can press the \inlineCode{Regenerate Signal} button to regenerate the signal.  For random signals, the \inlineCode{Regenerate Signal} button can also be used to visualize different realizations of the random signal.
    
    \section{Running a Modal Test}
    The Modal Testing Environment is then run on the \inlineCode{Run Test} tab of the controller.  With the data acquisition system armed, the \ac{GUI} initially looks like Figure \ref{fig:modalenvironmentruntestempty}.  This screen looks rather empty, but users can populate it with functions of their choice, as shown in Figure \ref{fig:modalenvironmentruntestfull}.

    \begin{figure}
        \centering
        \includegraphics[width=\linewidth]{figures/modal_environment_run_test_empty}
        \caption{Empty \ac{GUI} to run a modal test in Rattlesnake.}
        \label{fig:modalenvironmentruntestempty}
    \end{figure}
    
    
    \begin{figure}
        \centering
        \includegraphics[width=\linewidth]{figures/modal_environment_run_test_full}
        \caption{\Ac{GUI} to run a modal test in Rattlesnake populated with several data plots.}
        \label{fig:modalenvironmentruntestfull}
    \end{figure}
    
    \subsection{Acquisition}
    The \inlineCode{Acquisition} portion of the window contains controls for starting and stopping the measurement, as well as saving modal data to a file.  See Section \ref{sec:rattlesnake_environments_modal_output_files} for a description of the file format.
    
    \begin{description}
        \item[Preview Acquisition] Clicking the \inlineCode{Preview Acquisition} button will start the measurement running continuously, so the user can preview it and make sure it looks correct.  Data will not be saved to disk, nor will the acquisition automatically stop after the specified number of measurement frames have been acquired.
        \item[Start Acquisition]  Clicking the \inlineCode{Start Acquisition} button will start the measurement running.  Data will be stored to disk in the file specified in the \inlineCode{Data File:} box, and the acquisition will automatically stop when the specified number of averages has been acquired.
        \item[Stop Acquisition] Clicking the \inlineCode{Stop Acquisition} button will stop the current measurement.  Data already written to disk will remain on disk.
        \item[Data File:] The \inlineCode{Data File:} box displays the file name that the modal data will be written to.  See Section \ref{sec:rattlesnake_environments_modal_output_files} for a description of the file format.
        \item[Select File...] The \inlineCode{Select File...} button allows the user to select a file to use to save the modal data to disk.
        \item[Auto-Increment] If this checkbox is checked, a four-digit number will be appended to the file name, and will be incremented for each run.  This can be used to prevent accidentally overwriting data that has already been acquired.
    \end{description}
    
    \subsection{Averaging}
    The \inlineCode{Averaging} portion of the window displays the current number of measurement frames that have been acquired, as well as the total number of measurement frames that are to be acquired.  If \inlineCode{Acceptance} is set to \inlineCode{Manual}, then the \inlineCode{Accept} and \inlineCode{Reject} buttons will become available after each measurement frame is acquired, which will allow the users to manually accept or reject a measurement frame.
    
    If the measurement was started using the \inlineCode{Start Acquisition} button, the measurement will stop automatically when the current number of measurement frames is equal to the total number of measurement frames.  If \inlineCode{Preview Acquisition} was used, then the measurement will continue until stopped with the \inlineCode{Stop Acquisition} button.
    
    \subsection{Windows}
    The \inlineCode{Windows} portion of the screen allows the user to customize the data that they visualize during the test.
    
    \begin{description}
        \item[New] Clicking the \inlineCode{New} button will create a new window on the right-hand side of the screen.
        \item[New from Template...] The \inlineCode{New from Template...} drop down menu allows users to visualize common data types for modal testing.
        \begin{description}
            \item[Drive Point (Magnitude)] Rattlesnake will parse the channel names in the channel table to identify drive point \Ac{FRF}s, of which it will display the magnitude.  This will create one window for each drive point \Ac{FRF}.
            \item[Drive Point (Imaginary)] Rattlesnake will parse the channel names in the channel table to identify drive point \Ac{FRF}s, of which it will display the imaginary part.  This will create one window for each drive point \Ac{FRF}.  For an acceleration to force \ac{FRF}, the imaginary part of the drive point \ac{FRF} should always be positive if both channels reference and response channels have polarity of + or -, or negative if one of the reference or response channels has a polarity of + and the other has a polarity of -.  If the imaginary part of the drive point \inlineCode{FRF} crosses zero there could be an issue with the test setup.
            \item[Drive Point Coherence] Rattlesnake will parse the channel names in the channel table to identify drive point \Ac{FRF}s which it will display with the coherence overlaid.  This will create one window for each drive point \Ac{FRF}.
            \item[Reciprocity] Rattlesnake will parse the channel names in the channel table to identify reciprocal \ac{FRF}s to overlay.  It will find, for example the \ac{FRF} with input at degree of freedom \inlineCode{A} and response at degree of freedom \inlineCode{B} and overlay it with the \ac{FRF} with input at \inlineCode{B} and response at \inlineCode{A}.  It will generally create $\frac{n^2-n}{2}$ windows where $n$ is the number of drive points in the measurement.
            \item[Reference Autospectrum] Rattlesnake will plot the autospectrum for each reference channel in the test.  Users can use this plot to ensure adequate excitation over the bandwidth of interest.  It will generate one window per each reference channel in the test.
            \item[3x3 Channel Grid] Rattlesnake will create 9 windows which by default display time histories for the first 9 channels in the test.  This selection, combined with setting the \inlineCode{Increment Channels} selector to 9, can be used to quickly look through all of the channels in the test.
        \end{description}
        \item[Tile] Clicking the \inlineCode{Tile} button will tile all data display windows across the space available in the Rattlesnake \ac{GUI}.
        \item[Close All] Clicking this button will close all data display windows that are currently open.
        \item[Increment Channels] If a value is set in the \inlineCode{Increment Channels} selector, then when the \inlineCode{<<} or \inlineCode{>>} buttons are clicked, all of the response channels in all of the windows that are not locked will decrement or increment by that number.  This can be used to quickly look through all channels in a test without having to change the channel in each window manually.  For display windows with multiple channels, this will increment only the response channel.  The reference channel must be set manually.
    \end{description}
    
    \subsubsection{Window Types}
    Windows created in the \inlineCode{Run Test} tab of a Modal Testing environment are flexible in that they can show multiple different types of data in various formats for any channel in the test.  All windows have a \inlineCode{Lock} checkbox that when checked does not allow the channel to change via the \inlineCode{Increment Channels} arrow buttons.  A locked channel can still be changed manually, however.
    
    This section will walk through the various window types in a modal test.
    
    \subsubsection{Time Window}
    The \inlineCode{Time} window displays the time trace for a single measurement frame for the selected channel.  Figure \ref{fig:modaltimewindow} shows an example of this window.  The channel to visualize can be chosen from the drop-down menu.  The \inlineCode{Lock} checkbox can be checked to ensure the channel does not change when the \inlineCode{Increment Channels} functionality is utilized.
    
    
    \begin{figure}
        \centering
        \includegraphics[width=0.7\linewidth]{figures/modal_time_window}
        \caption{Modal Testing data window showing a channel's time signal.}
        \label{fig:modaltimewindow}
    \end{figure}
    
    \subsubsection{Windowed Time Window}
    The \inlineCode{Windowed Time} window displays the time trace for a single measurement with the window function applied.  This allows users to visualize the affect that the window has on the data.  Figure \ref{fig:modalwindowedtimewindow} shows an example of this window.  The channel to visualize can be chosen from the drop-down menu.  The \inlineCode{Lock} checkbox can be checked to ensure the channel does not change when the \inlineCode{Increment Channels} functionality is utilized.
    
    
    \begin{figure}
        \centering
        \includegraphics[width=0.7\linewidth]{figures/modal_windowed_time_window}
        \caption{Modal Testing data window showing a channel's windowed time signal.}
        \label{fig:modalwindowedtimewindow}
    \end{figure}
    
    \subsubsection{Spectrum Window}
    The \inlineCode{Spectrum} window displays the magnitude of the FFT of the (windowed) time trace for a single measurement.  Figure \ref{fig:modalspectrumwindow} shows an example of this window.  The channel to visualize can be chosen from the drop-down menu.  The \inlineCode{Lock} checkbox can be checked to ensure the channel does not change when the \inlineCode{Increment Channels} functionality is utilized.
    
    
    \begin{figure}
        \centering
        \includegraphics[width=0.7\linewidth]{figures/modal_spectrum_window}
        \caption{Modal Testing data window showing a channel's spectrum.}
        \label{fig:modalspectrumwindow}
    \end{figure}
    
    \subsubsection{Autospectrum Window}
    The \inlineCode{Autospectrum} window displays the magnitude of the autospectrum for a single measurement.  This is an averaged quantity, so it will generally improve as more averages are acquired.  Figure \ref{fig:modalspectrumwindow} shows an example of this window.  The channel to visualize can be chosen from the drop-down menu.  The \inlineCode{Lock} checkbox can be checked to ensure the channel does not change when the \inlineCode{Increment Channels} functionality is utilized.
    
    
    \begin{figure}
        \centering
        \includegraphics[width=0.7\linewidth]{figures/modal_autospectrum_window}
        \caption{Modal Testing data window showing a channel's autospectrum.}
        \label{fig:modalautospectrumwindow}
    \end{figure}
    
    \subsubsection{FRF Window}
    The \inlineCode{FRF} window displays a \Ac{FRF} for a reference/response channel combination.  This is an averaged quantity, so it will generally improve as more averages are acquired.  Figure \ref{fig:modalfrfwindow} shows an example of this window.  Two channel selection menus exist to select which \ac{FRF} to visualize.  The first corresponds to the \inlineCode{Response} channel, and the second corresponds to the \inlineCode{Reference} channel.  Channels selected as a Reference on the \inlineCode{Environment Definition} tab will show up in the reference channel selection menu, and channels that were not selected will show up in the response channel selection menu.  The \ac{FRF} can be visualized by looking at Real, Imaginary, Magnitude, or Phase parts, or it can be split into two plots to visualize Magnitude and Phase or Real and Imaginary parts simultaneously.  The \inlineCode{Lock} checkbox can be checked to ensure the response channel does not change when the \inlineCode{Increment Channels} functionality is utilized.
    
    
    \begin{figure}
        \centering
        \includegraphics[width=0.7\linewidth]{figures/modal_frf_window}
        \caption{Modal Testing data window showing the FRF between a reference and response channel.}
        \label{fig:modalfrfwindow}
    \end{figure}
    
    \subsubsection{Coherence Window}
    The \inlineCode{Coherence} window displays the coherence for a given channel.  Figure \ref{fig:modalcoherencewindow} shows an example of this window.  If only a single reference is used, this will be the regular coherence.  If multiple references are used, this will then be the Multiple Coherence function.  The channel to visualize can be chosen from the drop-down menu.  The \inlineCode{Lock} checkbox can be checked to ensure the channel does not change when the \inlineCode{Increment Channels} functionality is utilized.
    
    
    \begin{figure}
        \centering
        \includegraphics[width=0.7\linewidth]{figures/modal_coherence_window}
        \caption{Modal Testing data window showing a channel's coherence.}
        \label{fig:modalcoherencewindow}
    \end{figure}
    
    \subsubsection{FRF Coherence Window}
    The \inlineCode{FRF Coherence} window displays a \ac{FRF} overlaid with the Coherence plot.  This allows users to align drops in coherence with features of the \inlineCode{FRF} to determine if they occur at modes of the structure, which could suggest an issue with the data.  Figure \ref{fig:modalfrfcoherencewindow} shows an example of this window.  Two channel selection menus exist to select which \ac{FRF} to visualize.  The first corresponds to the \inlineCode{Response} channel, and the second corresponds to the \inlineCode{Reference} channel.  The coherence will generally correspond to the Response channel: if multiple references exist, then the multiple coherence that is plotted will be a function of the response degree of freedom; if only a single reference exists, then the regular coherence will be with respect to the only reference in the test, and will also then only change when the response is updated.  The \ac{FRF} can be visualized by looking at Real, Imaginary, Magnitude, or Phase parts, or it can be split into two plots to visualize Magnitude and Phase or Real and Imaginary parts simultaneously.  If these latter options are selected, the coherence will be overlaid on the second plot, which will be the Magnitude or Imaginary part.  The \inlineCode{Lock} checkbox can be checked to ensure the response channel does not change when the \inlineCode{Increment Channels} functionality is utilized.
    
    
    \begin{figure}
        \centering
        \includegraphics[width=0.7\linewidth]{figures/modal_frfcoherence_window}
        \caption{Modal Testing data window showing a channel's FRF overlaid with the coherence.}
        \label{fig:modalfrfcoherencewindow}
    \end{figure}
    
    \subsubsection{Reciprocity Window}
    The \inlineCode{Reciprocity} window displays two reciprocal \ac{FRF}s overlaid.  For a linear system, these \ac{FRF}s should be identical to one another.  Figure \ref{fig:modalreciprocitywindow} shows an example of this window.  Two channel selection menus exist to select which \ac{FRF} to visualize.  The first corresponds to the \inlineCode{Response} channel, and the second corresponds to the \inlineCode{Reference} channel.  Rattlesnake will automatically go through and find the reciprocal \ac{FRF} for that measurement and overlay it.  The \ac{FRF} can be visualized by looking at Real, Imaginary, Magnitude, or Phase parts, or it can be split into two plots to visualize Magnitude and Phase or Real and Imaginary parts simultaneously.  The \inlineCode{Lock} checkbox can be checked to ensure the response channel does not change when the \inlineCode{Increment Channels} functionality is utilized.
    
    
    \begin{figure}
        \centering
        \includegraphics[width=0.7\linewidth]{figures/modal_reciprocity_window}
        \caption{Modal Testing window showing two reciprocal FRFs.}
        \label{fig:modalreciprocitywindow}
    \end{figure}
    
    \FloatBarrier
    
    \subsection{Degree of Freedom Override}
    This section of the Modal Testing \inlineCode{Run Test} tab allows users to override the channel metadata for a given measurement.  This is particularly useful for roving hammer or roving accelerometer testing strategies.  The node identification number and direction are defined on the channel table on the first tab of the software, and it would be very tedious to need to move back to the first tab to redefine the node number, reinitialize the data acquisition, redefine the environment definition, and finally, rearm the test to take data.  Instead, the user can override the degree of freedom information without leaving the \inlineCode{Run Test} tab.
    
    To add an override channel, users can press the \inlineCode{Add Override} button.  This will create a new row in the override table.  These can be removed by clicking the \inlineCode{Remove Override} button with the particular row in the table selected.  Figure \ref{fig:modaldofoverride} shows an example of this functionality where the node information on the force channel has been updated.  Note that the windows will also be updated to display this new degree of freedom information.
    
    \begin{figure}
        \centering
        \includegraphics[width=\linewidth]{figures/modal_dof_override}
        \caption{Overriding the 16934 X+ degree of freedom information with 16950 X+ degree of freedom information.}
        \label{fig:modaldofoverride}
    \end{figure}
    
    
    \FloatBarrier
    
    \section{Output NetCDF File Structure}\label{sec:rattlesnake_environments_modal_output_files}
    When Rattlesnake saves data to a netCDF file, environment-specific parameters are stored in a netCDF group with the same name as the environment name.  Similar to the root netCDF structure described in Section \ref{sec:using_rattlesnake_output_files}, this group will have its own attributes, dimensions, and variables, which are described here.
    
    \subsection{NetCDF Dimensions}
    \begin{description}
        \item[reference\_channels] The number of reference channels in the measurement
        \item[response\_channels] The number of response channels in the measurement
    \end{description}
    
    \subsection{NetCDF Attributes}
    \begin{description}
        \item[samples\_per\_frame] The number of samples per measurement frame.
        \item[averaging\_type] A string specifying if linear or exponential averaging was performed in the measurement.
        \item[num\_averages] The number of averages used to compute spectral quantities like \ac{FRF}s and autospectra.
        \item[averaging\_coefficient] If Exponential Averaging is used, this is the weighting of the most recent frame compared to the weighting of the previous frames.  If the averaging coefficient is $\alpha$, then the most recent frame will be weighted $\alpha$, the frame before that will be weighted $\alpha(1-\alpha)$, the frame before that will be $\alpha(1-\alpha)^2$, etc.
        \item[frf\_technique] A string representing the estimator used to compute \ac{FRF}s, e.g. H1, H2, Hv.
        \item[frf\_window] A string representing the window function used when computing \ac{FRF}s.
        \item[overlap] The overlap fraction between measurement frames.
        \item[trigger\_type] A string representing the type of triggering used in the modal test (e.g. Free Run, First Frame, Every Frame).
        \item[accept\_type] A string representing the acceptance type used in the modal test (e.g. Accept All, Manual).
        \item[wait\_for\_steady\_state] The amount of time the measurement needed to wait for steady state to occur before acquiring data.
        \item[trigger\_channel] An integer representing the channel used to trigger the measurement.
        \item[pretrigger] The amount of the measurement frame that occurs prior to the trigger.
        \item[trigger\_slope\_positive] An integer equal to 1 if the slope was positive and 0 if the slope was negative.
        \item[trigger\_level] The level that was used to trigger the system as a fraction of the total range of the channel.
        \item[hysteresis\_level] The level that the data must return back to prior to being able to accept another trigger as a fraction of the total range of the channel.
        \item[hysteresis\_length] The length of time that a channel must return to the hysteresis level prior to being able to accept another trigger as a fraction of the measurement frame length.
        \item[signal\_generator\_type] The type of signal being generated by the modal test
        \item[signal\_generator\_level] The signal level, which could be a peak value for signals such as a sine wave, or an \ac{RMS} value for random signals
        \item[signal\_generator\_min\_frequency] The minimum frequency of the excitation signal
        \item[signal\_generator\_max\_frequency] The maximum frequency of the excitation signal
        \item[signal\_generator\_on\_fraction] The percent of time that a square wave is ``up'' or a burst random is ``active''.
        \item[exponential\_window\_value\_at\_frame\_end] The value of the exponential window function at the end of the measurement frame, used to set the strength of the exponential window.
        \item[acceptance\_function] If the system is set to automatically reject a measurement, this is the script and function name used to evaluate that signal.
    \end{description}
    
    \subsection{NetCDF Variables}
    \begin{description}
        \item[reference\_channel\_indices] The indices into the channel table that correspond to the reference channels.  Type: 32-bit int; Dimensions: \inlineCode{reference_channels}
        \item[response\_channel\_indices] The indices into the channel table that correspond to the response channels.  Type: 32-bit int; Dimensions: \inlineCode{response_channels}       
    \end{description}

    \subsection{Saving Modal Data}
    In addition to time streaming, Rattlesnake's Modal Testing environment will also save bespoke modal data files directly to the disk when a test is started with the \inlineCode{Start Acquisition} button.  The computed spectral quantities such as \ac{FRF} and coherence data are in a NetCDF file, and are updated at each average.  Time data is also stored to this file, but only the time data that is used for spectral computations.  For example, if a hammer test is being performed and the user takes a minute to evaluate the data before accepting or rejecting the measurement frame, that time data measured during that pause will not be stored to the modal data file; however, it would be stored to the streaming data file described in Section \ref{sec:using_rattlesnake_output_files}, which is essentially an open-tape measurement.
    
    The additional dimension is:
    
    \begin{description}
        \item[fft\_lines] The number of frequency lines in the spectral quantities.
    \end{description}
    
    There are also several additional variables to store the spectral data:
    
    \begin{description}
        \item[frf\_data\_real] The real part of the most recently computed value for the transfer functions between the reference signals and the response signals.  NetCDF files cannot handle complex data types so real and imaginary parts are split into two variables.  Type: 64-bit float; Dimensions: \inlineCode{fft_lines} $\times$ \inlineCode{response_channels} $\times$ \inlineCode{reference_channels}
        
        \item[frf\_data\_imag] The imaginary part of the most recently computed value for the transfer functions between the excitation signals and the control response signals.  NetCDF files cannot handle complex data types so real and imaginary parts are split into two variables.  Type: 64-bit float; Dimensions: \inlineCode{fft_lines} $\times$ \inlineCode{response_channels} $\times$ \inlineCode{reference_channels}
        
        \item[coherence] The multiple or regular coherence of the channels computed during the test, depending on if there are multiple references or a single reference, respectively.  Type: 64-bit float; Dimensions: \inlineCode{fft_lines} $\times$ \inlineCode{response_channels}
        
    \end{description}

    \chapter{Combined Environments}\label{sec:rattlesnake_environments_combining_environments}
    Rattlesnake offers functionality to run multiple events simultaneously.  This can be used to control multiple types of excitation devices (for example a centrifuge along with random vibration) as well as multiple environments using common excitation devices (for example a random vibration and shock environment using a single shaker table).
    
    A combined environments test is selected by selecting the \inlineCode{Combined Environments...} option in the window shown in Figure \ref{fig:environmentselection}.  After selection, a \ac{GUI} window will appear that allows the user to define which environments will be combined for the given test.  The user adds or removes environments from the table of environments displayed in this window by clicking the \inlineCode{Add Environment} or \inlineCode{Remove Selected Environment} buttons.  In the table itself, the type of environment can be chosen using the dropdown menu in the \inlineCode{Environment Type} column.  Each environment must also be given a unique name in the \inlineCode{Environment Name} column.  An example of this window is shown in Figure \ref{fig:combinedenvironmentselector}.  In this example, two environments have been chosen, a \Ac{MIMO} Random Vibration environment called ``Environment 1'' and a Time History Generator environment called ``Environment 2''.
    
    \begin{figure}
        \centering
        \includegraphics[width=0.5\linewidth]{figures/combined_environment_selector}
        \caption{Dialog used to select environments for a combined environments test.}
        \label{fig:combinedenvironmentselector}
    \end{figure}

    Combined environments tests have the potential to be very complex, as parameters for multiple environments need to be defined.  To make the setup of combined environments tests less tedious, Rattlesnake gives the user an option to load a combined environments profile using the \inlineCode{Load Profile...} button.  This profile is an Excel workbook with multiple worksheets defining the channel setup, global data acquisition parameters, environment parameters, as well as the test timeline.  Rattlesnake allows the user to save a profile template to disk, which can be opened and modified to specify the entire combined environments test by clicking the \inlineCode{Save Profile Template...} button.  Note that the template will be created with placeholder environments corresponding to the environments that are displayed in the Environment table, so it is in the user's best interest to add the required environments to the table using the \inlineCode{Add Environment} button and define each environment's type and name prior to saving the template.
    
    Once the environments in the combined environments test are selected, the \ac{GUI} is set up accordingly.  Users will notice the addition of an \inlineCode{Environment Table} on the \inlineCode{Data Acquisition Setup} tab which can be used to determine which excitation and response channels are to be used by each environment.  Excitation signals can be used by one or more environments.  If an excitation signal is used by multiple environments that will be run simultaneously, the excitation signal output to the hardware will end up being a summation of the individual environment's output signals.  It is up to the user to ensure that the combination of environment signals does not exceed the output that the hardware is capable of.  Similarly, response signals can be used by single environments, multiple environments, or no environments.  Only response signals belonging to a specific environment will be passed to that environment for analysis.  Note however that the part may be responding to excitation from multiple environments, so even if a given channel is only used by one environment, it may have response contributions from all active environments.  It is up to the user to design each environment such that this ``cross-talk'' is handled correctly, Rattlesnake will not attempt to separate a given response into contributions from each environment.  If a signal is used by no environments, it will still be recorded and saved to disk; however, data from it will not be delivered to any environment.
    
    With multiple environments defined, the \inlineCode{Environment Definition} and \inlineCode{Run Test} tabs will have a sub-tab to define or control each environment.  On the \inlineCode{Environment Definition} tab, each environment's sub-tab will need to be defined prior to selecting the \inlineCode{Initialize Environments} button.  Additionally, certain environments will have \inlineCode{System Identification} or \inlineCode{Test Prediction} tabs.  System identification will need to be performed for each environment that requires it prior to proceeding with the test setup.
    
    The last major change to the \inlineCode{GUI} for a combined environments test is the inclusion of the \inlineCode{Test Profile} tab which allows users to create a timeline of events that will occur in the combined environments test.  This is described more thoroughly in Section \ref{sec:using_rattlesnake_test_profiles}.
    
    \chapter{Implementing New Environment Types}\label{sec:rattlesnake_environments_implementation}
    
    \textbf{WARNING:} This chapter contains advanced Rattlesnake operations that require a very good understanding of the Rattlesnake software architecture.  Understanding this chapter is not required to run the software successfully, and therefore this chapter can be skipped if the user is not interested in implementing new environments.  Implementing an environment is a serious undertaking that will involve potentially thousands of lines of code depending on its complexity (the \Ac{MIMO} Random Vibration environment is over 3,000 lines of code), so this chapter will only present the basic concepts and minimum required coding, leaving the reader to dig into the code to understand the actual details required to implement the environment.
    
    \section{Overview of Environments}
    In Rattlesnake, an environment is the portion of the code that generates the signal that will be output, perhaps using previously acquired data to inform the next output signal.  In order to use environments interchangeably or simultaneously, Rattlesnake abstracts its environments so the interface between an environment and the main controller is common regardless of the type of environment used.  In general, an environment will consist of three main parts, and each of these parts inherits from an abstract base class defined in the \inlineCode{components/abstract_environment.py} source file in the \GitlabProject{} repository.
    
    \begin{description}
    \item[Environment Metadata or Parameters]Each environment requires defining an object that contains all of the parameters that completely define that environment, e.g. frame lengths, window functions, averaging, control targets, etc.
    \item[Environment-specific \ac{UI}] Each environment requires defining a \ac{UI} into its functionality.  This includes both the actual \ac{GUI} that appears on the screen, as well as the code to handle the operations performed when the user interacts with the \ac{GUI}.
    \item[Environment Implementation] Each environment requires defining its operations.  These operations will involve, for example, computing transfer functions between excitation and responses on the part, computing various metrics, and must at the end determine the next signal that will be output to the excitation devices in the test.
    \end{description}
    
    \section{Overview of Environment Interactions with the Controller}\label{sec:rattlesnake_environments_interactions_with_controller}
    
    Rattlesnake runs most of its tasks on separate processes in order to perform operations simultaneously.  Figure \ref{fig:controllerflowchart} shows a high-level overview of the controller processes.
    
    \begin{figure}
    \centering
    \includegraphics[width=\linewidth]{figures/controller_flowchart}
    \caption{High-level flow chart of the Rattlesnake software}
    \label{fig:controllerflowchart}
    \end{figure}
    
    The \ac{UI} of the controller exists on the main Rattlesnake process.  Therefore, any environment-specific \ac{UI} components are also run on this process.  At run-time after the environments in a given test have been selected, the controller assembles the \ac{GUI}, combining it with the environment-specific interfaces of each environment used in the test.  Typically, the constructor for each environment's \ac{UI} class will handle adding the environment-specific portion of the \ac{GUI} to the main \ac{GUI}.  The class implementing the \ac{UI} should inherit from the \inlineCode{AbstractUI} class in the \inlineCode{components/abstract_environment.py} file.
    
    The class implementing the environment computations should inherit from the \inlineCode{AbstractEnvironment} class in the \inlineCode{components/abstract_environment.py} file.  Each environment is run on a separate process to enable environments to run simultaneously with the acquisition and output processes.  Because of this, environments communicate with the rest of the controller primarily via \inlineCode{queue} objects from the Python \inlineCode{multiprocessing} module.  Messages put to and received from these queues are generally of the form \inlineCode{(message,data)} where \inlineCode{message} specifies the operation that the receiving process should perform, and the \inlineCode{data} is any parameters or input data required by that operation.  For example, a \inlineCode{message} might be \inlineCode{START_CONTROL} and the \inlineCode{data} might be a test level such as \inlineCode{-6} dB.  It is up to the environment to define what operations to perform when it receives the \inlineCode{START_CONTROL} message.
    
    Each environment will interact with many data streams, and therefore will interact with many \inlineCode{queue}s.  Each environment will have an \inlineCode{input_queue} or \inlineCode{command_queue} that it will pull these instructions and data from.  To enable the environments to pass messages back to the controller, a \inlineCode{controller_communication_queue} is provided where the environment can put its own messages and data to be read by the controller.  A \inlineCode{gui_update_queue} also exists for environments to put updates to the \ac{GUI}; this would be used to tell the controller to update a plot or set the value of some widget.
    
    The environment will receive acquired data through a \inlineCode{data_in_queue}, so it can use previously measured data to compute control decisions.  Similarly, the environment will put its next output data set to a \inlineCode{data_out_queue}, which will be read by the output process and put to the data acquisition system through the output hardware interface.
    
    Finally, the environment can put messages to a \inlineCode{log_file_queue} which will be written to the Rattlesnake log file, which is useful for debugging environment behavior.
    
    In general, the environment parameters will be loaded, selected, or otherwise defined in the environment-specific \ac{GUI} in the main process, but need to be passed to the environment implementation running on a separate process.  Rattlesnake requires that these parameters be packaged into an object that contains all the parameters to define the environment.  The class defining this object should inherit from the \inlineCode{AbstractMetadata} class in the \inlineCode{components/abstract_environment.py} file.
    
    \section{Writing an AbstractMetadata Subclass}
    In general, the first step to defining an environment is to define the parameters that define the environment.  These will be specified in a subclass of the \inlineCode{AbstractMetadata} class.
    
    This class is for the most part just a container for the parameters, but it does require the definition of a \inlineCode{store_to_netcdf} method, as the controller needs to know how to take the different parameters defining an environment and store them to a netCDF file as output from Rattlesnake.  For example, in the \Ac{MIMO} Random Vibration environment, this function defines the output structure described in Section \ref{sec:rattlesnake_environments_store_to_netcdf}.  The Streaming process will call this function when it creates a netCDF file for output from Rattlesnake.
    
    In many cases it can be helpful to specify other properties that are derived from the environment parameters as part of this class.  For example, if one environment specifies a ramp time between test levels, one common computation would be to compute the number of time samples this ramping event takes from the ramp time and the sample rate of the data acquisition system.
    
    \section{Designing the Envionment-Specific \ac{GUI}}
    With the parameters defining the environment defined in an \inlineCode{AbstractMetadata} subclass, \ac{GUI} components can be designed that allows a user of the software to specify these parameters and control the environment.  At run-time, these components will be added to the main \ac{GUI} as shown in Figure \ref{fig:addingenvironmentstomainui}.
        
    \begin{figure}
    \centering
    \includegraphics[width=\linewidth]{figures/adding_environments_to_main_ui}
    \caption{Graphical representation of adding a \ac{GUI} element defining the \Ac{MIMO} Random Vibration environment into the main controller \ac{GUI}.}
    \label{fig:addingenvironmentstomainui}
    \end{figure}
    
    These \ac{GUI} elements can be designed using a software such as QtDesigner and loaded at run-time, or assembled programmatically in the constructor for the environment.  Up to four \ac{GUI} components can designed, though two of the components are optional.
    
    \subsection{Environment Definition \ac{GUI}}
    The first \ac{GUI} component that should be created is the interface in the Environment Definition tab.  Essentially, the author of the environment should define a \inlineCode{GUI} element that contains all the parameters to define a specific environment, including signal processing parameters and the control targets.  This is a required portion of the \ac{GUI} as all the environment parameters must be defined before the controller can be run.
    
    \subsection{System Identification and Test Predictions \ac{GUI}}
    The second and third \ac{GUI} components are system identification and test predictions.  These are optional, as not all environments will require system identification (e.g. the Time History Generator environment described in Chapter \ref{sec:rattlesnake_environments_time_generator}).
    
    The system identification portion of the \inlineCode{GUI} is generally common between different controller types, as techniques for developing a transfer function between excitation and response are fairly universal.  The user should be able to start and stop acquisition of transfer function data, as well as visualize the transfer functions as they are computed.  Both the \Ac{MIMO} Random and \Ac{MIMO} Transient environments utilize a common system identification framework, which is defined in the \inlineCode{abstract_sysid_environment.py} and \inlineCode{abstract_sysid_data_analysis.py}.  If users wish to utilize this system identification framework for their own environments, they are welcome to make their environment inherit from this class, which will save a significant amount of coding to write a new system identification framework.
    
    When the system identification is performed, generally an environment can compute a prediction of the responses that could be achieved.  Thought should be given as to the best way to present the predicted responses on the Test Predictions \ac{GUI} component.  The user must be able to determine if a given test is feasible, so they should be aware of the required excitation as well as the accuracy the environment expects to achieve.
    
    \subsection{Environment Run \ac{GUI}}
    
    The final \ac{GUI} component controls running the test.  Here the user should be able to start and stop the environment, as well receive a visual of how well the environment is controlling the response of the test article.  Generally this will entail some type of overview plot or table.  This might also provide the functionality to visualize individual response or excitation channels.
    
    \subsection{Writing an AbstractUI Subclass}
    
    With the \ac{GUI} components defined, the code containing the \inlineCode{UI} logic and callbacks can be written.  This will generally inherit from the \inlineCode{AbstractUI} class in the \inlineCode{components/abstract_environment.py} file.
    
    \subsubsection{Abstract Methods in the AbstractUI class}
    
    Because it inherits from the \inlineCode{AbstractUI} base class, classes defining environment-specific \ac{UI}s in Rattlesnake require definition of a number abstract methods:
    
    \begin{description}
    \item[\_\_init\_\_] The class constructor that calls the superclass constructor, adds the \ac{GUI} elements to the main controller's \ac{GUI} tabs, and connects callback functions to the \ac{GUI} widgets.  This function will also set up the command map as described below.
   
    \item[collect\_environment\_definition\_parameters] This function gathers the environment parameters from the current values of the \ac{GUI} widgets and constructs the corresponding \inlineCode{AbstractMetadata} subclass object.
    
    \item[initialize\_data\_acquisition] This function is called when the Data Acquisition parameters are initialized (channel table, sample rate, etc.).  It should set up the environment user interface accordingly.
    
    \item[initialize\_environment] This function is called when the Environment Parameters are initialized.  It should set up the user interface and environment appropriately based off parameters specified in the environment's definition \ac{GUI} components.  It must also return an object of the \inlineCode{AbstractMetadata} subclass corresponding to the present environment, so the controller can send it to the environment.
    
    \item[retrieve\_metadata] This function retrieves parameters from a netCDF dataset that was written by the controller during streaming.  It must populate the widgets in the user interface with the proper information.  Where the \inlineCode{store_to_netcdf} function in the \inlineCode{AbstractMetadata} subclass writes parameters to the netCDF file, this function must be able to read that data back into the controller.
    
    \item[create\_environment\_template] This function creates a template worksheet in an Excel workbook where users can specify environment parameters that are loaded in via the \inlineCode{Load Profile...} button on the \inlineCode{Environment Selector} dialog box.  This function should be defined using the \inlineCode{@staticmethod} decorator.
    
    \item[set\_parameters\_from\_template] This function reads the template generated by the \inlineCode{create_environment_template} function after a user has filled in the required information.
    
    \item[start\_control] The function containing the logic to start the environment controlling.
    
    \item[stop\_control] The function containing the logic to gracefully shut down the environment.
    
    \item[update\_gui] This function receives data from the \inlineCode{gui_update_queue} that specifies how the user interface should be updated.  Data is typically recieved as \inlineCode{(instruction,data)} pairs stored as a tuple, where the \inlineCode{instruction} specifies which operation to perform or which widget to modify, and \inlineCode{data} provides additional information to use in the operation (a signal to plot, the value to update a widget to, etc.).
    \end{description}
    
    \subsubsection{Command Map for the \Ac{UI}}
    
    As described in section \ref{sec:using_rattlesnake_test_profiles}, Rattlesnake can automate a sequence of operations for combined environments tests.  To define such a sequence, the user creates a table of operations that occur at a given time in the test.  Each environment must therefore define which operations it can perform, as well as defining what happens when each of the operations is triggered.  To define these profile operations, each environment-specific \ac{UI} class contains a \inlineCode{command_map} dictionary that maps strings to functions in the class.  The main controller \ac{UI} reads the strings in this \inlineCode{command_map} object to populate the drop-down lists in the \inlineCode{Test Profile} table.  Then, when one of the operations is triggered, the controller calls the function mapped to that string in the \inlineCode{command_map} object.  In this way, the main controller \ac{UI} needs only to access the \inlineCode{command_map} object to know which operations can be automated using the test profile.  Figure \ref{fig:testprofilecommandmap} shows how specific entries in the \inlineCode{command_map} translate to operations in the \inlineCode{Test Profile} table.
    
    Note that the \inlineCode{AbstractUI} class defines two entries in the \inlineCode{command_map}, ``Start Control'' and ``Stop Control'', which map to the \inlineCode{start_control} and \inlineCode{stop_control} functions that are required to be defined.  Because new environment \ac{UI} classes will inherit from \inlineCode{AbstractUI}, they will also have these operations defined in their \inlineCode{command_map}.  Therefore, every environment will be able to be started and stopped via the \inlineCode{Test Profile} functionality.  Authors of new environments are then able to add additional operations to their subclasses if required.
    
    \begin{figure}
    \centering
    \includegraphics[width=\linewidth]{figures/test_profile_command_map}
    \caption{The \inlineCode{command_map} dictionary in each \inlineCode{AbstractUI} subclass maps functions in the subclass to strings that will appear in the \inlineCode{Test Profile} tab as available operations that can be automated in a combined environments test.}
    \label{fig:testprofilecommandmap}
    \end{figure}
    
    \section{Writing an AbstractEnvironment Subclass}
    The environment implementation is what performs the control calculations to generate the next output signal required by the controller.  It generally is started as a separate process from the main controller and communicates to the main controller using queues.  This enables multiple environments to be run simultaneously for combined environments tests.
    
    The environment implementation will generally inherit from the \inlineCode{AbstractEnvironment} abstract base class.  The \inlineCode{AbstractEnvironment} defines the minimum functionality required for an environment to interact with the controller; the author of a new environment is expected to significantly expand on this minimum functionality to produce a useful and functional environment in Rattlesnake.
    
    The \inlineCode{AbstractEnvironment} class defines the following methods:
    
    \begin{description}
    \item[log] Writes a message to the \inlineCode{log_file_queue} so it will eventually be put to the log file along with a timestamp of the message and the name of the environment putting the message to the queue.  The parallel nature of Rattlesnake means that messages written to the log file might be out of order depending on what is happening when the messages are written, so it is important to include the source and time of the message so operations can be reconstructed.
    \item[run] The main looping function that defines the environment's subprocess.  It continually waits for \inlineCode{(message,data)} instructions from the \inlineCode{command_queue} and then uses a command map to map the \inlineCode{message}s to functions defined in the class, passing the \inlineCode{data} as the argument to the function: \inlineCode{function = command_map[message]; output = function(data)}.  This function will continue to run until the \inlineCode{output} from one of the \inlineCode{command_map} functions returns \inlineCode{True}.
    \item[quit] Returns \inlineCode{True} to kick the class out of the loop in the \inlineCode{run} function so the process eventually exits.  This function can be overwritten by an inheriting subclass to perform extra cleanup or shutdown operations.
    \item[map\_command] A function that maps a \inlineCode{message} that will be received by the \inlineCode{command_queue} to a function that will be called in the class.
    \end{description}
    
    The \inlineCode{AbstractEnvironment} class also specifies the following abstract methods that must be defined by any subclasses:
    
    \begin{description}
    \item[initialize\_data\_acquisition\_parameters] Gets the data acquisition parameters (sample rate, channel table) from the controller and sets the environment up appropriately.
    \item[initialize\_environment\_test\_parameters] Gets the environment parameters (as an object of a subclass of \inlineCode{AbstractMetadata}) and sets up the enviornment appropriately
    \item[stop\_environment] Defines operations to gracefully shut down the environment
    \end{description}    
    
    Note that none of the functions above have anything to do with the actual generation of the signal in the environment.  As described previously, the above is the minimum requirements to interact with the controller.  It is up to the author to determine and implement the best approach to actually generate the signals and put them to the \inlineCode{data_out_queue}.  It is recommended that authors of new environments look at existing environments in Rattlesnake to understand the flow of information in those environments.
    
    As described above, the environment will sit in a loop in the \inlineCode{run} function waiting for \inlineCode{(message,data)} pairs to come down the \inlineCode{command_queue}.  Therefore, in order for any of the other functions listed above to be called, they should be mapped to a specific instruction \inlineCode{message}, or otherwise called from within a function that is mapped to a specific instruction \inlineCode{message}.  The environment classes that inherit from \inlineCode{AbstractEnvironment} utilize a similar \inlineCode{command_map} as the environment-specific \ac{UI} classes that inherit from \inlineCode{AbstractUI}.  However, instead of mapping strings to functions, the \inlineCode{command_map} for the actual environment implementation maps \inlineCode{message}s to functions.  Therefore, when a specific \inlineCode{message} is received, the \inlineCode{run} function queries the \inlineCode{command_map} to get the function mapped to that \inlineCode{message}.  That function is then called with \inlineCode{data} as the argument.
    
    Note that when creating a new process, a function is specified that defines the operations that will be performed on the new process.  Therefore a function must be defined that instantiates an object from the environment class and calls the \inlineCode{run} function of that object.  Any environment-specific set-up or clean-up operations should be performed in this function as well.
    
    \section{Connecting the Environment to the Controller}
    To get Rattlesnake to recognize a new environment, it must be added to a number of data structures in the \inlineCode{components/environments.py} file.  A ``short name'' for the environment should be added to the \inlineCode{ControlTypes} enumeration.  A ``long name'' for the environment can be added to the \inlineCode{environment_long_names} dictionary.  If the environment should be usable in combined environments, the corresponding \inlineCode{ControlTypes} enum should be added to the \inlineCode{combined_environments_capable} list.  If \inlineCode{*.ui} files from QtDesigner are used to define environment \ac{UI}s, the paths should be added to the corresponding \inlineCode{environment_definition_ui_paths}, \inlineCode{environment_prediction_ui_paths}, and \inlineCode{environment_run_ui_paths} dictionaries.  Finally the \ac{UI} class and target function that will be called by processes utilizing the given environment should be imported and assigned to the respective data containers.

    \bookmarksetup{startatroot}

    % ---------------------------------------------------------------------- %
    % References
    %
    \clearpage
    % If hyperref is included, then \phantomsection is already defined.
    % If not, we need to define it.
    \providecommand*{\phantomsection}{}
    \phantomsection
%    \addcontentsline{toc}{chapter}{References}
    \bibliographystyle{plain}
    \bibliography{Z:/Optical_Modal_PhD_Work/00_Literature_Review/uur_bibliography}


    % ---------------------------------------------------------------------- %
    %
    \appendix
%    \renewcommand{\partname}{Appendices Part}
%    \setcounter{part}{0}
%    
%    \part{Example Problems}\label{sec:example_problem}
    
    \chapter{Introduction to Example Problems}
    There is a significant amount of information contained in this document, which can be incredibly difficult to digest for a new user of the software.  These appendices will therefore contain example problems that users can implement either completely synthetically or using minimal hardware resources to try out Rattlesnake.  These can be used as a kind of quick-start guide to using Rattlesnake.
    
    Appendix \ref{sec:example_nidaqmx} demonstrates how a test can be run using relatively inexpensive NI hardware.
    
    Appendix \ref{sec:example_sdynpy} demonstrates how the same test from Appendix \ref{sec:example_nidaqmx} can be simulated synthetically using SDynPy \inlineCode{System} objects.  This example problem would be a good place to start for users without hardware, as it only relies on software.
    
    Appendix \ref{sec:example_state_space} demonstrates how the same test from Appendix \ref{sec:example_nidaqmx} can be simulated synthetically using State Space matrices.  State space matrices are more flexible than SDynPy \inlineCode{System} objects to represent generic linear time-invariant systems.
    
    \chapter{Example Problem using NI cDAQ Hardware and the NI DAQmx Interface}\label{sec:example_nidaqmx}
    
    This appendix will demonstrate a example problem using real hardware on a simple beam structure.  This example problem will use a simple cDAQ hardware.  A NI cDAQ-9185 chassis is used with two NI-9232 acquisition modules and two NI-9260 output modules.  Figure \ref{fig:examplecdaqhardware} shows this hardware setup.
    
    \begin{figure}[H]
        \centering
        \includegraphics[width=0.7\linewidth]{figures/example_cdaq_hardware}
        \caption{NI cDAQ hardware used for this demonstration.}
        \label{fig:examplecdaqhardware}
    \end{figure}
    
    This example problem will start with a modal test using an impact hammer, then demonstrate a modal test with a shaker, and finally perform a MIMO Random and MIMO Transient test.
    
    The beam being tested is 24 inches long by 1 inch high by 0.75 inches deep and made from steel. It is resting on soft foam to simulate a free-free boundary condition. We have six acquisition channels, so we will use four of those channels to measure response, and the other two to measure force and voltage signals as required by the testing configuration.  Figure \ref{fig:examplecdaqbeamstructure} shows the beam, as well as a block with the coordinate axes labeled.
    
    \begin{figure}[H]
        \centering
        \includegraphics[width=0.7\linewidth]{figures/example_cdaq_beam_structure}
        \caption{Beam structure that will be tested in this example problem.}
        \label{fig:examplecdaqbeamstructure}
    \end{figure}
    
    \section{Modal Testing with an Impact Hammer}
    
    The first demonstration will be a simple impact test using a modal hammer.  This will require no output from Rattlesnake.
    
    \subsection{Data Acquisition Setup}
    
    We will excite the beam in the thin direction, so we will place four uniaxial accelerometers on the beam at the 1, 12, 18, and 23 inch distances from the left end. We will name these locations 1, 12, 18, and 23, respectively.
    
    As we install instrumentation, we will fill out a channel table. Rather than filling out the channel table in the GUI, we will load it from an Excel spreadsheet, which allows us to more easily save our work. We can get a template channel table to fill out by opening Rattlesnake and selecting the Modal Testing environment (Figure \ref{fig:examplecdaqmodaltestingselection}).
    
    \begin{figure}[H]
        \centering
        \includegraphics[width=0.5\linewidth]{figures/example_cdaq_modal_testing_selection}
        \caption{Selecting the Modal Environment in Rattlesnake}
        \label{fig:examplecdaqmodaltestingselection}
    \end{figure}
    
    We can then click the \inlineCode{Save Channel Table} button to save out the empty channel table to an Excel spreadsheet (Figure \ref{fig:examplecdaqsavechanneltable}).
    
    \begin{figure}[H]
        \centering
        \includegraphics[width=0.7\linewidth]{figures/example_cdaq_save_channel_table}
        \caption{Saving out an empty channel table to fill out for the test.}
        \label{fig:examplecdaqsavechanneltable}
    \end{figure}
    
    As we add instrumentation and plug it into our data acquisition system, we will record that information in the channel table. In order to tell Rattlesnake the correct device, we open NI Measurement and Automation Explorer (NI MAX). In the Devices and Interfaces section, we can see our network devices, which includes our cDAQ system and its associated modules (Figure \ref{fig:examplecdaqnimax}).
    
    \begin{figure}[H]
        \centering
        \includegraphics[width=0.5\linewidth]{figures/example_cdaq_nimax}
        \caption{NI MAX dialog showing the cDAQ devices.}
        \label{fig:examplecdaqnimax}
    \end{figure}
    
    Our acquisition channels will be plugged into devices \inlineCode{cDAQ9185-217ED78Mod1} and \inlineCode{cDAQ9185-217ED78Mod2} which have acquisition channels \inlineCode{ai0}, \inlineCode{ai1}, and \inlineCode{ai2}, and our output devices are \inlineCode{cDAQ9185-217ED78Mod3} and \inlineCode{cDAQ9185-217ED78Mod4} which have output channels \inlineCode{ao0} and \inlineCode{ao1}.
    
    After we add our accelerometers, we will also add a channel for the modal hammer. Note that we will vary the hammer excitation location, so we will simply put in an arbitrary node number that we will override later. We will hit the beam on the surface opposite the side that the accelerometers are mounted on, so the direction of the impact is the identical to the directions of the accelerometers.
    
    When we finish our channel table it looks like Figure \ref{fig:examplecdaqchanneltableexcel}.  Note that there is no output from Rattlesnake, so there are no entries in the Output Feedback in the channel table.
    
    \begin{figure}[H]
        \centering
        \includegraphics[width=\linewidth]{figures/example_cdaq_channel_table_excel}
        \caption{Channel Table in Excel ready to load into Rattlesnake.}
        \label{fig:examplecdaqchanneltableexcel}
    \end{figure}
    
    The test setup after attaching the accelerometers and plugging in the instrumentation is shown in Figure \ref{fig:examplecdaqtestsetupmodalimpact}.
    
    \begin{figure}[H]
        \centering
        \includegraphics[width=0.7\linewidth]{figures/example_cdaq_test_setup_modal_impact}
        \caption{Modal test setup using the impact hammer.}
        \label{fig:examplecdaqtestsetupmodalimpact}
    \end{figure}
    
    We can now load our channel table into Rattlesnake by clicking the Load Channel Table button shown in Figure \ref{fig:examplecdaqloadchanneltable}.
    
    \begin{figure}[H]
        \centering
        \includegraphics[width=0.7\linewidth]{figures/example_cdaq_load_channel_table}
        \caption{Loading in the channel table for the test.}
        \label{fig:examplecdaqloadchanneltable}
    \end{figure}
    
    We should see that the information from the Excel spreadsheet is populated in the Rattlesnake GUI.
    
    In the \inlineCode{Data Acquisition Parameter} section we will leave \inlineCode{NI DAQmx} selected. We will set the \inlineCode{Sample Rate} to \inlineCode{5120} Hz, as this is a valid output rate for both the cDAQ device types we are using. We can leave the \inlineCode{Time per Read} and \inlineCode{Time per Write} to their defaults.  Figure \ref{fig:examplecdaqmodaltestimpactdaqsettings} shows the loaded channel table and data acquisition settings.
    
    \begin{figure}[H]
        \centering
        \includegraphics[width=\linewidth]{figures/example_cdaq_modal_test_impact_daq_settings}
        \caption{Data acquisition settings for the Modal Test using an impact hammer.}
        \label{fig:examplecdaqmodaltestimpactdaqsettings}
    \end{figure}
    
    At this point, we can click the \inlineCode{Initialize Data Acquisition} button shown in Figure \ref{fig:examplecdaqinitializedataacquisition} to proceed to the next stage of the test.
    
    \begin{figure}[H]
        \centering
        \includegraphics[width=0.5\linewidth]{figures/example_cdaq_initialize_data_acquisition}
        \caption{Click the \inlineCode{Initialize Data Acquisition} button to proceed.}
        \label{fig:examplecdaqinitializedataacquisition}
    \end{figure}
    
    \subsection{Environment Setup}
    
    At this point, we now need to set up our modal testing environment on the Environment Definition tab. As this is a simple beam, it will likely take a while to decay, so we will set up a 2 second measurement frame by setting the \inlineCode{Samples per Frame} to \inlineCode{10240}.
    
    We will lower the \inlineCode{Number of Averages} to \inlineCode{5}.
    
    For \inlineCode{Triggering}, we will set the \inlineCode{Triggering Type} to \inlineCode{Every Frame}, and set the \inlineCode{Acceptance} to \inlineCode{Manual}. We set the \inlineCode{Trigger Channel} to the Force channel with a \inlineCode{2%} \inlineCode{Pretrigger}. The \inlineCode{Trigger Level} we will set to \inlineCode{1%} and the \inlineCode{Hysteresis Level} to half of that value. Finally, we will set the \inlineCode{Hysteresis Frame %} to \inlineCode{50%}.
    
    In the \inlineCode{Channel Selections} table, we will check the \inlineCode{Reference} column for the Force channel.
    
    When all of these parameters are adjusted, we can click the \inlineCode{Initialize Environments} button to proceed.  Figure \ref{fig:examplecdaqmodalimpactsettings} shows the window with all settings set.
    
    \begin{figure}[H]
        \centering
        \includegraphics[width=\linewidth]{figures/example_cdaq_modal_impact_settings}
        \caption{Modal Test parameters for the impact test.}
        \label{fig:examplecdaqmodalimpactsettings}
    \end{figure}
    
    \subsection{Test Profile}
    We will leave the \inlineCode{Test Profile} tab empty. Simply click the \inlineCode{Initialize Profile} button to proceed on this tab.
    
    \subsection{Run Test}
    
    We will save the standard modal data file from this test, so we won't worry about setting up streaming.
    
    To enable the Modal Environment, first click the \inlineCode{Arm Data Acquisition} button. At this point, the modal environment buttons will become enabled, as shown in \ref{fig:examplecdaqmodaltestingready}
    
    \begin{figure}[H]
        \centering
        \includegraphics[width=\linewidth]{figures/example_cdaq_modal_testing_ready}
        \caption{Run Test tab with data acquisition armed}
        \label{fig:examplecdaqmodaltestingready}
    \end{figure}
    
    To start, we will use the \inlineCode{Degree of Freedom Override} section to choose our actual excitation location.  This allows us to more quickly adjust the impact location than returning to the \inlineCode{Data Acquisition Setup} page and reinitializing everything in the channel table. Click the \inlineCode{Add Override} button, and in the row that appears, change the \inlineCode{Channel} column to the force channel. In the \inlineCode{New Node} column, put \inlineCode{23}, which will be the position we excite at. For the \inlineCode{New Direction}, put \inlineCode{Y-}.  This is shown in Figure \ref{fig:examplecdaqchanneloverride}.
    
    \begin{figure}[H]
        \centering
        \includegraphics[width=0.5\linewidth]{figures/example_cdaq_channel_override}
        \caption{Using the Degree of Freedom Override to adjust the impact location metadata.}
        \label{fig:examplecdaqchanneloverride}
    \end{figure}
    
    Now with our channels correct, we can set up data display windows. We will use the \inlineCode{New from Template...} dropdown menu to create \inlineCode{Drive Point (Imaginary)}, \inlineCode{Drive Point Coherence}, and \inlineCode{Reference Autospectrum} channels. We will also manually create 5 windows by clicking the \inlineCode{New} button, and set them to visualize the time histories for each channel; these will help us determine if we need to change the trigger settings, frame length, or add a window. Clicking the \inlineCode{Tile} button will expand the windows to fill up the available space.  Figure \ref{fig:examplecdaqmodalimpactemptywindows} shows the result of these operations.
    
    \begin{figure}[H]
        \centering
        \includegraphics[width=\linewidth]{figures/example_cdaq_modal_impact_empty_windows}
        \caption{Run Test tab with modal windows created.}
        \label{fig:examplecdaqmodalimpactemptywindows}
    \end{figure}
    
    We will first preview the measurement to ensure that all of the parameters we have selected are adequate. Click the \inlineCode{Preview Acquisition} button and impact the structure behind the 23Y- accelerometer. If the data acquisition triggers, data should appear in the windows as shown in Figure \ref{fig:examplecdaqtrigger}.
    
    \begin{figure}[H]
        \centering
        \includegraphics[width=\linewidth]{figures/example_cdaq_trigger}
        \caption{Data acquisition triggered, waiting for acceptance or rejection of the measurement.}
        \label{fig:examplecdaqtrigger}
    \end{figure}
    
    After the data appears, the \inlineCode{Accept} and \inlineCode{Reject} buttons should become active. If the signal decays within the window and the force looks reasonable, we can click \inlineCode{Accept} which will then proceed to the next measurement. After a measurement is accepted, averaged quantities such as coherence, frequency response functions, and autopower spectra will also be computed as shown in Figure \ref{fig:examplecdaqfirstaccept}. Only Spectra (FFTs) and time histories can be visualized prior to accepting a measurement.
    
    \begin{figure}[H]
        \centering
        \includegraphics[width=\linewidth]{figures/example_cdaq_first_accept}
        \caption{Spectral quantities appear after acceptance.}
        \label{fig:examplecdaqfirstaccept}
    \end{figure}
    
    After a handful of measurements are taken, we should be able to see the coherence and frequency response functions stabilize. We can evaluate the measurement by investigating the drive point FRF, coherence, reference autospectra, or any other quantities of interest.  Figure \ref{fig:examplecdaqallaccept} shows these results.
    
    \begin{figure}[H]
        \centering
        \includegraphics[width=\linewidth]{figures/example_cdaq_all_accept}
        \caption{Quantities of interest after 5 averages.}
        \label{fig:examplecdaqallaccept}
    \end{figure}
    
    If this looks satisfactory, we can proceed with the measurement. Clicking the \inlineCode{Stop Acquisition} button will stop the preview, and allow us to start a real measurement. We must assign a \inlineCode{Data File} to save to by clicking the \inlineCode{Select File...} button. We will call our file \inlineCode{beam_hammer_impact.nc4}. To ensure we do not overwrite our data as we rove the hammer, we will click the \inlineCode{Auto-Increment} checkbox to automatically increment the filename.  This is shown in Figure \ref{fig:examplecdaqfilename}.
    
    \begin{figure}[H]
        \centering
        \includegraphics[width=0.5\linewidth]{figures/example_cdaq_filename}
        \caption{Setting a file name for the modal data.}
        \label{fig:examplecdaqfilename}
    \end{figure}
    
    With the data file selected, we can click the \inlineCode{Start Acquisition} button. We perform the measurement identically to the preview case. When the specified number of averages have been acquired, the measurement will stop automatically as shown in Figure \ref{fig:examplecdaqmodalimpactsacquired}.
    
    \begin{figure}[H]
        \centering
        \includegraphics[width=\linewidth]{figures/example_cdaq_modal_impacts_acquired}
        \caption{Acquisition stops automatically after the desired number of impacts.}
        \label{fig:examplecdaqmodalimpactsacquired}
    \end{figure}
    
    Now we will acquire the next measurement at node 18. We will update the \inlineCode{New Node} column of the \inlineCode{Degree of Freedom Override} table to be \inlineCode{18}. Note that this will also change the labels on the data windows. We will need to update the response channels on those windows to make sure we are still visualizing the drive point data, as shown in Figure \ref{fig:examplecdaqupdatewindows}.
    
    \begin{figure}[H]
        \centering
        \includegraphics[width=0.5\linewidth]{figures/example_cdaq_update_windows}
        \caption{Updating the signal being visualized in the data visualization window.}
        \label{fig:examplecdaqupdatewindows}
    \end{figure}
    
    We can simply press the \inlineCode{Start Acquisition} button, and a new file will be created. After that measurment has been obtained, we will then repeat again, updating the \inlineCode{Degree of Freedom Override} table to node \inlineCode{12} and start a measurement exciting at that location.
    
    When these measurements are complete, we can click \inlineCode{Disarm Data Acquisition} to stop the measurement, and close the Rattlesnake software.
    
    \FloatBarrier
    \subsection{Analyzing Modal Data}
    
    To analyze these data, we will use SDynPy\footnote{SDynPy can be installed via pip using \inlineCode{pip install sdynpy}, or otherwise downloaded from the Github repository here: \url{https://github.com/sandialabs/sdynpy}}, which is an open-source Python-based structural dynamics toolset. SDynPy is well-integrated with Rattlesnake, making it trivial to load in data from a Rattlesnake test.
    
    For example, reading in a modal dataset with SDynPy looks like:
    \begin{lstlisting}[language=Python]
# Import SDynPy and Numpy
import sdynpy as sdpy
import numpy as np

# Read a rattlesnake file using `read_modal_data`
time_data, frfs, coherence, channel_table = sdpy.rattlesnake.read_modal_data(
'beam_hammer_impact_0000.nc4')\end{lstlisting}
    
    We can plot the data easily using SDynPy's plotting functionality:
    \begin{lstlisting}[language=Python]
frfs.plot(one_axis=False,subplots_kwargs={'layout':'constrained'})\end{lstlisting}
    
    The results of this command are shown in Figure \ref{fig:examplecdaqplottedfrfs}.
    
    \begin{figure}[H]
        \centering
        \includegraphics[width=0.7\linewidth]{figures/example_cdaq_plotted_frfs}
        \caption{Frequency response functions plotted in SDynPy from the modal impact test.}
        \label{fig:examplecdaqplottedfrfs}
    \end{figure}
    
    We can combine and plot the frequency response functions from our multiple datasets:
    \begin{lstlisting}[language=Python]
all_frfs = []
for i in range(3):
    _, frfs, _, _ = sdpy.rattlesnake.read_modal_data('beam_hammer_impact_{:04}.nc4'.format(i))
all_frfs.append(frfs)
all_frfs = np.concatenate(all_frfs)
all_frfs.plot(one_axis=False,subplots_kwargs={'layout':'constrained','figsize':(10,6)})\end{lstlisting}
    
    The results of this command are shown in Figure \ref{fig:examplecdaqallfrfs}.
    
    \begin{figure}[H]
        \centering
        \includegraphics[width=\linewidth]{figures/example_cdaq_all_frfs}
        \caption{All FRFs from the three impacts}
        \label{fig:examplecdaqallfrfs}
    \end{figure}
    
    Ideally we would like to plot mode shapes of this test, so as a first step towards that goal we will construct a test geometry. This contains nodes, coordinate systems, and optionally tracelines or elements.  We can construct nodes with the ID numbers given in the test and coordinates specified by each node's position along the beam.  Here we construct and print a \inlineCode{node_array} in SDynPy
    \begin{lstlisting}[language=Python]
nodes = sdpy.node_array(id = [1,12,18,23],
                        coordinate = [[1,0,0],
                                      [12,0,0],
                                      [18,0,0],
                                      [23,0,0]])
print(nodes)\end{lstlisting}
    which has output
    \begin{lstlisting}
    Index,     ID,        X,        Y,        Z, DefCS, DisCS
     (0,),      1,    1.000,    0.000,    0.000,     1,     1
     (1,),     12,   12.000,    0.000,    0.000,     1,     1
     (2,),     18,   18.000,    0.000,    0.000,     1,     1
     (3,),     23,   23.000,    0.000,    0.000,     1,     1\end{lstlisting}
    
    We will do similarly with a global coordinate system
    \begin{lstlisting}[language=Python]
css = sdpy.coordinate_system_array(id=1,name='global')
print(css)\end{lstlisting}
    which has output
    \begin{lstlisting}
   Index,     ID,                 Name, Color,       Type
      (),      1,               global,     1,  Cartesian\end{lstlisting}
    
    We can then construct a Geometry object. We will then add a traceline to connect the nodes together.  We will then plot the geometry.
    \begin{lstlisting}[language=Python]
geometry = sdpy.Geometry(nodes,css)
geometry.add_traceline([1,12,18,23])
geometry.plot()\end{lstlisting}
    
    The plotted geometry is shown in Figure \ref{fig:examplecdaqgeometry}.
    
    \begin{figure}[H]
        \centering
        \includegraphics[width=0.7\linewidth]{figures/example_cdaq_geometry}
        \caption{Geometry used to plot mode shapes in the modal test.}
        \label{fig:examplecdaqgeometry}
    \end{figure}
    
    Now that we have a geometry, we can fit modes and plot mode shapes. We will only fit up to approximately 2000 Hz, which is where the antialiasing filter of the cDAQ device is clearly active.  We will use the \inlineCode{PolyPy_GUI} in SDynPy to do the mode fits.  We can add our geometry directly to the fitter by assigning to the \inlineCode{geometry} attribute.
    \begin{lstlisting}[language=Python]
ppgui = sdpy.PolyPy_GUI(all_frfs)
ppgui.geometry = geometry\end{lstlisting}
    
    We can select frequency ranges and polynomial orders to solve for, then click \inlineCode{Compute Stabilization} to proceed, as shown in Figure \ref{fig:examplecdaqpolypyranges}.
    
    \begin{figure}[H]
        \centering
        \includegraphics[width=0.7\linewidth]{figures/example_cdaq_polypy_ranges}
        \caption{Setting up frequency ranges and polynomial orders in PolyPy.}
        \label{fig:examplecdaqpolypyranges}
    \end{figure}
    
    We can then select stable poles from the diagram and plot resynthesized FRFs or mode indicator functions.  This operation and the results are shown in Figure \ref{fig:examplecdaqpolypyselectpoles}.
    
    \begin{figure}[H]
        \centering
        \includegraphics[width=0.7\linewidth]{figures/example_cdaq_polypy_select_poles}
        \includegraphics[width=\linewidth]{figures/example_cdaq_polypy_resynthesis}
        \caption{Selecting stable poles and visualizing resynthesized FRFs.}
        \label{fig:examplecdaqpolypyselectpoles}
    \end{figure}
    
    We can then plot mode shapes from the test.  Note that with only four sensors, our mode shapes will not be very highly resolved; however, they clearly correspond to the first three bending modes of the beam.  The first mode shape is shown in Figure \ref{fig:examplecdaqmodalfits}.
    
    \begin{figure}[H]
        \centering
        \includegraphics[width=0.7\linewidth]{figures/example_cdaq_modal_fits}
        \caption{First mode shape fit to the modal data from the impact test.}
        \label{fig:examplecdaqmodalfits}
    \end{figure}
    
    \section{Modal Testing with a Modal Shaker}\label{sec:example_nidaqmx_modal_shaker}
    
    Next we will demonstrate how to run a modal test with a shaker. A cap is placed over the accelerometer at node 1, and a force sensor is threaded into the cap. A receptacle is mounted to the force gauge to accept the shaker stinger.
    
    \subsection{Data Acquisition Settings}
    
    Recall that because Rattlesnake is required to measure the voltage associated with any drive signals it outputs, adding a shaker to a modal test requires two channels: one for the force sensor to measure the reference force, and one to measure the shaker voltage. Therefore, if we aim to measure the force for this test, we can only accommodate a single shaker input, as four of the six available channels channels are already taken by the accelerometers.
    
    When we set up the shaker, we put a tee into the output channel from the data acquisition system, and we tee off the shaker drive signal into the last acquisition channel and to the shaker amplifier, as shown in Figure \ref{fig:examplecdaqmodalshakerhardwaresetup}.  Figure \ref{fig:examplecdaqmodalshakersetup} shows a wider view with the modal shaker visible.
    
    \begin{figure}[H]
        \centering
        \includegraphics[width=0.7\linewidth]{figures/example_cdaq_modal_shaker_hardware_setup}
        \caption{Data acquisition setup for the shaker test showing the shaker signal teed into the last acquisition channel.}
        \label{fig:examplecdaqmodalshakerhardwaresetup}
    \end{figure}
    
    \begin{figure}[H]
        \centering
        \includegraphics[width=0.7\linewidth]{figures/example_cdaq_modal_shaker_setup}
        \caption{Modal shaker setup for running a modal test.}
        \label{fig:examplecdaqmodalshakersetup}
    \end{figure}
    
    As we fill out the channel table, we enter the physical device and channel that the shaker voltage is teed into in the \inlineCode{Channel Definition} section of the channel table, and enter the physical device and channel that the shaker voltage originates from in the \inlineCode{Output Feedback} section. In our case, the shaker signal comes from the \inlineCode{ao0} channel on the third module \inlineCode{cDAQ9185-217ED78Mod3}. The shaker signal is teed into acquisition channel \inlineCode{ai2} on the second module \inlineCode{cDAQ9185-217ED78Mod2}. We also enter the force sensor into the channel table.  note that the direction associated with the force is now in the \inlineCode{Y+} direction, as the push direction of the shaker is opposite the measurement direction of the accelerometer that is covered by the cap.  Figure \ref{fig:examplecdaqmodalshakerchanneltable} shows the channel table set up in Excel.
    
    \begin{figure}[H]
        \centering
        \includegraphics[width=\linewidth]{figures/example_cdaq_modal_shaker_channel_table}
        \caption{Channel table set up for the modal shaker test.}
        \label{fig:examplecdaqmodalshakerchanneltable}
    \end{figure}
    
    Note that when assigning a node and direction to the voltage signal, you do not want to use the same values as for the drive point, otherwise Rattlesnake will not be able to determine which frequency response function is the drive point FRF. The author's practice is to use the same node number as the force sensor, but offset by some value (e.g. 100 in this case). In the case of multiple shakers, it will be clear which voltage corresponds to which shaker if such a convention is followed.
    
    With the channel table set up, we can now start Rattlesnake and load the channel table. We will again select the Modal Testing environment, click the \inlineCode{Load Channel Table} button and load our file. We will again set the \inlineCode{Sample Rate} to \inlineCode{5120} Hz. With these parameters set, we can click the \inlineCode{Initialize Data Acquisition} button to proceed.  Figure \ref{fig:examplecdaqmodalshakerdaqsettings} shows these settings in the Rattlesnake \ac{GUI}.
    
    \begin{figure}[H]
        \centering
        \includegraphics[width=\linewidth]{figures/example_cdaq_modal_shaker_daq_settings}
        \caption{Data acquisition settings for the Modal Shaker Test}
        \label{fig:examplecdaqmodalshakerdaqsettings}
    \end{figure}
    
    \subsection{Environment Definition for Pseudorandom Excitation}
    
    Now we set up the Modal Testing environment parameters. Because we only have one shaker, we can use the \inlineCode{Pseudorandom} excitation, so we will plan for that.
    
    We will leave the \inlineCode{Samples per Frame} at its default value. We will change the \inlineCode{Number of Averages} to 10. We will leave \inlineCode{Triggering Type} set to \inlineCode{Free Run}, but we will add a delay of \inlineCode{5} in the \inlineCode{Wait for Steady State} box, to allow the system to come to steady state, as assumed by the Pseudorandom excitation.
    
    In the \inlineCode{Channel Selections} box, we select the Force channel as the reference.
    
    In the \inlineCode{Signal Generator} portion of the window, we select the \inlineCode{Pseudorandom} tab to set up a Pseudorandom signal. We can truncate the \inlineCode{Frequency Range} from \inlineCode{10 -- 2000} to remove the large rigid body motion that can occur at low frequency, as well as the portion of the excitation that is in the antialiasing filters of the system. We will set the \inlineCode{RMS} value to \inlineCode{0.1} V.
    
    With these parameters set, we can click the \inlineCode{Initialize Environments} button to proceed.  Figure \ref{fig:examplecdaqmodalpseudorandomsettings} shows these settings in the Rattlesnake \ac{GUI}.
    
    \begin{figure}[H]
        \centering
        \includegraphics[width=\linewidth]{figures/example_cdaq_modal_pseudorandom_settings}
        \caption{Parameters set for the Pseudorandom shaker excitation.}
        \label{fig:examplecdaqmodalpseudorandomsettings}
    \end{figure}
    
    \subsection{Test Profile}
    We will leave the \inlineCode{Test Profile} tab empty. Simply click the \inlineCode{Initialize Profile} button to proceed on this tab.
    
    \subsection{Run Test}
    
    \textbf{Warning:} cDAQ devices can output a pulse when plugged in, turned on, or activated for the first time, so make sure the shaker amplifier is off or the gain is set to zero before pressing the \inlineCode{Arm Data Acquisition} button for the first time. Once we \inlineCode{Arm Data Acquisition}, then we can unmute the shaker amplifier.
    
    To run the test we can again \inlineCode{Arm Data Acquisition} to enable the modal controls. Because we will not be roving the shaker, we have entered the proper node number into the channel table, and therefore do not need to utilize the \inlineCode{Degree of Freedom Override} table. We will again display the standard diagnostics of drive point FRF and coherence, reference autospectra, and time response from all sensors and the voltage signal.  This is shown in Figure \ref{fig:examplecdaqmodalshakerempty}.
    
    \begin{figure}[H]
        \centering
        \includegraphics[width=\linewidth]{figures/example_cdaq_modal_shaker_empty}
        \caption{Shaker testing with data acquisition armed and data visualization windows created.}
        \label{fig:examplecdaqmodalshakerempty}
    \end{figure}
    
    We can then \inlineCode{Preview Acquisition} to see how the test parameters look. Data will take a few seconds to appear after the shaker starts due to the \inlineCode{Wait for Steady State} parameter that was set.  Ideally if the pseudorandom excitation has reached steady state, we should see no appreciable difference between the time response of measurement frames, as the exact same signal is being applied and exact same responses are being measured.  This is shown in Figure \ref{fig:examplecdaqspeudorandompreview}.
    
    \begin{figure}[H]
        \centering
        \includegraphics[width=0.7\linewidth]{figures/example_cdaq_speudorandom_preview}
        \caption{Previewing the data acquisition with the pseudorandom excitation.}
        \label{fig:examplecdaqspeudorandompreview}
    \end{figure}
    
    If we are happy with the preview, we can move to \inlineCode{Start Acquisition}. We define a file to save the data to called \inlineCode{beam_shaker_pseudorandom.nc4}.  This is shown in Figure \ref{fig:examplecdaqpseudorandomacquire}.  The measurement and shakers should stop automatically when the requested number of averages have been acquired.
    
    \begin{figure}[H]
        \centering
        \includegraphics[width=\linewidth]{figures/example_cdaq_pseudorandom_acquire}
        \caption{Acquiring the pseudorandom modal data.}
        \label{fig:examplecdaqpseudorandomacquire}
    \end{figure}
    
    Once the test completes, we can mute our amplifier and Disarm Data Acquisition.
    
    \subsection{Other Shaker Excitation Types}
    We will return to the \inlineCode{Environment Definition} tab and set up a \inlineCode{Burst Random} test. We saw that the hammer ring-down took approximately 2 seconds, so we will set the \inlineCode{Samples Per Frame} value to \inlineCode{15360} to achieve a 3 second frame. We will set the \inlineCode{Number of Averages} to \inlineCode{30}. We will set the \inlineCode{Triggering Type} to \inlineCode{First Frame}, remove the \inlineCode{Wait for Steady State} value, and set the \inlineCode{Trigger Channel} to the voltage channel with a \inlineCode{Pretrigger} of \inlineCode{1%}. We will set the \inlineCode{Trigger Level} to \inlineCode{1%}, the \inlineCode{Hysteresis Level} to half that value, and the \inlineCode{Hysteresis Frame %} to \inlineCode{25%}.
    
    In the \inlineCode{Signal Generator} section, we will select a \inlineCode{Burst Random} signal and set the \inlineCode{On Percentage} to \inlineCode{33%} to provide a 1 second burst and 2 second ring down. We will keep the other parameters the same as our previous Pseudorandom test.
    
    Figure \ref{fig:examplecdaqburstrandomparameters} shows the parameters in the Rattlesnake \ac{GUI}.  Click the \inlineCode{Initialize Environments} button to proceed, and return to the \inlineCode{Run Test} tab.  
    
    \begin{figure}[H]
        \centering
        \includegraphics[width=\linewidth]{figures/example_cdaq_burst_random_parameters}
        \caption{Test parameters for the Burst Random excitation}
        \label{fig:examplecdaqburstrandomparameters}
    \end{figure}
    
    We can preview the signal by clicking the \inlineCode{Preview Acquisition} button to ensure that the triggering has been set up correctly and the signals ring down appropriately.  Figure \ref{fig:examplecdaqburstrandompreview} shows the previewed data.
    
    \begin{figure}[H]
        \centering
        \includegraphics[width=0.7\linewidth]{figures/example_cdaq_burst_random_preview}
        \caption{Previewing the measurement with the initial parameter set.}
        \label{fig:examplecdaqburstrandompreview}
    \end{figure}
    
    Apparently the shaker adds enough damping that the structure stops ringing by approximately 1/2 second after the burst finishes. Let's therefore \inlineCode{Stop Acquisition}, \inlineCode{Disarm Data Acquisition}, and return to the \inlineCode{Environment Definition} tab.  We will update the \inlineCode{Samples Per Frame} value to \inlineCode{5120} to achieve a 1 second measurement frame.  We will also change the burst \inlineCode{On Percentage} to \inlineCode{50%}.
    
    Figure \ref{fig:examplecdaqburstrandomparametersv2} shows the updated parameters in the Rattlesnake \ac{GUI}.  Click the \inlineCode{Initialize Environments} button to confirm these settings, and return to the \inlineCode{Run Test} tab.  
    
    \begin{figure}[H]
        \centering
        \includegraphics[width=\linewidth]{figures/example_cdaq_burst_random_parameters_v2}
        \caption{Updated burst random parameters after previewing the initial guess at the parameters.}
        \label{fig:examplecdaqburstrandomparametersv2}
    \end{figure}
    
    We can then preview this new configuration by \inlineCode{Arm Data Acquisition} and \inlineCode{Start Preview}.  The results are shown in Figure \ref{fig:examplecdaqburstrandompreviewv2}.  This configuration looks good and runs 3x faster, so we will keep it.
    
    \begin{figure}[H]
        \centering
        \includegraphics[width=0.7\linewidth]{figures/example_cdaq_burst_random_preview_v2}
        \caption{Previewing the measurement with the updated parameter set.}
        \label{fig:examplecdaqburstrandompreviewv2}
    \end{figure}
    
    Press the \inlineCode{Stop Acquisition} button to stop the preview, set a file to store the results to (\inlineCode{beam_shaker_burst.nc4}) and \inlineCode{Start Acquisition}. The acquisition should stop automatically when 30 averages are obtained.  Figure \ref{fig:examplecdaqburstrandomacquire} shows the measurement midway through acquisition.
    
    \begin{figure}[H]
        \centering
        \includegraphics[width=\linewidth]{figures/example_cdaq_burst_random_acquire}
        \caption{Acquiring the burst random data}
        \label{fig:examplecdaqburstrandomacquire}
    \end{figure}
    
    \subsection{Analyzing Modal Shaker Data}
    
    We can now analyze the shaker data similarly to the hammer data using SDynPy.  We will also load in the impact data from the tip.  While we did not excite the structure at node 1 with the hammer, we can use symmetry of the beam to compare between the two measurement types.
    
    \begin{lstlisting}[language=Python]
# Read a rattlesnake file using `read_modal_data`
time_data_pr, frfs_pr, coherence_pr, channel_table_pr = sdpy.rattlesnake.read_modal_data(
    'beam_shaker_pseudorandom.nc4')
time_data_br, frfs_br, coherence_br, channel_table_br = sdpy.rattlesnake.read_modal_data(
    'beam_shaker_burst.nc4')
# Compare to the impact
time_data_im, frfs_im, coherence_im, channel_table_im = sdpy.rattlesnake.read_modal_data(
    'beam_hammer_impact_0000.nc4')\end{lstlisting}
    
    We can use SDynPy to extract common degrees of freedom from each test and plot the FRFs for each.  
    
    \begin{lstlisting}[language=Python]
import matplotlib.pyplot as plt

ref_resp_dofs = sdpy.coordinate_array([1,23],'Y+')
compare_frfs = [frfs_pr[ref_resp_dofs[::-1]][np.newaxis],
frfs_br[ref_resp_dofs[::-1]][np.newaxis],
frfs_im[ref_resp_dofs][np.newaxis]]
fig,ax = plt.subplots(2)
for frf in compare_frfs:
frf.plot(ax)
ax[-1].legend(['Psuedorandom','Burst','Hammer'])
ax[-1].set_yscale('log')\end{lstlisting}
    
    This gives the following plot:
    
    \begin{figure}[H]
        \centering
        \includegraphics[width=0.5\linewidth]{figures/example_cdaq_modal_comparison}
        \caption{Comparison of hammer impact testing compared to shaker excitation with two types of signals.}
        \label{fig:examplecdaqmodalcomparison}
    \end{figure}
    
    We can note that the peaks are shifted slightly lower in frequency for the shaker testing, likely due to the small amount of mass loading of the shaker.  Additionally, there appears to be some shaker-structure interactions going on near the third mode, but otherwise there is reasonably good agreement between the three tests.
    
    \section{Running a Vibration Test with Rattlesnake}
    
    Now we will demonstrate how we can run a controlled vibration test with Rattlesnake.  We will demonstrate both Random and Transient vibration control in this section.
    
    \subsection{Constructing Vibration Specifications}\label{sec:example_nidaqmx_spec}
    
    Before we can run a test, we will first generate some environment data.  This is the response that we will control to during the vibration testing. We will do this using the modal hammer to randomly impact the structure, and then we will try to get the shaker to control to that response.
    
    To acquire environment data, we will simply run the modal environment in preview mode, but now use Rattlesnake's \inlineCode{Streaming} functionality to record time data, from which we will construct specifications for vibration testing.
    
    We will disconnect the shaker from our test article. Then we return to the \inlineCode{Environment Definition} page and set the \inlineCode{Triggering Type} to \inlineCode{Free Run}.
    
    On the \inlineCode{Run Test} tab, before selecting \inlineCode{Arm Data Acquisition} we select \inlineCode{Start Streaming Immediately} and \inlineCode{Select Streaming File...} to save a time history file called \inlineCode{environment_streaming.nc4}.  This is shown in Figure \ref{fig:examplecdaqsetupstreaming}.
    
    \begin{figure}[H]
        \centering
        \includegraphics[width=\linewidth]{figures/example_cdaq_set_up_streaming}
        \caption{Run Test tab with Streaming enabled}
        \label{fig:examplecdaqsetupstreaming}
    \end{figure}
    
    Now we can \inlineCode{Arm Data Acquisition} and then \inlineCode{Preview Acquisition}. Once the data acquisition is running, we will tap the structure with the modal hammer repeatedly along the length of the beam.  Don't worry that the hammer is not plugged in, as we only care about the accelerometer response. Once we have recorded a few seconds of impacts, we can select \inlineCode{Stop Acquisition} and \inlineCode{Disarm Data Acquisition}. Note that the streaming file will not be available until the data acquisition is disarmed.  We can now close Rattlesnake because we are done with the Modal environment.
    
    We will now load the streaming file with SDynPy to create an environment.
    
    \begin{lstlisting}[language=Python]
time_history, channel_table = sdpy.rattlesnake.read_rattlesnake_output(
    'environment_streaming.nc4')\end{lstlisting}
    
    We will extract just the first four acceleration channels for our environment, which are shown in Figure \ref{fig:examplecdaqspecificationsignal}
    
    \begin{lstlisting}[language=Python]
accelerations = time_history[:4]
accelerations.plot()\end{lstlisting}

    \begin{figure}[H]
        \centering
        \includegraphics[width=0.5\linewidth]{figures/example_cdaq_specification_signal}
        \caption{Measured acceleration data that we will turn into vibration specifications.}
        \label{fig:examplecdaqspecificationsignal}
    \end{figure}
    
    Looking at the data, we can truncate it to just the time range of interest, shown in Figure \ref{fig:examplecdaqspecificationsignaltruncated}.
    
    \begin{lstlisting}[language=Python]
truncated_accelerations = accelerations.extract_elements_by_abscissa(2.3,16.3)
truncated_accelerations.plot()\end{lstlisting}
    
    \begin{figure}[H]
        \centering
        \includegraphics[width=0.5\linewidth]{figures/example_cdaq_specification_signal_truncated}
        \caption{Truncated acceleration data that we will turn into vibration specifications.}
        \label{fig:examplecdaqspecificationsignaltruncated}
    \end{figure}
    
    To construct a random vibration environment, we will compute \ac{CPSD}s from this signal. While this isn't really an appropriate signal for a random vibration specification, it is good enough to demonstrate Rattlesnake's capabilities.  Figure \ref{fig:examplecdaqcpsdspecification} shows the \ac{APSD}s, which are the diagonal entries of the \ac{CPSD} specification matrix.
    
    \begin{lstlisting}[language=Python]
random_spec = truncated_accelerations.cpsd(samples_per_frame = 5120, overlap=0.5, window = 'hann')
# Remove low frequency and antialiasing filters
random_spec = random_spec.extract_elements_by_abscissa(20,2000)
# Scale it down so we don't break shakers trying to hit it
random_spec = random_spec / 100

ax = random_spec.plot_asds()
fig = ax[0,0].figure
fig.tight_layout()\end{lstlisting}
    
    \begin{figure}[H]
        \centering
        \includegraphics[width=0.7\linewidth]{figures/example_cdaq_cpsd_specification}
        \caption{Diagonal entries of the random vibration specification \ac{CPSD} matrix.}
        \label{fig:examplecdaqcpsdspecification}
    \end{figure}
    
    We will also construct a transient specification that we will use to test the transient vibration control, shown in Figure \ref{fig:examplecdaqtransientspecification}.
    
    \begin{lstlisting}[language=Python]
# We will make it shorter and again scale it down to make it so we don't break anything
transient_spec = truncated_accelerations.extract_elements_by_abscissa(11.48,16.3) / 5
transient_spec.plot()\end{lstlisting}
    
    \begin{figure}[H]
        \centering
        \includegraphics[width=0.7\linewidth]{figures/example_cdaq_transient_specification}
        \caption{Transient specification used for vibration control.}
        \label{fig:examplecdaqtransientspecification}
    \end{figure}
    
    We can now save these specifications to disk.
    
    \begin{lstlisting}[language=Python]
random_spec.to_rattlesnake_specification('random_specification.npz')
transient_spec.to_rattlesnake_specification('transient_specification.npz')\end{lstlisting}
    
    \subsection{Setting up the Shakers and Channel Table}
    
    We will now set up our vibration test.  We add a second shaker to the other end of the beam using a second drive cap over the accelerometer at node 23, shown in Figure \ref{fig:examplecdaqtwoshakersetup}.
    
    \begin{figure}[H]
        \centering
        \includegraphics[width=0.7\linewidth]{figures/example_cdaq_two_shaker_setup}
        \caption{Vibration test setup showing two shakers attached to the test article.}
        \label{fig:examplecdaqtwoshakersetup}
    \end{figure}
    
    Because we only have 6 acquisition channels, we must remove the force gauge to make room for the second shaker drive signal. We have now teed the first shaker's drive from \inlineCode{cDAQ9185-217ED78Mod3}'s \inlineCode{ao0} channel to the \inlineCode{ai1} channel of \inlineCode{cDAQ9185-217ED78Mod2}. The second shaker will be driven by \inlineCode{cDAQ9185-217ED78Mod3}'s \inlineCode{ao1} channel and is teed to the \inlineCode{ai2} channel of \inlineCode{cDAQ9185-217ED78Mod2}.  This is shown in Figure \ref{fig:examplecdaqtwoshakerhardwaresetup}.
    
    \begin{figure}[H]
        \centering
        \includegraphics[width=0.7\linewidth]{figures/example_cdaq_two_shaker_hardware_setup}
        \caption{cDAQ hardware setup using two shakers.}
        \label{fig:examplecdaqtwoshakerhardwaresetup}
    \end{figure}
    
    Our channel table now looks like Figure \ref{fig:examplecdaqtwoshakerchanneltable}.
    
    \begin{figure}[H]
        \centering
        \includegraphics[width=\linewidth]{figures/example_cdaq_two_shaker_channel_table}
        \caption{Channel table for the random vibration test setup}
        \label{fig:examplecdaqtwoshakerchanneltable}
    \end{figure}
    
    \subsection{Setting up the Random Vibration Test}
    Now that the test hardware and channel table are set up, we will set up our random vibration test. We will open Rattlesnake and select the \inlineCode{MIMO Random Vibration} environment as shown in Figure \ref{fig:examplecdaqmimorandomselection}.
    
    \begin{figure}[H]
        \centering
        \includegraphics[width=0.5\linewidth]{figures/example_cdaq_mimo_random_selection}
        \caption{Selecting the MIMO Random Vibration environment}
        \label{fig:examplecdaqmimorandomselection}
    \end{figure}
    
    On the \inlineCode{Data Acquisition Setup} tab, we will \inlineCode{Load Channel Table} to import our Excel channel spreadsheet. We will also set the \inlineCode{Sample Rate} to \inlineCode{5120}. The other parameters can be left at their default values.  Figure \ref{fig:examplecdaqvibrationdataacquisitionsetup} shows these parameters.
    
    \begin{figure}[H]
        \centering
        \includegraphics[width=\linewidth]{figures/example_cdaq_vibration_data_acquisition_setup}
        \caption{Data acquisition setup for the vibration tests.}
        \label{fig:examplecdaqvibrationdataacquisitionsetup}
    \end{figure}
    
    Click the \inlineCode{Initialize Data Acquisition} button to proceed.
    
    On the \inlineCode{Environment Definition} tab, we will keep the \inlineCode{Samples Per Frame} set to \inlineCode{5120}, which is consistent with the parameters used to construct the specification. In the \inlineCode{Control Python Script} section, we will click \inlineCode{Load} to select a control law. We will use the \inlineCode{control_laws.py} file, which is located in the \inlineCode{control_laws} folder of the main Rattlesnake directory. In the \inlineCode{Control Python Function} drop-down, we will select \inlineCode{pseudoinverse_control}.
    
    In the \inlineCode{Control Channels} section, we will select all four acceleration channels, which we will use to control to the environment.
    
    Finally we will click the \inlineCode{Load Spec} button and navigate to our random specification file \inlineCode{random_specification.npz} from SDynPy. The specification should then appear in the plots.
    
    With these parameters selected, we can click the Initialize Environments button to proceed.  Figure \ref{fig:examplecdaqrandomvibrationparameters} shows the parameters in the Rattlesnake \ac{GUI}.
    
    \begin{figure}[H]
        \centering
        \includegraphics[width=\linewidth]{figures/example_cdaq_random_vibration_parameters}
        \caption{Random vibration parameters specified in the Rattlesnake \ac{GUI}.}
        \label{fig:examplecdaqrandomvibrationparameters}
    \end{figure}
    
    At this point, we must run the \inlineCode{System Identification} phase of the controller. We will set the Level to \inlineCode{0.1 V RMS}, then click \inlineCode{Start} to start the system identification. A noise floor check will first be run, then the system identification follows. Rattlesnake will determine the transfer functions between the two shaker drive signals and the four accelerometer channels. When the specified number of averages have been obtained, the System Identification will stop automatically.  Figure \ref{fig:examplecdaqrandomvibrationsystemid} shows the system identification phase of the controller.
    
    \begin{figure}[H]
        \centering
        \includegraphics[width=\linewidth]{figures/example_cdaq_random_vibration_system_id}
        \caption{System identification for the Random Vibration environment.}
        \label{fig:examplecdaqrandomvibrationsystemid}
    \end{figure}
    
    On the \inlineCode{Test Prediction} tab, Rattlesnake will make predictions of the voltage that it will need to output to achieve the control specified, as well as the responses predicted to be obtained. Here it is important to note that the output signals are not too large for the shakers.  Figure \ref{fig:examplecdaqrandomvibrationpredictions} shows the predictions for this test.
    
    \begin{figure}[H]
        \centering
        \includegraphics[width=\linewidth]{figures/example_cdaq_random_vibration_predictions}
        \caption{Test predictions for the random vibration environment}
        \label{fig:examplecdaqrandomvibrationpredictions}
    \end{figure}
    
    If the predictions look good, we can proceed to the Test Profile tab, which we will again leave blank and click the Initialize Test Profile button to proceed to the Run Test tab.
    
    \subsection{Running a Random Vibration Test}
    On the \inlineCode{Run Test} tab, we see a different interface than the one we saw for the modal test. We can click the \inlineCode{Arm Data Acquisition} button to enable the Random environment.
    
    To start, we will set the \inlineCode{Current Test Level} to \inlineCode{-12} dB to ensure when we start the test we don't break anything. We will then step up slowly to the full \inlineCode{0} dB level test.
    
    We can then click the \inlineCode{Start Environment} button to start the test.  Figure \ref{fig:examplecdaqrandomvibrationstart} shows the \ac{GUI} prior to starting the test with the test level turned down to -12 dB.
    
    \begin{figure}[H]
        \centering
        \includegraphics[width=\linewidth]{figures/example_cdaq_random_vibration_start}
        \caption{Ready to start a random vibration test in Rattlesnake}
        \label{fig:examplecdaqrandomvibrationstart}
    \end{figure}
    
    As the test is running, the main Rattlesnake interface will show the sum of \ac{APSD}s, which is the trace of the \ac{CPSD} matrix and can be thought of as an average response of the test (Figure \ref{fig:examplecdaqrandomvibrationsumasds}). To visualize individual channels, windows can be created in the \inlineCode{Data Display} section of the window (figure \ref{fig:examplecdaqrandomvibrationindividualchannels}).
    
    \begin{figure}[H]
        \centering
        \includegraphics[width=\linewidth]{figures/example_cdaq_random_vibration_sum_asds}
        \caption{Rattlesnake \ac{GUI} showing the average test response}
        \label{fig:examplecdaqrandomvibrationsumasds}
    \end{figure}
    
    \begin{figure}[H]
        \centering
        \includegraphics[width=\linewidth]{figures/example_cdaq_random_vibration_individual_channels}
        \caption{Individual channel \ac{APSD}s showing how well each channel is matching the test level.}
        \label{fig:examplecdaqrandomvibrationindividualchannels}
    \end{figure}
    
    If the controller is controlling sufficiently well, we can step up the level to \inlineCode{0} dB as shown in Figure \ref{fig:examplecdaqrandomvibrationsumasdsfulllevel}.
    
    \begin{figure}
        \centering
        \includegraphics[width=\linewidth]{figures/example_cdaq_random_vibration_sum_asds_full_level}
        \caption{Vibration control at full specification level.}
        \label{fig:examplecdaqrandomvibrationsumasdsfulllevel}
    \end{figure}
    
    Feel free to investigate other control laws in the \inlineCode{control_laws.py} file, as they can result in a better control than the standard pseudoinverse control. Figure \ref{fig:examplecdaqrandomvibrationsumasdsmatchtrace} shows the results using the \inlineCode{match_trace_pseudoinverse} approach.  In particular, we see the regions between the peaks improve in this case.
    
    \begin{figure}[H]
        \centering
        \includegraphics[width=\linewidth]{figures/example_cdaq_random_vibration_sum_asds_match_trace}
        \caption{Improved control using a closed loop control law found in \inlineCode{match_trace_pseudoinverse}.}
        \label{fig:examplecdaqrandomvibrationsumasdsmatchtrace}
    \end{figure}
    
    Clearly, Rattlesnake has obtained relatively good control for this test.
    
    To save data from the test, we can either stream time data using the same approach used previously, or we can simply use the \inlineCode{Save Current Spectral Data} to save out the current realizations of the control \ac{CPSD}s, which we will call \inlineCode{random_vibration_spectral_results.nc4}. This is the quickest way to do comparisons between the specification and the control achieved. We can load this file with SDynPy and use SDynPy to make comparisons between the test and specification.  The results are shown in Figure \ref{fig:examplecdaqrandomvibrationerrorsummary}.
    
    \begin{lstlisting}[language=Python]
control, spec, drives = sdpy.rattlesnake.read_random_spectral_data(
    'random_vibration_spectral_results.nc4')
spec = spec.extract_elements_by_abscissa(20,2000)
control = control.extract_elements_by_abscissa(20,2000)
spec.error_summary(figure_kwargs={'figsize':(10,6)},Control=control)\end{lstlisting}
    
    \begin{figure}[H]
        \centering
        \includegraphics[width=0.7\linewidth]{figures/example_cdaq_random_vibration_error_summary}
        \caption{Error summary between the test and specification for the random vibration control.}
        \label{fig:examplecdaqrandomvibrationerrorsummary}
    \end{figure}
    
    We can then \inlineCode{Stop Environment}, \inlineCode{Disarm Data Acquisition} and close the software.
    
    \subsection{Setting up the Transient Vibration Test}
    
    The last demonstration will show how we can perform a MIMO Transient environment test. Open Rattlesnake and select the \inlineCode{MIMO Transient} environment as shown in Figure \ref{fig:examplecdaqmimotransientselection}.
    
    \begin{figure}[H]
        \centering
        \includegraphics[width=0.5\linewidth]{figures/example_cdaq_mimo_transient_selection}
        \caption{Selecting the \ac{MIMO} Transient environment}
        \label{fig:examplecdaqmimotransientselection}
    \end{figure}
    
    We will set the \inlineCode{Data Acquisition Setup} tab identical to what was used before, as shown in Figure \ref{fig:examplecdaqvibrationdataacquisitionsetup}. Click the \inlineCode{Initialize Data Acquisition} button to proceed.
    
    On the \inlineCode{Environment Definition tab}, we will again \inlineCode{Load} a control law file. This time we will load the \inlineCode{transient_control_laws.py} file, which contains transient control laws for Rattlesnake. We will select the \inlineCode{pseudoinverse_control} function.
    
    In the \inlineCode{Control Channels} section, we will again select all four accelerometer channels as the control channels.
    
    Finally, we will select the \inlineCode{Load Signal} button to load in our specification file \inlineCode{transient_specification.npz}.  Figure \ref{fig:examplecdaqtransientparameters} shows these parameters set in the Rattlesnake \ac{GUI}.  With these parameters set, we will click the \inlineCode{Initialize Environments} button to proceed.
    
    \begin{figure}[H]
        \centering
        \includegraphics[width=\linewidth]{figures/example_cdaq_transient_parameters}
        \caption{Parameters set for the Transient environment in the Rattlesnake \ac{GUI}.}
        \label{fig:examplecdaqtransientparameters}
    \end{figure}
    
    The Transient environment will also need to perform a \inlineCode{System Identification}. Again set the \inlineCode{Level} to \inlineCode{0.1 V RMS} and click \inlineCode{Start} to start the noise floor and system identification. Here it may be useful to look at the \inlineCode{Impulse Response}, as that is useful to understand how transient control will be performed.  This is shown in Figure \ref{fig:examplecdaqtransientsystemid}.
    
    \begin{figure}[H]
        \centering
        \includegraphics[width=\linewidth]{figures/example_cdaq_transient_system_id}
        \caption{System Identification for the Transient environment with Impulse Response function shown.}
        \label{fig:examplecdaqtransientsystemid}
    \end{figure}
    
    Once the system identification has been performed, Rattlesnake will again make a prediction as to the results of the test and show it on the \inlineCode{Test Predictions} tab. It will show the peak voltage it aims to output, as well as the response it predicts that it will get compared to the desired response.  Figure \ref{fig:examplecdaqtransientpredictions} shows these predictions.
    
    \begin{figure}[H]
        \centering
        \includegraphics[width=\linewidth]{figures/example_cdaq_transient_predictions}
        \caption{Test predictions and voltage that will be output for the transient control.}
        \label{fig:examplecdaqtransientpredictions}
    \end{figure}
    
    We can then proceed to the \inlineCode{Test Profile} tab, click the \inlineCode{Initialize Profile} button, and we are ready to run the test.
    
    \subsection{Running a Transient Test}
    
    We can now click the \inlineCode{Arm Data Acquisition} button to enable the Transient control functionality. We can set the \inlineCode{Signal Level} to \inlineCode{-12} dB to start so we don't break anything when we run. If it runs successfully, then we can step back up to full level.
    
    Rattlesnake will identify the control signal as it comes through the data acquisition and surround it with a box, as shown in Figure \ref{fig:examplecdaqtransientcontrolresults}.
    
    \begin{figure}[H]
        \centering
        \includegraphics[width=\linewidth]{figures/example_cdaq_transient_control_results}
        \caption{Running the transient control with the control signal identified.}
        \label{fig:examplecdaqtransientcontrolresults}
    \end{figure}
    
    To visualize individual channels, new windows can be created in the \inlineCode{Data Display} portion of the window.  The four control channels can be seen in Figure \ref{fig:examplecdaqtransientcontrolchannels}
    
    \begin{figure}[H]
        \centering
        \includegraphics[width=\linewidth]{figures/example_cdaq_transient_control_channels}
        \caption{Control achieved on each control channel in the Transient vibration test.}
        \label{fig:examplecdaqtransientcontrolchannels}
    \end{figure}
    
    To save out data, we can click the \inlineCode{Save Current Control Data} button. We will save the data to the file \inlineCode{transient_vibration_control_results.nc4}. This can be loaded into SDynPy to do comparisons to the desired signal.  It can be useful to zoom in to the starts of each impulse to see that while the ring down does match quite well, the initial impulse cannot be achieved as well with the modal shakers as it was with the impact hammer.  These results are shown in Figure \ref{fig:examplecdaqtransientzoom}.
    
    \begin{lstlisting}[language=Python]
control, specification, drives = sdpy.rattlesnake.read_transient_control_data(
'transient_vibration_control_results.nc4')

ax = specification.plot(
    one_axis=False, plot_kwargs = {'linewidth':0.5},
    subplots_kwargs = {'figsize':(10,6),'sharex':True,'sharey':True})
for a,c in zip(ax.flatten(),control):
    c.plot(a,plot_kwargs = {'linewidth':0.5})
# Zoom in on a specific portion
a.set_xlim(0.45,0.55)\end{lstlisting}
    
    \begin{figure}[H]
        \centering
        \includegraphics[width=0.7\linewidth]{figures/example_cdaq_transient_zoom}
        \caption{Zoom into an impact during the transient control showing good results achieved after the impact.}
        \label{fig:examplecdaqtransientzoom}
    \end{figure}
    
    \section{Summary}
    This appendix has walked through several example tests using real hardware with Rattlesnake.  Users are encouraged to use this appendix as a quick-start guide to using Rattlesnake.  A relatively inexpensive NI cDAQ device was used to perform modal and vibration testing on a simple beam structure with minimal instrumentation.
    
    Modal testing was performed using both hammer impacts and shaker excitation.  Modes were fit to the data using the SDynPy Python package.
    
    Vibration testing was performed using modal shakers.  An impact hammer was used to create an environment that the vibration test would aim to replicate.  Both \ac{MIMO} Random and \ac{MIMO} transient vibration tests were run, where the \ac{MIMO} Random attempted replicate a prescribed \ac{CPSD} matrix, and the \ac{MIMO} Transient attempted to replicate the time history directly.  Both control schemes achieved good results.
    
    \chapter{Synthetic Example Problem with a SDynPy System}\label{sec:example_sdynpy}
    This example problem will utilize the Python package SDynPy\footnote{SDynPy can be installed via pip using \inlineCode{pip install sdynpy}, or otherwise downloaded from the Github repository here: \url{https://github.com/sandialabs/sdynpy}} to construct a \inlineCode{System} object that can be integrated over time to produce a synthetic control problem in Rattlesnake.  This example problem can run entirely on a desktop or laptop computer and does not require any dedicated data acquisition, instrumentation, or shaker hardware.  Synthetic tests are therefore very useful for developing an understanding of how Rattlesnake works.  Rattlesnake deliberately tries to make its synthetic operation look as close as possible to running a real test, which makes it useful for learning or trying out new things without worrying about breaking an expensive piece of equipment if a mistake is made.  For more information in using SDynPy objects with Rattlesnake, see Chapter \ref{sec:rattlesnake_hardware_sdynpy_system}.
    
    This test problem will largely mirror the example problem from Appendix \ref{sec:example_nidaqmx}, but will be performed synthetically.
    
    The general process for this control problem is as follows:
    \begin{itemize}
        \item Import the SDynPy package.
        \item Create a \inlineCode{System} object for the structure, consisting of a Mass, Stiffness, and Damping matrix.
        \item Select degrees of freedom to control
        \item Create a specification to make Rattlesnake try to achieve
        \item Create Rattlesnake tests for \Ac{MIMO} Random Vibration, \Ac{MIMO} Transient, and Modal Testing environments.
        \item Load the Rattlesnake data back into SDynPy to make comparison plots.
    \end{itemize}
    
    Being a Python package, we run SDynPy from within a Python environment.  The author is using Python 3.11 for this example.  The full source code for this example is attached to this document as a \inlineCode{.txt} file\footnote{Many PDF readers will block files with executable extensions for security reasons, which includes \inlineCode{.py} files, which is why we are renaming the extension of the code that is attached to this document as a benign text file.}, which can be transformed to a Python file by renaming the \inlineCode{.txt} to \inlineCode{.py}: \attachfile{attachments/sdynpy_example_setup.txt}
    
    \section{Creating a SDynPy System and Specifications for Random and Transient Vibration Tests}
    
    This section will contain the SDynPy commands that will be used to set up the test.
    
    \subsection{Importing Relevant Python Packages}
    
    In order to use SDynPy and other Python Packages, they need to be imported 
    
    \begin{lstlisting}[language=Python]
# Import SDynPy module to get structural dynamics functionality
import sdynpy as sdpy
# We'll also import some common Python packages
import numpy as np # NumPy for numeric calculations
import matplotlib.pyplot as plt # Matplotlib for plotting\end{lstlisting}
    
    \subsection{Creating the Demonstration Objects}
    
    Next we will construct the demonstration objects.  We will construct \inlineCode{Geometry} and \inlineCode{System} objects from a beam finite element model that we create with SDynPy.  The \inlineCode{Geometry} contains all of the spatial information about the object: locations and orientations of various degrees of freedom on the structure, as well as connections between them.  The \inlineCode{System} contains the dynamic information: mass, stiffness, and damping matrices and the degrees of freedom associated with each row and column of those matrices.
    
    We will plot the degrees of freedom on the geometry to help us understand where we would like to perform measurements and place excitation devices.  The results are shown in Figure \ref{fig:examplesdynpybeamgeometry}.
    
    \begin{lstlisting}[language=Python]
# Create a system and geometry object using the beam functionality in SDynPy
system,geometry = sdpy.System.beam(
    length = 24 * 0.0254, # meters
    width = 0.75 * 0.0254, # meters
    height = 1.0 * 0.0254, # meters
    num_nodes = 25,
    material = 'steel')
# We want to see which degrees of freedom we have to work with, so we will
# plot the geometry with coordinate labels
geometry.plot_coordinate(label_dofs=True,arrow_scale=0.02)
\end{lstlisting}
    
    \begin{figure}[H]
        \centering
        \includegraphics[width=0.7\linewidth]{figures/example_sdynpy_beam_geometry}
        \caption{Geometry and degrees of freedom from the beam finite element model.}
        \label{fig:examplesdynpybeamgeometry}
    \end{figure}
    
    Because Rattlesnake will try to integrate equations of motion real-time, it is useful if we can get a reduced-order model of the structure to integrate.  Transforming to a modal model is a natural choice, as it will reduce the number of degrees of freedom as well as uncouple them.  It also affords us the ability to add modal damping to the structure, which can make the test more realistic.  SDynPy will keep track of the modal transformation (mode shape matrix) that transforms the modal degrees of freedom back to physical degrees of freedom.
    
    \begin{lstlisting}[language=Python]
# The finite element model system of equations is too large to integrate
# real-time, so we will create a reduced modal system to integrate based on the
# desired bandwidth of the test
test_bandwidth = 2560 # Hz

# Solve for modes up to 1.5x the bandwidth to lessen modal truncation effects
modes = system.eigensolution(maximum_frequency = test_bandwidth*1.5)

# This also gives us the option to add some damping to the model
modes.damping = 0.005

# Transform to modal system: modal mass, modal stiffness, modal damping
modal_system = modes.system()

# Save it to a file that we will load in Rattlesnake
modal_system.save('sdynpy_system.npz')\end{lstlisting}

    \subsection{Selecting Degrees of Freedom}
    
    We will now select the degrees of freedom for this test.  We will compute a vibration specification at these degrees of freedom that we will later use to control the test.  In this case, we will select the same response degrees of freedom as from Appendix \ref{sec:example_nidaqmx}.  Note, however, that the node numbers will be offset by 1 from that test, as the node that is 1 inch from the end of the beam is actually called node 2 in our finite element model.  Figure \ref{fig:examplesdynpycontroldofs} shows these degrees of freedom.
    
    We will also select some excitation degrees of freedom that we will use to hit repeatedly with a hammer impact, similar to what was done in Section \ref{sec:example_nidaqmx_spec} to create the vibration specification.  We will integrate equations of motion where the excitation is random strikes at each of these nodes in the same way we measured the response to random tapping on the beam in the real test.  Note that these are not the nodes where shakers will be placed in the actual modal and vibration tests, but are only used to construct the specification.  Figure \ref{fig:examplesdynpyexcitationdofs} shows these degrees of freedom.
    
    \begin{lstlisting}[language=Python]
# Now let's select some degrees of freedom for our test.  We will measure all
# points in the vertical (z) direction.
response_dofs = sdpy.coordinate_array([2,13,19,24],'Y-')

# Let's also select some excitation degrees of freedom for the test
excitation_dofs = sdpy.coordinate_array(geometry.node.id[::3],'Y-')

geometry.plot_coordinate(response_dofs,label_dofs=True,arrow_scale=0.05)
geometry.plot_coordinate(excitation_dofs,label_dofs=True,arrow_scale=0.05)\end{lstlisting}

    \begin{figure}[H]
        \centering
        \includegraphics[width=0.5\linewidth]{figures/example_sdynpy_control_dofs}
        \caption{Degrees of freedom measured and used to control the vibration test.}
        \label{fig:examplesdynpycontroldofs}
    \end{figure}
    
    \begin{figure}[H]
        \centering
        \includegraphics[width=0.5\linewidth]{figures/example_sdynpy_excitation_dofs}
        \caption{Degrees of freedom at which random hammer impacts will be applied to construct a vibration environment.}
        \label{fig:examplesdynpyexcitationdofs}
    \end{figure}
    
    \subsection{Simulating an Environment}\label{sec:example_sdynpy_spec}
    
    With the system set up and degrees of freedom selected, we will now integrate the equations of motion to generate specifications that we can use to control the test.  We will apply a multi-hammer impact to each of the nodes in Figure \ref{fig:examplesdynpyexcitationdofs} and extract responses at the nodes in Figure \ref{fig:examplesdynpycontroldofs}.
    
    \begin{lstlisting}[language=Python]
# Let's simulate the multi-hammer impact that was performed on the test article
# in the example problem
        
frame_length = test_bandwidth*2 # Samples

responses, references = modal_system.simulate_test(
    bandwidth = test_bandwidth,
    frame_length = frame_length,
    num_averages = 30,
    excitation = 'multi-hammer',
    references = excitation_dofs,
    responses = response_dofs,
    excitation_level = 0.1,
    excitation_max_frequency = test_bandwidth*1.5)\end{lstlisting}
    
    This will give us 30 averages of the multi-hammer excitation that we will use to construct specifications.  From the response data we will construct a CPSD matrix.  The \ac{APSD}s are shown in Figure \ref{fig:examplesdynpycpsdspec}.
    
    \begin{lstlisting}[language=Python]
# Now we will convert that into a CPSD and a transient specification that we
# will try to replicate.
cpsd = responses.cpsd(frame_length, overlap = 0.0, window = 'boxcar')

# We'll truncate the low frequency to make sure that we don't get huge rigid
# body response as well as the high frequency above the antialiasing filters of
# the system
cpsd.ordinate[(cpsd.abscissa < 20) | (cpsd.abscissa > 2000)] = 0

# Plot the diagonal so we can see levels
cpsd.plot_asds({'layout':'constrained','sharex':True,'sharey':True})\end{lstlisting}
    
    \begin{figure}[H]
        \centering
        \includegraphics[width=0.5\linewidth]{figures/example_sdynpy_cpsd_spec}
        \caption{Specification for the \ac{MIMO} Random Vibration test.}
        \label{fig:examplesdynpycpsdspec}
    \end{figure}
    
    We will also pull off the first frame of data to use as a specification for a transient test, shown in Figure \ref{fig:examplesdynpytransientspec}.
    
    \begin{lstlisting}[language=Python]
# We'll also grab just the first average to make into a rattlesnake transient
# specification
transient = responses.idx_by_el[:frame_length]

transient.plot(
    one_axis=False,
    subplots_kwargs={'layout':'constrained'},
    plot_kwargs={'linewidth':0.5})\end{lstlisting}

    \begin{figure}[H]
        \centering
        \includegraphics[width=0.5\linewidth]{figures/example_sdynpy_transient_spec}
        \caption{Specification for the \ac{MIMO} Transient Vibration test.}
        \label{fig:examplesdynpytransientspec}
    \end{figure}

    We can then save these files to disk.
    
    \begin{lstlisting}[language=Python]
# Now we can save these to specifications for Rattlesnake
cpsd.to_rattlesnake_specification('random_specification.npz')
transient.to_rattlesnake_specification('transient_specification.npz')\end{lstlisting}

    \section{Running a Synthetic Modal Test}\label{sec:example_sdynpy_modal}
    
    Now that we have the pieces of information required, we can run our modal and vibration tests.  The primary difference between this test and the one run in Appendix \ref{sec:example_nidaqmx} is that there is no voltage signal being measured, only forces.  The SDynPy object we will be integrating does not know anything about any shaker we might attach to the structure, it only knows about the test article's dynamics.  This means ``inputs'' to the system are forces rather than voltages.
    
    Another major difference between this test and the previous on from Appendix \ref{sec:example_nidaqmx} is that it is more difficult to perform a synthetic hammer impact test.  One can achieve this by using the combined environments capabilities of Rattlesnake, specifying perhaps a \inlineCode{Time History Generator} environment to provide a impact-like pulse, which could be used to trigger the modal acquisition.  However, this advanced usage is out of the scope of this beginner-level tutorial, so it will not be discussed here.
    
    \subsection{Data Acquisition Setup}
    
    Similarly to Appendix \ref{sec:example_nidaqmx}, we will open Rattlesnake in the Modal Environment to save out a channel table template that we will use for this test.  Rather than filling out the channel table in the GUI, we will load it from an Excel spreadsheet, which allows us to more easily save our work. We can get a template channel table to fill out by opening Rattlesnake and selecting the Modal Testing environment (Figure \ref{fig:examplesdynpymodaltestingselection}).
    
    \begin{figure}[H]
        \centering
        \includegraphics[width=0.5\linewidth]{figures/example_cdaq_modal_testing_selection}
        \caption{Selecting the Modal Environment in Rattlesnake}
        \label{fig:examplesdynpymodaltestingselection}
    \end{figure}
    
    We can then click the \inlineCode{Save Channel Table} button to save out the empty channel table to an Excel spreadsheet (Figure \ref{fig:examplesdynpysavechanneltable}).
    
    \begin{figure}[H]
        \centering
        \includegraphics[width=0.7\linewidth]{figures/example_cdaq_save_channel_table}
        \caption{Saving out an empty channel table to fill out for the test.}
        \label{fig:examplesdynpysavechanneltable}
    \end{figure}
    
    In a real test, we would apply instrumentation to the test article and then record that information in the channel table.  In this case, the degrees of freedom names are already stored in the SDynPy \inlineCode{System} object, we only need to tell Rattlesnake which degrees of freedom to use for excitations and responses.  Recall in our code, we specified our responses as nodes 2, 13, 19, and 24 in the Y- direction.  We will fill our our channel table with those nodes and directions.  We will also put the excitation signal (i.e. force) at node 2, which mirrors what was done in Section \ref{sec:example_nidaqmx_modal_shaker}.  Note that Rattlesnake will interrogate the SDynPy \inlineCode{System} object to handle the bookkeeping to determine which \inlineCode{System} degree of freedom is equivalent to each channel in the Rattlesnake test.  Rattlesnake will handle polarity changes is for example, \inlineCode{Y-} direction is specified in the test whereas \inlineCode{Y+} direction was specified in the \inlineCode{System} object.
    
    We can look back at Section \ref{sec:rattlesnake_hardware_sdynpy_system} to see what needs to be set up in the channel table for a test using a SDynPy \inlineCode{System} object.  We will fill in the \inlineCode{Node Number} and \inlineCode{Node Direction} with the node numbers and directions.  We will put \inlineCode{Virtual} in the \inlineCode{Physical Device} column to specify that a given channel is active.  We will put \inlineCode{Input} in the \inlineCode{Feedback Device} column.  The \inlineCode{Type} must also be specified, as it will determine the derivative applied to an acquired response channel.  We will put the word \inlineCode{Acceleration} in that column for the response channels.  We will put \inlineCode{Force} in that channel for the excitation signal.
    
    Other columns of the channel table may be filled out to provide additional documentation, but they are not required.  For example, we could put units on our data to tell users of the dataset that the data will be in SI units, or we could put comments to describe various aspects of the test.  The channel table is shown in Figure \ref{fig:examplesdynpychanneltable}.
    
    \begin{figure}[H]
        \centering
        \includegraphics[width=\linewidth]{figures/example_sdynpy_channel_table}
        \caption{Channel Table used in the SDynPy Modal Test.}
        \label{fig:examplesdynpychanneltable}
    \end{figure}
    
    Note that there is no voltage signal in the channel table because the SDynPy system has no knowledge of a shaker model to use in an integration.  However, Rattlesnake still needs to know which signals are responses and which signals are inputs, so we must specify those signals in the \inlineCode{Feedback Device} column of the channel table.  It may sound strange to think of ``feedback'' in terms of a synthetic test, as we are not teeing any signal back to a separate channel on the data acquisition.  However, one can think of this as a kind of ``virtual'' feedback: the signal being applied to the structure as an input when integrating equations of motions is also ``fed back'' to the user as a response that is measured as an output of the integration.
    
    With the channel table set up, we can now start Rattlesnake and load the channel table. We will again select the Modal Testing environment, click the \inlineCode{Load Channel Table} button and load our file.  Note that we must switch the \inlineCode{Hardware Selector} from \inlineCode{NI DAQmx} to \inlineCode{SDynPy System Integration...} as shown in Figure \ref{fig:examplesdynpysystemintegrationselection}.
    
    \begin{figure}[H]
        \centering
        \includegraphics[width=0.5\linewidth]{figures/example_sdynpy_system_integration_selection}
        \caption{Selecting \inlineCode{SDynPy System Integration...} as the hardware device used in the test.}
        \label{fig:examplesdynpysystemintegrationselection}
    \end{figure}
    
    When this option is selected, a file dialog will appear, and the user must select the file in which the \inlineCode{System} object was saved.  In this case, it is \inlineCode{sdynpy_system.npz}.  An extra parameter \inlineCode{Integration Oversample} will also appear, which specifies the Sample Rate used by the time integration compared to the sample rate of the test.  We will again set the \inlineCode{Sample Rate} to \inlineCode{5120} Hz, and we will leave the default \inlineCode{Integration Oversample} factor at \inlineCode{10}, which will mean that the time integration will take place at \inlineCode{51200} time steps per second, but the result will be downsampled back to \inlineCode{5120} Hz to provide commonality with other hardware devices.  Figure \ref{fig:examplesdynpymodalshakerdaqsettings} shows these settings in the Rattlesnake \ac{GUI}.

    \begin{figure}[H]
        \centering
        \includegraphics[width=\linewidth]{figures/example_sdynpy_modal_shaker_daq_settings}
        \caption{Data Acquisition Settings used to perform the Synthetic SDynPy Modal Test.}
        \label{fig:examplesdynpymodalshakerdaqsettings}
    \end{figure}
    
    At this point, we can click the \inlineCode{Initialize Data Acquisition} button shown in Figure \ref{fig:examplesdynpyinitializedataacquisition} to proceed to the next stage of the test.
    
    \begin{figure}[H]
        \centering
        \includegraphics[width=0.5\linewidth]{figures/example_cdaq_initialize_data_acquisition}
        \caption{Click the \inlineCode{Initialize Data Acquisition} button to proceed.}
        \label{fig:examplesdynpyinitializedataacquisition}
    \end{figure}

    \subsection{Environment Definition for Pseudorandom Excitation}
    
    Now we set up the Modal Testing environment parameters. Because we only have one shaker, we can use the \inlineCode{Pseudorandom} excitation, so we will plan for that.
    
    We will leave the \inlineCode{Samples per Frame} at its default value. We will change the \inlineCode{Number of Averages} to 10. We will leave \inlineCode{Triggering Type} set to \inlineCode{Free Run}, but we will add a delay of \inlineCode{5} in the \inlineCode{Wait for Steady State} box, to allow the system to come to steady state, as assumed by the Pseudorandom excitation.
    
    In the \inlineCode{Channel Selections} box, we select the Force channel as the reference.
    
    In the \inlineCode{Signal Generator} portion of the window, we select the \inlineCode{Pseudorandom} tab to set up a Pseudorandom signal. We can truncate the \inlineCode{Frequency Range} from \inlineCode{10 -- 2000} to remove the large rigid body motion that can occur at low frequency, as well as the portion of the excitation that is in the anti-aliasing filters of the system (note that there are no anti-aliasing filters in the synthetic integration, but it is good practice for when moving to real data acquisition systems). We will set the \inlineCode{RMS} value to \inlineCode{0.1} V.  Note, however, that this will actually be the force level in Newtons due to how the synthetic test is set up.
    
    With these parameters set, we can click the \inlineCode{Initialize Environments} button to proceed.  Figure \ref{fig:examplesdynpymodalpseudorandomsettings} shows these settings in the Rattlesnake \ac{GUI}.
    
    \begin{figure}[H]
        \centering
        \includegraphics[width=\linewidth]{figures/example_sdynpy_modal_pseudorandom_settings}
        \caption{Parameters set for the Pseudorandom shaker excitation.}
        \label{fig:examplesdynpymodalpseudorandomsettings}
    \end{figure}
    
    \subsection{Test Profile}
    We will leave the \inlineCode{Test Profile} tab empty. Simply click the \inlineCode{Initialize Profile} button to proceed on this tab.
    
    \subsection{Run Test}
    We will save the standard modal data file from this test, so we won't worry about setting up streaming.
    
    To enable the Modal Environment, first click the \inlineCode{Arm Data Acquisition} button. At this point, the modal environment buttons will become enabled, as shown in \ref{fig:examplesdynpymodaltestingready}
    
    \begin{figure}[H]
        \centering
        \includegraphics[width=\linewidth]{figures/example_sdynpy_modal_testing_ready}
        \caption{Run Test tab with data acquisition armed}
        \label{fig:examplesdynpymodaltestingready}
    \end{figure}
    
     Next we can set up data display windows. We will use the \inlineCode{New from Template...} dropdown menu to create \inlineCode{Drive Point (Imaginary)}, \inlineCode{Drive Point Coherence}, and \inlineCode{Reference Autospectrum} channels. We will also manually create 5 windows by clicking the \inlineCode{New} button, and set them to visualize the time histories for each channel; these will help us determine if we need to change the trigger settings, frame length, or add a window. Clicking the \inlineCode{Tile} button will expand the windows to fill up the available space.  Figure \ref{fig:examplesdynpymodalshakeremptywindows} shows the result of these operations.
    
    \begin{figure}[H]
        \centering
        \includegraphics[width=\linewidth]{figures/example_sdynpy_modal_shaker_empty_windows}
        \caption{Run Test tab with modal windows created.}
        \label{fig:examplesdynpymodalshakeremptywindows}
    \end{figure}
    
    We can then \inlineCode{Preview Acquisition} to see how the test parameters look. Data will take a few seconds to appear after the shaker starts due to the \inlineCode{Wait for Steady State} parameter that was set.  Ideally if the pseudorandom excitation has reached steady state, we should see no appreciable difference between the time response of measurement frames, as the exact same signal is being applied and exact same responses are being measured.  This is shown in Figure \ref{fig:examplesdynpypseudorandompreview}.  Note that due to there being no excitation signal past 2000 Hz and below 10 Hz, the \ac{FRF} values may be quite large in those frequency ranges, because the computation is essentially dividing by zero at those frequency lines.  You may need to adjust the limits of the Y-axis appropriately to visualize the data from 10--2000 Hz.
    
    \begin{figure}[H]
        \centering
        \includegraphics[width=0.7\linewidth]{figures/example_sdynpy_pseudorandom_preview}
        \caption{Previewing the data acquisition with the pseudorandom excitation.}
        \label{fig:examplesdynpypseudorandompreview}
    \end{figure}
    
    If we are happy with the preview, we can move to \inlineCode{Start Acquisition}. We must first define a file to save the data, which we will call \inlineCode{beam_shaker_pseudorandom.nc4}.  This is shown in Figure \ref{fig:examplesdynpypseudorandomacquire}.  The measurement and shakers should stop automatically when the requested number of averages have been acquired.
    
    \begin{figure}[H]
        \centering
        \includegraphics[width=\linewidth]{figures/example_sdynpy_pseudorandom_acquire}
        \caption{Acquiring the pseudorandom modal data.}
        \label{fig:examplesdynpypseudorandomacquire}
    \end{figure}
    
    \subsection{Other Shaker Excitation Types}
    We will return to the \inlineCode{Environment Definition} tab and set up a \inlineCode{Burst Random} test. We don't know exactly how long the structure will take to ring down, so we will set the \inlineCode{Samples Per Frame} value to \inlineCode{15360} to achieve a 3 second frame. We will set the \inlineCode{Number of Averages} to \inlineCode{30}. We will set the \inlineCode{Triggering Type} to \inlineCode{First Frame}, remove the \inlineCode{Wait for Steady State} value, and set the \inlineCode{Trigger Channel} to the force channel with a \inlineCode{Pretrigger} of \inlineCode{1%}. We will set the \inlineCode{Trigger Level} to \inlineCode{0.1%}, the \inlineCode{Hysteresis Level} to half that value, and the \inlineCode{Hysteresis Frame %} to \inlineCode{25%}.  Note that because there is no ``sensitivity'' value that translates voltage to engineering unit in the synthetic case, the two values are treated equally.
    
    In the \inlineCode{Signal Generator} section, we will select a \inlineCode{Burst Random} signal and set the \inlineCode{On Percentage} to \inlineCode{33%} to provide a 1 second burst and 2 second ring down. We will keep the other parameters the same as our previous Pseudorandom test.
    
    Figure \ref{fig:examplesdynpyburstrandomparameters} shows the parameters in the Rattlesnake \ac{GUI}.  Click the \inlineCode{Initialize Environments} button to proceed, and return to the \inlineCode{Run Test} tab.  
    
    \begin{figure}[H]
        \centering
        \includegraphics[width=\linewidth]{figures/example_sdynpy_burst_random_parameters}
        \caption{Test parameters for the Burst Random excitation}
        \label{fig:examplesdynpyburstrandomparameters}
    \end{figure}
    
    We can preview the signal by clicking the \inlineCode{Preview Acquisition} button to ensure that the triggering has been set up correctly and the signals ring down appropriately.  Figure \ref{fig:examplesdynpyburstrandompreview} shows the previewed data.
    
    \begin{figure}[H]
        \centering
        \includegraphics[width=0.7\linewidth]{figures/example_sdynpy_burst_random_preview}
        \caption{Previewing the measurement with the initial parameter set.}
        \label{fig:examplesdynpyburstrandompreview}
    \end{figure}
    
    Looking at the preview data, the damping is such that the that the structure stops ringing by approximately 1/2 second after the burst finishes. Let's therefore \inlineCode{Stop Acquisition}, \inlineCode{Disarm Data Acquisition}, and return to the \inlineCode{Environment Definition} tab.  We will update the \inlineCode{Samples Per Frame} value to \inlineCode{5120} to achieve a 1 second measurement frame.  We will also change the burst \inlineCode{On Percentage} to \inlineCode{50%}.
    
    Figure \ref{fig:examplesdynpyburstrandomparametersv2} shows the updated parameters in the Rattlesnake \ac{GUI}.  Click the \inlineCode{Initialize Environments} button to confirm these settings, and return to the \inlineCode{Run Test} tab.  
    
    \begin{figure}[H]
        \centering
        \includegraphics[width=\linewidth]{figures/example_sdynpy_burst_random_parameters_v2}
        \caption{Updated burst random parameters after previewing the initial guess at the parameters.}
        \label{fig:examplesdynpyburstrandomparametersv2}
    \end{figure}
    
    We can then preview this new configuration by \inlineCode{Arm Data Acquisition} and \inlineCode{Start Preview}.  The results are shown in Figure \ref{fig:examplesdynpyburstrandompreviewv2}.  This configuration looks good and runs 3x faster, so we will keep it.
    
    \begin{figure}[H]
        \centering
        \includegraphics[width=0.7\linewidth]{figures/example_sdynpy_burst_random_preview_v2}
        \caption{Previewing the measurement with the updated parameter set.}
        \label{fig:examplesdynpyburstrandompreviewv2}
    \end{figure}
    
    Press the \inlineCode{Stop Acquisition} button to stop the preview, set a file to store the results to (\inlineCode{beam_shaker_burst.nc4}) and \inlineCode{Start Acquisition}. The acquisition should stop automatically when 30 averages are obtained.  Figure \ref{fig:examplesdynpyburstrandomacquire} shows the measurement midway through acquisition.
    
    \begin{figure}[H]
        \centering
        \includegraphics[width=\linewidth]{figures/example_sdynpy_burst_random_acquire}
        \caption{Acquiring the burst random data}
        \label{fig:examplesdynpyburstrandomacquire}
    \end{figure}
    
    \subsection{Analyzing Modal Data}
    
    To analyze these data, we will again use SDynPy\footnote{SDynPy can be installed via pip using \inlineCode{pip install sdynpy}, or otherwise downloaded from the Github repository here: \url{https://github.com/sandialabs/sdynpy}}, which is an open-source Python-based structural dynamics toolset.  We will focus on the Burst Random test data, but analyzing the Pseudorandom data should be very similar.  A Python script containing the commands is attached to this document as a \inlineCode{.txt} file\footnote{Many PDF readers will block files with executable extensions for security reasons, which includes \inlineCode{.py} files, which is why we are renaming the extension of the code that is attached to this document as a benign text file.}, which can be transformed to a Python file by renaming the \inlineCode{.txt} to \inlineCode{.py}: \attachfile{attachments/sdynpy_example_analysis.txt}
    
    SDynPy is well-integrated with Rattlesnake, making it trivial to load in data from a Rattlesnake test.  For example, reading in a modal dataset with SDynPy looks like:
    \begin{lstlisting}[language=Python]
# Import SDynPy and Numpy
import sdynpy as sdpy
import numpy as np

# Read a rattlesnake file using `read_modal_data`
time_data, frfs, coherence, channel_table = sdpy.rattlesnake.read_modal_data(
    'beam_hammer_impact_0000.nc4')\end{lstlisting}
    
    We can plot the data easily using SDynPy's plotting functionality:
    \begin{lstlisting}[language=Python]
frfs.plot(one_axis=False,subplots_kwargs={'layout':'constrained'})\end{lstlisting}
    
    The results of this command are shown in Figure \ref{fig:examplesdynpyplottedfrfs}.
    
    \begin{figure}[H]
        \centering
        \includegraphics[width=0.7\linewidth]{figures/example_sdynpy_plotted_frfs}
        \caption{Frequency response functions plotted in SDynPy from the modal impact test.}
        \label{fig:examplesdynpyplottedfrfs}
    \end{figure}
    
    Ideally we would like to plot mode shapes of this test, so as a first step towards that goal we will construct a test geometry. This contains nodes, coordinate systems, and optionally tracelines or elements.  We can construct nodes with the ID numbers given in the test and coordinates specified by each node's position along the beam.  Here we construct and print a \inlineCode{node_array} in SDynPy
    \begin{lstlisting}[language=Python]
nodes = sdpy.node_array(id = [2,13,19,24],
                        coordinate = [[1,0,0],
                                      [12,0,0],
                                      [18,0,0],
                                      [23,0,0]])
print(nodes)\end{lstlisting}
    which has output
    \begin{lstlisting}
Index,     ID,        X,        Y,        Z, DefCS, DisCS
 (0,),      2,    1.000,    0.000,    0.000,     1,     1
 (1,),     13,   12.000,    0.000,    0.000,     1,     1
 (2,),     19,   18.000,    0.000,    0.000,     1,     1
 (3,),     24,   23.000,    0.000,    0.000,     1,     1\end{lstlisting}
    
    We will do similarly with a global coordinate system
    \begin{lstlisting}[language=Python]
css = sdpy.coordinate_system_array(id=1,name='global')
print(css)\end{lstlisting}
    which has output
    \begin{lstlisting}
Index,     ID,                 Name, Color,       Type
   (),      1,               global,     1,  Cartesian\end{lstlisting}
    
    We can then construct a Geometry object. We will then add a traceline to connect the nodes together.  We will then plot the geometry.
    \begin{lstlisting}[language=Python]
geometry = sdpy.Geometry(nodes,css)
geometry.add_traceline([2,13,19,24])
geometry.plot()\end{lstlisting}
    
    The plotted geometry is shown in Figure \ref{fig:examplesdynpygeometry}.
    
    \begin{figure}[H]
        \centering
        \includegraphics[width=0.7\linewidth]{figures/example_cdaq_geometry}
        \caption{Geometry used to plot mode shapes in the modal test.}
        \label{fig:examplesdynpygeometry}
    \end{figure}
    
    Note that an alternative way to construct the geometry would have been to use the geometry output directly from the \inlineCode{sdynpy.System.beam} function, and remove any nodes that were not defined in the test.
    
    Now that we have a geometry, we can fit modes and plot mode shapes. We will only fit up to approximately 2000 Hz, which was the maximum frequency of the excitation.  We will use the \inlineCode{PolyPy_GUI} in SDynPy to do the mode fits.  We can add our geometry directly to the fitter by assigning to the \inlineCode{geometry} attribute.
    \begin{lstlisting}[language=Python]
ppgui = sdpy.PolyPy_GUI(frfs)
ppgui.geometry = geometry\end{lstlisting}
    
    We can select frequency ranges and polynomial orders to solve for, then click \inlineCode{Compute Stabilization} to proceed, as shown in Figure \ref{fig:examplesdynpypolypyranges}.
    
    \begin{figure}[H]
        \centering
        \includegraphics[width=0.7\linewidth]{figures/example_sdynpy_polypy_ranges}
        \caption{Setting up frequency ranges and polynomial orders in PolyPy.}
        \label{fig:examplesdynpypolypyranges}
    \end{figure}
    
    We can then select stable poles from the diagram and plot resynthesized FRFs or mode indicator functions, shown in Figure \ref{fig:examplesdynpypolypyselectpoles}.
    
    \begin{figure}[H]
        \centering
        \includegraphics[width=0.7\linewidth]{figures/example_sdynpy_polypy_select_poles}
        \includegraphics[width=\linewidth]{figures/example_cdaq_polypy_resynthesis}
        \caption{Selecting stable poles and visualizing resynthesized FRFs.}
        \label{fig:examplesdynpypolypyselectpoles}
    \end{figure}
    
    We can then plot mode shapes from the test.  Note that with only four sensors, our mode shapes will not be very highly resolved; however, they clearly correspond to the first three bending modes of the beam.  The first mode shape is shown in Figure \ref{fig:examplesdynpymodalfits}.
    
    \begin{figure}[H]
        \centering
        \includegraphics[width=0.7\linewidth]{figures/example_sdynpy_modal_fits}
        \caption{First mode shape fit to the modal data from the impact test.}
        \label{fig:examplesdynpymodalfits}
    \end{figure}
    
    \section{Running a Synthetic Vibration Test with Rattlesnake}\label{sec:example_sdynpy_random_vibration}
    
    Now we will demonstrate how we can simulate a controlled vibration test with Rattlesnake.  We will demonstrate both Random and Transient vibration control in this section.
    
    \subsection{Setting up the Channel Table}
    
    We will now set up our vibration test.  We add a second shaker to the other end of the beam using a second drive at node 24, in addition to the drive at node 2 used in the modal test. We will now have two force channels applied to the structure.  Note that the control channels in the test must be in the same order as the channels in the specification that we constructed in Section \ref{sec:example_sdynpy_spec}.
    
    Our channel table now looks like Figure \ref{fig:examplesdynpytwoshakerchanneltable}.
    
    \begin{figure}[H]
        \centering
        \includegraphics[width=\linewidth]{figures/example_sdynpy_two_shaker_channel_table}
        \caption{Channel table for the random vibration test setup}
        \label{fig:examplesdynpytwoshakerchanneltable}
    \end{figure}
    
    \subsection{Setting up the Random Vibration Test}
    Now that the specification and channel table are set up, we will set up our random vibration test. We will open Rattlesnake and select the \inlineCode{MIMO Random Vibration} environment as shown in Figure \ref{fig:examplesdynpymimorandomselection}.
    
    \begin{figure}[H]
        \centering
        \includegraphics[width=0.5\linewidth]{figures/example_cdaq_mimo_random_selection}
        \caption{Selecting the MIMO Random Vibration environment}
        \label{fig:examplesdynpymimorandomselection}
    \end{figure}
    
    On the \inlineCode{Data Acquisition Setup} tab, we will \inlineCode{Load Channel Table} to import our Excel channel spreadsheet. We will set the \inlineCode{Hardware Selector} to \inlineCode{SDynPy System Integration...} and select our system file \inlineCode{sdynpy_system.npz}, as well as set the sample \inlineCode{Sample Rate} to \inlineCode{5120}. The other parameters can be left at their default values.  Figure \ref{fig:examplesdynpyvibrationdataacquisitionsetup} shows these parameters.
    
    \begin{figure}[H]
        \centering
        \includegraphics[width=\linewidth]{figures/example_sdynpy_vibration_data_acquisition_setup}
        \caption{Data acquisition setup for the vibration tests.}
        \label{fig:examplesdynpyvibrationdataacquisitionsetup}
    \end{figure}
    
    Click the \inlineCode{Initialize Data Acquisition} button to proceed.
    
    On the \inlineCode{Environment Definition} tab, we will keep the \inlineCode{Samples Per Frame} set to \inlineCode{5120}, which is consistent with the parameters used to construct the specification. In the \inlineCode{Control Python Script} section, we will click \inlineCode{Load} to select a control law. We will use the \inlineCode{control_laws.py} file, which is located in the \inlineCode{control_laws} folder of the main Rattlesnake directory. In the \inlineCode{Control Python Function} drop-down, we will select \inlineCode{pseudoinverse_control}.
    
    In the \inlineCode{Control Channels} section, we will select all four acceleration channels, which we will use to control to the environment.
    
    Finally we will click the \inlineCode{Load Spec} button and navigate to our random specification file \inlineCode{random_specification.npz} from SDynPy. The specification should then appear in the plots.
    
    With these parameters selected, we can click the Initialize Environments button to proceed.  Figure \ref{fig:examplesdynpyrandomvibrationparameters} shows the parameters in the Rattlesnake \ac{GUI}.
    
    \begin{figure}[H]
        \centering
        \includegraphics[width=\linewidth]{figures/example_sdynpy_random_vibration_parameters}
        \caption{Random vibration parameters specified in the Rattlesnake \ac{GUI}.}
        \label{fig:examplesdynpyrandomvibrationparameters}
    \end{figure}
    
    At this point, we must run the \inlineCode{System Identification} phase of the controller. We will set the Level to \inlineCode{0.1 V RMS}, then click \inlineCode{Start} to start the system identification. A noise floor check will first be run (noise will be zero for a synthetic test), then the system identification follows. Rattlesnake will determine the transfer functions between the two shaker drive forces and the four accelerometer channels. When the specified number of averages have been obtained, the System Identification will stop automatically.  Figure \ref{fig:examplesdynpyrandomvibrationsystemid} shows the system identification phase of the controller.
    
    \begin{figure}[H]
        \centering
        \includegraphics[width=\linewidth]{figures/example_sdynpy_random_vibration_system_id}
        \caption{System identification for the Random Vibration environment.}
        \label{fig:examplesdynpyrandomvibrationsystemid}
    \end{figure}
    
    On the \inlineCode{Test Prediction} tab, Rattlesnake will make predictions of the voltage that it will need to output to achieve the control specified, as well as the responses predicted to be obtained. Here it is important to note that the output signals are not too large for the shakers, though with synthetic data, there is no real risk of damaging anything if they are too large.  Figure \ref{fig:examplesdynpyrandomvibrationpredictions} shows the predictions for this test.
    
    \begin{figure}[H]
        \centering
        \includegraphics[width=\linewidth]{figures/example_sdynpy_random_vibration_predictions}
        \caption{Test predictions for the random vibration environment}
        \label{fig:examplesdynpyrandomvibrationpredictions}
    \end{figure}
    
    If the predictions look good, we can proceed to the Test Profile tab, which we will again leave blank and click the Initialize Test Profile button to proceed to the Run Test tab.
    
    \subsection{Running a Random Vibration Test}
    On the \inlineCode{Run Test} tab, we see a different interface than the one we saw for the modal test. We can click the \inlineCode{Arm Data Acquisition} button to enable the Random environment.
    
    To start, we will set the \inlineCode{Current Test Level} to \inlineCode{-12} dB to ensure when we start the test we don't break anything (again, a best practice; there is no risk to break anything with a synthetic test). We will then step up slowly to the full \inlineCode{0} dB level test.
    
    We can then click the \inlineCode{Start Environment} button to start the test.

    As the test is running, the main Rattlesnake interface will show the sum of \ac{APSD}s, which is the trace of the \ac{CPSD} matrix and can be thought of as an average response of the test (Figure \ref{fig:examplesdynpyrandomvibrationsumasds}). To visualize individual channels, windows can be created in the \inlineCode{Data Display} section of the window (figure \ref{fig:examplesdynpyrandomvibrationindividualchannels}).
    
    \begin{figure}[H]
        \centering
        \includegraphics[width=\linewidth]{figures/example_sdynpy_random_vibration_sum_asds}
        \caption{Rattlesnake \ac{GUI} showing the average test response}
        \label{fig:examplesdynpyrandomvibrationsumasds}
    \end{figure}
    
    \begin{figure}[H]
        \centering
        \includegraphics[width=\linewidth]{figures/example_sdynpy_random_vibration_individual_channels}
        \caption{Individual channel \ac{APSD}s showing how well each channel is matching the test level.}
        \label{fig:examplesdynpyrandomvibrationindividualchannels}
    \end{figure}
    
    If the controller is controlling sufficiently well, we can step up the level to \inlineCode{0} dB for the amount of time specified by the test.  Clearly, Rattlesnake has obtained relatively good control for this test.
    
    To save data from the test, we can either stream time data using the same approach used previously, or we can simply use the \inlineCode{Save Current Spectral Data} to save out the current realizations of the control \ac{CPSD}s, which we will call \inlineCode{random_vibration_spectral_results.nc4}. This is the quickest way to do comparisons between the specification and the control achieved. We can load this file with SDynPy and use SDynPy to make comparisons between the test and specification.  The results are shown in Figure \ref{fig:examplesdynpyrandomvibrationerrorsummary}.
    
    \begin{lstlisting}[language=Python]
control, spec, drives = sdpy.rattlesnake.read_random_spectral_data(
    'random_vibration_spectral_results.nc4')
spec = spec.extract_elements_by_abscissa(20,2000)
control = control.extract_elements_by_abscissa(20,2000)
spec.error_summary(figure_kwargs={'figsize':(10,6)},Control=control)\end{lstlisting}
    
    \begin{figure}[H]
        \centering
        \includegraphics[width=0.7\linewidth]{figures/example_sdynpy_random_vibration_error_summary}
        \caption{Error summary between the test and specification for the random vibration control.}
        \label{fig:examplesdynpyrandomvibrationerrorsummary}
    \end{figure}
    
    We can then \inlineCode{Stop Environment}, \inlineCode{Disarm Data Acquisition} and close the software.
    
    \subsection{Setting up the Transient Vibration Test}
    
    The last demonstration will show how we can perform a MIMO Transient environment test. Open Rattlesnake and select the \inlineCode{MIMO Transient} environment as shown in Figure \ref{fig:examplesdynpymimotransientselection}.
    
    \begin{figure}[H]
        \centering
        \includegraphics[width=0.5\linewidth]{figures/example_cdaq_mimo_transient_selection}
        \caption{Selecting the \ac{MIMO} Transient environment}
        \label{fig:examplesdynpymimotransientselection}
    \end{figure}
    
    We will set the \inlineCode{Data Acquisition Setup} tab identical to what was used before, as shown in Figure \ref{fig:examplesdynpyvibrationdataacquisitionsetup}. Click the \inlineCode{Initialize Data Acquisition} button to proceed.
    
    On the \inlineCode{Environment Definition tab}, we will again \inlineCode{Load} a control law file. This time we will load the \inlineCode{transient_control_laws.py} file, which contains transient control laws for Rattlesnake. We will select the \inlineCode{pseudoinverse_control} function.
    
    In the \inlineCode{Control Channels} section, we will again select all four accelerometer channels as the control channels.
    
    Finally, we will select the \inlineCode{Load Signal} button to load in our specification file \inlineCode{transient_specification.npz}.  Figure \ref{fig:examplesdynpytransientparameters} shows these parameters set in the Rattlesnake \ac{GUI}.  With these parameters set, we will click the \inlineCode{Initialize Environments} button to proceed.
    
    \begin{figure}[H]
        \centering
        \includegraphics[width=\linewidth]{figures/example_sdynpy_transient_parameters}
        \caption{Parameters set for the Transient environment in the Rattlesnake \ac{GUI}.}
        \label{fig:examplesdynpytransientparameters}
    \end{figure}
    
    The Transient environment will also need to perform a \inlineCode{System Identification}. Again set the \inlineCode{Level} to \inlineCode{0.1 V RMS} and click \inlineCode{Start} to start the noise floor and system identification. Here it may be useful to look at the \inlineCode{Impulse Response}, as that is useful to understand how transient control will be performed.  This is shown in Figure \ref{fig:examplesdynpytransientsystemid}.
    
    \begin{figure}[H]
        \centering
        \includegraphics[width=\linewidth]{figures/example_sdynpy_transient_system_id}
        \caption{System Identification for the Transient environment with Impulse Response function shown.}
        \label{fig:examplesdynpytransientsystemid}
    \end{figure}
    
    Once the system identification has been performed, Rattlesnake will again make a prediction as to the results of the test and show it on the \inlineCode{Test Predictions} tab. It will show the peak signal it aims to output, as well as the response it predicts that it will get compared to the desired response.  Figure \ref{fig:examplesdynpytransientpredictions} shows these predictions.
    
    \begin{figure}[H]
        \centering
        \includegraphics[width=\linewidth]{figures/example_sdynpy_transient_prediction}
        \caption{Test predictions and voltage that will be output for the transient control.}
        \label{fig:examplesdynpytransientpredictions}
    \end{figure}
    
    We can then proceed to the \inlineCode{Test Profile} tab, click the \inlineCode{Initialize Profile} button, and we are ready to run the test.
    
    \subsection{Running a Transient Test}
    
    We can now click the \inlineCode{Arm Data Acquisition} button to enable the Transient control functionality. We can set the \inlineCode{Signal Level} to \inlineCode{-12} dB to start so we don't break anything when we run (again, a best practice; there is no risk to break anything with a synthetic test). If it runs successfully, then we can step back up to full level.
    
    Rattlesnake will identify the control signal as it comes through the data acquisition and surround it with a box, as shown in Figure \ref{fig:examplesdynpytransientcontrolresults}.
    
    \begin{figure}[H]
        \centering
        \includegraphics[width=\linewidth]{figures/example_sdynpy_transient_control_results}
        \caption{Running the transient control with the control signal identified.}
        \label{fig:examplesdynpytransientcontrolresults}
    \end{figure}
    
    To visualize individual channels, new windows can be created in the \inlineCode{Data Display} portion of the window.  The four control channels can be seen in Figure \ref{fig:examplesdynpytransientcontrolchannels}
    
    \begin{figure}[H]
        \centering
        \includegraphics[width=\linewidth]{figures/example_sdynpy_transient_control_channels}
        \caption{Control achieved on each control channel in the Transient vibration test.}
        \label{fig:examplesdynpytransientcontrolchannels}
    \end{figure}
    
    To save out data, we can click the \inlineCode{Save Current Control Data} button. We will save the data to the file \inlineCode{transient_vibration_control_results.nc4}. This can be loaded into SDynPy to do comparisons to the desired signal.  These results are shown in Figure \ref{fig:examplesdynpytransientzoom}.
    
    \begin{lstlisting}[language=Python]
control, specification, drives = sdpy.rattlesnake.read_transient_control_data(
    'transient_vibration_control_results.nc4')
        
ax = specification.plot(
    one_axis=False, plot_kwargs = {'linewidth':0.5},
    subplots_kwargs = {'figsize':(10,6),'sharex':True,'sharey':True})
for a,c in zip(ax.flatten(),control):
    c.plot(a,plot_kwargs = {'linewidth':0.5})
# Zoom in on a specific portion
a.set_xlim(0.45,0.48)\end{lstlisting}
    
    \begin{figure}[H]
        \centering
        \includegraphics[width=0.7\linewidth]{figures/example_sdynpy_transient_zoom}
        \caption{Zoom into an impact during the transient control showing good results achieved after the impact.}
        \label{fig:examplesdynpytransientzoom}
    \end{figure}
    
    \section{Summary}
    This appendix has walked through several example tests using virtual hardware with Rattlesnake.  Users are encouraged to use this appendix as a quick-start guide to using Rattlesnake.  Rattlesnake constructed equations of motion from a \inlineCode{System} object from the open-source SDynPy Python package.  These equations of motion were integrated over time by Rattlesnake to simulate a test being performed on that structure.
    
    Modal testing was performed using shaker excitation.  Modes were fit to the data using the SDynPy Python package.  Impact hammer testing could be performed using the combined environments capabilities of Rattlesnake, but that advanced usage would only confuse this beginner-level tutorial, so it was not included in this demonstration.
    
    Vibration testing was performed using simulated modal shakers.  A series of impacts was used to create an environment that the vibration test would aim to replicate.  Both \ac{MIMO} Random and \ac{MIMO} transient vibration tests were run, where the \ac{MIMO} Random attempted to replicate a prescribed \ac{CPSD} matrix, and the \ac{MIMO} Transient attempted to replicate the time history directly.  Both control schemes achieved good results.

    \chapter{Synthetic Example Problem using State Space Matrices}\label{sec:example_state_space}

    One final example will be provided which demonstrates the use of State Space matrices to define a virtual system from which a system of equations can be created and integrated.  The State Space formulation provides the most flexibility to implement arbitrary systems of equations in Rattlesnake.  For example, if the user wanted to substructure a shaker model to a typical mass, stiffness, and damping matrix, this formulation would provide for that flexibility.  However, because of this flexibility, it requires a bit more work and care to set up correctly.
    
    In the current example, we will not do anything as complicated as substructuring, but rather construct a state space model from the typical mass, stiffness, and damping matrices directly.  However, the user should be aware that state-space models can be used to represent more complex configurations.
    
    \section{Creating a State Space Model}
    
    Chapter \ref{sec:rattlesnake_hardware_state_space} provides an overview of the State Space Matrices in equations \eqref{eq:state_eoms} and \eqref{eq:output_eoms}.  They are reproduced here for convenience.
    
    \begin{equation}
        \dot{\mathbf{x}} = \mathbf{A}\mathbf{x} + \mathbf{B}\mathbf{u}
    \end{equation}
    \begin{equation}
        \mathbf{y} = \mathbf{C}\mathbf{x} + \mathbf{D}\mathbf{u}
    \end{equation}
    
    Here $\mathbf{A}$ is the state matrix, $\mathbf{B}$ is the input matrix, $\mathbf{C}$ is the output matrix, $\mathbf{D}$ is the feedthrough matrix, $\mathbf{x}$ is the state vector, $\mathbf{u}$ is the input vector, and $\mathbf{y}$ is the output vector. 
    
    The signals that Rattlesnake outputs are provided to those equations as the inputs $\mathbf{u}$ to the model, and the signals Rattlesnake measures are the output degrees of freedom $\mathbf{y}$.  Users must therefore construct the state space matrices $\mathbf{A}$, $\mathbf{B}$, $\mathbf{C}$, and $\mathbf{D}$ with the correct shape, row order, and column order.  The rows of $\mathbf{u}$ (and therefore the columns of $\mathbf{B}$ and $\mathbf{D}$) must be in the same order as the excitation rows in the channel table (these will be rows with an entry in the \inlineCode{Feedback Device} column).  Similarly, the rows of $\mathbf{y}$ (and therefore the rows of $\mathbf{A}$ and $\mathbf{C}$) must be in the same order as the acquisition rows in the channel table (these will be rows with an entry in the \inlineCode{Physical Device} column).  Note that because Rattlesnake requires its excitation signals to be also measured by the acquisition (e.g., shaker signals must be teed back to acquisition channels for real hardware), there will generally be a row partition of the $\mathbf{C}$ matrix containing zeros and a row partition of the $\mathbf{D}$ matrix containing the identity matrix.  This would result in the input signals $\mathbf{u}$ being directly returned as output degrees of freedom $\mathbf{y}$.  Equation \eqref{eq:excitation_signal_passthrough} shows this, where a partition of zeros in the $\mathbf{C}$ matrix and a partition of identity matrix $\mathbf{I}$ in the $\mathbf{D}$ matrix results in a partition of the output degrees of freedom $\mathbf{y}_i$ being equivalent to the input signals $\mathbf{u}$.
    
    \begin{equation}\label{eq:excitation_signal_passthrough}
        \begin{bmatrix}
            \vdots \\
            \mathbf{y}_{i} \\
            \vdots \\
        \end{bmatrix} = \begin{bmatrix}
            \vdots & \vdots \\
            \mathbf{0} & \mathbf{0} \\
            \vdots & \vdots
        \end{bmatrix}\mathbf{x} + 
        \begin{bmatrix}
            \vdots \\ \mathbf{I} \\ \vdots
        \end{bmatrix}\mathbf{u}
    \end{equation}
    
    \section{Creating the State Space Model}
    Prior to running virtual control using a state space model, that model must be created along with the channel table corresponding to the output matrices $\mathbf{C}$ and $\mathbf{D}$ and the specification that will be controlled.
    
    This example will create a state space model from a beam mass, stiffness, and damping matrix.  In this example, we will use SDynPy to generate the beam mass, stiffness, and damping matrices, from which state space matrices will be derived.  However, one should note that this is not required.  State space matrices can be computed from many different sources.  Note that we will use a reduced order model of the system for faster integration.  This will require incorporating the reduction transformation that the beam uses.
    
    The first step of this process is to create the beam finite element model.  For simplicity, we will use the same code and model used in Appendix \ref{sec:example_sdynpy}.
    
    \begin{lstlisting}[language=Python]
# Import SDynPy module to get structural dynamics functionality
import sdynpy as sdpy
# We'll also import some common Python packages
import numpy as np # NumPy for numeric calculations
import matplotlib.pyplot as plt # Matplotlib for plotting

# Create a system and geometry object using the beam functionality in SDynPy
system,geometry = sdpy.System.beam(
    length = 24 * 0.0254, # meters
    width = 0.75 * 0.0254, # meters
    height = 1.0 * 0.0254, # meters
    num_nodes = 25,
    material = 'steel')
# We want to see which degrees of freedom we have to work with, so we will
# plot the geometry with coordinate labels
geometry.plot_coordinate(label_dofs=True,arrow_scale=0.02)
    
# The finite element model system of equations is too large to integrate
# real-time, so we will create a reduced modal system to integrate based on the
# desired bandwidth of the test
test_bandwidth = 2560 # Hz

# Solve for modes up to 1.5x the bandwidth to lessen modal truncation effects
modes = system.eigensolution(maximum_frequency = test_bandwidth*1.5)

# This also gives us the option to add some damping to the model
modes.damping = 0.005

# Transform to modal system: modal mass, modal stiffness, modal damping
modal_system = modes.system()\end{lstlisting}

    While SDynPy's \inlineCode{System} object has a \inlineCode{to_state_space} method to automatically construct state space matrices from its internal data, we will instead perform these operations ourselves to demonstrate the approach that can be used if a user is constructing these matrices from some other source.  Prior to constructing the state space matrices, we must determine which physical degrees of freedom we wish to use as outputs, and which we wish to use as inputs, as this will determine which rows and columns of the matrices are used.  We will use the same degrees of freedom in our channel table as were used in Chapter \ref{sec:example_sdynpy}. These will include Accelerations as response degrees of freedom and Forces as excitation degrees of freedom.
    
    We will first extract the mass matrix $\mathbf{M}$, stiffness matrix $\mathbf{K}$, damping matrix $\mathbf{C}$, and reduction transformation matrix $\mathbf{\Phi}$.  We will split up the transformation matrix into $\mathbf{\Phi}_{r}$, which is the rows of $\mathbf{\Phi}$ corresponding to the physical response degrees of freedom, and $\mathbf{\Phi}_{i}$, which is the rows of $\mathbf{\Phi}$ corresponding to the input degrees of freedom.  Note that if the system in question is a physical system without any reduction, then the following equations and code still apply, except that the transformation matrix $\mathbf{\Phi}$ will equal the identity matrix.
    
    \begin{lstlisting}[language=Python]
# Select degrees of freedom to use in the test, corresponds to the channel table    
response_coordinates = sdpy.coordinate_array([2,13,19,24],'Y-')
input_coordinates = sdpy.coordinate_array([2,24],'Y+')

# Extract the System matrices, including transformation matrices.
M = modal_system.M
K = modal_system.K
C = modal_system.C
phi_response = modal_system.transformation_matrix_at_coordinates(response_coordinates)
phi_input = modal_system.transformation_matrix_at_coordinates(input_coordinates)

ndofs = modal_system.ndof
tdofs_response = phi_response.shape[0]
tdofs_input = phi_input.shape[0]
    \end{lstlisting}
    
    We can then compute the state space matrices from the system matrices.
    
    The state matrix $\mathbf{A}$ has the equation and code:
    
    \begin{equation}
        \mathbf{A} = \begin{bmatrix}
        \mathbf{0} & \mathbf{I} \\
        -\mathbf{M}^{-1}\mathbf{K} & -\mathbf{M}^{-1}\mathbf{C}
        \end{bmatrix}
    \end{equation}
    
    \begin{lstlisting}[language=Python]
A_state = np.block([[np.zeros((ndofs, ndofs)), np.eye(ndofs)],
                    [-np.linalg.solve(M, K), -np.linalg.solve(M, C)]])\end{lstlisting}
    
    The input matrix $\mathbf{B}$ has the equation and code:
    
    \begin{equation}
        \mathbf{B} = \begin{bmatrix}
        \mathbf{0} \\ \mathbf{M}^{-1}{\mathbf{\Phi}_i}^T
        \end{bmatrix}
    \end{equation}
    
    \begin{lstlisting}[language=Python]
B_state = np.block([[np.zeros((ndofs, tdofs_input))],
                    [np.linalg.solve(M, phi_input.T)]])\end{lstlisting}
                        
    Note the inclusion of the transformation matrix ${\mathbf{\Phi}_i}^T$ which transforms physical inputs $\mathbf{u}$ into the reduced space so it can be multiplied by the reduced mass matrix $\mathbf{M}$.
    
    Our outputs will include accelerations and forces, so we will include those partitions of the output matrix $\mathbf{C}$ and feedforward matrix $\mathbf{D}$.  The equations and code are:
    
    \begin{equation}
        \mathbf{C} = \begin{bmatrix}
        -\mathbf{\Phi}_r\mathbf{M}^{-1}\mathbf{K} & -\mathbf{\Phi}_r\mathbf{M}^{-1}\mathbf{C} \\
        \mathbf{0} & \mathbf{0}
        \end{bmatrix}
    \end{equation}
    \begin{equation}
        \mathbf{D} = \begin{bmatrix}
        \mathbf{\Phi}_r \mathbf{M}^{-1}{\mathbf{\Phi}_i}^T \\ \mathbf{I}
        \end{bmatrix}
    \end{equation}
    
    \begin{lstlisting}[language=Python]
C_state = np.block([[-phi_response@np.linalg.solve(M, K), -phi_response@np.linalg.solve(M, C)],
                    [np.zeros((tdofs_input, ndofs)), np.zeros((tdofs_input, ndofs))]])

D_state = np.block([[phi_response @ np.linalg.solve(M, phi_input.T)],
                    [np.eye(tdofs_input)]])\end{lstlisting}
    
    Note the multiplication by the response transformation $\mathbf{\Phi}_r$ to transform reduced degrees of freedom to physical degrees of freedom for the output for the top acceleration partition.  For the bottom force partition, we have zeros in the $\mathbf{C}$ matrix and identity in the $\mathbf{D}$ matrix, so the physical inputs are passed through directly to outputs without any modifications, as required by Rattlesnake.  This is the virtual equivalent of teeing the excitation signal into a response channel on a data acquisition system.  
    
    We must finally save these matrices to a \inlineCode{.mat} file or \inlineCode{.npz} file with fields \inlineCode{A}, \inlineCode{B}, \inlineCode{C}, and \inlineCode{D}, which can be read by Rattlesnake.
    
    Note that this system is set up with channels in the same order as the channel table shown in Figure \ref{fig:examplesdynpytwoshakerchanneltable}, so we can simply use that channel table for this test.  Be aware that the input and response locations are hard-coded into the state space model, therefore if we wish to change the number of shakers or their locations, or change the number of responses or their locations, we must generate a new state space model to use in Rattlesnake.  The advantage of the SDynPy approach described in Appendix \ref{sec:example_sdynpy} is that that hardware will parse the channel table and construct the system equations of motion based off that data, so we need not worry about the channel bookkeeping or generating separate models.  While the state space approach offers more flexibility, we must be aware that it is our responsibility to ensure the channel table in Rattlesnake matches the state space system we have created.

    \begin{lstlisting}[language=Python]
np.savez('state_space.npz',A = A_state, B = B_state, C = C_state, D = D_state)\end{lstlisting}

    \section{Running a Synthetic Vibration Test}
    Because we have set up our state space system equivalently to the example problem in Appendix \ref{sec:example_sdynpy}, we will utilize the channel table and specifications from that example problem for this example problem as well.  Readers are encouraged to fill out a channel table and run the code from that Appendix to generate the specifications if they have not already.
    
    It is worth noting again that attempting to run the Modal test as described in Section \ref{sec:example_sdynpy_modal} will not work without modification, because the channel table for that test only utilized one shaker, while our current state space object has two shakers.  After completion of this example problem, a good exercise for the reader would be to modify the system object to remove the second shaker, or to add a second shaker to the modal test and run a burst random test with two shakers.
    
    \subsection{Data Acquisition Setup}
    
    Now that the state space matrices are constructed, we will set up a random or transient vibration test using the specifications and channel tables developed in Appendix \ref{sec:example_sdynpy}.  We will open Rattlesnake and select either the \inlineCode{MIMO Random Vibration} environment or \inlineCode{MIMO Transient Environment}.
        
    On the \inlineCode{Data Acquisition Setup} tab, we will \inlineCode{Load Channel Table} to import our Excel channel spreadsheet. We will set the \inlineCode{Hardware Selector} to \inlineCode{State Space Integration...} and select our state space file \inlineCode{state_space.npz}, as well as set the sample \inlineCode{Sample Rate} to \inlineCode{5120}. The other parameters can be left at their default values.  Figure \ref{fig:examplestatespacevibrationdataacquisitionsetup} shows these parameters.
       
   \begin{figure}[H]
       \centering
       \includegraphics[width=\linewidth]{figures/example_state_space_vibration_data_acquisition_setup}
       \caption{Data acquisition setup for the State Space example problem.}
       \label{fig:examplestatespacevibrationdataacquisitionsetup}
   \end{figure}

    Click the \inlineCode{Initialize Data Acquisition} button to proceed.
    
    \subsection{Environment Setup and Running a Synthetic Test}
    
    The remainder of the operations are identical to those described in Section \ref{sec:example_sdynpy_random_vibration}, so readers are directed to that Section to complete the test.  Rattlesnake is designed such that it utilizes largely the same workflow regardless of which data acquisition hardware is being used.  Figures \ref{fig:examplestatespacevibrationenvironmentsetup}--\ref{fig:examplestatespacevibrationruntest} show screenshots from the Rattlesnake software at different stages of the test.  These should look nominally identical to Figures \ref{fig:examplesdynpyrandomvibrationparameters}--\ref{fig:examplesdynpyrandomvibrationsumasds} from Section \ref{sec:example_sdynpy_random_vibration}.

    \begin{figure}[H]
        \centering
        \includegraphics[width=\linewidth]{figures/example_state_space_vibration_environment_setup}
        \caption{Environment Definition for the State Space example problem.}
        \label{fig:examplestatespacevibrationenvironmentsetup}
    \end{figure}
    
    \begin{figure}[H]
        \centering
        \includegraphics[width=\linewidth]{figures/example_state_space_vibration_system_identification}
        \caption{System Identification for the State Space example problem}
        \label{fig:examplestatespacevibrationsystemidentification}
    \end{figure}

    \begin{figure}[H]
        \centering
        \includegraphics[width=\linewidth]{figures/example_state_space_vibration_predictions}
        \caption{Test Predictions for the State Space Example Problem}
        \label{fig:examplestatespacevibrationpredictions}
    \end{figure}

    \begin{figure}[H]
    \centering
    \includegraphics[width=\linewidth]{figures/example_state_space_vibration_run_test}
    \caption{Running a MIMO Random environment with the State Space example problem.}
    \label{fig:examplestatespacevibrationruntest}
    \end{figure}
    
    \section{Summary}
    This appendix has walked through an example test using virtual hardware with Rattlesnake.  Users are encouraged to use this Appendix as a quick-start guide to using Rattlesnake.  Rattlesnake constructed equations of motion from state space matrices.  These equations of motion were integrated over time by Rattlesnake to simulate a test being performed on that structure.
    
    This Appendix only showed the data acquisition setup portion of the controller, as the remaining controller activities are identical to those from Appendix \ref{sec:example_sdynpy}, highlighting Rattlesnake's common workflow regardless of which data acquistion hardware is used.
    % \printindex

    \begin{SANDdistribution}[NM]% or [CA]
%    % \SANDdistCRADA    % If this report is about CRADA work
%    % \SANDdistPatent    % If this report has a Patent Caution or Patent Interest
%    % \SANDdistLDRD    % If this report is about LDRD work
%
%    % External Address Format: {num copies}{Address}
%    \SANDdistExternal{}{}
%    \bigskip
%
%    % The following MUST BE between the external and internal distributions!
%    % \SANDdistClassified % If this report is classified
%
%    % Internal Address Format: {num copies}{Mail stop}{Name}{Org}
%    \SANDdistInternal{}{}{}{}
%
%    % Mail Channel Address Format: {num copies}{Mail Channel}{Name}{Org}
%    \SANDdistInternalM{}{}{}{}
    \end{SANDdistribution}


    % The second printing
    %\begin{SANDreDistribution}
    %    \SANDdistExternal{}{}
    %    \bigskip
    %    \SANDdistInternal{}{}{}{}
    %    \SANDdistInternalM{}{}{}{}
    %\end{SANDreDistribution}

\end{document}
